\documentclass[11pt,a4paper,openany]{book}

% ---------- Packages ----------
\usepackage[utf8]{inputenc}
\usepackage[T1]{fontenc}
\usepackage{lmodern}
\usepackage[margin=2.5cm]{geometry}
\usepackage{amsmath,amssymb,amsthm,mathtools}
\usepackage{bm}
\usepackage{graphicx}
\usepackage{booktabs,array,tabularx}
\usepackage{enumitem}
\usepackage{xcolor}
\usepackage{tikz}
\usetikzlibrary{arrows.meta,calc,decorations.pathmorphing,positioning,3d,shapes.geometric}
\usepackage{pgfplots}
\pgfplotsset{compat=1.17}
\usepackage{listings}
\usepackage[most]{tcolorbox}
\usepackage[hidelinks]{hyperref}
\usepackage[nameinlink,capitalise]{cleveref}

% ---------- Theorems ----------
\theoremstyle{plain}
\newtheorem{theorem}{Theorem}[chapter]
\newtheorem{proposition}[theorem]{Proposition}
\newtheorem{lemma}[theorem]{Lemma}
\newtheorem{corollary}[theorem]{Corollary}
\theoremstyle{definition}
\newtheorem{definition}[theorem]{Definition}
\newtheorem{example}[theorem]{Example}
\newtheorem{exercise}{Exercise}[chapter]
\theoremstyle{remark}
\newtheorem{remark}[theorem]{Remark}

\newenvironment{solution}{\begin{proof}[Solution]}{\end{proof}}

% ---------- Boxes ----------
\tcbset{enhanced,breakable,boxrule=0.6pt,arc=2pt,left=6pt,right=6pt,top=4pt,bottom=4pt}
\newtcolorbox{origin}[1][]{colback=blue!3,colframe=blue!50!black,
  title={Physical origin#1},fonttitle=\bfseries\small}
\newtcolorbox{keypoint}[1][]{colback=green!3,colframe=green!40!black,
  title={Key point#1},fonttitle=\bfseries\small}
\newtcolorbox{caveat}[1][]{colback=red!3,colframe=red!50!black,
  title={Caveat#1},fonttitle=\bfseries\small}

% ---------- Code ----------
\lstset{language=Python,basicstyle=\ttfamily\small,keywordstyle=\color{blue!60!black}\bfseries,
  commentstyle=\color{green!40!black}\itshape,stringstyle=\color{red!60!black},
  frame=single,framerule=0.3pt,rulecolor=\color{gray!60},backgroundcolor=\color{gray!4},
  breaklines=true,columns=fullflexible,showstringspaces=false,upquote=true}

% ---------- Macros ----------
\newcommand{\ket}[1]{\left|#1\right\rangle}
\newcommand{\bra}[1]{\left\langle#1\right|}
\newcommand{\braket}[2]{\left\langle#1\middle|#2\right\rangle}
\newcommand{\ketbra}[2]{\left|#1\middle\rangle\!\middle\langle#2\right|}
\newcommand{\proj}[1]{\ketbra{#1}{#1}}
\newcommand{\mel}[3]{\left\langle#1\middle|#2\middle|#3\right\rangle}
\newcommand{\expval}[1]{\left\langle#1\right\rangle}
\newcommand{\tr}{\operatorname{Tr}}
\newcommand{\Tr}{\operatorname{Tr}}
\newcommand{\id}{I}
\newcommand{\Ha}{\mathcal H}
\newcommand{\R}{\mathbb R}
\newcommand{\C}{\mathbb C}
\newcommand{\N}{\mathbb N}
\newcommand{\Z}{\mathbb Z}
\newcommand{\dd}{\mathrm d}
\newcommand{\e}{\mathrm e}
\newcommand{\ii}{\mathrm i}
\newcommand{\vsig}{\vec\sigma}
\newcommand{\comm}[2]{\left[#1,#2\right]}
\newcommand{\acomm}[2]{\left\{#1,#2\right\}}
\newcommand{\abs}[1]{\left|#1\right|}
\newcommand{\norm}[1]{\left\|#1\right\|}
\newcommand{\CNOT}{\mathrm{CNOT}}
\newcommand{\gcdf}{\operatorname{gcd}}
\newcommand{\diag}{\operatorname{diag}}
\newcommand{\Dis}{\mathcal D}
\newcommand{\order}{\mathcal O}
\newcommand{\starred}{$^\star$}

% Quantum circuits
\usepackage{quantikz}

\numberwithin{equation}{chapter}
\setlist{itemsep=2pt,topsep=4pt}
\emergencystretch=2em


\title{\textbf{Quantum Computing: the Physics}\\[0.6em]
\large Expanded lecture notes, with derivations, proofs and physical models\\[0.3em]
\normalsize Based on the lecture course of D.\,D.~Solnyshkov (IMAPP master, UCA)}
\author{Expanded edition}
\date{\today}

\begin{document}
\frontmatter
\maketitle

\chapter*{How to read these notes}
\addcontentsline{toc}{chapter}{How to read these notes}
These notes follow the structure of the physics part of the \emph{Quantum Computing} course
(D.\,D.~Solnyshkov). Every topic of the original notes is kept. Each one is expanded in four ways:
\begin{enumerate}
  \item \textbf{Where it comes from.} When a Hamiltonian, an equation of motion or a physical system appears,
  we say where it comes from and derive it from simpler principles whenever this can be done in a few pages.
  These parts are marked by a box \emph{``Physical origin''}.
  \item \textbf{Proofs and derivations.} Statements that the original notes give without proof
  (for example the projector $P_\pm=(I\pm\vec n\cdot\vec\sigma)/2$, the Tsirelson bound, the steady state of the
  damped oscillator, the Grover rotation angle, the QFT circuit, the Uhlmann theorem) are proved here.
  \item \textbf{Missing material.} Sections that were empty or only announced in the original notes
  (Ekert's protocol, fidelity and Uhlmann's theorem, universal gate sets, the quantum Fourier transform
  and phase estimation) are written out. Topics needed to make the story complete are added:
  superdense coding, Simon's algorithm, Kraus operators, quantum channels, stabilizer codes, the
  transmon, the Jaynes--Cummings model, and others.
  \item \textbf{Worked examples and exercises.} Each chapter ends with exercises, most with solutions or hints.
\end{enumerate}
Material that goes beyond what is needed for the exam is marked with a star ($\star$) in the section title.

\paragraph{Conventions.}
We write $\hbar$ explicitly in physical models. In the quantum-information chapters we often set $\hbar=1$.
The computational basis of one qubit is $\ket0=(1,0)^T$, $\ket1=(0,1)^T$. Multi-qubit states are written in
the \emph{big-endian} order $\ket{q_0q_1\dots q_{n-1}}=\ket{q_0}\otimes\ket{q_1}\otimes\dots$ (the first qubit is the most significant bit),
except when we explicitly discuss Qiskit, which uses the opposite convention (\cref{sec:endian}).
Pauli matrices are written $X,Y,Z$ or $\sigma_x,\sigma_y,\sigma_z$ interchangeably.
The mathematical background (Hilbert spaces, tensor products, operator exponentials, the Trotter product formula,
density matrices) is treated rigorously in the companion mathematics notes (C.~Kriegler); here we use
it at the level of a physicist.

\tableofcontents

\mainmatter
\part{One qubit}
\chapter{Introduction}\label{ch:intro}

\section{What quantum information is good for}
Quantum informatics, the use of quantum systems to store, transmit and process information, is
useful and important for
\begin{itemize}
  \item \textbf{basic science}: a better understanding of quantum physics, and the solution of quantum many-body problems;
  \item \textbf{communications}: quantum key distribution, quantum networks, quantum repeaters;
  \item \textbf{computing}: algorithms that are faster than any known classical algorithm for specific problems;
  \item \textbf{simulation}: using a controllable quantum system to emulate another one (this was the original motivation of Feynman in 1982);
  \item \textbf{sensing and metrology}: using entanglement and coherence to beat classical limits of precision.
\end{itemize}
Private companies are mostly interested in communications and computing, while academia can use
quantum computing and quantum simulation for research. Quantum machine learning, for example, can be useful for both.

A specialist in quantum computing in a company or in a research team should be able to answer the following questions:
\begin{itemize}
  \item Do we need a quantum computer at all? Is there a known quantum algorithm with an advantage for this problem?
  \item Can it work in our case? This requires an estimate of the error rate of the algorithm and of its implementation on a particular quantum computer
  (number of gates $\times$ error per gate, coherence time, connectivity).
  \item How do we map an optimization problem onto the hardware (QUBO, Ising model, embedding)?
\end{itemize}
In research, quantum computing is of interest for high-energy physics, nanophysics and photonics, i.e.\
the research fields of the laboratories of UCA (LPCA, Institut Pascal).

\section{Why quantum mechanics changes computation}
A classical bit is either $0$ or $1$. A \emph{qubit} is a quantum two-level system, whose state is a vector
\begin{equation}
  \ket\psi=\alpha\ket0+\beta\ket1,\qquad |\alpha|^2+|\beta|^2=1,
\end{equation}
in a two-dimensional complex Hilbert space. Three features of quantum mechanics, absent from classical
physics, are used by quantum computers:
\begin{enumerate}
  \item \textbf{Superposition and interference.} Amplitudes, not probabilities, are added. Amplitudes are complex
  numbers and can cancel each other (destructive interference) or reinforce each other (constructive interference).
  Quantum algorithms are clever arrangements of interference.
  \item \textbf{Entanglement.} The state of $n$ qubits lives in a space of dimension $2^n$, and a generic state cannot be
  described qubit by qubit. Describing it classically requires $2^n$ complex numbers: about $10^{15}$ bytes for $n=46$.
  This exponential size is both the reason why quantum systems are hard to simulate and the resource a quantum computer exploits.
  \item \textbf{Measurement.} The outcome of a measurement is random, and the measurement changes the state.
  At the end of a computation we obtain only $n$ classical bits out of the $2^n$ amplitudes. This severely limits
  what can be gained, and is the reason why only a few problems have a quantum speed-up.
\end{enumerate}

\section{Where we are}
Do not expect too much from this course: it is a general introduction to the topic. And do not expect too much from
quantum computing in the short term either. In spite of the ``Python interface'', we are at the age of ENIAC
(programming by hard-wiring registers), even before the age of the assembly language. At the same time the programming
paradigm is very different from classical programming, which makes things even harder. The positive point is that you can
become the new Knuth or Dijkstra.

The present era is often called NISQ (Noisy Intermediate-Scale Quantum, a term coined by J.~Preskill in 2018):
devices with $10^2$--$10^3$ physical qubits, error rates of $10^{-3}$--$10^{-2}$ per two-qubit gate, and no or only
rudimentary error correction. Error-corrected (``fault-tolerant'') quantum computers capable of running Shor's
algorithm on cryptographically relevant numbers would require of the order of $10^6$ physical qubits.

\section{Plan of the notes}
\begin{itemize}
  \item \Cref{ch:qm}--\cref{ch:applications1}: one qubit. Postulates of quantum mechanics, the two-level system and
  its physical origins, the Bloch sphere, measurement, and first applications (random numbers, no-cloning, BB84).
  \item \Cref{ch:entanglement}--\cref{ch:teleportation}: many qubits. Entanglement, Bell inequalities, the density
  matrix, decoherence and master equations, entanglement criteria, teleportation.
  \item \Cref{ch:gates}--\cref{ch:qec}: gates and algorithms. Universal gates, Qiskit, Deutsch--Jozsa, Grover, QFT and Shor,
  complexity, error correction.
  \item \Cref{ch:annealing}--\cref{ch:qml}: optimization, annealing, adiabatic computing and quantum machine learning.
  \item \Cref{ch:implementations}: physical implementations, with derivations of the Hamiltonians of superconducting
  qubits, trapped ions, cavity QED, NMR, spins in quantum dots, neutral atoms and photons.
\end{itemize}

\chapter{Basics of quantum mechanics}\label{ch:qm}

\section{The postulates}
Quantum computing uses only a small part of quantum mechanics, but it uses it in a very precise way.
It is therefore worth stating the rules explicitly. We consider systems with a finite number of levels.

\begin{enumerate}[label=\textbf{P\arabic*.}]
  \item \textbf{States.} The state of an isolated system is a unit vector $\ket\psi$ of a complex Hilbert space $\Ha$
  (here $\Ha=\C^d$). Two vectors that differ by a global phase, $\ket\psi$ and $\e^{\ii\gamma}\ket\psi$, describe the same physical state.
  \item \textbf{Observables.} A measurable quantity is represented by a Hermitian operator $A=A^\dagger$. Its eigenvalues are real
  and are the possible outcomes of the measurement.
  \item \textbf{Born rule.} If $A=\sum_m a_m P_m$ is the spectral decomposition, the probability of obtaining $a_m$ is
  $p(m)=\mel\psi{P_m}\psi$ and the state just after the measurement is $P_m\ket\psi/\sqrt{p(m)}$ (\cref{ch:measurement}).
  \item \textbf{Dynamics.} The time evolution of an isolated system is given by the Schrödinger equation
  \begin{equation}\label{eq:schrodinger}
    \ii\hbar\frac{\partial}{\partial t}\ket{\psi(t)}=\hat H\ket{\psi(t)},
  \end{equation}
  where the Hamiltonian $\hat H$ is the observable associated with the energy.
  \item \textbf{Composite systems.} The state space of a system made of two parts $A$ and $B$ is the tensor product $\Ha_A\otimes\Ha_B$
  (\cref{ch:entanglement}).
\end{enumerate}

\begin{origin}[: why the Schrödinger equation?]
The Schrödinger equation can be derived from two requirements.
\begin{enumerate}
  \item \emph{Conservation of probability.} The norm of the state must be preserved, so the evolution from time $0$ to $t$
  is a linear map $U(t)$ with $U^\dagger U=I$: a unitary operator. (Linearity is the superposition principle.)
  \item \emph{Time homogeneity and continuity.} For a system with time-independent properties, evolving for $t_1$ and then $t_2$
  is the same as evolving for $t_1+t_2$: $U(t_1+t_2)=U(t_2)U(t_1)$, with $U(0)=I$, and $U(t)$ depends continuously on $t$.
\end{enumerate}
A continuous one-parameter group of unitaries has the form $U(t)=\e^{-\ii Kt}$ with $K$ Hermitian (Stone's theorem; in finite
dimension this is elementary: differentiate at $t=0$, $U(\dd t)\approx I-\ii K\dd t$, and $U^\dagger U=I$ to first order gives $K^\dagger=K$).
Differentiating, $\ii\,\dd\ket{\psi}/\dd t=K\ket\psi$. Identifying the generator of time translations with the energy divided by $\hbar$,
$K=\hat H/\hbar$, as suggested by the Planck--Einstein relation $E=\hbar\omega$ for a plane wave $\e^{-\ii\omega t}$, gives
\cref{eq:schrodinger}. In wave mechanics, with $\hat H=\hat p^2/2m+V(\hat x)$ and $\hat p=-\ii\hbar\,\partial_x$, this is the
familiar $\ii\hbar\,\partial_t\psi=-\frac{\hbar^2}{2m}\partial_x^2\psi+V\psi$.
\end{origin}

\section{Stationary states and the time-evolution operator}
If the Hamiltonian does not depend on time, the variables can be separated. Writing $\psi(x,t)=\phi(x)f(t)$ in
$\ii\hbar\,\partial_t\psi=\hat H\psi$ and dividing by $\phi f$ gives $\ii\hbar\,\dot f/f=(\hat H\phi)/\phi$. The left side depends only
on $t$, the right side only on $x$, so both equal a constant $E$. Hence $f(t)=\e^{-\ii Et/\hbar}$ and $\phi$ satisfies the stationary
Schrödinger equation
\begin{equation}\label{eq:stationary}
  \hat H\psi_n(x)=E_n\psi_n(x).
\end{equation}
Since $\hat H$ is Hermitian, its eigenvectors form an orthonormal basis, and the general solution is the superposition
\begin{equation}\label{eq:superposition}
  \psi(x,t)=\sum_n c_n\psi_n(x)\,\e^{-\ii E_nt/\hbar},\qquad c_n=\braket{\psi_n}{\psi(0)},
\end{equation}
where the coefficients are fixed by the initial condition. The probabilities $|c_n|^2$ do not change in time, but the relative
phases $\e^{-\ii(E_n-E_m)t/\hbar}$ between components do. This is the origin of all quantum oscillations: \emph{the observable dynamics
is governed by the energy differences $E_n-E_m$, the Bohr frequencies.}

An equivalent description uses the time-evolution operator. Since \cref{eq:schrodinger} is a linear first-order equation with constant
coefficients, $\partial_t\psi=-(\ii\hat H/\hbar)\psi$, its solution is
\begin{equation}\label{eq:U}
  \ket{\psi(t)}=\hat U(t)\ket{\psi(0)},\qquad \hat U(t)=\e^{-\ii\hat Ht/\hbar}.
\end{equation}

\subsection{The matrix exponential}
The exponential of a square matrix $X$ is defined by the series
\begin{equation}
  \e^X=\sum_{k=0}^\infty\frac{X^k}{k!},
\end{equation}
which converges for any $X$ since $\norm{X^k}\le\norm X^k$. It satisfies $\frac{\dd}{\dd t}\e^{tX}=X\e^{tX}$, which is why \cref{eq:U}
solves the Schrödinger equation. Its main properties are:
\begin{itemize}
  \item $\e^{X}\e^{Y}=\e^{X+Y}$ \emph{if} $[X,Y]=0$. In general this is false; see the Baker--Campbell--Hausdorff formula below.
  \item $(\e^X)^{-1}=\e^{-X}$, and $(\e^{X})^\dagger=\e^{X^\dagger}$. Hence if $H$ is Hermitian, $\e^{-\ii Ht}$ is unitary:
  $(\e^{-\ii Ht})^\dagger\e^{-\ii Ht}=\e^{\ii Ht}\e^{-\ii Ht}=I$.
  \item $\det\e^X=\e^{\tr X}$ (prove it with the diagonalization below, or with the Jordan form).
  \item If $X=PDP^{-1}$, then $X^k=PD^kP^{-1}$ (all inner $P^{-1}P$ cancel), and therefore $\e^X=P\e^DP^{-1}$.
\end{itemize}
The series is useful for proofs, but a better way to compute the exponential is to diagonalize $X$. For a Hermitian Hamiltonian,
$H=PDP^{-1}$ with $D=\diag(E_1,E_2,\dots)$ and $P$ unitary ($P^{-1}=P^\dagger$); the columns of $P$ are the normalized eigenvectors.
Then
\begin{equation}\label{eq:Udiag}
  U(t)=P\begin{pmatrix}\e^{-\ii E_1t/\hbar}&0&\cdots\\0&\e^{-\ii E_2t/\hbar}&\cdots\\ \vdots&\vdots&\ddots\end{pmatrix}P^{-1},
  \qquad \psi(t)=U(t)\psi(0).
\end{equation}
Reading it from right to left: $P^{-1}$ expresses the initial state in the eigenbasis (it computes the $c_n$), the diagonal matrix
attaches the phases, and $P$ returns to the original basis. This is exactly \cref{eq:superposition}.

Another useful formula is the \emph{spectral} form
\begin{equation}\label{eq:spectralU}
  U(t)=\sum_n\e^{-\ii E_nt/\hbar}\proj{\psi_n},
\end{equation}
which follows from $H=\sum_nE_n\proj{\psi_n}$ and $f(H)=\sum_n f(E_n)\proj{\psi_n}$ for any function $f$.

\subsection{Useful identities}
\begin{proposition}[Exponential of an involution]\label{prop:involution}
If $A^2=I$ (for example any Pauli matrix, or $\vec n\cdot\vec\sigma$ with $|\vec n|=1$), then
\begin{equation}
  \e^{-\ii\theta A}=\cos\theta\,I-\ii\sin\theta\,A.
\end{equation}
\end{proposition}
\begin{proof}
Split the series into even and odd powers, using $A^{2k}=I$ and $A^{2k+1}=A$:
$\e^{-\ii\theta A}=\sum_k\frac{(-\ii\theta)^{2k}}{(2k)!}I+\sum_k\frac{(-\ii\theta)^{2k+1}}{(2k+1)!}A
=\cos\theta\,I-\ii\sin\theta\,A$.
\end{proof}

\begin{proposition}[Hadamard lemma and BCH]\label{prop:bch}
For any matrices $A,B$,
\begin{equation}
  \e^{A}B\e^{-A}=B+[A,B]+\frac1{2!}[A,[A,B]]+\frac1{3!}[A,[A,[A,B]]]+\cdots
\end{equation}
and, if $[A,B]$ commutes with both $A$ and $B$, $\e^{A}\e^{B}=\e^{A+B+\frac12[A,B]}$.
\end{proposition}
\begin{proof}
Let $f(s)=\e^{sA}B\e^{-sA}$. Then $f'(s)=\e^{sA}[A,B]\e^{-sA}$, $f''(s)=\e^{sA}[A,[A,B]]\e^{-sA}$, etc. The Taylor series of $f$ at $s=0$,
evaluated at $s=1$, gives the first formula. For the second, let $g(s)=\e^{sA}\e^{sB}$. Then
$g'(s)=(A+\e^{sA}B\e^{-sA})g(s)=(A+B+s[A,B])g(s)$, because the series stops at the first commutator. Since all operators in the
bracket commute with each other, $g(s)=\e^{s(A+B)+\frac{s^2}2[A,B]}$.
\end{proof}

When $A$ and $B$ do not commute, $\e^{(A+B)t}\neq\e^{At}\e^{Bt}$, but the Lie--Trotter product formula holds:
\begin{equation}\label{eq:trotter}
  \e^{A+B}=\lim_{n\to\infty}\left(\e^{A/n}\e^{B/n}\right)^n,\qquad
  \e^{(A+B)\delta}=\e^{A\delta}\e^{B\delta}+\order(\delta^2).
\end{equation}
This is the basis of \emph{digital quantum simulation}: the evolution under a Hamiltonian $H=\sum_k H_k$, where each $H_k$ acts on only a
few qubits, is approximated by a product of short evolutions under each $H_k$, each implemented by a few gates (Lloyd 1996). The rigorous
proof is in the mathematics notes.

\section{Time-dependent Hamiltonians and the pictures of quantum mechanics}
Quantum gates are implemented by time-dependent fields, so we also need the case $H=H(t)$. The evolution operator still exists, defined by
$\ii\hbar\,\partial_tU(t,t_0)=H(t)U(t,t_0)$, $U(t_0,t_0)=I$, and is unitary, but in general it is \emph{not} $\e^{-\ii\int H\dd t/\hbar}$
because $H(t)$ and $H(t')$ need not commute. Splitting $[t_0,t]$ into small steps,
\begin{equation}
  U(t,t_0)=\lim_{N\to\infty}\e^{-\ii H(t_N)\delta t/\hbar}\cdots\e^{-\ii H(t_1)\delta t/\hbar}\equiv
  \mathcal T\exp\left(-\frac\ii\hbar\int_{t_0}^tH(t')\dd t'\right),
\end{equation}
the \emph{time-ordered exponential}, with later times to the left.

\paragraph{Heisenberg picture.} Mean values $\expval{A}(t)=\mel{\psi(0)}{U^\dagger AU}{\psi(0)}$ can be computed either by evolving the
state (Schrödinger picture) or by evolving the operator, $A_H(t)=U^\dagger(t)AU(t)$ (Heisenberg picture). Differentiating,
\begin{equation}\label{eq:heisenberg}
  \frac{\dd A_H}{\dd t}=\frac\ii\hbar[H,A_H]\qquad\text{(time-independent $H$, no explicit $t$-dependence of $A$)}.
\end{equation}
Taking the mean value gives the \emph{Ehrenfest theorem} $\dd\expval A/\dd t=\frac\ii\hbar\expval{[H,A]}$. In particular, an observable that
commutes with $H$ is conserved. This equation is the starting point of the Bloch-vector equation of \cref{ch:bloch}.

\paragraph{Interaction picture.} When $H=H_0+V(t)$ with $H_0$ simple, one removes the ``trivial'' evolution under $H_0$ by writing
$\ket{\psi(t)}=\e^{-\ii H_0t/\hbar}\ket{\psi_I(t)}$. Substituting into \cref{eq:schrodinger},
\begin{equation}\label{eq:interaction}
  \ii\hbar\,\partial_t\ket{\psi_I}=V_I(t)\ket{\psi_I},\qquad V_I(t)=\e^{\ii H_0t/\hbar}V(t)\e^{-\ii H_0t/\hbar}.
\end{equation}
More generally, any time-dependent unitary change of frame $\ket{\tilde\psi}=R(t)\ket\psi$ gives a new Hamiltonian
\begin{equation}\label{eq:frame}
  \tilde H=RHR^\dagger+\ii\hbar\,\dot RR^\dagger.
\end{equation}
(Proof: $\ii\hbar\,\partial_t(R\psi)=\ii\hbar\dot R\psi+R H\psi=(\ii\hbar\dot RR^\dagger+RHR^\dagger)R\psi$.) The extra term
$\ii\hbar\dot RR^\dagger$ is the quantum analogue of the fictitious forces in a rotating frame. We use \cref{eq:frame} to derive the
rotating-wave approximation in \cref{sec:rabi}.

\section{Commutators and the uncertainty relation}
\begin{theorem}[Robertson uncertainty relation]
For any state $\ket\psi$ and observables $A$, $B$, with $\Delta A^2=\expval{A^2}-\expval{A}^2$,
\begin{equation}
  \Delta A\,\Delta B\ \ge\ \tfrac12\abs{\expval{[A,B]}}.
\end{equation}
\end{theorem}
\begin{proof}
Let $a=A-\expval A$, $b=B-\expval B$, $\ket f=a\ket\psi$, $\ket g=b\ket\psi$. By Cauchy--Schwarz,
$\Delta A^2\Delta B^2=\braket ff\braket gg\ge|\braket fg|^2\ge(\operatorname{Im}\braket fg)^2$. Now
$\braket fg-\braket gf=\expval{ab-ba}=\expval{[A,B]}$, so $\operatorname{Im}\braket fg=\expval{[A,B]}/(2\ii)$.
\end{proof}
For the position and momentum, $[x,p]=\ii\hbar$ gives $\Delta x\Delta p\ge\hbar/2$. For a spin $1/2$, $[\sigma_x,\sigma_y]=2\ii\sigma_z$ gives
$\Delta\sigma_x\Delta\sigma_y\ge|\expval{\sigma_z}|$: a qubit cannot have well-defined values of two orthogonal spin components.
This is precisely what makes the two bases of the BB84 protocol (\cref{sec:bb84}) incompatible.

\section{Exercises}
\begin{exercise}
Show that $\det\e^X=\e^{\tr X}$ for a diagonalizable $X$. Deduce that $\e^{-\ii Ht}$ has determinant of modulus one when $H$ is Hermitian.
\end{exercise}
\begin{exercise}
Using \cref{prop:bch}, show that $\e^{\ii\theta Z/2}X\e^{-\ii\theta Z/2}=\cos\theta\,X-\sin\theta\,Y$. Interpret it as a rotation of the
$x$ axis about the $z$ axis.
\end{exercise}
\begin{exercise}
Show that the error of the first-order Trotter formula is $\e^{(A+B)\delta}-\e^{A\delta}\e^{B\delta}=-\tfrac{\delta^2}{2}[A,B]+\order(\delta^3)$,
and that the symmetric formula $\e^{A\delta/2}\e^{B\delta}\e^{A\delta/2}$ has an error $\order(\delta^3)$.
\end{exercise}

\chapter{A qubit: the two-level system}\label{ch:twolevel}

A qubit is any quantum system whose dynamics can be restricted to two states. Before treating the general case,
we solve the simplest example of the original notes and then explain where such Hamiltonians come from in nature.

\section{The simplest coupled two-level system}\label{sec:simplest}
Consider the Hamiltonian
\begin{equation}\label{eq:Hsimple}
  H=\begin{pmatrix}0&\hbar\omega/2\\ \hbar\omega/2&0\end{pmatrix}=\frac{\hbar\omega}2\sigma_x,
\end{equation}
which describes two states of equal energy (set to zero) coupled by a matrix element $\hbar\omega/2$.

\paragraph{Eigenvalues.} The characteristic equation is $\det(H-E)=E^2-(\hbar\omega/2)^2=0$, so $E_\pm=\pm\hbar\omega/2$.
\paragraph{Eigenvectors.} For $E_+=\hbar\omega/2$,
\begin{equation}
  \begin{pmatrix}0&\hbar\omega/2\\ \hbar\omega/2&0\end{pmatrix}\begin{pmatrix}\psi_1\\ \psi_2\end{pmatrix}
  =\frac{\hbar\omega}2\begin{pmatrix}\psi_1\\ \psi_2\end{pmatrix}
  \ \Longrightarrow\ \psi_2=\psi_1\ \Longrightarrow\ \psi_+=\frac1{\sqrt2}\begin{pmatrix}1\\1\end{pmatrix},
\end{equation}
and the eigenvector for $E_-=-\hbar\omega/2$ is orthogonal to it, $\psi_-=(1,-1)^T/\sqrt2$. Hence
\begin{equation}
  P=\frac1{\sqrt2}\begin{pmatrix}1&1\\1&-1\end{pmatrix}=P^{-1},
\end{equation}
(the Hadamard matrix, which is its own inverse).
\paragraph{Evolution operator.} By \cref{eq:Udiag},
\begin{equation}\label{eq:Usimple}
  U(t)=P\begin{pmatrix}\e^{-\ii\omega t/2}&0\\0&\e^{\ii\omega t/2}\end{pmatrix}P^{-1}
  =\begin{pmatrix}\cos\frac{\omega t}2&-\ii\sin\frac{\omega t}2\\-\ii\sin\frac{\omega t}2&\cos\frac{\omega t}2\end{pmatrix}
  =\cos\tfrac{\omega t}2\,I-\ii\sin\tfrac{\omega t}2\,\sigma_x .
\end{equation}
Let us check the matrix product explicitly for the $(1,1)$ element: $\frac12(\e^{-\ii\omega t/2}+\e^{\ii\omega t/2})=\cos\frac{\omega t}2$;
for the $(1,2)$ element: $\frac12(\e^{-\ii\omega t/2}-\e^{\ii\omega t/2})=-\ii\sin\frac{\omega t}2$. The last form also follows directly
from \cref{prop:involution} since $\sigma_x^2=I$.

\paragraph{Dynamics.} For $\psi(0)=(1,0)^T$,
\begin{equation}
  \psi(t)=\begin{pmatrix}\cos\frac{\omega t}2\\-\ii\sin\frac{\omega t}2\end{pmatrix},\qquad
  P_1(t)=\abs{\begin{pmatrix}1&0\end{pmatrix}\psi(t)}^2=\cos^2\frac{\omega t}2=\frac{1+\cos\omega t}2.
\end{equation}
The system oscillates between the two states at the frequency $\omega=(E_+-E_-)/\hbar$, the Bohr frequency of the
two eigenstates (\cref{fig:rabi}). The probability is periodic with period $2\pi/\omega$, while the state vector itself has period $4\pi/\omega$:
$U(2\pi/\omega)=-I$. This sign is a global phase and is not observable here, but it becomes observable when the two-level system is
part of an interferometer (the famous $4\pi$ periodicity of spinors, observed with neutrons in 1975).

\begin{figure}[h]\centering
\begin{tikzpicture}
\begin{axis}[width=10cm,height=4.2cm,xlabel={$\omega t$},ylabel={$P_1(t)$},xmin=0,xmax=12.57,ymin=0,ymax=1.05,
  xtick={0,3.1416,6.2832,9.4248,12.566},xticklabels={$0$,$\pi$,$2\pi$,$3\pi$,$4\pi$},ytick={0,0.5,1}]
\addplot[blue,thick,domain=0:12.57,samples=200]{cos(deg(x/2))^2};
\addplot[red,dashed,thick,domain=0:12.57,samples=200]{0.5*sin(deg(sqrt(2)*x/2))^2};
\end{axis}\end{tikzpicture}
\caption{Solid: resonant oscillation $P_1=\cos^2(\omega t/2)$. Dashed: probability to reach the other state for a detuned
system with $\Delta=\Omega$, \cref{eq:rabiformula}: the oscillation is faster and incomplete.}\label{fig:rabi}
\end{figure}

\paragraph{A quantum gate.} At $\omega t=\pi$, $U=-\ii\sigma_x$: up to the global phase $-\ii$, the evolution has implemented the NOT gate
$X=\sigma_x$, exchanging $\ket0$ and $\ket1$. A ``$\pi$-pulse'' is a NOT gate. At $\omega t=\pi/2$ we get
$U=(I-\ii\sigma_x)/\sqrt2$, which transforms $\ket0$ into the equal superposition $(\ket0-\ii\ket1)/\sqrt2$: a ``$\pi/2$-pulse'' creates
superpositions. All single-qubit gates are realised in this way, by letting the qubit evolve under a chosen Hamiltonian for a chosen
time. In the Bloch-sphere language of \cref{ch:bloch}, $U(t)$ is a rotation of the pseudospin about the $x$ axis by the angle $\omega t$,
which brings $\vec S=S_z\vec e_z$ to $-S_z\vec e_z$ at $\omega t=\pi$.

\section{Physical origin of two-level Hamiltonians}
Where do Hamiltonians like \cref{eq:Hsimple} come from? Real systems have infinitely many levels. A two-level description is an
\emph{effective model}, valid when only two levels are involved in the dynamics. Here are three fundamental examples.

\subsection{A particle in a double well: the ammonia molecule}\label{sec:doublewell}
\begin{origin}[: tunnelling in a double well]
Consider a particle in a symmetric double-well potential $V(x)$, with two minima at $\pm a$ separated by a barrier. If the barrier
is high, the low-energy states are well described by the two states $\ket L$ and $\ket R$ localized in the left and right wells
(for example the ground states of each well considered alone, of energy $E_0$). We look for eigenstates in the form
$\ket\psi=c_L\ket L+c_R\ket R$ (a variational, or ``tight-binding'', ansatz). Projecting $H\ket\psi=E\ket\psi$ onto $\bra L$ and $\bra R$,
and neglecting the small overlap $\braket LR$,
\begin{equation}
  H_{\rm eff}=\begin{pmatrix}\mel LHL&\mel LHR\\ \mel RHL&\mel RHR\end{pmatrix}=\begin{pmatrix}E_0&-A\\-A&E_0\end{pmatrix}.
\end{equation}
The diagonal elements are equal by symmetry. The off-diagonal element $-A$ is the \emph{tunnelling amplitude}: it is non-zero because the
wave function of $\ket L$ has an exponentially small tail under the barrier and in the right well. A WKB estimate gives
$A\simeq\frac{\hbar\omega_0}{\pi}\exp\left(-\frac1\hbar\int_{-b}^{b}\sqrt{2m(V(x)-E_0)}\,\dd x\right)$, where $\omega_0$ is the oscillation
frequency in one well and $\pm b$ are the classical turning points: the tunnelling is exponentially sensitive to the height and width of the barrier.

Dropping $E_0$ (a global phase), $H_{\rm eff}=-A\sigma_x$, i.e.\ \cref{eq:Hsimple} with $\hbar\omega/2=-A$. The eigenstates are the
symmetric and antisymmetric combinations $(\ket L\pm\ket R)/\sqrt2$, split by $2A$. A particle initially in the left well oscillates
between the two wells at the frequency $2A/\hbar$.

\emph{The ammonia molecule} NH$_3$ is the textbook example (Feynman Lectures, vol.~III, chap.~8--9). The nitrogen atom can be above or
below the plane of the three hydrogen atoms: these are $\ket L$ and $\ket R$. Tunnelling of the nitrogen through the plane splits the
ground state into two levels separated by $2A/h\approx23.87$\,GHz. This transition was used in the first maser (Townes, 1954). Feynman,
Vernon and Hellwarth introduced the geometrical (pseudospin) representation of the next chapter precisely to describe masers \cite{FVH1957}.
\end{origin}

The same model describes an electron in a double quantum dot (charge qubit), a Cooper pair box (\cref{ch:implementations}),
a flux qubit (clockwise vs anticlockwise current), or a photon in two coupled waveguides (directional coupler).

\subsection{A spin 1/2 in a magnetic field}\label{sec:spin}
\begin{origin}[: the Zeeman Hamiltonian]
A particle with spin $\vec S=\frac\hbar2\vec\sigma$ has a magnetic moment $\vec\mu=\gamma\vec S$, where $\gamma$ is the gyromagnetic ratio
($\gamma=-g_e e/2m_e$ with $g_e\approx2$ for the electron, $\gamma/2\pi=42.58$\,MHz/T for the proton). As for a classical magnetic dipole, its
energy in a magnetic field $\vec B$ is
\begin{equation}
  H=-\vec\mu\cdot\vec B=-\frac{\hbar\gamma}2\,\vec B\cdot\vec\sigma.
\end{equation}
This is the most general traceless $2\times2$ Hermitian matrix (\cref{sec:pauli}): any two-level Hamiltonian is formally that of a spin in an
effective magnetic field. A field along $z$ splits the levels by $\hbar\gamma B$ (Zeeman effect); a field along $x$ couples them as in
\cref{eq:Hsimple}. The spin precesses about $\vec B$ at the Larmor frequency $\omega_L=\gamma B$. Nuclear magnetic resonance
(\cref{sec:nmr}) and electron spin qubits in quantum dots use exactly this Hamiltonian.
\end{origin}

\subsection{An atom driven by a laser: Rabi oscillations}\label{sec:rabi}
\begin{origin}[: the electric-dipole Hamiltonian and the rotating-wave approximation]
Consider two levels $\ket g=\ket0$ and $\ket e=\ket1$ of an atom, with energy difference $\hbar\omega_0$, and an electric field
$\vec E(t)=\vec E_0\cos\omega t$ of a laser. Since the wavelength (${\sim}500$\,nm) is much larger than the atom (${\sim}0.1$\,nm), the field is
uniform over the atom (\emph{dipole approximation}) and the coupling is $V=-\hat{\vec d}\cdot\vec E(t)$, where $\hat{\vec d}=-e\hat{\vec r}$ is the
electric dipole operator. (It is obtained from the minimal-coupling Hamiltonian $(\vec p+e\vec A)^2/2m$ by a gauge transformation, the
Göppert-Mayer transformation.) For atomic states of well-defined parity, $\mel g{\vec r}g=\mel e{\vec r}e=0$, so only the off-diagonal elements
$\vec d_{eg}=\mel e{\hat{\vec d}}g$ survive. Choosing the phases of the states so that $\vec d_{eg}\cdot\vec E_0$ is real, and setting
$\hbar\Omega=-\vec d_{eg}\cdot\vec E_0$,
\begin{equation}\label{eq:Hdriven}
  H(t)=-\frac{\hbar\omega_0}2Z+\hbar\Omega\cos(\omega t)\,X,
\end{equation}
where $Z=\proj0-\proj1$, so that $\ket0=\ket g$ has energy $-\hbar\omega_0/2$.

This Hamiltonian is time dependent. We go to a frame rotating at the laser frequency, $\ket{\tilde\psi}=R(t)\ket\psi$ with
$R=\e^{-\ii\omega tZ/2}$. By \cref{eq:frame}, $\ii\hbar\dot RR^\dagger=\frac{\hbar\omega}2Z$, and
$RXR^\dagger=\cos(\omega t)X+\sin(\omega t)Y$ (\cref{prop:bch}). Hence
\begin{equation}
  \tilde H=-\frac{\hbar\Delta}2Z+\frac{\hbar\Omega}2\Big[X\big(1+\cos2\omega t\big)+Y\sin2\omega t\Big],\qquad\Delta=\omega_0-\omega.
\end{equation}
Near resonance, $|\Delta|,\Omega\ll\omega$, the terms oscillating at $2\omega$ average to zero over the slow dynamics, which takes place on the
time scale $1/\Omega$. Neglecting them is the \emph{rotating-wave approximation} (RWA), and gives a time-independent Hamiltonian:
\begin{equation}\label{eq:HRWA}
  \tilde H_{\rm RWA}=\frac\hbar2\left(-\Delta\,Z+\Omega\,X\right).
\end{equation}
The neglected terms give a small shift of the resonance, the Bloch--Siegert shift $\Omega^2/4\omega$.
\end{origin}

\paragraph{Rabi formula.} $\tilde H_{\rm RWA}=\frac{\hbar\Omega_R}2\,\vec n\cdot\vec\sigma$ with the \emph{generalized Rabi frequency}
$\Omega_R=\sqrt{\Omega^2+\Delta^2}$ and the unit vector $\vec n=(\Omega,0,-\Delta)/\Omega_R$. By \cref{prop:involution},
\begin{equation}
  \tilde U(t)=\cos\frac{\Omega_Rt}2\,I-\ii\sin\frac{\Omega_Rt}2\,\vec n\cdot\vec\sigma,
\end{equation}
and the probability to find the atom excited, starting from $\ket g$, is $|\mel1{\tilde U}0|^2$ (the frame change $R$ only adds phases):
\begin{equation}\label{eq:rabiformula}
  P_e(t)=\frac{\Omega^2}{\Omega^2+\Delta^2}\,\sin^2\frac{\sqrt{\Omega^2+\Delta^2}\,t}2.
\end{equation}
This is Rabi's formula (I.~I.~Rabi, 1937, for spins in magnetic fields). At resonance ($\Delta=0$) the transfer is complete; off
resonance it is faster and incomplete (\cref{fig:rabi}). The Lorentzian prefactor $\Omega^2/(\Omega^2+\Delta^2)$ is the line shape of
the transition: it has a width $\sim\Omega$ (power broadening). For a qubit, the lesson is that gates are made by \emph{resonant pulses}:
the frequency $\omega$ selects which transition is driven, and the pulse area $\Omega t$ selects the rotation angle. The phase of the
drive ($\cos\omega t$ vs $\sin\omega t$) selects the rotation axis ($X$ or $Y$), as you can check by repeating the derivation with
$\sin\omega t$.

\begin{keypoint}
All single-qubit gates in the laboratory are rotations generated by Hamiltonians of the form \cref{eq:HRWA}. The original example
\cref{eq:Hsimple} is the resonant case $\Delta=0$, $\Omega=\omega$.
\end{keypoint}

\section{Exercises}
\begin{exercise}
Derive \cref{eq:rabiformula} directly by solving the coupled equations for the amplitudes $c_g(t)$, $c_e(t)$ in the rotating frame.
\end{exercise}
\begin{exercise}
For the ammonia molecule, estimate $A$ in eV from $2A/h=23.87$\,GHz. Compare with the thermal energy at room temperature. What fraction
of molecules is in the upper state at $300$\,K? (Answer: $A\approx4.9\times10^{-5}$\,eV, nearly one half.) Why does this make a maser
possible only after a state selection (Townes used an electrostatic focuser)?
\end{exercise}
\begin{exercise}
Repeat the RWA derivation with the drive $\hbar\Omega\sin(\omega t+\phi)X$. Show that the effective Hamiltonian is
$\frac\hbar2(-\Delta Z+\Omega(\sin\phi\,X+\cos\phi\,Y))$, so that the phase of the microwave or laser selects the rotation axis in the $xy$ plane.
\end{exercise}

\chapter{The Bloch sphere and single-qubit gates}\label{ch:bloch}

The geometrical representation of a two-level system as a vector on a sphere appeared as the \emph{pseudospin} representation of
Feynman, Vernon and Hellwarth \cite{FVH1957}, building on F.~Bloch's equations for nuclear spins (1946). For light, the same
construction is the Poincaré sphere (H.~Poincaré, 1892) of polarization states.

\section{Pauli matrices}\label{sec:pauli}
\subsection{Algebra}
The identity and the Pauli matrices are
\begin{equation}
  \sigma_0=\begin{pmatrix}1&0\\0&1\end{pmatrix},\quad
  \sigma_x=\begin{pmatrix}0&1\\1&0\end{pmatrix},\quad
  \sigma_y=\begin{pmatrix}0&-\ii\\ \ii&0\end{pmatrix},\quad
  \sigma_z=\begin{pmatrix}1&0\\0&-1\end{pmatrix}.
\end{equation}
They are Hermitian, unitary and traceless (except $\sigma_0$), and square to the identity. Their products follow from a single identity,
\begin{equation}\label{eq:pauliproduct}
  \sigma_j\sigma_k=\delta_{jk}\,I+\ii\sum_l\varepsilon_{jkl}\sigma_l,\qquad j,k,l\in\{x,y,z\},
\end{equation}
where $\varepsilon_{jkl}$ is the Levi-Civita symbol. For example $\sigma_x\sigma_y=\ii\sigma_z$; check it by multiplying the matrices.
Consequences:
\begin{equation}
  [\sigma_j,\sigma_k]=2\ii\varepsilon_{jkl}\sigma_l,\qquad\{\sigma_j,\sigma_k\}=2\delta_{jk}I,\qquad
  (\vec a\cdot\vec\sigma)(\vec b\cdot\vec\sigma)=(\vec a\cdot\vec b)\,I+\ii(\vec a\times\vec b)\cdot\vec\sigma.
\end{equation}
The last identity, valid for real vectors, follows by multiplying \cref{eq:pauliproduct} by $a_jb_k$ and summing. In particular
$(\vec n\cdot\vec\sigma)^2=I$ for a unit vector $\vec n$, so the eigenvalues of $\vec n\cdot\vec\sigma$ are $\pm1$.

\subsection{Decomposition of a $2\times2$ Hamiltonian}
The four matrices are orthogonal for the Hilbert--Schmidt scalar product $\langle A,B\rangle=\tr(A^\dagger B)$:
\begin{equation}
  \tr(\sigma_\alpha\sigma_\beta)=2\delta_{\alpha\beta},\qquad\alpha,\beta\in\{0,x,y,z\},
\end{equation}
which follows from \cref{eq:pauliproduct} and $\tr\sigma_l=0$. Four orthogonal (hence linearly independent) elements of the
four-dimensional space of $2\times2$ complex matrices form a basis. Any matrix can therefore be written
\begin{equation}
  H=\begin{pmatrix}a&b\\c&d\end{pmatrix}=h_0\sigma_0+h_x\sigma_x+h_y\sigma_y+h_z\sigma_z
  =\begin{pmatrix}h_0+h_z&h_x-\ii h_y\\h_x+\ii h_y&h_0-h_z\end{pmatrix},\qquad h_\alpha=\tfrac12\tr(H\sigma_\alpha),
\end{equation}
where the formula for the coefficients is obtained by multiplying by $\sigma_\beta$ and taking the trace. Identifying the elements,
\begin{equation}
  h_0=\frac{a+d}2,\quad h_z=\frac{a-d}2,\quad h_x=\frac{b+c}2,\quad h_y=\frac{\ii(b-c)}2.
\end{equation}
For a Hermitian matrix, $a,d\in\R$ and $c=b^*$, so all four coefficients are real: $h_x=\operatorname{Re}b$, $h_y=-\operatorname{Im}b$.
Writing $\vec h=(h_x,h_y,h_z)$,
\begin{equation}\label{eq:Hpauli}
  H=h_0\sigma_0+\vec h\cdot\vec\sigma,\qquad
  E_\pm=h_0\pm|\vec h|=\frac{a+d}2\pm\sqrt{\left(\frac{a-d}2\right)^2+|b|^2}.
\end{equation}
The eigenvalues follow from $(\vec h\cdot\vec\sigma)^2=|\vec h|^2I$. The term $h_0\sigma_0$ adds only a global phase $\e^{-\ii h_0t/\hbar}$ and does not
affect the dynamics. \emph{The physics of a two-level system is entirely contained in a three-dimensional real vector $\vec h$}: its length gives
the level splitting $2|\vec h|$, its direction gives the eigenstates. For example, for $a=d=0$, $b=c=\hbar\omega/2$, we get
$\vec h=(\hbar\omega/2,0,0)$, $E_\pm=\pm\hbar\omega/2$, i.e.\ $H=\frac{\hbar\omega}2\sigma_x$.

\section{The Bloch sphere}
\subsection{Parametrization of a pure state}
A qubit state $\ket\psi=\alpha\ket0+\beta\ket1$ has four real parameters. Normalization removes one and the global phase another, leaving two,
which can be taken as angles on a sphere. Choosing the global phase so that $\alpha$ is real and non-negative,
\begin{equation}\label{eq:blochstate}
  \ket\psi=\cos\frac\theta2\ket0+\e^{\ii\varphi}\sin\frac\theta2\ket1,\qquad\theta\in[0,\pi],\ \varphi\in[0,2\pi).
\end{equation}
The original notes write the equivalent form $(\cos\frac\theta2\e^{-\ii\varphi},\sin\frac\theta2)^T$, which differs by the global phase
$\e^{-\ii\varphi}$ (and describes the same physical state). The point of the unit sphere with polar angle $\theta$ and azimuth $\varphi$,
\begin{equation}
  \vec n=(\sin\theta\cos\varphi,\ \sin\theta\sin\varphi,\ \cos\theta),
\end{equation}
is the \emph{Bloch vector}. The north pole is $\ket0$, the south pole $\ket1$, the equator contains the equal superpositions
$\ket\pm=(\ket0\pm\ket1)/\sqrt2$ ($\pm x$) and $\ket{\pm\ii}=(\ket0\pm\ii\ket1)/\sqrt2$ ($\pm y$).

\begin{proposition}\label{prop:blochvector}
The Bloch vector is the vector of mean values of the Pauli matrices, $\vec n=\mel\psi{\vec\sigma}\psi$, and $\ket\psi$ is the eigenvector of
$\vec n\cdot\vec\sigma$ with eigenvalue $+1$.
\end{proposition}
\begin{proof}
With $c=\cos\frac\theta2$, $s=\sin\frac\theta2$:
$\expval{\sigma_z}=c^2-s^2=\cos\theta$;
$\expval{\sigma_x}=\alpha^*\beta+\beta^*\alpha=2cs\cos\varphi=\sin\theta\cos\varphi$;
$\expval{\sigma_y}=-\ii\alpha^*\beta+\ii\beta^*\alpha=2cs\sin\varphi=\sin\theta\sin\varphi$.
For the second statement,
$\vec n\cdot\vec\sigma=\begin{pmatrix}\cos\theta&\sin\theta\e^{-\ii\varphi}\\ \sin\theta\e^{\ii\varphi}&-\cos\theta\end{pmatrix}$, and
the first component of $(\vec n\cdot\vec\sigma)\ket\psi$ is $\cos\theta\,c+\sin\theta\,s=c(\cos\theta+2s^2)=c$, using
$\sin\theta=2sc$ and $\cos\theta=1-2s^2$; the second component is checked in the same way.
\end{proof}

\paragraph{Why half angles?} Orthogonal states, $\braket{\psi}{\psi'}=0$, are at \emph{antipodal} points of the sphere ($\theta'=\pi-\theta$,
$\varphi'=\varphi+\pi$): the angle on the sphere is twice the ``angle'' $\arccos|\braket\psi{\psi'}|$ in Hilbert space. Indeed
\begin{equation}\label{eq:overlapbloch}
  |\braket\psi{\psi'}|^2=\frac{1+\vec n\cdot\vec n'}2=\cos^2\frac{\Theta}2,
\end{equation}
where $\Theta$ is the angle between the Bloch vectors (proof: use \cref{prop:blochvector} and the projector \cref{eq:projector} below).
For photons this explains the factor two between polarization angles and angles on the Poincaré sphere (\cref{sec:photonic}).

\begin{figure}[h]\centering
\begin{tikzpicture}[scale=2.1]
  \draw[thick] (0,0) circle (1);
  \draw[dashed] (1,0) arc (0:180:1 and 0.3);
  \draw (-1,0) arc (180:360:1 and 0.3);
  \draw[->] (0,0)--(0,1.25) node[above]{$z$};
  \draw[->] (0,0)--(-0.6,-0.5) node[below left]{$x$};
  \draw[->] (0,0)--(1.25,0) node[right]{$y$};
  \node[right] at (0.03,1.05) {$\ket0$};
  \node[right] at (0.03,-1.08) {$\ket1$};
  \fill (0,1) circle (0.7pt); \fill (0,-1) circle (0.7pt);
  \draw[->,very thick,red] (0,0)--(0.45,0.62) node[above right]{$\vec n$};
  \draw[dotted] (0.45,0.62)--(0.45,-0.12);
  \draw[dotted] (0,0)--(0.45,-0.12);
  \draw (0,0.35) arc (90:54:0.35); \node at (0.12,0.45) {$\theta$};
  \draw (-0.18,-0.15) arc (220:345:0.25 and 0.1); \node at (0.1,-0.28) {$\varphi$};
  \node[left] at (-0.62,-0.52) {$\ket+$};
\end{tikzpicture}
\caption{The Bloch sphere. A pure qubit state $\cos\frac\theta2\ket0+\e^{\ii\varphi}\sin\frac\theta2\ket1$ corresponds to the unit vector $\vec n$.}
\label{fig:bloch}
\end{figure}

\subsection{Dynamics: the Bloch equation}
The pseudospin representation is based on the fact that the Schrödinger equation $\ii\hbar\partial_t\psi=H\psi$ for a two-level system is
equivalent to the classical equation of a precessing vector,
\begin{equation}\label{eq:precession}
  \dot{\vec S}=\vec S\times\vec\Omega,\qquad\text{with}\qquad
  H=-\frac\hbar2(\Omega_x,\Omega_y,\Omega_z)\cdot\begin{pmatrix}\sigma_x\\ \sigma_y\\ \sigma_z\end{pmatrix}=-\vec\Omega\cdot\hat{\vec S},
\end{equation}
where $\hat{\vec S}=\frac\hbar2\vec\sigma$ is the spin operator and $S_\alpha=\mel\psi{\sigma_\alpha}\psi$ are the pseudospin components
(with or without the factor $\hbar/2$; often just $1/2$, or $1$ as here).
\begin{proof}
By the Ehrenfest theorem (\cref{eq:heisenberg}),
\begin{equation*}
  \dot S_k=\frac\ii\hbar\expval{[H,\sigma_k]}=-\frac\ii2\sum_j\Omega_j\expval{[\sigma_j,\sigma_k]}
  =-\frac\ii2\sum_{j,l}\Omega_j\,2\ii\varepsilon_{jkl}S_l=\sum_{j,l}\varepsilon_{klj}S_l\Omega_j=(\vec S\times\vec\Omega)_k .
\end{equation*}
Conversely, the two real parameters of the Bloch vector determine the state up to a global phase, so the Bloch equation contains all the dynamics.
\end{proof}

\subsection{What precesses around what?}\label{sec:precession}
Equation \eqref{eq:precession} is easy to misread, so let us state precisely what moves and around which axis.
\begin{itemize}
  \item \textbf{The vector that moves} is the Bloch vector (pseudospin) $\vec S(t)=\mel{\psi(t)}{\vec\sigma}{\psi(t)}$, i.e.\ the \emph{state}.
  \item \textbf{The axis} is the vector $\vec\Omega$ (equivalently $\vec h$), which is fixed by the \emph{Hamiltonian}: for a spin it is the direction of the
  magnetic field. It does not move (for a time-independent $H$).
\end{itemize}
\emph{The state precesses around the Hamiltonian axis}, never the other way round. Comparing $H=h_0\sigma_0+\vec h\cdot\vec\sigma$ with $H=-\frac\hbar2\vec\Omega\cdot\vec\sigma$
gives $\vec\Omega=-2\vec h/\hbar$, and \cref{eq:precession} can be rewritten in the two equivalent forms
\begin{equation}\label{eq:precession2}
  \dot{\vec S}=-\vec\Omega\times\vec S=\vec\omega_p\times\vec S,\qquad\vec\omega_p=-\vec\Omega=\frac2\hbar\vec h .
\end{equation}
The equation $\dot{\vec S}=\vec\omega_p\times\vec S$ is the kinematics of a rigid rotation with \emph{angular velocity vector} $\vec\omega_p$ (as $\dot{\vec r}=\vec\omega\times\vec r$
for a point of a rotating solid). Therefore:
\begin{itemize}
  \item the Bloch vector rotates about the axis $\vec h$, \emph{counterclockwise} (right-hand rule: thumb along $\vec h$) or, equivalently, clockwise about $\vec\Omega$;
  \item with the angular frequency $|\vec\omega_p|=2|\vec h|/\hbar=(E_+-E_-)/\hbar$, the Bohr frequency of the two eigenstates.
\end{itemize}

\paragraph{The motion is on a cone.} Two quantities are conserved:
\begin{equation}
  \frac{\dd}{\dd t}|\vec S|^2=2\vec S\cdot(\vec\omega_p\times\vec S)=0,\qquad
  \frac{\dd}{\dd t}(\vec S\cdot\vec h)=(\vec\omega_p\times\vec S)\cdot\vec h=0 .
\end{equation}
The length of the Bloch vector is constant (pure states stay pure), and so is its projection on $\vec h$, i.e.\ the mean energy $\expval H=h_0+\vec h\cdot\vec S$. The tip of $\vec S$
therefore moves on a circle: $\vec S$ sweeps a \emph{cone} of fixed opening angle $\beta_0$ around $\vec h$ (\cref{fig:precession}). In quantum language, the populations of the two
eigenstates, $\cos^2\frac{\beta_0}2$ and $\sin^2\frac{\beta_0}2$, are constant (\cref{eq:overlapbloch}); only their relative phase evolves, at the Bohr frequency, and this
phase \emph{is} the azimuthal angle around the cone.
\begin{itemize}
  \item The two eigenstates of $H$ are the points $\pm\vec h/|\vec h|$ ($\beta_0=0$ or $\pi$): a cone of zero opening, so they do not move. Stationary states are the fixed points of the rotation.
  \item A state perpendicular to $\vec h$ ($\beta_0=\pi/2$) moves on a great circle: this is the case of maximal oscillation of the observables.
\end{itemize}

\begin{figure}[h]\centering
\begin{tikzpicture}[scale=2.2]
  \draw[gray] (0,0) circle (1);
  \draw[gray,dashed] (1,0) arc (0:180:1 and 0.28);
  \draw[gray] (-1,0) arc (180:360:1 and 0.28);
  \draw[->,very thick,blue] (0,0)--(0,1.35) node[above]{$\vec h$ (axis fixed by $H$, e.g.\ $\vec B$)};
  \draw[thick,red!70!black] (0,0.71) ellipse (0.71 and 0.2);
  \draw[->,very thick,red] (0,0)--(0.54,0.58) node[right]{$\vec S(t)$};
  \draw[dashed,red!60] (0,0)--(-0.71,0.71);
  \draw[->,thick,red!70!black] (0.25,0.53) arc (-70:-20:0.5 and 0.15);
  \node[red!70!black,font=\small] at (1.15,0.75) {$\vec\omega_p=\tfrac2\hbar\vec h$};
  \draw (0,0.3) arc (90:47:0.3); \node at (0.1,0.4) {\small$\beta_0$};
  \fill[blue] (0,1) circle (0.8pt) node[left,font=\small]{eigenstate};
  \fill[blue] (0,-1) circle (0.8pt) node[left,font=\small]{eigenstate};
\end{tikzpicture}
\caption{Precession on the Bloch sphere. The axis $\vec h$ is set by the Hamiltonian and stays fixed. The Bloch vector $\vec S$ of the state rotates around it on a cone,
counterclockwise when viewed from the tip of $\vec h$, at the Bohr frequency $(E_+-E_-)/\hbar$. The opening angle $\beta_0$ (i.e.\ the populations of the eigenstates)
is conserved. The eigenstates $\pm\vec h/|\vec h|$ do not move.}\label{fig:precession}
\end{figure}

\paragraph{Explicit solution for an axis along $z$.} Take $\vec h=h_z\vec e_z$, i.e.\ $H=h_z\sigma_z$. Then \cref{eq:precession2} reads $\dot S_x=-\omega_pS_y$,
$\dot S_y=\omega_pS_x$, $\dot S_z=0$, with $\omega_p=2h_z/\hbar$, so
\begin{equation}
  S_x+\ii S_y=\sin\theta\,\e^{\ii(\varphi_0+\omega_pt)},\qquad S_z=\cos\theta=\text{const}.
\end{equation}
The polar angle is fixed and the azimuth grows linearly. The same result follows directly from the Schrödinger equation: the component $\ket0$ (energy $+h_z$) acquires the
phase $\e^{-\ii h_zt/\hbar}$ and $\ket1$ the phase $\e^{+\ii h_zt/\hbar}$, so the relative phase $\varphi$ in \cref{eq:blochstate} grows as $\varphi_0+2h_zt/\hbar$.
\emph{Precession is nothing but the evolution of the relative phase between the two energy eigenstates.}

\paragraph{Example 1: the gate of \cref{sec:simplest}.} For $H=\frac{\hbar\omega}2\sigma_x$, $\vec h=\frac{\hbar\omega}2\vec e_x$ and $\vec\omega_p=\omega\vec e_x$. Starting from
$\ket0$ (north pole), $\vec S$ is perpendicular to the axis and moves on the great circle in the $yz$ plane:
$\vec S(t)=(0,-\sin\omega t,\cos\omega t)$. (Check with $\psi(t)=(\cos\frac{\omega t}2,-\ii\sin\frac{\omega t}2)$: $S_z=\cos^2-\sin^2=\cos\omega t$,
$S_y=2\operatorname{Im}(\alpha^*\beta)=-\sin\omega t$.) The right-hand rotation about $+x$ takes $+z$ towards $-y$, as found. At $\omega t=\pi$ the spin is at the south pole:
the NOT gate is a half-turn around $x$.

\paragraph{Example 2: detuned Rabi oscillations.} For the driven atom in the rotating frame, \cref{eq:HRWA}, $\vec h=\frac\hbar2(\Omega,0,-\Delta)$: the axis is tilted from the
equator by the detuning. Starting at the north pole, $\vec S$ moves on a cone around this tilted axis, with angular frequency $\Omega_R=\sqrt{\Omega^2+\Delta^2}$. If $\beta$
is the angle between the axis and $\vec e_z$, $\cos\beta=-\Delta/\Omega_R$, the cone reaches at most the polar angle $2\beta$, so the maximal excited population is
$\frac{1-\cos2\beta}2=\sin^2\beta=\Omega^2/\Omega_R^2$. The Rabi formula \cref{eq:rabiformula} is pure geometry: only at resonance ($\Delta=0$) is the axis in the equatorial
plane and the cone a great circle through both poles.

\paragraph{A real spin in a magnetic field.} With $H=-\vec\mu\cdot\vec B=-\frac{\hbar\gamma}2\vec B\cdot\vec\sigma$ (\cref{sec:spin}), $\vec\Omega=\gamma\vec B$ and
\cref{eq:precession} is the Larmor equation
\begin{equation}
  \frac{\dd\vec\mu}{\dd t}=\gamma\,\vec\mu\times\vec B=-\gamma\vec B\times\vec\mu .
\end{equation}
The magnetic moment (and the spin) precesses around $\vec B$ at the Larmor frequency $\omega_L=|\gamma|B$. The sense depends on the sign of $\gamma$: for the proton
($\gamma>0$) the precession is clockwise seen from the tip of $\vec B$; for the electron ($\gamma<0$, the moment is antiparallel to the spin) it is counterclockwise.
Classically, this is a gyroscope: the field exerts the torque $\vec\tau=\vec\mu\times\vec B$ on a body with angular momentum $\vec L=\vec\mu/\gamma$, and since $\dd\vec L/\dd t=\vec\tau$
is perpendicular to $\vec L$, the torque does not align $\vec L$ with $\vec B$ (as a compass needle with friction would) but makes it turn around $\vec B$, exactly like a spinning
top in gravity. \emph{A spin never aligns with the field by precession alone}: alignment, i.e.\ relaxation to the ground state, requires energy exchange with an environment,
described by $T_1$ (\cref{ch:open}). Precession conserves energy; relaxation does not.

\paragraph{Rotating frame.} In NMR or ESR, a static field $B_0\vec e_z$ makes the spin precess quickly ($\omega_0=\gamma B_0\sim$ hundreds of MHz), and a weak field $B_1$
rotating in the $xy$ plane at frequency $\omega$ drives transitions. Viewed from a frame rotating with $B_1$, the spin feels the \emph{effective field}
$\vec B_{\rm eff}=(B_0-\omega/\gamma)\vec e_z+B_1\vec e_x$: the rotation of the frame subtracts $\omega/\gamma$ from the static field (this is the Bloch-vector version of
\cref{eq:frame}). At resonance the $z$ component vanishes and the spin precesses slowly around $B_1$ in the equatorial plane, from the north to the south pole: this slow
precession, seen in the lab as a spiral down the sphere, is the Rabi oscillation (``nutation'') of \cref{sec:rabi}, with $\Omega=\gamma B_1$ and $\Delta=\omega_0-\omega$.

\begin{keypoint}
It is very useful to write both the Hamiltonian and the spinor in polar coordinates: the action of any single-qubit gate is then a rotation
of the Bloch sphere, which can be visualized.
\end{keypoint}

\section{Rotations and single-qubit gates}
\subsection{Rotation operators}
Define the rotation by an angle $\alpha$ about the unit axis $\vec n$,
\begin{equation}\label{eq:rotation}
  R_{\vec n}(\alpha)=\e^{-\ii\alpha\,\vec n\cdot\vec\sigma/2}=\cos\frac\alpha2\,I-\ii\sin\frac\alpha2\,\vec n\cdot\vec\sigma,
\end{equation}
where the second form is \cref{prop:involution}. It is the evolution operator of the Hamiltonian $H=\frac{\hbar\omega}2\vec n\cdot\vec\sigma$ for
the time $t=\alpha/\omega$. Explicitly,
\begin{equation}
  R_x(\alpha)=\begin{pmatrix}\cos\frac\alpha2&-\ii\sin\frac\alpha2\\-\ii\sin\frac\alpha2&\cos\frac\alpha2\end{pmatrix},\quad
  R_y(\alpha)=\begin{pmatrix}\cos\frac\alpha2&-\sin\frac\alpha2\\ \sin\frac\alpha2&\cos\frac\alpha2\end{pmatrix},\quad
  R_z(\alpha)=\begin{pmatrix}\e^{-\ii\alpha/2}&0\\0&\e^{\ii\alpha/2}\end{pmatrix}.
\end{equation}
\begin{proposition}
$R_{\vec n}(\alpha)$ rotates the Bloch vector by the angle $\alpha$ about $\vec n$: $R\,(\vec m\cdot\vec\sigma)R^\dagger=(\mathcal R\vec m)\cdot\vec\sigma$,
where $\mathcal R$ is the $3\times3$ rotation matrix.
\end{proposition}
\begin{proof}
It suffices to check it for an infinitesimal rotation, $R\approx I-\ii\frac{\delta\alpha}2\vec n\cdot\vec\sigma$:
$R(\vec m\cdot\vec\sigma)R^\dagger\approx\vec m\cdot\vec\sigma-\frac{\ii\delta\alpha}2[\vec n\cdot\vec\sigma,\vec m\cdot\vec\sigma]
=\vec m\cdot\vec\sigma+\delta\alpha\,(\vec n\times\vec m)\cdot\vec\sigma$, using $[\vec n\cdot\vec\sigma,\vec m\cdot\vec\sigma]=2\ii(\vec n\times\vec m)\cdot\vec\sigma$.
This is the infinitesimal rotation $\vec m\to\vec m+\delta\alpha\,\vec n\times\vec m$, and finite rotations are products of infinitesimal ones.
\end{proof}
Since $R_{\vec n}(2\pi)=-I$ but the rotation of the sphere is the identity, each rotation of $SO(3)$ corresponds to two matrices $\pm R$ of $SU(2)$:
$SU(2)$ is the \emph{double cover} of $SO(3)$. Every $2\times2$ unitary matrix is $\e^{\ii\gamma}R_{\vec n}(\alpha)$ for some global phase $\gamma$, axis and angle.

\subsection{Standard gates}
\begin{center}
\begin{tabular}{@{}llll@{}}
\toprule
Gate & Matrix & Rotation (up to phase) & Action\\ \midrule
$X$ (NOT) & $\begin{pmatrix}0&1\\1&0\end{pmatrix}$ & $\pi$ about $x$ & $\ket0\leftrightarrow\ket1$\\[2ex]
$Y$ & $\begin{pmatrix}0&-\ii\\ \ii&0\end{pmatrix}$ & $\pi$ about $y$ & bit and phase flip\\[2ex]
$Z$ & $\begin{pmatrix}1&0\\0&-1\end{pmatrix}$ & $\pi$ about $z$ & phase flip $\ket1\to-\ket1$\\[2ex]
$H$ (Hadamard) & $\frac1{\sqrt2}\begin{pmatrix}1&1\\1&-1\end{pmatrix}$ & $\pi$ about $(\vec e_x+\vec e_z)/\sqrt2$ & $\ket0\to\ket+$, $\ket1\to\ket-$\\[2ex]
$S$ (phase) & $\begin{pmatrix}1&0\\0&\ii\end{pmatrix}$ & $\pi/2$ about $z$ & $\ket+\to\ket{+\ii}$\\[2ex]
$T$ ($\pi/8$) & $\begin{pmatrix}1&0\\0&\e^{\ii\pi/4}\end{pmatrix}$ & $\pi/4$ about $z$ & \\ \bottomrule
\end{tabular}
\end{center}
For example $H=\frac{X+Z}{\sqrt2}=\ii R_{\vec n}(\pi)$ with $\vec n=(1,0,1)/\sqrt2$: the Hadamard gate exchanges the $x$ and $z$ axes. Hence
$HXH=Z$, $HZH=X$, $HYH=-Y$. Similarly $T=\e^{\ii\pi/8}R_z(\pi/4)$, which explains its historical name ``$\pi/8$ gate''. Note $T^2=S$ and $S^2=Z$.

\begin{theorem}[Euler decomposition]\label{thm:euler}
Any single-qubit unitary can be written as
\begin{equation*}U=\e^{\ii\gamma}R_z(\beta)R_y(\theta)R_z(\delta).\end{equation*}
\end{theorem}
\begin{proof}
A unitary $2\times2$ matrix has orthonormal columns, so it can be written
\begin{equation*}
U=\e^{\ii\gamma}\begin{pmatrix}\e^{-\ii(\beta+\delta)/2}\cos\frac\theta2&-\e^{-\ii(\beta-\delta)/2}\sin\frac\theta2\\
\e^{\ii(\beta-\delta)/2}\sin\frac\theta2&\e^{\ii(\beta+\delta)/2}\cos\frac\theta2\end{pmatrix}
\end{equation*}
for suitable real
$\gamma,\beta,\theta,\delta$ (the moduli of the entries are fixed by the first column and unitarity; four phases remain, of which only three
are independent after the determinant is factored out). Multiplying out $R_z(\beta)R_y(\theta)R_z(\delta)$ gives exactly this matrix.
\end{proof}
Geometrically: any rotation of the sphere can be performed as a rotation about $z$, then about $y$, then about $z$ again (Euler angles).
In hardware, $R_z$ is often ``free'' (a change of reference phase of later pulses, the ``virtual $Z$ gate''), and the IBM basis gates
$\{\texttt{rz},\texttt{sx},\texttt{x}\}$, with $\texttt{sx}=\sqrt X$, suffice for any single-qubit gate.

\section{Exercises}
\begin{exercise}
Prove \cref{eq:overlapbloch}. Deduce that two states are orthogonal if and only if their Bloch vectors are opposite.
\end{exercise}
\begin{exercise}
Find the Bloch vectors of the eigenstates of $H=\hbar(\sigma_x+\sigma_z)$ and the precession frequency of a generic state.
\end{exercise}
\begin{exercise}
Show that $HR_x(\alpha)H=R_z(\alpha)$ and that $SXS^\dagger=Y$. Show that $\{H,S\}$ generates a finite group (the Clifford group of one qubit,
with 24 elements up to phase, the rotation group of the cube), and therefore that $T$ is needed for universality.
\end{exercise}

\chapter{Measurement and probability}\label{ch:measurement}

\section{Projective measurements}
An observable $M$ has the spectral decomposition
\begin{equation}
  M=\sum_m m\,P_m,
\end{equation}
where $P_m$ is the projector onto the eigenspace of eigenvalue $m$. The possible outcomes are the eigenvalues. The probability to get $m$ is
\begin{equation}\label{eq:born}
  p(m)=\mel\psi{P_m}\psi,
\end{equation}
and if the value $m$ is obtained, the state becomes
\begin{equation}\label{eq:collapse}
  \ket{\psi'}=\frac{P_m\ket\psi}{\sqrt{p(m)}}.
\end{equation}
The projectors satisfy
\begin{equation}
  \sum_mP_m=I\qquad\text{and}\qquad P_mP_{m'}=\delta_{mm'}P_m,
\end{equation}
which guarantee $\sum_mp(m)=1$ and $p(m)\ge0$ ($\mel\psi{P_m}\psi=\norm{P_m\psi}^2$). Often one defines only the projectors; there is then an implicit
observable $M=\sum_m m P_m$ with arbitrary distinct labels $m$. The projectors are often defined by a basis (``measure in the basis $\{\ket m\}$''):
$P_m=\proj m$, and then $p(m)=|\braket m\psi|^2$.

\paragraph{Properties.}
\begin{itemize}
  \item \emph{Mean value:} $\expval M=\sum_m m\,p(m)=\mel\psi M\psi$.
  \item \emph{Repeatability:} measuring $M$ again immediately gives $m$ with certainty, since $P_m\ket{\psi'}=\ket{\psi'}$.
  \item \emph{Global phases are irrelevant:} $\ket\psi\to\e^{\ii\gamma}\ket\psi$ does not change \cref{eq:born}. Relative phases, however, matter as soon
  as one measures in another basis.
\end{itemize}

\begin{example}
$\ket\psi=(\ket0+\ket1)/\sqrt2$, $P_0=\proj0$:
$p(0)=\frac{\bra0+\bra1}{\sqrt2}\proj0\frac{\ket0+\ket1}{\sqrt2}=\frac12$. The state after the observation is
$P_0\ket\psi/\sqrt{p(0)}=\frac{\ket0/\sqrt2}{\sqrt{1/2}}=\ket0$.
\end{example}

\section{Measurement along an arbitrary axis}
Let $\vec n$ be a unit vector. Then $\vec n\cdot\vec\sigma=n_1\sigma_1+n_2\sigma_2+n_3\sigma_3$ is the observable ``spin along $\vec n$''
(a Stern--Gerlach apparatus oriented along $\vec n$). Its eigenvalues are $\pm1$ and its projectors are (Problem 2.60 in \cite{NielsenChuang})
\begin{equation}\label{eq:projector}
  P_\pm=\frac{I\pm\vec n\cdot\vec\sigma}2.
\end{equation}
\begin{proof}
Since $(\vec n\cdot\vec\sigma)^2=I$, $P_\pm^2=\frac{I\pm2\vec n\cdot\vec\sigma+I}{4}=P_\pm$; $P_+P_-=\frac{I-I}4=0$; $P_++P_-=I$; and
$P_+-P_-=\vec n\cdot\vec\sigma$. These are exactly the properties of the spectral projectors of $\vec n\cdot\vec\sigma$ with eigenvalues $\pm1$.
For $\vec n=\vec e_z$, $P_+=\proj0$ and $P_-=\proj1$: indeed $\sigma_z=\proj0-\proj1$ and $I=\proj0+\proj1$.
\end{proof}
For a qubit prepared in $\ket0$,
\begin{equation}
  p(+)=\mel0{P_+}0=\frac{1+n_3}2=\frac{1+\cos\theta}2=\cos^2\frac\theta2,
\end{equation}
where $\theta$ is the polar angle of $\vec n$ on the Bloch sphere (equivalently, the Poincaré sphere for photons): $\theta=0\Rightarrow p=1$,
$\theta=\pi/2\Rightarrow p=1/2$, $\theta=\pi\Rightarrow p=0$. More generally, for a state with Bloch vector $\vec m$,
$p(+)=\frac{1+\vec m\cdot\vec n}2$, which proves \cref{eq:overlapbloch}.

\paragraph{Measuring in any basis with a computational-basis detector.} Quantum computers usually measure only $Z$. To measure along $\vec n$,
one first applies the rotation $U$ that brings $\vec n$ to $\vec e_z$, then measures $Z$, since
$\mel\psi{U^\dagger P_\pm^{(z)}U}\psi=\mel\psi{P_\pm^{(\vec n)}}\psi$. For example, $X$ is measured by applying $H$ before the measurement.

\section{Statistics of repeated measurements}
A single measurement gives one random outcome; quantum mechanics predicts probabilities, which are estimated by repeating the experiment
(``shots''). With $N$ shots, the number of outcomes $+$ is binomial, and the estimated probability $\hat p=N_+/N$ has a standard deviation
\begin{equation}
  \delta\hat p=\sqrt{\frac{p(1-p)}N}\le\frac1{2\sqrt N}.
\end{equation}
This \emph{projection noise} (or shot noise) is the reason why a quantum computer must run each circuit thousands of times: the precision of
an estimated expectation value improves only as $1/\sqrt N$. For variational algorithms (\cref{ch:qml}), this is a major cost.

\section{Generalized measurements (POVM)\starred}\label{sec:povm}
Real detectors are often not projective: a photon detector destroys the photon, and some measurements have more outcomes than the dimension of
the space. The general description uses \emph{measurement operators} $M_m$ with $\sum_mM_m^\dagger M_m=I$:
\begin{equation}
  p(m)=\mel\psi{M_m^\dagger M_m}\psi,\qquad\ket{\psi'}=\frac{M_m\ket\psi}{\sqrt{p(m)}}.
\end{equation}
If we care only about the probabilities, it suffices to know the positive operators $E_m=M_m^\dagger M_m$, $\sum_mE_m=I$, the
\emph{positive operator-valued measure} (POVM). Projective measurements are the special case $E_m=P_m$. By Naimark's theorem, any POVM can be
realised as a projective measurement on a larger system (the qubit plus an ancilla) after a unitary interaction.

\begin{example}[Unambiguous discrimination]
Alice sends either $\ket0$ or $\ket+$. No measurement can distinguish them with certainty since they are not orthogonal. The three-outcome POVM
$E_1=c\proj1$, $E_2=c\proj-$, $E_3=I-E_1-E_2$, with $c=\sqrt2/(1+\sqrt2)$, never errs: outcome $1$ excludes $\ket0$ (so the state was $\ket+$),
outcome $2$ excludes $\ket+$ (so it was $\ket0$), and outcome $3$ means ``don't know''. It succeeds with probability $1-1/\sqrt2\approx0.29$, which
is optimal.
\end{example}

\section{Exercises}
\begin{exercise}
A qubit is in the state $\ket\psi=\cos\frac\theta2\ket0+\e^{\ii\varphi}\sin\frac\theta2\ket1$. Compute the probabilities of the outcomes of
a measurement of $X$, of $Y$ and of $Z$. Show that the three measurements together determine the state (quantum state tomography).
\end{exercise}
\begin{exercise}
Check that $E_3$ in the unambiguous discrimination example is positive, and compute the success probability.
\end{exercise}
\begin{exercise}[Quantum Zeno effect]
A qubit evolves under $H=\frac{\hbar\omega}2\sigma_x$ starting from $\ket0$, and is measured in the $Z$ basis $N$ times during the interval
$[0,T]$, with $\omega T=\pi$. Show that the probability to find it in $\ket0$ at all measurements is $[\cos^2(\pi/2N)]^N\to1$ as $N\to\infty$.
Frequent measurement freezes the evolution.
\end{exercise}

\chapter{What can be done with one qubit}\label{ch:applications1}

\section{Photonic qubits: polarization}\label{sec:photonic}
Light is quantized; its particles are called photons. A photon in a given spatial mode has two polarization states (horizontal and vertical
linear polarizations, as in 3D glasses; the circular basis can be used as well), so each photon is a two-level system, a qubit. Quantum gates
acting on photons are implemented by waveplates (quarter- and half-waveplates), made of birefringent materials, which have different
refractive indices for the two polarizations and therefore accumulate different phases. Projective measurements are implemented by
polarizers and polarizing beam splitters.

\subsection{Jones vectors and the Poincaré sphere}
\begin{origin}[: from Maxwell's equations to a qubit]
A monochromatic plane wave propagating along $z$ has the electric field
$\vec E(z,t)=\operatorname{Re}\left[(E_H\vec e_x+E_V\vec e_y)\e^{\ii(kz-\omega t)}\right]$. The complex two-component \emph{Jones vector}
$(E_H,E_V)$, normalized to unit intensity, describes the polarization. Its global phase is irrelevant to the polarization (it is just a
shift of the origin of time). The polarization of a classical wave is therefore described by exactly the same mathematics as a qubit:
a normalized vector of $\C^2$ modulo a global phase. For a single photon, the Jones vector becomes the quantum state
$\ket\psi=\alpha\ket H+\beta\ket V$ and $|\alpha|^2$ becomes a probability.
\end{origin}
We identify $\ket H=\ket0$, $\ket V=\ket1$. Then
\begin{center}
\begin{tabular}{@{}lll@{}}\toprule
Polarization & State & Point on the Poincaré sphere\\ \midrule
horizontal, vertical & $\ket H$, $\ket V$ & $\pm\vec e_z$\\
diagonal, antidiagonal ($\pm45^\circ$) & $\ket{D},\ket{A}=(\ket H\pm\ket V)/\sqrt2$ & $\pm\vec e_x$\\
right, left circular & $\ket{R},\ket L=(\ket H\pm\ii\ket V)/\sqrt2$ & $\pm\vec e_y$\\ \bottomrule
\end{tabular}
\end{center}
(The handedness convention for circular polarization varies between textbooks.) The components of the Bloch vector are the normalized
\emph{Stokes parameters}, measured classically with polarizers: $S_1=I_H-I_V$, $S_2=I_D-I_A$, $S_3=I_R-I_L$, divided by the total intensity.
A linear polarization at angle $\eta$ from the horizontal, $\cos\eta\ket H+\sin\eta\ket V$, corresponds by \cref{eq:blochstate} to
$\theta=2\eta$, $\varphi=0$: a point on the equator-crossing great circle in the $xz$ plane at polar angle $2\eta$.

The pseudospin is thus rotated by a \emph{double angle} with respect to the polarization: orthogonal polarizations ($H$ and $V$) differ by
$90^\circ$ in real space, while orthogonal qubit states are $180^\circ$ apart on the sphere. One needs to get used to it when working with photons.

\subsection{Waveplates as quantum gates}
A plate of thickness $L$ of a birefringent crystal with indices $n_f$ (fast axis) and $n_s$ (slow axis) delays the slow component by the phase
\begin{equation}
  \delta=\frac{2\pi(n_s-n_f)L}{\lambda}.
\end{equation}
With the fast axis horizontal, the plate acts as $J_0(\delta)=\diag(1,\e^{\ii\delta})=\e^{\ii\delta/2}R_z(\delta)$: a rotation of the Poincaré
sphere about the $z$ axis by $\delta$. A half-wave plate (HWP) has $\delta=\pi$ ($J_0=Z$), a quarter-wave plate (QWP) has $\delta=\pi/2$ ($J_0=S$).

A plate rotated by an angle $\alpha$ in the laboratory is $J(\alpha)=\mathcal R(\alpha)J_0\mathcal R(-\alpha)$ with the real rotation matrix
$\mathcal R(\alpha)=\begin{pmatrix}\cos\alpha&-\sin\alpha\\ \sin\alpha&\cos\alpha\end{pmatrix}=R_y(2\alpha)$: a physical rotation by $\alpha$ is a rotation of the
sphere by $2\alpha$ about $y$. For a half-wave plate,
\begin{equation}
  J_{\rm HWP}(\alpha)=\begin{pmatrix}\cos2\alpha&\sin2\alpha\\ \sin2\alpha&-\cos2\alpha\end{pmatrix}.
\end{equation}
It maps $\ket H$ to the linear polarization at $2\alpha$: a HWP at angle $\alpha$ \emph{mirrors} the polarization about its axis. Special cases:
at $\alpha=0$ it is the $Z$ gate; at $\alpha=45^\circ$ it is the $X$ gate ($H\leftrightarrow V$); at $\alpha=22.5^\circ$ it is exactly the
\emph{Hadamard gate}. A QWP at $45^\circ$ transforms linear into circular polarization. With a QWP--HWP--QWP sequence, any $SU(2)$ rotation
can be realised (this is the optical version of the Euler decomposition, \cref{thm:euler}).

\subsection{Measurement: Malus's law}
A polarizer along the direction $\eta_0$ realizes the projective measurement onto $\ket{\eta_0}=\cos\eta_0\ket H+\sin\eta_0\ket V$. The
equivalent of \cref{eq:projector} is \emph{Malus's law}: the probability of transmission is
\begin{equation}
  p=|\braket{\eta_0}\psi|^2=\cos^2\eta=\cos^2\frac\Theta2,
\end{equation}
where $\eta$ is the angle between the polarization of light and the transmission axis of the polarizer, and $\Theta=2\eta$ is the angle between
the pseudospin of the light and the pseudospin eigenvector of the polarizer. For a classical beam, Malus's law (1809) gives the transmitted
intensity; for a single photon it gives a probability.

\section{Random number generation}
\subsection{Why true randomness matters}
Random numbers are an essential resource in many areas: cryptographic keys and nonces, quantum key distribution, Monte Carlo simulations,
randomized algorithms, statistical sampling and lotteries. Most random numbers used in practice come from \emph{pseudo-random number generators}.
These are deterministic algorithms that expand a short seed into a long sequence that looks random. Anyone who knows the algorithm and the
seed can reproduce the whole sequence, which is a serious weakness in cryptography, and the sequence is eventually periodic.

Classical physical generators, based for instance on thermal noise or chaotic circuits, avoid the explicit seed. Their unpredictability,
however, comes only from our ignorance of the microscopic details: the underlying dynamics is deterministic in principle, and it is hard to
certify that no hidden correlations, drifts or external influences are present.

Quantum mechanics offers a fundamentally different source. According to the Born rule, the outcome of a measurement on a superposition state
is intrinsically unpredictable. No additional information about the system, even in principle, can determine it in advance (this statement is
made precise, and testable, by Bell's theorem, \cref{ch:bell}). This is the basis of \emph{quantum random number generators} (QRNGs) \cite{Jacak2021}.

\subsection{Example: a photon and a polarizer}
Consider a single photon linearly polarized at an angle $\theta$ with respect to the horizontal axis:
\begin{equation}
  \ket\psi=\cos\theta\ket H+\sin\theta\ket V.
\end{equation}
The photon is sent onto a polarizer that transmits horizontal polarization, followed by a single-photon detector. This realizes the
projective measurement with projectors $P_H=\proj H$ (transmitted, detector clicks) and $P_V=\proj V$ (absorbed, no click). The transmission
probability is
\begin{equation}
  p(H)=\mel\psi{P_H}\psi=|\braket H\psi|^2=\cos^2\theta,
\end{equation}
which is Malus's law at the level of a single photon. For $\theta=45^\circ$, $\ket\psi=(\ket H+\ket V)/\sqrt2$ and $p(H)=p(V)=1/2$. Each photon is
transmitted or absorbed with equal probability. Assigning the bit value $1$ to a click and $0$ to no click, a stream of identically prepared
photons produces a sequence of independent, unbiased random bits. After the measurement the transmitted photon is in state $\ket H$, so its
polarization no longer carries any memory of the random choice.

\subsection{Practical remarks}
A real device deviates from this ideal picture. If the angle is not exactly $45^\circ$, the bits are biased, with $p(1)=\cos^2\theta\neq\frac12$.
Photon losses and the finite detector efficiency also look like ``no click'' events and add to the bias. A common improvement replaces the
polarizer by a \emph{polarizing beam splitter} with a detector in each output port, so that both outcomes are actually detected and events with no
photon are discarded. Most commercial QRNGs actually use a beam splitter, which is the same measurement for a photon in a superposition of two
spatial modes.

Residual bias is removed by classical post-processing. The simplest method is \emph{von Neumann's procedure} \cite{vonNeumann1951}: group the bits
in pairs, map $01\to0$ and $10\to1$, and discard $00$ and $11$.
\begin{proposition}
For independent bits with any fixed bias $P(1)=p$, the output of the von Neumann extractor is exactly unbiased; on average it produces
$p(1-p)$ output bits per input bit.
\end{proposition}
\begin{proof}
For independent bits, $P(01)=(1-p)p=P(10)$, so conditioned on keeping the pair, both outputs have probability $1/2$. A pair is kept with
probability $2p(1-p)$ and yields one bit, i.e.\ $p(1-p)$ bits per input bit (at most $1/4$, for $p=1/2$).
\end{proof}
The method relies on independence; correlated bits require stronger \emph{randomness extractors} (e.g.\ hashing with Toeplitz matrices), whose
output length is set by the min-entropy $H_{\min}=-\log_2\max_xP(x)$ of the raw data.

Finally, the randomness of such a device rests on trusting its physical model (state preparation, polarizer and detectors). \emph{Device-independent}
QRNGs remove this assumption by certifying the randomness through the violation of a Bell inequality, at the cost of a much lower bit rate.

\section{The no-cloning theorem}\label{sec:nocloning}
A natural question is whether an unknown quantum state can be copied. Suppose there exists a linear operator $C$ that clones any state $\ket\alpha$
onto a second system initially prepared in a reference state $\ket0$:
\begin{equation}\label{eq:clone}
  C\ket\alpha\ket0=\ket\alpha\ket\alpha\qquad\text{for all }\ket\alpha.
\end{equation}
\begin{theorem}[No-cloning \cite{WoottersZurek1982,Dieks1982}]
No linear (in particular, no unitary) operator satisfies \cref{eq:clone} for all states.
\end{theorem}
\begin{proof}[Proof 1 (linearity)]
Let $\ket a$, $\ket b$ be two orthogonal states. By \cref{eq:clone}, $C\ket a\ket0=\ket a\ket a$ and $C\ket b\ket0=\ket b\ket b$. For the
superposition $\ket x=(\ket a+\ket b)/\sqrt2$, linearity gives
\begin{equation}
  C\ket x\ket0=\frac1{\sqrt2}\left(\ket a\ket a+\ket b\ket b\right),
\end{equation}
whereas a true clone would be
\begin{equation}
  \ket x\ket x=\frac12\left(\ket a\ket a+\ket a\ket b+\ket b\ket a+\ket b\ket b\right).
\end{equation}
The two differ: the first does not contain the cross terms, and even its coefficients differ. Contradiction.
\end{proof}
\begin{proof}[Proof 2 (unitarity)]
If $U\ket a\ket0=\ket a\ket a$ and $U\ket b\ket0=\ket b\ket b$, take the scalar product of the two equations. Since $U^\dagger U=I$, the left-hand side gives
$\braket ab\braket00=\braket ab$ and the right-hand side gives $\braket ab^2$. Thus $\braket ab=\braket ab^2$, so $\braket ab=0$ or $1$:
only identical or orthogonal states can be cloned by the same machine. (Allowing the machine to end in a state $\ket{M_a}$ that depends on the input does
not help: one gets $|\braket ab|=|\braket ab|^2|\braket{M_a}{M_b}|\le|\braket ab|^2$.)
\end{proof}
Cloning is therefore not forbidden for states known to belong to a fixed orthogonal set $\{\ket a,\ket b\}$: classical information can be copied (a CNOT
copies $\ket0$ and $\ket1$). It is the requirement to copy their superpositions as well that leads to the contradiction.

\begin{remark}[History and relation to relativity]
The theorem was found in 1982 in response to N.~Herbert's ``FLASH'' proposal for faster-than-light communication. Alice and Bob share an entangled pair
(\cref{ch:entanglement}); Alice measures her particle either in the $Z$ or in the $X$ basis, which leaves Bob's particle in a state of the $Z$ basis or of
the $X$ basis. If Bob could make many copies of his particle, he could measure them and find out which basis Alice chose, instantaneously. No-cloning
forbids it. More generally, perfect cloning would allow signalling and perfect state discrimination.
\end{remark}
\begin{remark}[Approximate cloning\starred]
Imperfect copies are allowed. The best universal cloner for qubits (Bužek and Hillery, 1996) produces two copies, each with fidelity
$F=\mel\psi{\rho_{\rm copy}}\psi=5/6$ for every input state. Its Bloch vector is shrunk by a factor $2/3$.
\end{remark}

\section{Quantum key distribution: BB84}\label{sec:bb84}
\subsection{The problem}
The only provably secure classical cipher is the \emph{one-time pad} (Vernam, 1917): the message $m$ is combined with a random key $k$ of the same length,
$c=m\oplus k$, and decrypted by $m=c\oplus k$. The ciphertext is statistically independent of the message (Shannon, 1949), provided that the key is
truly random, used only once, and \emph{known only to the two parties}. The hard part is to distribute the key. Public-key cryptography (RSA) solves
this problem only under computational assumptions, which Shor's algorithm (\cref{ch:shor}) would break.

The BB84 protocol, proposed by Bennett and Brassard in 1984 \cite{BB84}, was the first quantum key distribution (QKD) scheme. Its goal is to establish a
shared secret key between two parties, traditionally called Alice and Bob. Its security does not rely on computational assumptions. It follows from the
laws of quantum mechanics: an eavesdropper cannot copy an unknown quantum state in order to measure it later (no-cloning), and cannot measure it without
disturbing it if she does not know the basis.

The protocol uses two channels: a \emph{quantum channel}, along which Alice sends single photons to Bob, and a public but \emph{authenticated} classical
channel, along which they compare their choices of basis. (Authentication requires a short pre-shared secret; QKD is therefore really a key
\emph{expansion} protocol.) It was originally formulated for single photons. Ekert later proposed a version based on entangled pairs \cite{Ekert1991}
(\cref{sec:e91}).

\subsection{Encoding}
Alice encodes each bit in the polarization of a photon, using one of two bases:
\begin{align}
  \text{rectilinear basis }(+):&\quad 0\mapsto\ket\uparrow,\quad 1\mapsto\ket\rightarrow,\\
  \text{diagonal basis }(\times):&\quad 0\mapsto\ket\nearrow,\quad 1\mapsto\ket\searrow,
\end{align}
with $\ket\nearrow,\ket\searrow=(\ket\rightarrow\pm\ket\uparrow)/\sqrt2$. The states within one basis are orthogonal, while the two bases are
\emph{mutually unbiased}: $|\braket\uparrow\nearrow|^2=|\braket\uparrow\searrow|^2=1/2$. On the Poincaré sphere, the four states lie at the four points
$\pm\vec e_z$, $\pm\vec e_x$ of a great circle. If a photon is measured in the basis in which it was prepared, the result is certain. If it is measured in the
other basis, the result is completely random.

\subsection{The protocol}
\begin{enumerate}
  \item Alice generates a random bit string and, for each bit, a random choice of basis ($+$ or $\times$). She sends Bob one photon per bit, polarized accordingly.
  \item Bob independently chooses a random basis for each photon and measures its polarization in that basis.
  \item Over the public channel, Alice and Bob announce the bases they used, but not the bit values. They keep only the bits for which the bases coincide
  (on average half of them). This is the \emph{sifted key}.
  \item They publicly compare a random subset of the sifted key to estimate the \emph{quantum bit error rate} (QBER) $Q$, and discard these bits.
  \item If $Q$ is below a threshold, they perform classical \emph{error correction} (so that both have the same string) and \emph{privacy amplification}
  (hashing the string into a shorter one about which Eve's information is negligible).
\end{enumerate}
An example is shown in \cref{tab:bb84}. In columns 2, 4, 5 and 7 the bases differ, and Bob's result is random. In the other columns his result coincides with Alice's bit.

\begin{table}[h]\centering
\begin{tabular}{@{}lcccccccc@{}}\toprule
Alice's random bit & 0&1&1&0&1&0&0&1\\
Alice's basis & $+$&$+$&$\times$&$+$&$\times$&$\times$&$\times$&$+$\\
Polarization sent & $\uparrow$&$\rightarrow$&$\searrow$&$\uparrow$&$\searrow$&$\nearrow$&$\nearrow$&$\rightarrow$\\
Bob's basis & $+$&$\times$&$\times$&$\times$&$+$&$\times$&$+$&$+$\\
Polarization measured & $\uparrow$&$\nearrow$&$\searrow$&$\nearrow$&$\rightarrow$&$\nearrow$&$\rightarrow$&$\rightarrow$\\
Same basis? & \checkmark&--&\checkmark&--&--&\checkmark&--&\checkmark\\
Shared key & 0&&1&&&0&&1\\ \bottomrule
\end{tabular}
\caption{An example of the BB84 protocol with eight photons.}\label{tab:bb84}
\end{table}

\subsection{Security against intercept--resend}
Suppose an eavesdropper, Eve, intercepts each photon, measures it in a randomly chosen basis, and resends to Bob a photon in the state she found.
Consider a bit that survives sifting (Alice and Bob used the same basis).
\begin{itemize}
  \item With probability $1/2$ Eve chose the right basis: she learns the bit and resends the correct state. No error.
  \item With probability $1/2$ she chose the wrong basis: her result is random, and she resends a state of the wrong basis. Bob, measuring in Alice's basis,
  obtains a random result, wrong with probability $1/2$.
\end{itemize}
Eve therefore introduces errors in $\frac12\times\frac12=25\%$ of the sifted bits, while learning $50\%$ of them. Comparing $m$ bits, Alice and Bob
detect her with probability $1-(3/4)^m$, i.e.\ $>99.99\%$ for $m=32$. Because of the no-cloning theorem, Eve cannot avoid this disturbance by copying
the photons. More subtle attacks trade information for disturbance, but there is always a trade-off: \emph{information gain implies disturbance}.

\begin{caveat}[: how much error is tolerable?]
Real channels have errors even without Eve, and Alice and Bob must attribute \emph{all} errors to Eve. The Shor--Preskill security proof (2000)
shows that, against the most general attacks, a secret key can be distilled at the rate (secret bits per sifted bit)
\begin{equation}
  r=1-2h(Q),\qquad h(Q)=-Q\log_2Q-(1-Q)\log_2(1-Q),
\end{equation}
where one $h(Q)$ is the cost of error correction and the other the cost of privacy amplification. The rate vanishes at $Q\approx11\%$. Above this threshold
the key is discarded.
\end{caveat}

\subsection{Practical QKD}
Single-photon sources are difficult, so most systems use attenuated laser pulses, coherent states with a mean photon number $\mu\approx0.1$--$0.5$. The number of
photons per pulse is Poisson distributed, $P(n)=\e^{-\mu}\mu^n/n!$, so a fraction $P(n\ge2)/P(n\ge1)\approx\mu/2$ of the non-empty pulses contains several photons.
Eve can keep one of them and let the others go to Bob (\emph{photon-number-splitting attack}). The \emph{decoy-state} method (Hwang 2003; Lo, Ma, Chen 2005)
defeats this attack by randomly varying $\mu$, which lets Alice and Bob estimate the single-photon contribution. Detector side channels (``blinding'' attacks)
motivate measurement-device-independent QKD. Commercial QKD systems reach a few hundred kilometres in fibre; the Chinese satellite Micius distributed keys
over $1200$\,km (2017). Beyond that, \emph{quantum repeaters}, based on entanglement swapping (\cref{ch:teleportation}) and quantum memories, are required, since
the no-cloning theorem also forbids optical amplification of quantum signals.

\section{Exercises}
\begin{exercise}
Show that a half-wave plate at $22.5^\circ$ implements the Hadamard gate on polarization qubits, and that a QWP at $45^\circ$ maps $\ket H$ to a circular polarization.
\end{exercise}
\begin{exercise}
In BB84, Eve measures every photon in the \emph{Breidbart basis}, halfway between the two bases (at $22.5^\circ$ in real space). Show that she guesses
each sifted bit correctly with probability $\cos^2(\pi/8)\approx0.85$, and that she induces a QBER of $25\%$.
\end{exercise}
\begin{exercise}
Compute the probability that Eve, measuring every photon as in the intercept--resend attack, remains undetected when Alice and Bob compare $m$ bits. How many bits must
be compared to detect her with probability $1-10^{-9}$?
\end{exercise}


\part{Many qubits, mixed states and decoherence}
\chapter{Many qubits and entanglement}\label{ch:entanglement}

So far we have seen that a lot of useful things can be done already with a single qubit: random numbers and quantum cryptography are part of the
larger field of quantum informatics. Quantum computing with gates really starts when we go beyond a single qubit.

\section{Composite systems and the tensor product}\label{sec:tensor}
\subsection{Why a tensor product?}
\begin{origin}[: the tensor-product postulate]
If system $A$ can be in any of the states $\ket{i}_A$ ($i=1..d_A$) and system $B$ in any of the states $\ket\mu_B$ ($\mu=1..d_B$), then the
pair can be in any of the $d_Ad_B$ joint configurations $\ket i_A\ket\mu_B$, and, by the superposition principle, in any linear combination of
them. The state space is therefore of dimension $d_Ad_B$ (not $d_A+d_B$, as it would be for a classical description by a list of properties).
The tensor product $\Ha_A\otimes\Ha_B$ is precisely the space spanned by the products, with the scalar product
$(\bra i\bra\mu)(\ket j\ket\nu)=\braket ij\braket\mu\nu$. Probabilities of independent events multiply, which is consistent with
$|\braket{i\mu}{\psi_A\psi_B}|^2=|\braket i{\psi_A}|^2|\braket\mu{\psi_B}|^2$. For $n$ qubits the dimension is $2^n$.
\end{origin}

\subsection{Kronecker products of matrices and vectors}
Let $A$ be an $m\times n$ matrix and $B$ a $p\times q$ matrix. Their tensor (Kronecker) product $A\otimes B$ is the $mp\times nq$ matrix with elements
\begin{equation}
  (A\otimes B)_{p(r-1)+v,\ q(s-1)+w}=a_{rs}b_{vw},
\end{equation}
that is, the block matrix
\begin{equation}
  A\otimes B=\begin{pmatrix}a_{11}B&\cdots&a_{1n}B\\ \vdots&\ddots&\vdots\\ a_{m1}B&\cdots&a_{mn}B\end{pmatrix}.
\end{equation}
For example,
\begin{equation}
  \sigma_z\otimes\sigma_x=\begin{pmatrix}1\cdot\sigma_x&0\cdot\sigma_x\\0\cdot\sigma_x&-1\cdot\sigma_x\end{pmatrix}
  =\begin{pmatrix}0&1&0&0\\1&0&0&0\\0&0&0&-1\\0&0&-1&0\end{pmatrix}.
\end{equation}
The same rule applies to vectors, which fixes the ordering of the basis of the composite system:
\begin{equation}
  \ket0\otimes\ket0=\ket{00}=\begin{pmatrix}1\\0\end{pmatrix}\otimes\begin{pmatrix}1\\0\end{pmatrix}=\begin{pmatrix}1\\0\\0\\0\end{pmatrix},\quad
  \ket{01}=\begin{pmatrix}0\\1\\0\\0\end{pmatrix},\quad\ket{10}=\begin{pmatrix}0\\0\\1\\0\end{pmatrix},\quad\ket{11}=\begin{pmatrix}0\\0\\0\\1\end{pmatrix},
\end{equation}
and similarly $\ket{000}=(1,0,0,0,0,0,0,0)^T$. In general, the basis states of $n$ qubits are ordered as the binary numbers $00\dots0$, $00\dots1$, \dots, $11\dots1$
(big-endian convention, used throughout these notes). Useful rules:
\begin{equation}
  (A\otimes B)(C\otimes D)=AC\otimes BD,\quad(A\otimes B)^\dagger=A^\dagger\otimes B^\dagger,\quad\tr(A\otimes B)=\tr A\,\tr B,
\end{equation}
and the eigenvalues of $A\otimes B$ are the products $a_ib_j$ of the eigenvalues. An operator acting on qubit $k$ only, such as $\sigma^z_k$, means
$I\otimes\cdots\otimes\sigma^z\otimes\cdots\otimes I$ in the full $2^n$-dimensional space.

\section{Entanglement}
\begin{definition}
A pure state of $A$ and $B$ is a \emph{product} (separable) state if it can be written $\ket\psi=\ket{\psi_A}\otimes\ket{\psi_B}$. Otherwise it is \emph{entangled}.
For $N$ parts, a state is fully separable if it factorizes into a tensor product of $N$ single-part states.
\end{definition}
A state that can be factored is, for example,
\begin{equation}
  \Psi_{1+2}=\frac1{\sqrt2}\left(\ket0_1\ket1_2-\ket1_1\ket1_2\right)=\frac1{\sqrt2}\left(\ket0_1-\ket1_1\right)\ket1_2 .
\end{equation}
On the other hand, the following maximally entangled state cannot be factored:
\begin{equation}\label{eq:singlet}
  \Psi_{1+2}=\frac1{\sqrt2}\left(\ket0_1\ket1_2-\ket1_1\ket0_2\right).
\end{equation}
\begin{proof}
Suppose it can: $\frac1{\sqrt2}(\ket{01}-\ket{10})=(a\ket0+b\ket1)\otimes(c\ket0+d\ket1)=ac\ket{00}+ad\ket{01}+bc\ket{10}+bd\ket{11}$. Identifying coefficients,
\begin{equation}
  ac=0,\quad bd=0,\quad ad=\tfrac1{\sqrt2},\quad bc=-\tfrac1{\sqrt2}.
\end{equation}
The third and fourth equations require $a,b,c,d$ all non-zero, which contradicts the first two.
\end{proof}

\begin{proposition}[Entanglement criterion for two qubits]\label{prop:det}
The state $\ket\psi=\alpha\ket{00}+\beta\ket{01}+\gamma\ket{10}+\delta\ket{11}$ is a product state if and only if $\alpha\delta-\beta\gamma=0$.
\end{proposition}
\begin{proof}
Arrange the coefficients in the matrix $\Psi=\begin{pmatrix}\alpha&\beta\\ \gamma&\delta\end{pmatrix}$. A product state $(a,b)\otimes(c,d)$ has
$\Psi=\begin{pmatrix}a\\b\end{pmatrix}\begin{pmatrix}c&d\end{pmatrix}$, a matrix of rank one, hence $\det\Psi=0$. Conversely, a non-zero $2\times2$ matrix with
zero determinant has rank one, i.e.\ is an outer product of two vectors.
\end{proof}
The quantity $C=2|\alpha\delta-\beta\gamma|\in[0,1]$ is the \emph{concurrence} of the pure state; it is $0$ for product states and $1$ for maximally entangled
states such as \cref{eq:singlet}. The argument generalizes: a bipartite state is a product state iff its coefficient matrix has rank one. The rank is the Schmidt
number (\cref{app:schmidt}).

\paragraph{Consequence.} As soon as two subsystems 1 and 2 are entangled, there is no such thing as ``the wave function of subsystem 1'' alone. The whole is more
than the sum of its parts: all the information may be in the correlations, none in the parts. The description of a part requires the density matrix
(\cref{ch:density}). Entangled states are not rare: for random states of many qubits, product states form a set of measure zero.

\subsection{The Bell states}
The four maximally entangled states
\begin{equation}\label{eq:bellstates}
  \ket{\Phi^\pm}=\frac{\ket{00}\pm\ket{11}}{\sqrt2},\qquad\ket{\Psi^\pm}=\frac{\ket{01}\pm\ket{10}}{\sqrt2}
\end{equation}
form an orthonormal basis of two qubits, the \emph{Bell basis}. $\ket{\Psi^-}$ is the spin \emph{singlet} (total spin $0$); the other three form the
\emph{triplet} (total spin $1$). They are all related by local operations on one qubit: $\ket{\Phi^-}=(Z\otimes I)\ket{\Phi^+}$,
$\ket{\Psi^+}=(X\otimes I)\ket{\Phi^+}$, $\ket{\Psi^-}=(XZ\otimes I)\ket{\Phi^+}$ (up to sign). They are created from product states by the circuit
\begin{center}
\begin{quantikz}
\lstick{$\ket{x}$} & \gate{H} & \ctrl{1} & \qw \\
\lstick{$\ket{y}$} & \qw & \targ{} & \qw
\end{quantikz}
\qquad$\ket{00}\to\ket{\Phi^+},\ \ket{01}\to\ket{\Psi^+},\ \ket{10}\to\ket{\Phi^-},\ \ket{11}\to\ket{\Psi^-}$.
\end{center}
For example, $H\otimes I$ maps $\ket{00}$ to $(\ket{00}+\ket{10})/\sqrt2$, and the CNOT (\cref{sec:cnot}) flips the second qubit in the second term, giving
$(\ket{00}+\ket{11})/\sqrt2$. Running the circuit backwards (CNOT, then $H$) maps the Bell basis onto the computational basis: this is how a \emph{Bell measurement}
is performed on a quantum computer.

\subsection{Multipartite entanglement\starred}
For three qubits, two inequivalent kinds of genuine tripartite entanglement exist:
\begin{equation}
  \ket{\mathrm{GHZ}}=\frac{\ket{000}+\ket{111}}{\sqrt2},\qquad\ket{W}=\frac{\ket{001}+\ket{010}+\ket{100}}{\sqrt3}.
\end{equation}
Tracing out one qubit of GHZ leaves a separable (classically correlated) state, while for W the remaining two qubits are still entangled. No local operations,
even assisted by classical communication and with a finite probability of success, can convert one into the other (Dür, Vidal, Cirac 2000).

\section{Quantum parallelism and entanglement}\label{sec:parallelism}
Why is entanglement important for quantum computing? Because it makes the link between the input and the output.

The idea of \emph{quantum parallelism} is based on the possibility for a qubit to be in a superposition of two states. It is therefore possible to perform a
computation for all possible inputs (input qubits in their superpositions) and obtain all possible outputs in a single run. Given a classical function
$f:\{0,1\}^n\to\{0,1\}^m$, the reversible implementation $U_f\ket x\ket y=\ket x\ket{y\oplus f(x)}$ (\cref{sec:reversible}) gives
\begin{equation}
  U_f\left(\frac1{\sqrt{2^n}}\sum_x\ket x\right)\ket0=\frac1{\sqrt{2^n}}\sum_x\ket x\ket{f(x)}.
\end{equation}
A measurement at the end gives one single result, but quantum algorithms enhance the probability to measure the result ``of interest''. For example, one can look
for the solutions of an equation, the output being its right-hand side. All possible inputs give all possible values, but interference can increase the probability
to measure the correct output. How do we then know to which input it corresponds, i.e.\ to which values of the unknowns? Only thanks to entanglement: each output
is entangled with the corresponding inputs, so that when we observe a value of the output qubits, the input qubits collapse onto the states corresponding to that output.

\begin{example}[Multiplication of two bits]
We want
\begin{equation}
  \ket{00}\to\ket0,\quad\ket{01}\to\ket0,\quad\ket{10}\to\ket0,\quad\ket{11}\to\ket1 .
\end{equation}
The state of the system should be a superposition of the four 3-qubit states $\ket{000}$, $\ket{010}$, $\ket{100}$, $\ket{111}$ (2 input qubits and 1 output qubit)
and no other terms. This state is entangled: a product of single-qubit states containing $\ket{000}$ and $\ket{111}$ would also contain terms such as $\ket{001}$ or
$\ket{110}$. Measuring the output and finding $1$ projects the inputs onto $\ket{11}$; finding $0$ projects them onto $(\ket{00}+\ket{01}+\ket{10})/\sqrt3$.
The circuit is a Toffoli gate (\cref{sec:toffoli}) with the output qubit as target.
\end{example}

\begin{remark}
Entanglement is necessary but not sufficient for a quantum speed-up. By the Gottesman--Knill theorem, circuits made only of $H$, $S$ and CNOT gates can
create highly entangled states but are efficiently simulable classically. Conversely, a pure-state computation whose entanglement stays bounded (small Schmidt rank
across every cut) can be simulated efficiently with tensor networks (Vidal 2003). Since entanglement is crucial, there is a large effort to create and maximize it,
and this requires being able to measure it and describe its decay by decoherence (\cref{ch:density,ch:open}).
\end{remark}

\section{Exercises}
\begin{exercise}
Which of the following states are entangled? $\frac12(\ket{00}+\ket{01}+\ket{10}+\ket{11})$; $\frac12(\ket{00}+\ket{01}+\ket{10}-\ket{11})$;
$\frac1{\sqrt2}(\ket{00}+\ii\ket{11})$. Compute their concurrence.
\end{exercise}
\begin{exercise}
Show that $\ket{\Psi^-}$ has the same form in every basis: $\ket{\Psi^-}=\frac1{\sqrt2}(\ket{\vec n,+}\ket{\vec n,-}-\ket{\vec n,-}\ket{\vec n,+})$ for any axis $\vec n$.
(Hint: $(U\otimes U)\ket{\Psi^-}=\det U\,\ket{\Psi^-}$ for any $2\times2$ matrix $U$.)
\end{exercise}

\chapter{Bell inequalities and nonlocality}\label{ch:bell}

\section{The EPR argument}
In 1935 Einstein, Podolsky and Rosen (EPR) \cite{EPR1935} argued that quantum mechanics might be an incomplete description of reality. In the form
given by Bohm, the argument goes as follows. Two spin-$\frac12$ particles are prepared in the singlet state and fly apart. If Alice measures $\sigma_z$ on
her particle, she can predict with certainty the result of a measurement of $\sigma_z$ on Bob's particle (the opposite one), without disturbing it in any way,
if the two measurements are space-like separated. According to EPR, the value of $\sigma_z$ of Bob's particle is then an ``element of reality''. But Alice
could equally well have measured $\sigma_x$, so $\sigma_x$ of Bob's particle must also be an element of reality. Quantum mechanics does not allow both
values to be defined simultaneously ($[\sigma_x,\sigma_z]\neq0$), so, EPR concluded, it must be incomplete: it should be completed by additional
\emph{hidden variables}.

In 1964 John Bell \cite{Bell1964} took the EPR thought experiment one step further. He showed that the question is not merely philosophical: if the hidden
variables are \emph{local}, they lead to predictions that can be tested experimentally, and that contradict quantum mechanics.

\section{The setup and the quantum prediction}
A source emits pairs of spin-$\frac12$ particles in the singlet state
\begin{equation}
  \ket{\psi^-}=\frac1{\sqrt2}\left(\ket{\uparrow\downarrow}-\ket{\downarrow\uparrow}\right),
\end{equation}
which fly apart in opposite directions (\cref{fig:bellsetup}). One particle reaches Alice's analyzer, which measures the spin component along a unit
vector $\vec a$; the other reaches Bob's analyzer, oriented along $\vec b$. Each measurement gives the result $+1$ or $-1$ (in units of $\hbar/2$).

\begin{figure}[h]\centering
\begin{tikzpicture}
  \node[draw,circle,fill=yellow!30] (s) at (0,0) {source};
  \node[draw,rectangle,minimum width=1.3cm,minimum height=1cm] (a) at (-4.5,0) {Alice};
  \node[draw,rectangle,minimum width=1.3cm,minimum height=1cm] (b) at (4.5,0) {Bob};
  \draw[->,thick] (s)--(a); \draw[->,thick] (s)--(b);
  \draw[->,red,thick] (-4.5,0.6)--(-4.0,1.3) node[above]{$\vec a$};
  \draw[->,red,thick] (4.5,0.6)--(5.1,1.2) node[above]{$\vec b$};
\end{tikzpicture}
\caption{A Bell experiment. A pair of particles emitted by the source is analyzed along the directions $\vec a$ and $\vec b$, which can be chosen freely.}
\label{fig:bellsetup}
\end{figure}

Bell's simple but powerful idea was to consider the average of the product of the two results, the \emph{correlation function} $P(\vec a,\vec b)$.
\begin{proposition}
For the singlet state,
\begin{equation}\label{eq:singletcorr}
  P_{\rm QM}(\vec a,\vec b)=\mel{\psi^-}{(\vec a\cdot\vec\sigma)\otimes(\vec b\cdot\vec\sigma)}{\psi^-}=-\vec a\cdot\vec b.
\end{equation}
\end{proposition}
\begin{proof}
The singlet has total spin zero: $(\vec\sigma\otimes I+I\otimes\vec\sigma)\ket{\psi^-}=0$, which can be checked component by component (e.g.\
$(Z\otimes I+I\otimes Z)(\ket{01}-\ket{10})=(1-1)\ket{01}-(-1+1)\ket{10}=0$). Therefore $(I\otimes\vec b\cdot\vec\sigma)\ket{\psi^-}=-(\vec b\cdot\vec\sigma\otimes I)\ket{\psi^-}$ and
\begin{equation}
  P_{\rm QM}=-\mel{\psi^-}{(\vec a\cdot\vec\sigma)(\vec b\cdot\vec\sigma)\otimes I}{\psi^-}
  =-\vec a\cdot\vec b-\ii(\vec a\times\vec b)\cdot\mel{\psi^-}{\vec\sigma\otimes I}{\psi^-}=-\vec a\cdot\vec b,
\end{equation}
where we used the product rule of Pauli matrices and the fact that each particle alone is completely unpolarized,
$\mel{\psi^-}{\vec\sigma\otimes I}{\psi^-}=0$ (its reduced density matrix is $I/2$, \cref{ch:density}).
\end{proof}
If the angle between the analyzers is $\theta_{ab}$, $P=-\cos\theta_{ab}$, and the probability of equal results is $P(A=B)=\frac{1+P}2=\sin^2\frac{\theta_{ab}}2$.
The antiparallel and parallel configurations recover the EPR situation:
\begin{equation}
  P(\vec a,-\vec a)=+1,\qquad P(\vec a,\vec a)=-1.
\end{equation}
If both analyzers are along the same axis, the results are always opposite, so measuring one particle allows us to predict with certainty the result for the other.
Note, however, that each result alone is completely random: $P(A=+1)=1/2$ for any $\vec a$.

\section{Local hidden variables and Bell's inequality}
Assume now that the wave function is not the complete description, and that each pair also carries some unknown variables, collectively denoted $\lambda$,
distributed with a probability density $\rho(\lambda)\ge0$, $\int\rho(\lambda)\dd\lambda=1$. In such a complete theory the results of the measurements
are determined by $\lambda$:
\begin{equation}
  A(\vec a,\lambda)=\pm1,\qquad B(\vec b,\lambda)=\pm1.
\end{equation}
The essential assumption is \emph{locality}: Alice's result $A$ depends only on her own setting $\vec a$ and on $\lambda$, but not on Bob's setting $\vec b$,
and vice versa. A second, implicit, assumption is \emph{measurement independence} (``free will''): $\rho(\lambda)$ does not depend on the settings. The correlation
function is then
\begin{equation}\label{eq:lhvcorr}
  P(\vec a,\vec b)=\int\rho(\lambda)A(\vec a,\lambda)B(\vec b,\lambda)\dd\lambda.
\end{equation}
To reproduce the perfect EPR anticorrelations, the hidden variables must satisfy $A(\vec a,\lambda)=-B(\vec a,\lambda)$ for all $\vec a,\lambda$.

\paragraph{Derivation.} Substituting $B=-A$,
\begin{equation}
  P(\vec a,\vec b)=-\int\rho(\lambda)A(\vec a,\lambda)A(\vec b,\lambda)\dd\lambda .
\end{equation}
Introducing a third orientation $\vec c$,
\begin{equation}
  P(\vec a,\vec b)-P(\vec a,\vec c)=-\int\rho(\lambda)\left[A(\vec a,\lambda)A(\vec b,\lambda)-A(\vec a,\lambda)A(\vec c,\lambda)\right]\dd\lambda.
\end{equation}
Since $A(\vec b,\lambda)^2=1$, we can factor out $A(\vec a,\lambda)A(\vec b,\lambda)$:
\begin{equation}
  P(\vec a,\vec b)-P(\vec a,\vec c)=-\int\rho(\lambda)\left[1-A(\vec b,\lambda)A(\vec c,\lambda)\right]A(\vec a,\lambda)A(\vec b,\lambda)\dd\lambda .
\end{equation}
Now $-1\le A(\vec a,\lambda)A(\vec b,\lambda)\le1$ and $\rho(\lambda)[1-A(\vec b,\lambda)A(\vec c,\lambda)]\ge0$. It follows that
\begin{equation}
  \abs{P(\vec a,\vec b)-P(\vec a,\vec c)}\le\int\rho(\lambda)\left[1-A(\vec b,\lambda)A(\vec c,\lambda)\right]\dd\lambda.
\end{equation}
The first term of the integral is $1$ by normalization, and the second is $P(\vec b,\vec c)$ (using $A(\vec c)=-B(\vec c)$ again). We obtain \emph{Bell's inequality}
\begin{equation}\label{eq:bellineq}
  \abs{P(\vec a,\vec b)-P(\vec a,\vec c)}\le1+P(\vec b,\vec c).
\end{equation}

\paragraph{Violation by quantum mechanics.} Take $\vec a\perp\vec b$ in a plane, with $\vec c$ at $45^\circ$ from both. Then
$P_{\rm QM}(\vec a,\vec b)=0$ and $P_{\rm QM}(\vec a,\vec c)=P_{\rm QM}(\vec b,\vec c)=-\cos45^\circ\approx-0.707$, and \cref{eq:bellineq} would require
$0.707\le1-0.707=0.293$, which is false. \emph{No local hidden-variable theory can reproduce all the predictions of quantum mechanics.}

\section{The CHSH inequality}
Bell's original inequality relies on perfect anticorrelations, which are never observed exactly in experiments. Clauser, Horne, Shimony and Holt \cite{CHSH1969}
derived an inequality without this assumption. Each side chooses between two settings, $\vec a,\vec a'$ for Alice and $\vec b,\vec b'$ for Bob. Define
\begin{equation}
  S=E(\vec a,\vec b)+E(\vec a,\vec b')+E(\vec a',\vec b)-E(\vec a',\vec b'),
\end{equation}
where $E$ is the measured correlation.
\begin{theorem}[CHSH]
In a local hidden-variable theory, $|S|\le2$.
\end{theorem}
\begin{proof}
For each $\lambda$, write $A=A(\vec a,\lambda)$, $A'=A(\vec a',\lambda)$, $B=B(\vec b,\lambda)$, $B'=B(\vec b',\lambda)$, all equal to $\pm1$. Then
\begin{equation}
  AB+AB'+A'B-A'B'=A(B+B')+A'(B-B').
\end{equation}
Either $B=B'$, and then the first bracket is $\pm2$ and the second is $0$, or $B=-B'$, and the converse. In both cases the expression is $\pm2$.
Averaging over $\lambda$ with the positive weight $\rho(\lambda)$ gives $|S|\le2$.
\end{proof}
In quantum mechanics, with the singlet, $E=-\cos\theta$. Choosing the four directions in a plane at the angles $a=0$, $a'=90^\circ$, $b=45^\circ$, $b'=-45^\circ$,
\begin{equation}
  S=-\left[\cos45^\circ+\cos45^\circ+\cos45^\circ-\cos135^\circ\right]=-2\sqrt2,\qquad|S|=2\sqrt2\approx2.83>2.
\end{equation}

\begin{theorem}[Tsirelson bound]
In quantum mechanics, $|S|\le2\sqrt2$ for any state and any measurements.
\end{theorem}
\begin{proof}
Let $A,A',B,B'$ be Hermitian with eigenvalues $\pm1$ (so $A^2=\dots=I$), acting on Alice's and Bob's spaces. The CHSH operator
$\mathcal C=A\otimes(B+B')+A'\otimes(B-B')$ satisfies (expand, using $(B\pm B')^2=2I\pm\{B,B'\}$)
\begin{equation}
  \mathcal C^2=4\,I\otimes I-[A,A']\otimes[B,B'].
\end{equation}
Since $\norm{[A,A']}\le2\norm A\norm{A'}=2$ and similarly for $B$, $\norm{\mathcal C^2}\le8$, hence $|\expval{\mathcal C}|\le\norm{\mathcal C}\le2\sqrt2$.
\end{proof}
The bound $2\sqrt2$ is reached only if $[A,A']$ and $[B,B']$ are ``maximal'', i.e.\ for anticommuting observables like $\sigma_x$ and $\sigma_z$: incompatibility of
measurements is the resource behind nonlocality. Theories with $|S|$ up to $4$ (``PR boxes'') would still forbid signalling; why nature stops at $2\sqrt2$ is an open
foundational question.

\paragraph{Photons.} For polarization-entangled photons in $\ket{\Phi^+}=(\ket{HH}+\ket{VV})/\sqrt2$ and polarizers at angles $\alpha$ and $\beta$ in real space,
$E(\alpha,\beta)=\cos2(\alpha-\beta)$ (the factor 2 is the Poincaré-sphere double angle). The maximal violation is obtained at $\alpha=0$, $\alpha'=45^\circ$,
$\beta=22.5^\circ$, $\beta'=-22.5^\circ$.

\section{Experiments and loopholes}
The question can be settled by experiment. The first tests were done by Freedman and Clauser (1972). The experiments of Aspect, Grangier and Roger \cite{Aspect1982} with
entangled photons from an atomic cascade gave $S=2.697\pm0.015$, in agreement with quantum mechanics, and a later experiment by Aspect, Dalibard and Roger switched the
settings during the flight of the photons. Real experiments must close \emph{loopholes}:
\begin{itemize}
  \item the \emph{locality loophole}: the choice of setting on one side must be space-like separated from the measurement on the other side (closed by Weihs et al.\ in 1998
  with random settings and 400\,m separation);
  \item the \emph{detection loophole}: if many particles are lost, the detected subset may be unrepresentative (closed with ions in 2001, with superconducting detectors for photons later);
  \item the \emph{freedom-of-choice loophole}: the settings must be independent of $\lambda$ (``cosmic Bell tests'' used light from distant quasars to choose them).
\end{itemize}
In 2015, three groups (Delft, with electron spins separated by 1.3\,km; Vienna and NIST, with photons) performed the first loophole-free Bell tests. The 2022 Nobel prize
in physics was awarded to A.~Aspect, J.~Clauser and A.~Zeilinger for these experiments. The conclusion is that nature is not described by any local hidden-variable theory.
This is also the foundation of device-independent random number generation and cryptography: a violation of a Bell inequality certifies that the outcomes did not exist before
the measurement, whatever the inner workings of the devices.

\paragraph{No signalling.} Nonlocal correlations do not allow faster-than-light communication. Bob's marginal probabilities do not depend on Alice's choice of setting:
$\sum_AP(A,B|\vec a,\vec b)=P(B|\vec b)$, because Bob's reduced density matrix is not changed by operations on Alice's side (\cref{ch:density}). The correlations appear only
when the two lists of results are compared, which requires a classical channel.

\section{GHZ: nonlocality without inequalities\starred}
For three qubits in $\ket{\rm GHZ}=(\ket{000}+\ket{111})/\sqrt2$ one checks, using $X\ket0=\ket1$, $Y\ket0=\ii\ket1$, $Y\ket1=-\ii\ket0$,
\begin{equation}
  XXX\ket{\rm GHZ}=+\ket{\rm GHZ},\qquad XYY\ket{\rm GHZ}=YXY\ket{\rm GHZ}=YYX\ket{\rm GHZ}=-\ket{\rm GHZ}.
\end{equation}
A local hidden-variable model would assign values $x_i,y_i=\pm1$ to each qubit. The last three relations require $x_1y_2y_3=y_1x_2y_3=y_1y_2x_3=-1$; their product is
$x_1x_2x_3(y_1y_2y_3)^2=x_1x_2x_3=-1$, in contradiction with the first relation, which requires $+1$. A single run of the experiment contradicts local realism
(Greenberger, Horne, Zeilinger 1989; Mermin 1990).

\section{Entanglement-based QKD: Ekert's protocol (E91)}\label{sec:e91}
Ekert \cite{Ekert1991} proposed in 1991 a key distribution protocol in which the security is tested by a Bell inequality.
\begin{enumerate}
  \item A source (which may even be controlled by Eve) distributes pairs of qubits, ideally in the singlet state, one to Alice and one to Bob.
  \item Alice measures her qubit along one of three directions in the $xz$ plane of the Bloch sphere, chosen at random, at angles $a_1=0$, $a_2=45^\circ$, $a_3=90^\circ$.
  Bob measures along $b_1=45^\circ$, $b_2=90^\circ$, $b_3=135^\circ$. (For polarization photons, divide the angles by two.)
  \item They announce their choices of directions publicly. In the cases where they used the same direction ($a_2=b_1$ or $a_3=b_2$, a fraction $2/9$ of the pairs)
  their results are perfectly anticorrelated, and Bob flips his bit: this is the \emph{key}.
  \item With the results for $(a_1,a_3)\times(b_1,b_3)$, they compute the CHSH quantity
  \begin{equation}
    S=E(a_1,b_1)-E(a_1,b_3)+E(a_3,b_1)+E(a_3,b_3),
  \end{equation}
  which for the singlet equals $-2\sqrt2$ (check it with $E=-\cos(a-b)$). The remaining combinations are discarded.
\end{enumerate}
\paragraph{Why it is secure.} If Eve gains information about the outcomes, for instance by measuring the particles in transit, or by preparing the pairs in states with
predetermined values, the correlations become those of a local hidden-variable model with $\lambda$ = Eve's information, and then $|S|\le2$. Observing $|S|=2\sqrt2$
therefore guarantees that the outcomes did not exist before the measurements and are unknown to anybody. More quantitatively, the monogamy of entanglement implies that
if Alice and Bob are maximally entangled, Eve is uncorrelated with them. The key is created at the moment of measurement and is never ``sent'': it does not exist
before. This is the first device-independent idea in quantum cryptography. For identical devices and the singlet source, E91 is equivalent to BB84 (Bennett, Brassard,
Mermin 1992): Alice's measurement prepares Bob's qubit in one of the BB84 states.

\section{Exercises}
\begin{exercise}
Show that for $\ket{\Phi^+}$, $\mel{\Phi^+}{(\vec a\cdot\vec\sigma)\otimes(\vec b\cdot\vec\sigma)}{\Phi^+}=a_xb_x-a_yb_y+a_zb_z$. Find settings that give $|S|=2\sqrt2$.
\end{exercise}
\begin{exercise}
For the \emph{Werner state} $\rho=p\proj{\psi^-}+(1-p)I/4$, show that the maximal CHSH value is $2\sqrt2\,p$. Hence the state violates CHSH only if $p>1/\sqrt2$,
although it is entangled for $p>1/3$ (\cref{ex:werner}).
\end{exercise}
\begin{exercise}
Construct an explicit local hidden-variable model that reproduces the perfect anticorrelations $P(\vec a,\vec a)=-1$ and gives $P(\vec a,\vec b)=-1+2\theta_{ab}/\pi$
(the ``linear'' correlation). Compare with $-\cos\theta_{ab}$. (Hint: $\lambda$ is a random unit vector, and $A=\operatorname{sign}(\vec a\cdot\vec\lambda)$, $B=-\operatorname{sign}(\vec b\cdot\vec\lambda)$.)
\end{exercise}

\chapter{The density matrix}\label{ch:density}

\section{Why we need it}
For an entangled system, the wave function of an isolated subsystem does not exist. Consider two systems $A$ and $B$, initially independent:
\begin{equation}
  t=0:\qquad\ket{\psi(0)}=\ket{\psi_A(0)}\otimes\ket{\psi_B(0)}.
\end{equation}
Because of the interaction between $A$ and $B$, at later times in general $\ket{\psi(t)}\neq\ket{\psi_A(t)}\otimes\ket{\psi_B(t)}$, for example
\begin{equation}
  \ket{\psi(t)}=\frac1{\sqrt2}\left(\ket{01}-\ket{10}\right).
\end{equation}
We cannot write $\ket{\psi_A}=\ket0-\ket1$, nor $\ket{\psi_A}=\ket0+\ket1$, because we ``do not know'' where to put the minus sign. More formally, the relative phase between
$\ket0$ and $\ket1$ is not defined for $A$ alone, but only for the whole system $AB$. A second reason is classical ignorance: a source may emit $\ket0$ or $\ket1$ at
random (an unpolarized beam, a thermal atom). Both situations are described by the density matrix, which generalizes the state vector, $\ket\psi\to\hat\rho$.

\section{Definition}
\subsection{Pure states}
Start from a pure state $\ket\psi=\sum_nc_n\ket n$, $\sum_n|c_n|^2=1$. The mean value of an observable $\hat A$ is
\begin{equation}
  \mel\psi{\hat A}\psi=\sum_{n,m}c_nc_m^*\mel m{\hat A}n.
\end{equation}
We define the density operator
\begin{equation}
  \hat\rho=\proj\psi=\sum_{n,m}c_nc_m^*\ketbra nm,
\end{equation}
with matrix elements $\rho_{nm}=\mel n{\hat\rho}m=\braket n\psi\braket\psi m=c_nc_m^*$. The mean value becomes
\begin{equation}\label{eq:meantrace}
  \mel\psi{\hat A}\psi=\sum_{n,m}\rho_{nm}\mel m{\hat A}n=\sum_n\mel n{\hat\rho\hat A}n=\tr(\hat\rho\hat A)=\tr(\hat A\hat\rho).
\end{equation}
The diagonal terms $\rho_{nn}=|c_n|^2=p_n$ are the probabilities of the basis states (\emph{populations}). The off-diagonal terms $\rho_{nm}=c_nc_m^*$ describe the
phase relations between them (\emph{coherences}). The density matrix contains no unphysical global phase: $\e^{\ii\gamma}\ket\psi$ gives the same $\hat\rho$.

\subsection{Mixed states}
If the system is in the state $\ket{\psi_i}$ with probability $p_i$ (an \emph{ensemble} $\{p_i,\ket{\psi_i}\}$), the mean value of $\hat A$ is
$\sum_ip_i\mel{\psi_i}{\hat A}{\psi_i}=\tr(\hat\rho\hat A)$ with
\begin{equation}
  \hat\rho=\sum_ip_i\proj{\psi_i}.
\end{equation}
The probability of an outcome $m$ of a projective measurement is $p(m)=\tr(P_m\hat\rho)$, and the state after the measurement is $P_m\hat\rho P_m/p(m)$.
\emph{All} measurable predictions depend only on $\hat\rho$: two ensembles with the same density matrix cannot be distinguished by any experiment
(\cref{app:hjw}). For example, an equal mixture of $\ket0$ and $\ket1$ and an equal mixture of $\ket+$ and $\ket-$ both give $\hat\rho=I/2$.

\begin{theorem}[Characterization]
An operator $\hat\rho$ is the density operator of some ensemble if and only if (i) $\tr\hat\rho=1$ and (ii) $\hat\rho\ge0$ (positive semi-definite:
$\mel\phi{\hat\rho}\phi\ge0$ for all $\ket\phi$). In particular $\hat\rho$ is Hermitian, and its eigenvalues are probabilities.
\end{theorem}
\begin{proof}
If $\hat\rho=\sum p_i\proj{\psi_i}$, then $\tr\hat\rho=\sum p_i=1$ and $\mel\phi{\hat\rho}\phi=\sum p_i|\braket\phi{\psi_i}|^2\ge0$. Conversely, a positive operator is
Hermitian and has a spectral decomposition $\hat\rho=\sum_k\lambda_k\proj k$ with $\lambda_k\ge0$, $\sum\lambda_k=1$: it is the ensemble $\{\lambda_k,\ket k\}$.
\end{proof}

\subsection{Purity}
\begin{proposition}[Purity criterion]
For any density matrix $\tr(\hat\rho^2)\le1$, with equality if and only if $\hat\rho=\proj\psi$ is pure. The minimum, for dimension $d$, is $1/d$, reached by the maximally mixed state $I/d$.
\end{proposition}
\begin{proof}
In the eigenbasis, $\tr\hat\rho^2=\sum_k\lambda_k^2\le(\max_k\lambda_k)\sum_k\lambda_k\le1$, with equality iff one eigenvalue is $1$ and all others $0$, i.e.\ $\hat\rho$ is a
rank-one projector. The minimum of $\sum\lambda_k^2$ under $\sum\lambda_k=1$ is at $\lambda_k=1/d$ (Cauchy--Schwarz: $1=(\sum\lambda_k)^2\le d\sum\lambda_k^2$).
\end{proof}

\subsection{The Bloch ball}
For a qubit, the decomposition on Pauli matrices gives, with $\tr\hat\rho=1$ and $\hat\rho$ Hermitian,
\begin{equation}\label{eq:blochball}
  \hat\rho=\frac{I+\vec r\cdot\vec\sigma}2=\frac12\begin{pmatrix}1+r_z&r_x-\ii r_y\\ r_x+\ii r_y&1-r_z\end{pmatrix},\qquad\vec r=\tr(\hat\rho\vec\sigma)=\expval{\vec\sigma}.
\end{equation}
The eigenvalues are $(1\pm|\vec r|)/2$, so positivity requires $|\vec r|\le1$: density matrices fill the \emph{Bloch ball}. The purity is
$\tr\hat\rho^2=\frac{1+|\vec r|^2}2$: pure states are on the surface, mixed states inside, and $I/2$ is the centre. Decoherence (\cref{ch:open}) shrinks the Bloch vector.
The Bloch equation \cref{eq:precession} remains valid for mixed states, since it is linear in $\vec r$.

\section{Entanglement and the reduced density matrix}\label{sec:reduced}
Consider the entangled state of two qubits $A$ and $B$ ($A$ is the ``system of interest'' and $B$ the ``environment''):
\begin{equation}
  \ket\psi=\frac1{\sqrt2}\left(\ket{01}+\ket{10}\right),\qquad\hat\rho=\proj\psi=\frac12\left(\ket{01}+\ket{10}\right)\left(\bra{01}+\bra{10}\right).
\end{equation}
In the basis $\{\ket{00},\ket{01},\ket{10},\ket{11}\}$,
\begin{equation}
  \hat\rho=\frac12\begin{pmatrix}0&0&0&0\\0&1&1&0\\0&1&1&0\\0&0&0&0\end{pmatrix}.
\end{equation}
As discussed above, subsystem $A$ alone has no wave function. It does, however, have a density matrix, obtained by \emph{tracing out} subsystem $B$:
\begin{equation}
  \hat\rho_A=\tr_B\hat\rho=\mel{0_B}{\hat\rho}{0_B}+\mel{1_B}{\hat\rho}{1_B}=\frac12\begin{pmatrix}1&0\\0&1\end{pmatrix}
\end{equation}
in the basis $\{\ket0,\ket1\}$. The off-diagonal terms (coherences) of $\hat\rho_A$ have disappeared: they would involve the overlap $\braket{0_B}{1_B}=0$ of the orthogonal states of $B$
with which $\ket0_A$ and $\ket1_A$ are correlated. \emph{This is precisely the mechanism of decoherence.} The coherence is not destroyed, but transferred to the correlations between
the system and its environment, and it is no longer accessible from the system alone.

\subsection{The partial trace}
The partial trace is defined on product operators by
\begin{equation}
  \tr_B\left(\ketbra{a_1}{a_2}\otimes\ketbra{b_1}{b_2}\right)=\ketbra{a_1}{a_2}\,\tr\left(\ketbra{b_1}{b_2}\right),
\end{equation}
and extended by linearity. The trace of an outer product is
\begin{equation}
  \tr\ketbra xy=\sum_n\braket nx\braket yn=\sum_n\braket yn\braket nx=\braket yx,
\end{equation}
where we used the completeness relation $\sum_n\proj n=I$. Hence
\begin{equation}\label{eq:ptracerule}
  \tr_B\left(\ketbra{a_1}{a_2}\otimes\ketbra{b_1}{b_2}\right)=\ketbra{a_1}{a_2}\braket{b_2}{b_1}.
\end{equation}
In components, $(\rho_A)_{ij}=\sum_\mu\rho_{i\mu,j\mu}$: sum the diagonal of each block over the index of $B$.

\begin{proposition}[Why the partial trace]
$\hat\rho_A=\tr_B\hat\rho$ is the unique operator such that $\tr(\hat\rho_A\hat O_A)=\tr\left(\hat\rho\,(\hat O_A\otimes I_B)\right)$ for every observable $\hat O_A$ of $A$.
It therefore gives all the predictions for measurements on $A$ alone.
\end{proposition}
\begin{proof}
Write $\hat\rho=\sum\rho_{i\mu,j\nu}\ketbra ij\otimes\ketbra\mu\nu$. Then $\tr(\hat\rho\,\hat O_A\otimes I)=\sum\rho_{i\mu,j\nu}\mel j{\hat O_A}i\delta_{\mu\nu}=\tr(\hat\rho_A\hat O_A)$.
Uniqueness: if $\tr(X\hat O)=0$ for all Hermitian $\hat O$, then $X=0$.
\end{proof}
\begin{corollary}[No signalling]
A unitary $U_B$ or a measurement performed by Bob on $B$ does not change $\hat\rho_A$. For a unitary: $\tr_B[(I\otimes U_B)\hat\rho(I\otimes U_B^\dagger)]=\tr_B\hat\rho$, by
cyclicity of the trace over $B$. For a measurement with unknown outcome: $\sum_m\tr_B[(I\otimes P_m)\hat\rho(I\otimes P_m)]=\tr_B\hat\rho$ using $\sum_mP_m=I$.
Nothing Bob does can be detected by Alice without classical communication.
\end{corollary}

\begin{example}
For the Bell state $\ket\psi=(\ket{00}+\ket{11})/\sqrt2$,
\begin{equation}
  \hat\rho=\frac12\left(\ketbra{00}{00}+\ketbra{11}{11}+\ketbra{00}{11}+\ketbra{11}{00}\right).
\end{equation}
Tracing out the second qubit term by term with \cref{eq:ptracerule},
\begin{equation}
  \hat\rho'=\tr_B\hat\rho=\frac12\left(\proj0\braket00+\ketbra10\braket01+\ketbra01\braket10+\proj1\braket11\right)=\frac{\proj0+\proj1}2=\frac I2.
\end{equation}
This is a mixed state, since $\tr[(I/2)^2]=1/2<1$. A maximally entangled pure state of two qubits has maximally mixed reduced states.
\end{example}

\subsection{Entanglement entropy}
The \emph{von Neumann entropy} $S(\hat\rho)=-\tr(\hat\rho\log_2\hat\rho)=-\sum_k\lambda_k\log_2\lambda_k$ measures how mixed a state is: $0$ for pure states, $\log_2d$
for $I/d$. For a pure state of $AB$, $S(\hat\rho_A)=S(\hat\rho_B)$ (\cref{app:schmidt}) and this common value, the \emph{entanglement entropy}, is the standard measure of
entanglement: $0$ for product states, $1$ ebit for a Bell state. It counts how many Bell pairs can be distilled from many copies of the state by local operations and
classical communication (Bennett et al., 1996).

\section{Equation of motion}
In the pure case, the time evolution of the density matrix follows from the Schrödinger equation:
\begin{equation}
  \frac{\partial\hat\rho}{\partial t}=\left(\frac{\partial}{\partial t}\ket\psi\right)\bra\psi+\ket\psi\left(\frac{\partial}{\partial t}\bra\psi\right)
  =\frac1{\ii\hbar}\left(\hat H\proj\psi-\proj\psi\hat H\right)=\frac1{\ii\hbar}[\hat H,\hat\rho].
\end{equation}
By linearity the same \emph{von Neumann equation} holds for mixtures. Its solution is $\hat\rho(t)=U\hat\rho(0)U^\dagger$. In the eigenbasis of $\hat H$,
$\rho_{nm}(t)=\rho_{nm}(0)\e^{-\ii(E_n-E_m)t/\hbar}$: populations are constant and coherences rotate at the Bohr frequencies. Unitary evolution conserves $\tr\hat\rho^2$
and the entropy: pure states stay pure. Decoherence requires coupling to an environment (\cref{ch:open}).

\section{Thermal states}\label{sec:thermal}
A system in contact with a reservoir at temperature $T$ reaches thermal equilibrium. Its density matrix commutes with $\hat H$ and is diagonal in the energy basis:
\begin{equation}
  \rho_{nm}=p_n\delta_{nm},\qquad p_n=\frac{\e^{-\beta E_n}}Z,\quad Z=\sum_n\e^{-\beta E_n},\quad\beta=\frac1{k_BT},
\end{equation}
i.e.\ $\hat\rho=\e^{-\beta\hat H}/Z$.
\begin{origin}[: the Gibbs state from maximum entropy]
Among all states with a given mean energy $\tr(\hat\rho\hat H)=E$, the Gibbs state maximizes the entropy $S=-k_B\tr\hat\rho\ln\hat\rho$. With Lagrange multipliers for the
normalization and the energy, $\delta[S/k_B-\beta(\tr\hat\rho\hat H-E)-\mu(\tr\hat\rho-1)]=0$ gives $-\ln\hat\rho-1-\beta\hat H-\mu=0$, i.e.\ $\hat\rho\propto\e^{-\beta\hat H}$.
The multiplier $\beta$ is identified with $1/k_BT$ through $\partial S/\partial E=1/T$. The same state is obtained as the reduced state of a small system weakly coupled to a large
reservoir (canonical ensemble), and as the steady state of the master equations of \cref{ch:open}.
\end{origin}
For a qubit with $H=-\frac{\hbar\omega}2Z$ (ground state $\ket0$), $\expval Z=\tanh(\beta\hbar\omega/2)$, and the excited-state population is
$p_1=1/(1+\e^{\beta\hbar\omega})$. For a superconducting qubit at $\omega/2\pi=5$\,GHz and $T=20$\,mK, $\hbar\omega/k_BT=12$ and $p_1\approx6\times10^{-6}$: the qubit is
reset by simply waiting. At $T=300$\,K, $p_1\approx1/2$: this is why superconducting qubits need dilution refrigerators, and why room-temperature NMR uses
pseudo-pure states (\cref{sec:nmr}).

\subsection{The harmonic oscillator}\label{sec:oscillator}
\begin{origin}[: quantization of the harmonic oscillator]
The harmonic oscillator, $H=\frac{\hat p^2}{2m}+\frac12m\omega^2\hat x^2$ with $[\hat x,\hat p]=\ii\hbar$, describes any system close to a minimum of its potential
energy: a vibrating molecule or ion in a trap, an LC circuit (\cref{sec:lc}), a mode of the electromagnetic field. Define
\begin{equation}
  \hat a=\sqrt{\frac{m\omega}{2\hbar}}\left(\hat x+\frac{\ii\hat p}{m\omega}\right),\qquad \hat a^\dagger=\sqrt{\frac{m\omega}{2\hbar}}\left(\hat x-\frac{\ii\hat p}{m\omega}\right).
\end{equation}
Then $[\hat a,\hat a^\dagger]=1$ and $H=\hbar\omega(\hat a^\dagger\hat a+\frac12)$. With $\hat N=\hat a^\dagger\hat a$, $[\hat N,\hat a]=-\hat a$: if $\hat N\ket n=n\ket n$, then
$\hat a\ket n$ is an eigenvector with eigenvalue $n-1$. Since $\norm{\hat a\ket n}^2=n\ge0$, the lowering must stop at a state $\ket0$ with $\hat a\ket0=0$, so $n=0,1,2,\dots$, and
\begin{equation}
  E_n=\hbar\omega\left(n+\tfrac12\right),\qquad\hat a\ket n=\sqrt n\ket{n-1},\qquad\hat a^\dagger\ket n=\sqrt{n+1}\ket{n+1}.
\end{equation}
In dimensionless units ($\hbar=m=\omega=1$), $\hat x=(\hat a+\hat a^\dagger)/\sqrt2$ and $\hat p=-\ii(\hat a-\hat a^\dagger)/\sqrt2$. The levels are \emph{equally spaced}: this is
why a harmonic oscillator alone is not a qubit (\cref{ch:implementations}). The quantum $\ket n$ is called a phonon (vibration) or a photon (light, microwaves).
\end{origin}
For the harmonic oscillator, measuring the energies from the ground state ($E_n=\hbar\omega n$),
\begin{equation}\label{eq:thermalHO}
  \hat\rho=\sum_n\left(1-\e^{-\beta\hbar\omega}\right)\e^{-\beta\hbar\omega n}\proj n,
\end{equation}
where the prefactor ensures $\sum_np_n=1$ (geometric series). The mean occupation is the Bose--Einstein (Planck) distribution
\begin{equation}
  \bar n=\sum_nnp_n=\frac1{\e^{\beta\hbar\omega}-1},\qquad p_n=\frac{\bar n^n}{(\bar n+1)^{n+1}}.
\end{equation}
(Proof: $\sum_nnx^n(1-x)=x/(1-x)$ with $x=\e^{-\beta\hbar\omega}$; and $x=\bar n/(\bar n+1)$, $1-x=1/(\bar n+1)$.)

\section{Exercises}
\begin{exercise}
Show that $\hat\rho=\frac12\begin{pmatrix}1&1\\1&1\end{pmatrix}$ is pure, find its state vector and Bloch vector. Show that $\frac14\begin{pmatrix}3&1\\1&1\end{pmatrix}$ is
mixed, and compute its purity and entropy.
\end{exercise}
\begin{exercise}
Compute the reduced density matrices and the entanglement entropy of $\cos\alpha\ket{00}+\sin\alpha\ket{11}$. For which $\alpha$ is the entanglement maximal?
\end{exercise}
\begin{exercise}
Show that the thermal state of a qubit has Bloch vector $\vec r=\tanh(\beta|\vec h|)\,\vec h/|\vec h|$ for $H=-\vec h\cdot\vec\sigma$.
\end{exercise}

\chapter{Open quantum systems and decoherence}\label{ch:open}

A qubit is never perfectly isolated. Its coupling to uncontrolled degrees of freedom (phonons, photons, nuclear spins, charge fluctuators, \dots) makes its evolution
non-unitary: populations relax and coherences decay. This chapter derives the tools to describe it: quantum channels (Kraus operators) and master equations
(Lindblad form), and applies them to the damped harmonic oscillator and to the qubit relaxation times $T_1$ and $T_2$.

\section{Quantum channels and Kraus operators}
\subsection{Derivation from a system--environment model}
Let the system $S$ be in the state $\hat\rho$, initially uncorrelated with an environment $E$ in the pure state $\ket{e_0}$ (a mixed environment can always be
purified, \cref{app:purification}). The total system evolves unitarily, $U$, and we then look at $S$ alone:
\begin{equation}
  \mathcal E(\hat\rho)=\tr_E\left[U(\hat\rho\otimes\proj{e_0})U^\dagger\right]=\sum_k\mel{e_k}{U}{e_0}\,\hat\rho\,\mel{e_0}{U^\dagger}{e_k}=\sum_kK_k\hat\rho K_k^\dagger,
\end{equation}
where $\{\ket{e_k}\}$ is a basis of $E$ and $K_k=\mel{e_k}U{e_0}$ are operators on $S$, the \emph{Kraus operators}. Unitarity of $U$ gives the completeness relation
\begin{equation}
  \sum_kK_k^\dagger K_k=\mel{e_0}{U^\dagger U}{e_0}=I,
\end{equation}
which guarantees $\tr\mathcal E(\hat\rho)=1$. The map $\mathcal E$ is linear, trace preserving, and \emph{completely positive}: it maps density matrices to density matrices even
when it acts on a part of a larger entangled system. Conversely (Stinespring, Kraus), every completely positive trace-preserving map, called a \emph{quantum channel},
has this form. The Kraus representation is not unique: $K'_j=\sum_ku_{jk}K_k$ with a unitary $u$ gives the same channel (a change of basis of the environment).

\subsection{Standard qubit channels}
With probability $p$ an error occurs:
\begin{align}
  \text{bit flip:}&\quad\mathcal E(\hat\rho)=(1-p)\hat\rho+pX\hat\rho X,\\
  \text{phase flip (dephasing):}&\quad\mathcal E(\hat\rho)=(1-p)\hat\rho+pZ\hat\rho Z,\\
  \text{depolarizing:}&\quad\mathcal E(\hat\rho)=(1-p)\hat\rho+\tfrac p3\left(X\hat\rho X+Y\hat\rho Y+Z\hat\rho Z\right).
\end{align}
In the Bloch-ball picture \cref{eq:blochball}, using $Z\vec\sigma Z=(-\sigma_x,-\sigma_y,\sigma_z)$, the phase flip maps $(r_x,r_y,r_z)\to((1-2p)r_x,(1-2p)r_y,r_z)$: the ball
is squeezed into a cigar along $z$, and the coherences $\rho_{01}$ are multiplied by $1-2p$. The depolarizing channel shrinks the ball uniformly, $\vec r\to(1-\frac43p)\vec r$,
and can also be written $\mathcal E(\hat\rho)=(1-\frac43p)\hat\rho+\frac43p\frac I2$, using the identity $\hat\rho+X\hat\rho X+Y\hat\rho Y+Z\hat\rho Z=2I$ for any $\hat\rho$ of unit trace.

\paragraph{Amplitude damping (spontaneous emission).} An excited atom $\ket1$ emits a photon into the empty environment with probability $p$:
\begin{equation}
  U\ket0\ket{0_E}=\ket0\ket{0_E},\qquad U\ket1\ket{0_E}=\sqrt{1-p}\ket1\ket{0_E}+\sqrt p\ket0\ket{1_E}.
\end{equation}
The Kraus operators $K_k=\mel{k_E}U{0_E}$ are
\begin{equation}
  K_0=\begin{pmatrix}1&0\\0&\sqrt{1-p}\end{pmatrix},\qquad K_1=\begin{pmatrix}0&\sqrt p\\0&0\end{pmatrix},
\end{equation}
and
\begin{equation}
  \mathcal E\begin{pmatrix}\rho_{00}&\rho_{01}\\ \rho_{10}&\rho_{11}\end{pmatrix}=\begin{pmatrix}\rho_{00}+p\rho_{11}&\sqrt{1-p}\,\rho_{01}\\ \sqrt{1-p}\,\rho_{10}&(1-p)\rho_{11}\end{pmatrix}.
\end{equation}
For a continuous decay, $p=1-\e^{-\gamma t}$: the population decays as $\e^{-\gamma t}$ and the coherence as $\e^{-\gamma t/2}$. This factor $1/2$ is the origin of the relation
$T_2=2T_1$ for purely radiative decay, found below with the master equation. The same model is a beam splitter of transmission $1-p$ acting on a photonic qubit: \emph{losses
are amplitude damping}.

\section{The Lindblad master equation}
\subsection{Structure}
For a memoryless (Markovian) environment, the evolution over a short time $\dd t$ is a channel close to the identity. Its Kraus operators must be of the form
$K_0=I-(\frac\ii\hbar H+\frac12\sum_kL_k^\dagger L_k)\dd t$ and $K_k=\sqrt{\dd t}\,L_k$ ($k\ge1$), to satisfy $\sum K^\dagger K=I$ to first order in $\dd t$. Then
$\hat\rho(t+\dd t)=\sum_kK_k\hat\rho K_k^\dagger$ gives, to first order,
\begin{equation}\label{eq:lindblad}
  \dot{\hat\rho}=-\frac\ii\hbar[H,\hat\rho]+\sum_k\Dis[L_k]\hat\rho,\qquad\Dis[L]\hat\rho=L\hat\rho L^\dagger-\frac12\left\{L^\dagger L,\hat\rho\right\}.
\end{equation}
This is the \emph{Lindblad} (Gorini--Kossakowski--Sudarshan--Lindblad, 1976) master equation, the most general generator of a Markovian completely positive evolution.
The \emph{jump operators} $L_k$ describe the processes induced by the environment: $L\hat\rho L^\dagger$ is the ``jump'' (e.g.\ emission of a photon, detected or not), and the
anticommutator is the corresponding loss of norm of the no-jump evolution.

\subsection{Microscopic derivation for an oscillator coupled to a bath\starred}
\begin{origin}[: the Born--Markov approximation]
Consider a mode $\hat a$ of frequency $\omega$ (a cavity mode, a vibrational mode, or an LC resonator) coupled to a bath of oscillators $\hat b_k$:
\begin{equation}
  H=\hbar\omega\hat a^\dagger\hat a+\sum_k\hbar\omega_k\hat b_k^\dagger\hat b_k+\hbar\sum_kg_k\left(\hat a\hat b_k^\dagger+\hat a^\dagger\hat b_k\right).
\end{equation}
(This form follows from a linear coupling, e.g.\ capacitive or dipolar, in the rotating-wave approximation.) In the interaction picture, to second order in the coupling,
and assuming (i) \emph{Born}: the bath is large and stays in its thermal state, $\hat\rho_{\rm tot}(t)\approx\hat\rho(t)\otimes\hat\rho_B$; (ii) \emph{Markov}: the bath
correlation functions $\expval{\hat B(t)\hat B^\dagger(t-\tau)}$, with $\hat B=\sum_kg_k\hat b_k\e^{-\ii\omega_kt}$, decay on a time $\tau_c$ much shorter than the time
scale of the system dynamics, one finds
\begin{equation}
  \dot{\hat\rho}_I=-\int_0^\infty\dd\tau\,\tr_B\left[V_I(t),\left[V_I(t-\tau),\hat\rho_I(t)\otimes\hat\rho_B\right]\right]/\hbar^2 .
\end{equation}
The time integrals produce $\int_0^\infty\dd\tau\,\e^{\ii(\omega-\omega_k)\tau}=\pi\delta(\omega-\omega_k)+\ii\,\mathcal P\frac1{\omega-\omega_k}$. The principal part
gives a small frequency shift (Lamb shift), absorbed in $\omega$. The delta function gives the rates. With the thermal bath occupations
$\expval{\hat b_k^\dagger\hat b_k}=\bar n(\omega_k)$, the result is \cref{eq:lindblad} with the two jump operators
\begin{equation}
  L_1=\sqrt{\gamma(\bar n+1)}\,\hat a\quad\text{(emission)},\qquad L_2=\sqrt{\gamma\bar n}\,\hat a^\dagger\quad\text{(absorption)},
\end{equation}
where $\bar n=\bar n(\omega)=(\e^{\hbar\omega/k_BT}-1)^{-1}$ and
\begin{equation}
  \gamma=2\pi\sum_kg_k^2\delta(\omega-\omega_k)=2\pi g(\omega)^2D(\omega),
\end{equation}
with $D$ the density of bath modes. This is Fermi's golden rule. Emission is stimulated ($\bar n+1$) and absorption proportional to $\bar n$, as in Einstein's theory of
radiation (1917). For an atom in free space, the same calculation with the vacuum electromagnetic modes gives the Wigner--Weisskopf spontaneous emission rate
$\gamma=\omega^3|d_{eg}|^2/(3\pi\varepsilon_0\hbar c^3)$, which gives a lifetime of $1.6$\,ns for the $2p$ state of hydrogen. Engineering $D(\omega)$, for example with a
cavity, changes $\gamma$ (Purcell effect): this is used both to protect qubits and to read them out quickly.
\end{origin}

\section{The damped harmonic oscillator}
With $H=\hbar\omega\hat a^\dagger\hat a$ and a thermal bath,
\begin{equation}\label{eq:dampedHO}
  \dot{\hat\rho}=-\ii\omega[\hat a^\dagger\hat a,\hat\rho]+\gamma(\bar n+1)\Dis[\hat a]\hat\rho+\gamma\bar n\Dis[\hat a^\dagger]\hat\rho .
\end{equation}
Explicitly,
\begin{equation}
  \dot{\hat\rho}=-\ii\omega[\hat a^\dagger\hat a,\hat\rho]+\frac{\gamma(\bar n+1)}2\left(2\hat a\hat\rho\hat a^\dagger-\hat a^\dagger\hat a\hat\rho-\hat\rho\hat a^\dagger\hat a\right)
  +\frac{\gamma\bar n}2\left(2\hat a^\dagger\hat\rho\hat a-\hat a\hat a^\dagger\hat\rho-\hat\rho\hat a\hat a^\dagger\right).
\end{equation}

\paragraph{Matrix elements.} In the Fock basis, using $\hat a\ket n=\sqrt n\ket{n-1}$, $\hat a^\dagger\ket n=\sqrt{n+1}\ket{n+1}$:
$(\hat a\hat\rho\hat a^\dagger)_{mn}=\sqrt{(m+1)(n+1)}\,\rho_{m+1,n+1}$, $(\hat a^\dagger\hat a\hat\rho)_{mn}=m\rho_{mn}$, $(\hat\rho\hat a^\dagger\hat a)_{mn}=n\rho_{mn}$,
$(\hat a^\dagger\hat\rho\hat a)_{mn}=\sqrt{mn}\,\rho_{m-1,n-1}$, $(\hat a\hat a^\dagger\hat\rho)_{mn}=(m+1)\rho_{mn}$. Therefore
\begin{multline}
  \dot\rho_{mn}=-\ii\omega(m-n)\rho_{mn}+\gamma(\bar n+1)\left[\sqrt{(m+1)(n+1)}\,\rho_{m+1,n+1}-\frac{m+n}2\rho_{mn}\right]\\
  +\gamma\bar n\left[\sqrt{mn}\,\rho_{m-1,n-1}-\frac{m+n+2}2\rho_{mn}\right].
\end{multline}
The equations for populations and coherences decouple, and the coherence $\rho_{mn}$ only couples to $\rho_{m\pm1,n\pm1}$, which oscillate at the same frequency.

\paragraph{Populations.} For $m=n$, $p_n=\rho_{nn}$:
\begin{equation}\label{eq:rateeq}
  \dot p_n=\gamma(\bar n+1)\left[(n+1)p_{n+1}-np_n\right]+\gamma\bar n\left[np_{n-1}-(n+1)p_n\right].
\end{equation}
This is a classical rate equation: the level $n$ is fed from $n+1$ by emission and from $n-1$ by absorption, and depleted by both processes.
\begin{proposition}
The steady state of \cref{eq:rateeq} is the thermal distribution $p_n^{ss}=\bar n^n/(\bar n+1)^{n+1}$, and $\expval{\hat a^\dagger\hat a}$ relaxes to $\bar n$ at the rate $\gamma$:
\begin{equation}
  \frac{\dd}{\dd t}\expval{\hat n}=-\gamma\left(\expval{\hat n}-\bar n\right),\qquad\expval{\hat n}(t)=\bar n+\left(\expval{\hat n}(0)-\bar n\right)\e^{-\gamma t}.
\end{equation}
\end{proposition}
\begin{proof}
\emph{Steady state.} Write \cref{eq:rateeq} as $\dot p_n=J_{n+1}-J_n$ with the net downward flux $J_n=\gamma(\bar n+1)np_n-\gamma\bar n\,np_{n-1}$ between levels $n$ and $n-1$.
Since $J_0=0$, a stationary solution with $J_n=0$ for all $n$ (detailed balance) has $p_n/p_{n-1}=\bar n/(\bar n+1)=\e^{-\beta\hbar\omega}$, i.e.\ the thermal state \cref{eq:thermalHO}.
\emph{Mean occupation.} For an operator $A$, $\frac{\dd}{\dd t}\expval A=\tr(A\dot{\hat\rho})$; by cyclicity of the trace,
$\tr(A\,\Dis[L]\hat\rho)=\expval{L^\dagger AL-\frac12\{L^\dagger L,A\}}$. For $A=\hat n$ and $L=\hat a$: $\hat a^\dagger\hat n\hat a-\hat n^2=\hat a^\dagger\hat a^\dagger\hat a\hat a-\hat n^2=\hat n(\hat n-1)-\hat n^2=-\hat n$.
For $L=\hat a^\dagger$: $\hat a\hat n\hat a^\dagger-\frac12\{\hat a\hat a^\dagger,\hat n\}=(\hat n+1)^2-(\hat n+1)\hat n=\hat n+1$. The Hamiltonian commutes with $\hat n$. Hence
$\frac{\dd}{\dd t}\expval{\hat n}=-\gamma(\bar n+1)\expval{\hat n}+\gamma\bar n(\expval{\hat n}+1)=-\gamma\expval{\hat n}+\gamma\bar n$.
\end{proof}
Similarly $\frac{\dd}{\dd t}\expval{\hat a}=-(\ii\omega+\frac\gamma2)\expval{\hat a}$: the amplitude of a classical oscillation decays at the rate $\gamma/2$, the energy at the rate $\gamma$,
exactly as for a classical damped oscillator with quality factor $Q=\omega/\gamma$.

\section{Truncation to a qubit: $T_1$ and $T_2$}
If only the states $\ket0$ and $\ket1$ are involved (because the oscillator is anharmonic, or because it is an atom), we replace
\begin{equation}
  \hat a\to\sigma^-=\ketbra01,\qquad\hat a^\dagger\to\sigma^+=\ketbra10,\qquad\hat a^\dagger\hat a\to\sigma^+\sigma^-=\proj1=\frac{I-Z}2,
\end{equation}
where $Z=\proj0-\proj1$ as in \cref{ch:bloch}. (The original notes write $\proj1=\frac12(1+\sigma_z)$, which corresponds to the opposite, ``atomic'', sign convention
$\sigma_z=\proj1-\proj0$; the equations below are the same in both conventions when written with $\rho_{00},\rho_{11},\rho_{01},\rho_{10}$.) The master equation becomes
\begin{equation}
  \dot{\hat\rho}=-\ii\omega[\proj1,\hat\rho]+\gamma(\bar n+1)\Dis[\sigma^-]\hat\rho+\gamma\bar n\Dis[\sigma^+]\hat\rho .
\end{equation}
In components, with $\rho_{00}=1-\rho_{11}$ and $\rho_{01}=\rho_{10}^*$:
\begin{align}
  \dot\rho_{11}&=-\gamma(\bar n+1)\rho_{11}+\gamma\bar n\rho_{00}=-\gamma(2\bar n+1)\rho_{11}+\gamma\bar n,\\
  \dot\rho_{10}&=-\left(\ii\omega+\frac{\gamma(2\bar n+1)}2\right)\rho_{10},\qquad\dot\rho_{01}=\left(\ii\omega-\frac{\gamma(2\bar n+1)}2\right)\rho_{01}.
\end{align}
(Check: $(\sigma^-\hat\rho\sigma^+)_{11}=0$, $(\sigma^+\sigma^-\hat\rho)_{11}=\rho_{11}$, $(\sigma^+\hat\rho\sigma^-)_{11}=\rho_{00}$, $(\sigma^-\sigma^+\hat\rho)_{11}=0$; and for the
coherence each dissipator contributes $-\frac12\times$ its rate.) These are the \emph{Bloch equations}, with the parameters
\begin{gather}
  T_1^{-1}=\gamma(2\bar n+1),\qquad T_2=2T_1,\\
  \rho_{11}^{ss}=\frac{\bar n}{2\bar n+1},\qquad\frac{\rho_{11}^{ss}}{\rho_{00}^{ss}}=\frac{\bar n}{\bar n+1},\qquad
  \expval{\proj0-\proj1}_{ss}=\frac1{2\bar n+1},
\end{gather}
and the solutions
\begin{equation}
  \rho_{11}(t)=\rho_{11}^{ss}+\left[\rho_{11}(0)-\rho_{11}^{ss}\right]\e^{-t/T_1},\qquad\rho_{01}(t)=\rho_{01}(0)\,\e^{\ii\omega t}\e^{-t/T_2}.
\end{equation}
Note that the steady state $\rho_{11}/\rho_{00}=\e^{-\beta\hbar\omega}$ is again thermal (here for a two-level system, $\frac{1}{2\bar n+1}=\tanh\frac{\beta\hbar\omega}2$).
$T_1$ is the \emph{energy relaxation} (longitudinal) time; $T_2$ the \emph{coherence} (transverse) time. On the Bloch sphere, $r_z$ relaxes to its equilibrium value in $T_1$
and the transverse components spiral inwards in $T_2$.

\paragraph{Pure dephasing.} Fluctuations of the qubit frequency (e.g.\ from charge or flux noise, or nuclear spins) randomize the phase without exchanging energy. They are
described by the additional jump operator $L_\varphi=\sqrt{\gamma_\varphi/2}\,Z$, which does not change the populations and damps the coherences at the rate $\gamma_\varphi$
(since $\Dis[Z]\hat\rho=Z\hat\rho Z-\hat\rho$ multiplies $\rho_{01}$ by $-2$). In general
\begin{equation}
  \frac1{T_2}=\frac1{2T_1}+\frac1{T_\varphi},\qquad T_2\le2T_1 .
\end{equation}
For slow (quasi-static) noise the decay is not exponential: a static random detuning $\delta$ with Gaussian distribution of width $\sigma$ gives
$\overline{\e^{\ii\delta t}}=\e^{-\sigma^2t^2/2}$ ($T_2^*$ decay). Such inhomogeneous dephasing can be reversed by a $\pi$ pulse in the middle of the evolution (Hahn spin echo,
1950), which is the simplest example of \emph{dynamical decoupling}.

\paragraph{Driven qubit: optical Bloch equations.} Adding the Rabi drive \cref{eq:HRWA} in the rotating frame gives the optical Bloch equations, which describe Rabi oscillations
damped on the time scales $T_1,T_2$, and a steady state with excited population $\rho_{11}^{ss}=\frac{(\Omega^2/4)(2T_1/T_2)}{\Delta^2+T_2^{-2}+\Omega^2T_1/T_2}$ (for $\bar n=0$; for radiative decay, $T_2=2T_1$, this is the familiar $\frac{\Omega^2/4}{\Delta^2+\gamma^2/4+\Omega^2/2}$),
which saturates at $1/2$ for strong driving.

\section{Decoherence of an observable: an example}
Consider the position of a harmonic oscillator, $\hat x=(\hat a+\hat a^\dagger)/\sqrt2$, in the state $\ket\psi=(\ket0+\ket1)/\sqrt2$. The mean value $\mel\psi{\hat x}\psi$
oscillates in time. In the density matrix formalism,
\begin{equation}
  \expval{\hat x}=\tr(\hat\rho\hat x)=\sum_n\mel n{\hat\rho\frac{\hat a+\hat a^\dagger}{\sqrt2}}n
  =\frac1{\sqrt2}\sum_n\left(\sqrt{n+1}\,\rho_{n,n+1}+\sqrt n\,\rho_{n,n-1}\right).
\end{equation}
In our case only $n=0,1$ contribute: $\expval{\hat x}=\frac1{\sqrt2}(\rho_{01}+\rho_{10})=\sqrt2\operatorname{Re}\rho_{01}$, which involves only the off-diagonal terms. With
$\rho_{01}(t)=\frac12\e^{\ii\omega t}\e^{-t/T_2}$ we get $\expval{\hat x}=\frac1{\sqrt2}\cos(\omega t)\,\e^{-t/T_2}$: a damped oscillation. If the coherences have decayed, there are no
oscillations:
\begin{equation}
  \hat\rho=\begin{pmatrix}1/2&0\\0&1/2\end{pmatrix}\quad\Longrightarrow\quad\dot{\hat\rho}=0\ \text{(unitary part)},
\end{equation}
since $\hat\rho$ commutes with the Hamiltonian $\hat H=\hbar\omega(\hat a^\dagger\hat a+\frac12)$. (For $\bar n$ finite it still relaxes towards the thermal populations.)
\emph{Physical observables that ``move'' are carried by the coherences; decoherence freezes them.} This is why quantum superpositions of macroscopically distinct states
(Schrödinger cats) are never observed: their coherences decay at a rate proportional to the square of their ``distance'', which is enormous for macroscopic objects.

\section{Exercises}
\begin{exercise}
Show that the depolarizing channel with $p=3/4$ maps every state to $I/2$, and that the composition of two phase-flip channels with probabilities $p_1,p_2$ is a phase-flip
channel with $1-2p=(1-2p_1)(1-2p_2)$. Deduce that repeated dephasing gives an exponential decay of the coherence.
\end{exercise}
\begin{exercise}
Derive $\frac{\dd}{\dd t}\expval{\hat a}=-(\ii\omega+\frac\gamma2)\expval{\hat a}$ from \cref{eq:dampedHO}, and show that it does not depend on $\bar n$.
\end{exercise}
\begin{exercise}
A transmon has $T_1=100\,\mu$s and $T_2=80\,\mu$s. Compute $T_\varphi$. A two-qubit gate takes $50$\,ns: estimate the error per gate due to decoherence.
(Answer: $T_\varphi\approx133\,\mu$s; error $\sim t/T_2\sim6\times10^{-4}$.)
\end{exercise}

\chapter{Detecting entanglement: the Peres criterion}\label{ch:peres}

\section{Separable mixed states}
For pure states, entanglement means ``not a product'' (\cref{prop:det}). For mixed states the definition must allow for classical correlations.
\begin{definition}[Werner, 1989]
A density matrix of $AB$ is \emph{separable} if it is a convex combination of product states,
\begin{equation}\label{eq:separable}
  \hat\rho=\sum_ip_i\,\hat\rho_i^A\otimes\hat\rho_i^B,\qquad p_i\ge0,\ \sum_ip_i=1.
\end{equation}
Otherwise it is \emph{entangled}.
\end{definition}
Separable states are exactly those that Alice and Bob can prepare with local operations and classical communication (LOCC): they agree on $i$ with probability $p_i$
(e.g.\ by a phone call) and prepare $\hat\rho_i^A$ and $\hat\rho_i^B$ locally. For example $\frac12(\proj{00}+\proj{11})$ is correlated (both qubits are always equal in the
$Z$ basis) but separable. In contrast, $\proj{\Phi^+}$ has the same diagonal but also the coherences $\ketbra{00}{11}$, which make it entangled: it is correlated in
every basis. Deciding whether a given $\hat\rho$ is separable is hard in general (NP-hard in the dimension, Gurvits 2003), but there are simple necessary criteria.

\section{Partial transpose}
Let $\hat\rho$ be a density matrix on $\Ha_A\otimes\Ha_B$:
\begin{equation}
  \hat\rho=\sum_{ijkl}p^{ij}_{kl}\ketbra ij\otimes\ketbra kl.
\end{equation}
Its \emph{partial transpose} with respect to $B$ is
\begin{equation}
  \hat\rho^{T_B}=(I\otimes T)\hat\rho=\sum_{ijkl}p^{ij}_{kl}\ketbra ij\otimes\left(\ketbra kl\right)^T=\sum_{ijkl}p^{ij}_{kl}\ketbra ij\otimes\ketbra lk
  =\sum_{ijkl}p^{ij}_{lk}\ketbra ij\otimes\ketbra kl,
\end{equation}
where in the last step the summation indices $k$ and $l$ have been exchanged. In terms of matrix elements, only the indices of subsystem $B$ are exchanged:
\begin{equation}
  \left(\rho^{T_B}\right)_{m\mu,n\nu}=\rho_{m\nu,n\mu},
\end{equation}
where Latin indices refer to $A$ and Greek indices to $B$. With the big-endian ordering, $\hat\rho$ is a block matrix made of $\dim\Ha_A\times\dim\Ha_A$ blocks, each of size
$\dim\Ha_B\times\dim\Ha_B$:
\begin{equation}
  \hat\rho=\begin{pmatrix}A_{11}&A_{12}&\cdots\\A_{21}&A_{22}&\cdots\\ \vdots&\vdots&\ddots\end{pmatrix}\quad\Longrightarrow\quad
  \hat\rho^{T_B}=\begin{pmatrix}A_{11}^T&A_{12}^T&\cdots\\A_{21}^T&A_{22}^T&\cdots\\ \vdots&\vdots&\ddots\end{pmatrix}.
\end{equation}
Each block is transposed, but the order of the blocks is not changed. (The partial transpose depends on the basis, but its eigenvalues do not.)

\section{The Peres--Horodecki criterion}
\begin{theorem}[Peres \cite{Peres1996}; Horodecki \cite{Horodecki1996}]
If $\hat\rho$ is separable, then all the eigenvalues of $\hat\rho^{T_B}$ are non-negative (``PPT'': positive partial transpose). Equivalently, if $\hat\rho^{T_B}$ has a negative
eigenvalue, then $\hat\rho$ is entangled. For systems of dimension $2\times2$ and $2\times3$, the converse is also true: PPT implies separability.
\end{theorem}
\begin{proof}[Proof of the first part]
A separable state is a convex sum of product states, \cref{eq:separable}. Then
\begin{equation}
  \hat\rho^{T_B}=(I\otimes T)\hat\rho=\sum_ip_i\,\hat\rho_i^A\otimes\left(\hat\rho_i^B\right)^T.
\end{equation}
Transposition preserves the eigenvalues, so $(\hat\rho_i^B)^T$ has the same non-negative eigenvalues as $\hat\rho_i^B$, and is itself a valid density matrix (it is
also the complex conjugate of $\hat\rho_i^B$, since $\hat\rho_i^B$ is Hermitian). $\hat\rho^{T_B}$ is therefore a convex sum of positive semi-definite operators, and is positive
semi-definite. A negative eigenvalue violates this requirement.
\end{proof}
\begin{remark}[Why the transpose?]
The transpose $T$ is a \emph{positive} map (it maps density matrices to density matrices) which is not \emph{completely positive}: $I\otimes T$ applied to an entangled
state can give a non-positive operator. Physically, $\hat\rho^T=\hat\rho^*$ is the time-reversed state; reversing time on one half of an entangled system only is not a
physical operation. Every positive but not completely positive map gives an entanglement criterion, and the set of all of them characterizes separability (Horodecki 1996).
In higher dimensions there exist entangled PPT states (``bound entanglement''), which cannot be distilled into Bell pairs.
\end{remark}

\section{Examples}
\begin{example}[An entangled state]
For $\ket\psi=(\ket{01}+\ket{10})/\sqrt2$, in the basis $\{\ket{00},\ket{01},\ket{10},\ket{11}\}$,
\begin{equation}
  \hat\rho=\frac12\begin{pmatrix}0&0&0&0\\0&1&1&0\\0&1&1&0\\0&0&0&0\end{pmatrix},
\end{equation}
where, for example, $\rho_{10,01}=\braket{10}\psi\braket\psi{01}=\frac12$. The partial transpose moves the off-diagonal elements: $\rho_{01,10}\to(\rho^{T_B})_{00,11}$ and
$\rho_{10,01}\to(\rho^{T_B})_{11,00}$ (the upper-right block $\begin{pmatrix}0&0\\1&0\end{pmatrix}$ becomes $\begin{pmatrix}0&1\\0&0\end{pmatrix}$), so that
\begin{equation}
  \hat\rho^{T_B}=\frac12\begin{pmatrix}0&0&0&1\\0&1&0&0\\0&0&1&0\\1&0&0&0\end{pmatrix}.
\end{equation}
Its eigenvalues follow from (omitting the factor $\frac12$)
\begin{equation}
  \det\begin{pmatrix}-\lambda&0&0&1\\0&1-\lambda&0&0\\0&0&1-\lambda&0\\1&0&0&-\lambda\end{pmatrix}=\lambda^2(1-\lambda)^2-(1-\lambda)^2=(1-\lambda)^2(\lambda-1)(\lambda+1)=0.
\end{equation}
The eigenvalues of $\hat\rho^{T_B}$ are therefore $\frac12,\frac12,\frac12$ and $\lambda_4=-\frac12<0$: the state is entangled. The eigenvector of the negative eigenvalue is
$(\ket{00}-\ket{11})/\sqrt2$.
\end{example}

\begin{example}[A separable state]
For $\ket{\psi_2}=(\ket{10}+\ket{00})/\sqrt2=\frac1{\sqrt2}(\ket0+\ket1)\otimes\ket0$,
\begin{equation}
  \hat\rho=\frac12\begin{pmatrix}1&0&1&0\\0&0&0&0\\1&0&1&0\\0&0&0&0\end{pmatrix}\quad\Longrightarrow\quad\hat\rho^{T_B}=\hat\rho,
\end{equation}
because every block is symmetric. The partial transpose is the same matrix, with eigenvalues $1,0,0,0\ge0$, as expected for a separable state.
\end{example}

\begin{example}[The Werner state]\label{ex:werner}
Consider the mixture of a singlet with white noise, $\hat\rho_W=p\proj{\psi^-}+(1-p)\frac I4$, $0\le p\le1$. This models a Bell pair sent through a depolarizing channel.
Since $I^{T_B}=I$ and $\proj{\psi^-}^{T_B}$ has eigenvalues $\frac12,\frac12,\frac12,-\frac12$ (as in the first example, up to a local unitary), $\hat\rho_W^{T_B}$ has the eigenvalues
\begin{equation}
  \frac p2+\frac{1-p}4\ (\times3),\qquad-\frac p2+\frac{1-p}4=\frac{1-3p}4.
\end{equation}
The Werner state is entangled if and only if $p>1/3$. (It violates the CHSH inequality only for $p>1/\sqrt2$: entanglement and nonlocality are not the same thing.)
\end{example}

\section{Quantifying entanglement}
\begin{itemize}
  \item The \emph{negativity} $\mathcal N=\sum_{\lambda_i<0}|\lambda_i(\hat\rho^{T_B})|$ ($=\frac12$ for a Bell state, $\frac{3p-1}4$ for the Werner state with $p>1/3$) and the
  logarithmic negativity $E_N=\log_2(1+2\mathcal N)$ are easily computable entanglement measures.
  \item The \emph{concurrence} of Wootters (1998), $C(\hat\rho)=\max(0,\lambda_1-\lambda_2-\lambda_3-\lambda_4)$, where $\lambda_i$ are the square roots of the eigenvalues of
  $\hat\rho(Y\otimes Y)\hat\rho^*(Y\otimes Y)$ in decreasing order, gives the exact entanglement of formation of any two-qubit state. For pure states it reduces to
  $2|\alpha\delta-\beta\gamma|$ (\cref{prop:det}).
  \item An \emph{entanglement witness} is an observable $W$ with $\tr(W\hat\rho)\ge0$ for all separable states. Measuring $\tr(W\hat\rho)<0$ proves entanglement with a single
  observable, without full tomography. For a target state $\ket\psi$, $W=\alpha I-\proj\psi$ with $\alpha$ the maximal squared overlap of $\ket\psi$ with product states ($\frac12$
  for Bell states) is a witness: a fidelity with a Bell state larger than $1/2$ proves entanglement.
\end{itemize}

\section{Exercises}
\begin{exercise}
Show that the partial transpose with respect to $A$ has the same spectrum as the partial transpose with respect to $B$. (Hint: $\hat\rho^{T_A}=(\hat\rho^{T_B})^T$.)
\end{exercise}
\begin{exercise}
For which values of $p$ is $\hat\rho=p\proj{\Phi^+}+(1-p)\proj{01}$ entangled?
(Answer: for all $p>0$. Compute the eigenvalue of the block $\{\ket{01},\ket{10}\}$ of $\hat\rho^{T_B}$.)
\end{exercise}

\chapter{Quantum teleportation and dense coding}\label{ch:teleportation}

Entanglement is a resource that can be consumed to perform communication tasks impossible classically. The two fundamental protocols are teleportation
(1 Bell pair + 2 bits $\Rightarrow$ 1 qubit) and superdense coding (1 Bell pair + 1 qubit $\Rightarrow$ 2 bits).

\section{Quantum teleportation}
Quantum teleportation \cite{Bennett1993} transfers an unknown quantum state from one place to another. It uses a previously shared entangled pair and two bits of classical
communication; the quantum system itself is not transmitted. Suppose Alice wants to send Bob the state of qubit $C$:
\begin{equation}
  \ket\psi_C=\alpha\ket0_C+\beta\ket1_C.
\end{equation}
She does not know $\alpha$ and $\beta$. Even if she did, a single copy of the state would not allow her to measure them, and the no-cloning theorem forbids making copies.
Sending classical information about $\alpha,\beta$ would in any case require infinitely many bits.

\paragraph{Bell basis.} The protocol relies on the Bell states \cref{eq:bellstates},
\begin{equation}
  \ket{\Phi^\pm}_{AB}=\frac1{\sqrt2}\left(\ket0_A\ket0_B\pm\ket1_A\ket1_B\right),\qquad\ket{\Psi^\pm}_{AB}=\frac1{\sqrt2}\left(\ket0_A\ket1_B\pm\ket1_A\ket0_B\right).
\end{equation}
Conversely, the product basis states can be expressed in the Bell basis (add and subtract the definitions):
\begin{align}
  \ket0\ket0&=\tfrac1{\sqrt2}\left(\ket{\Phi^+}+\ket{\Phi^-}\right),&\ket1\ket1&=\tfrac1{\sqrt2}\left(\ket{\Phi^+}-\ket{\Phi^-}\right),\nonumber\\
  \ket0\ket1&=\tfrac1{\sqrt2}\left(\ket{\Psi^+}+\ket{\Psi^-}\right),&\ket1\ket0&=\tfrac1{\sqrt2}\left(\ket{\Psi^+}-\ket{\Psi^-}\right).\label{eq:prodtobell}
\end{align}

\paragraph{The protocol.} Alice and Bob share the pair $\ket{\Phi^+}_{AB}$. Alice holds qubits $A$ and $C$, Bob holds qubit $B$. The state of the three qubits is
\begin{align}
  \ket{\Phi^+}_{AB}\otimes\ket\psi_C&=\frac1{\sqrt2}\left(\ket0_A\ket0_B+\ket1_A\ket1_B\right)\otimes\left(\alpha\ket0_C+\beta\ket1_C\right)\nonumber\\
  &=\frac1{\sqrt2}\left(\alpha\ket{00}_{AC}\ket0_B+\beta\ket{01}_{AC}\ket0_B+\alpha\ket{10}_{AC}\ket1_B+\beta\ket{11}_{AC}\ket1_B\right).
\end{align}
In the second line we regrouped the qubits $A$ and $C$, which are both in Alice's hands. We now use \cref{eq:prodtobell} to rewrite the pair $AC$ in the Bell basis. Collecting the
terms multiplying each Bell state (for example, $\ket{\Phi^+}_{AC}$ appears in $\ket{00}_{AC}$ and $\ket{11}_{AC}$ with coefficient $\frac1{\sqrt2}$ each),
\begin{align}
  \ket{\Phi^+}_{AB}\otimes\ket\psi_C=\frac12\Big[&\ket{\Phi^+}_{AC}\otimes\left(\alpha\ket0_B+\beta\ket1_B\right)
  +\ket{\Phi^-}_{AC}\otimes\left(\alpha\ket0_B-\beta\ket1_B\right)\nonumber\\
  +&\ket{\Psi^+}_{AC}\otimes\left(\beta\ket0_B+\alpha\ket1_B\right)
  +\ket{\Psi^-}_{AC}\otimes\left(\beta\ket0_B-\alpha\ket1_B\right)\Big].\label{eq:teleport}
\end{align}
No physical operation has been performed yet: this is the same state, written in a different basis. Note, however, that the coefficients $\alpha$ and $\beta$ now appear in Bob's qubit.

Alice then performs a measurement of her two qubits $A$ and $C$ in the Bell basis. The state collapses onto one of the four terms of \cref{eq:teleport}, each with probability $1/4$
(independently of $\alpha,\beta$: Alice learns nothing about the state). For example, if she obtains $\ket{\Phi^+}_{AC}$, the result of the collapse is
$\ket{\Phi^+}_{AC}\otimes(\alpha\ket0_B+\beta\ket1_B)$, and Bob's qubit is already in the state $\ket\psi$. In the other three cases, Bob's qubit is in a state related to $\ket\psi$ by a
known unitary operation. Alice sends Bob the result of her measurement (two classical bits), and Bob applies the corresponding correction:
\begin{center}
\begin{tabular}{@{}lll@{}}\toprule
Alice's result & Bob's qubit & Bob's correction\\ \midrule
$\ket{\Phi^+}$ & $\alpha\ket0+\beta\ket1$ & $I$\\
$\ket{\Phi^-}$ & $\alpha\ket0-\beta\ket1$ & $Z$\\
$\ket{\Psi^+}$ & $\beta\ket0+\alpha\ket1$ & $X$\\
$\ket{\Psi^-}$ & $\beta\ket0-\alpha\ket1$ & $ZX$\\ \bottomrule
\end{tabular}
\end{center}
In the last case, $ZX(\beta\ket0-\alpha\ket1)=Z(\beta\ket1-\alpha\ket0)=-(\alpha\ket0+\beta\ket1)$, which differs from $\ket\psi$ only by an irrelevant global phase. In all four cases Bob ends up
with the state $\ket\psi$.

\paragraph{Circuit.} On a quantum computer, the Bell measurement is done by the inverse of the Bell-state preparation circuit (CNOT then $H$) followed by a computational-basis measurement,
and the corrections can be applied as classically controlled gates (qubits $A$ and $B$ start in $\ket{\Phi^+}$):
\begin{center}
\begin{quantikz}
\lstick{$C:\ \ket\psi$} & \ctrl{1} & \gate{H} & \meter{} & \cw & \cwbend{2} \\
\lstick{$A$} & \targ{} & \qw & \meter{} & \cwbend{1} & \\
\lstick{$B$} & \qw & \qw & \qw & \gate{X} & \gate{Z} & \qw\ \rstick{$\ket\psi$}
\end{quantikz}
\end{center}
The measured bits $(m_C,m_A)$ select the correction $Z^{m_C}X^{m_A}$. (Check: $\ket{\Psi^-}_{AC}$ is mapped by CNOT$_{C\to A}$ and $H_C$ to $m_C=m_A=1$.) The classical
communication can even be replaced by quantum gates (``deferred measurement''): a CNOT from $A$ to $B$ and a CZ from $C$ to $B$.

\paragraph{Remarks.}
\begin{itemize}
  \item \emph{No faster-than-light communication.} Until Bob receives Alice's two bits, his qubit is in an equal mixture of the four states of the table:
  $\frac14\sum_k\sigma_k\proj\psi\sigma_k=\frac I2$, the maximally mixed state, independently of $\ket\psi$ (this is the identity used for the depolarizing channel).
  \item \emph{No cloning.} After Alice's measurement, qubit $C$ is left in a Bell state with $A$ and no longer carries any information about $\alpha$ and $\beta$. The original state is destroyed.
  \item \emph{Experiments.} Teleportation of photon polarization was first demonstrated in 1997 \cite{Bouwmeester1997}, later between atoms, ions, superconducting qubits, over 143\,km
  between Canary islands, and from ground to the Micius satellite (1400\,km, 2017).
  \item \emph{Teleportation is a computing primitive.} Gate teleportation (Gottesman--Chuang 1999) applies a gate by teleporting through a suitably prepared entangled state; it is the
  basis of measurement-based quantum computing and of the fault-tolerant implementation of the $T$ gate (magic states).
\end{itemize}

\section{Entanglement swapping}
If Alice shares $\ket{\Phi^+}_{A_1B_1}$ with Charlie and Charlie shares $\ket{\Phi^+}_{C_2B_2}$ with Bob, a Bell measurement by Charlie on $B_1,C_2$ ``teleports'' the entanglement:
$A_1$ and $B_2$, which never interacted, end up in a Bell state (up to a known Pauli correction). The proof is \cref{eq:teleport} with $\ket\psi_C$ replaced by half of a Bell pair.
Entanglement swapping is the elementary step of \emph{quantum repeaters}, which distribute entanglement over long distances by swapping between short links, combined with
entanglement purification and quantum memories.

\section{Superdense coding}
Conversely, a shared Bell pair allows Alice to send \emph{two} classical bits by transmitting a \emph{single} qubit (Bennett and Wiesner, 1992). Alice and Bob share $\ket{\Phi^+}_{AB}$.
To send the bits $(b_1,b_2)$, Alice applies $Z^{b_1}X^{b_2}$ to her qubit:
\begin{equation}
  (I)\ket{\Phi^+}=\ket{\Phi^+},\quad(X)\ket{\Phi^+}=\ket{\Psi^+},\quad(Z)\ket{\Phi^+}=\ket{\Phi^-},\quad(ZX)\ket{\Phi^+}=-\ket{\Psi^-},
\end{equation}
and sends her qubit to Bob, who performs a Bell measurement on both qubits and recovers the two bits, since the four Bell states are orthogonal. This does not contradict the
Holevo bound (one qubit cannot carry more than one bit of accessible information): the second qubit, sent earlier, is also needed.

\section{Exercises}
\begin{exercise}
Derive \cref{eq:teleport} explicitly for the $\ket{\Psi^\pm}$ terms.
\end{exercise}
\begin{exercise}
Show that if the shared pair is the Werner state $\hat\rho_W$ instead of $\ket{\Phi^+}$ (with the singlet replaced by $\ket{\Phi^+}$), Bob receives $p\proj\psi+(1-p)I/2$, and that the
average fidelity $\bar F=\frac{1+p}2$ exceeds the best classical value $2/3$ (measure and send) only if $p>1/3$, which coincides with the entanglement threshold.
\end{exercise}


\part{Gates and algorithms}
\chapter{Multi-qubit gates and universality}\label{ch:gates}

\section{Reversible computation}\label{sec:reversible}
A quantum computation (before measurement) is a unitary transformation, and unitary transformations are invertible. Classical logic, however, is mostly irreversible:
the AND gate maps four inputs to two outputs and loses information. Landauer (1961) showed that erasing one bit dissipates at least $k_BT\ln2$ of heat; Bennett (1973) showed that
any classical computation can be made reversible, at the price of keeping extra (``garbage'') bits. The standard construction turns any function $f:\{0,1\}^n\to\{0,1\}^m$ into the
permutation of basis states
\begin{equation}\label{eq:Uf}
  U_f\ket x\ket y=\ket x\ket{y\oplus f(x)},
\end{equation}
where $\oplus$ is the bitwise addition modulo 2. It is a permutation (hence unitary), and its own inverse: $U_f^2\ket x\ket y=\ket x\ket{y\oplus f(x)\oplus f(x)}=\ket x\ket y$.
Every classical circuit for $f$ with $G$ gates can be turned into a reversible circuit for $U_f$ with $\order(G)$ Toffoli gates and ancillas, which are reset to $\ket0$ by
\emph{uncomputing} the intermediate results (compute, copy the result with CNOTs, run the circuit backwards). Uncomputation is essential in quantum computing: garbage left
entangled with the result would destroy the interference that algorithms rely on.

\section{The CNOT gate}\label{sec:cnot}
The controlled-NOT (CNOT) gate acts on two qubits, a \emph{control} and a \emph{target}. If the control qubit is in state $\ket0$, the target is left unchanged. If the control is
in state $\ket1$, the target is inverted ($\ket0\leftrightarrow\ket1$). The control qubit itself is never modified in the computational basis. Truth table:
\begin{center}
\begin{tabular}{@{}cc|cc@{}}\toprule
\multicolumn{2}{c|}{Before}&\multicolumn{2}{c}{After}\\
Control & Target & Control & Target\\ \midrule
0&0&0&0\\0&1&0&1\\1&0&1&1\\1&1&1&0\\ \bottomrule
\end{tabular}
\qquad
\begin{quantikz}
\lstick{$\ket x$} & \ctrl{1} & \qw\ \rstick{$\ket x$}\\
\lstick{$\ket y$} & \targ{} & \qw\ \rstick{$\ket{y\oplus x}$}
\end{quantikz}
\end{center}
The target output is the exclusive OR of the two inputs:
\begin{equation}
  \CNOT\ket x\ket y=\ket x\ket{y\oplus x}.
\end{equation}
Unlike the classical XOR gate, which has two inputs and one output, CNOT keeps the control bit in its output. It is therefore reversible, as every unitary quantum gate must be.

\paragraph{Matrix.} In the basis $\ket0=(1,0)^T$, $\ket1=(0,1)^T$ a qubit is $\ket\psi=\alpha\ket0+\beta\ket1=(\alpha,\beta)^T$, and the inverted qubit is $X\ket\psi=(\beta,\alpha)^T$.
In the basis $\{\ket{00},\ket{01},\ket{10},\ket{11}\}$, with the first qubit as control,
\begin{equation}\label{eq:cnotmatrix}
  \CNOT=\begin{pmatrix}1&0&0&0\\0&1&0&0\\0&0&0&1\\0&0&1&0\end{pmatrix}=\proj0\otimes I+\proj1\otimes X.
\end{equation}
The upper-left block is the identity acting on the target when the control is $\ket0$; the lower-right block is the NOT gate, acting when the control is $\ket1$.
Multiplying by the column vectors $\ket{00}=(1,0,0,0)^T,\dots,\ket{11}=(0,0,0,1)^T$ leaves the first two unchanged and exchanges the last two:
$\CNOT\ket{10}=\ket{11}$, $\CNOT\ket{11}=\ket{10}$, in agreement with the truth table.

\paragraph{What CNOT does to superpositions.}
\begin{itemize}
  \item It \emph{entangles}: $\CNOT\,\ket+\ket0=\ket{\Phi^+}$. A product state becomes maximally entangled.
  \item It ``copies'' basis states, $\ket x\ket0\to\ket x\ket x$, but not superpositions (no-cloning): $(\alpha\ket0+\beta\ket1)\ket0\to\alpha\ket{00}+\beta\ket{11}$.
  \item \emph{The control is not always unchanged.} In the $X$ basis the roles are reversed: $(H\otimes H)\CNOT_{1\to2}(H\otimes H)=\CNOT_{2\to1}$. For example
  $\CNOT\ket+\ket-=\ket-\ket-$: the phase is ``kicked back'' onto the control. Phase kickback is the mechanism of the Deutsch--Jozsa, Grover and phase-estimation algorithms.
\end{itemize}

\section{Other two-qubit gates}
\paragraph{Controlled-$Z$.} $\mathrm{CZ}=\diag(1,1,1,-1)=\proj0\otimes I+\proj1\otimes Z$. It is symmetric in the two qubits (it only puts a sign on $\ket{11}$). Since $HXH=Z$,
\begin{equation}
  (I\otimes H)\,\CNOT\,(I\otimes H)=\mathrm{CZ},
\end{equation}
which can be checked by multiplying the matrices, for instance with SymPy (\cref{ch:qiskit}).

\paragraph{SWAP.} $\mathrm{SWAP}\ket{x}\ket{y}=\ket y\ket x$. It is made of three CNOTs with alternating control:
\begin{equation}
  \mathrm{SWAP}=\CNOT_{1\to2}\,\CNOT_{2\to1}\,\CNOT_{1\to2}.
\end{equation}
Proof on basis states: $\ket{x,y}\to\ket{x,x\oplus y}\to\ket{x\oplus(x\oplus y),x\oplus y}=\ket{y,x\oplus y}\to\ket{y,x\oplus y\oplus y}=\ket{y,x}$. This is the reversible version of the
classic XOR-swap trick. On hardware with limited connectivity, SWAPs are used to bring distant qubits next to each other, which is one of the main overheads of transpilation.

\paragraph{Controlled-$U$.} For any single-qubit $U$, $C(U)=\proj0\otimes I+\proj1\otimes U$.
\begin{proposition}[\cite{NielsenChuang}, Cor.~4.2]
Write $U=\e^{\ii\alpha}AXBXC$ with $ABC=I$ (always possible by the Euler decomposition: $A=R_z(\beta)R_y(\theta/2)$, $B=R_y(-\theta/2)R_z(-(\delta+\beta)/2)$,
$C=R_z((\delta-\beta)/2)$). Then $C(U)$ is implemented by
\begin{center}
\begin{quantikz}
 & \qw & \ctrl{1} & \qw & \ctrl{1} & \qw & \gate{\diag(1,\e^{\ii\alpha})} & \qw\\
 & \gate{C} & \targ{} & \gate{B} & \targ{} & \gate{A} & \qw & \qw
\end{quantikz}
\end{center}
\end{proposition}
\begin{proof}
If the control is $\ket0$, the target undergoes $ABC=I$. If it is $\ket1$, it undergoes $AXBXC$ and the phase gate gives the factor $\e^{\ii\alpha}$, i.e.\ $U$.
The identity $XR_y(\theta)X=R_y(-\theta)$, $XR_z(\theta)X=R_z(-\theta)$ is used to check $AXBXC=R_z(\beta)R_y(\theta)R_z(\delta)$.
\end{proof}

\section{The Toffoli gate}\label{sec:toffoli}
The Toffoli gate, or controlled-controlled-NOT (CCNOT), flips the target if and only if both controls are $\ket1$: $\ket{x,y,z}\to\ket{x,y,z\oplus xy}$. With $z=0$ it computes the AND,
and with $x=y=1$ it is a NOT. Since $\{\text{AND},\text{NOT}\}$ is universal for Boolean logic, \emph{the Toffoli gate alone is universal for classical reversible computation} (with ancillas).
It is not natively available in hardware and is decomposed into two-qubit gates. With $V=\sqrt X$ (so that $V^2=X$),
\begin{center}
\begin{quantikz}
\lstick{$x$} & \ctrl{2} & \qw & \qw\\
\lstick{$y$} & \ctrl{1} & \qw & \qw\\
\lstick{$z$} & \targ{} & \qw & \qw
\end{quantikz}
\ $=$\
\begin{quantikz}
 & \qw & \ctrl{1} & \qw & \ctrl{1} & \ctrl{2} & \qw\\
 & \ctrl{1} & \targ{} & \ctrl{1} & \targ{} & \qw & \qw\\
 & \gate{V} & \qw & \gate{V^\dagger} & \qw & \gate{V} & \qw
\end{quantikz}
\end{center}
Check: the target receives $V^{y}$, then $V^{\dagger(x\oplus y)}$, then $V^{x}$; the total exponent is $y-(x\oplus y)+x=2xy$, so the target undergoes $V^{2xy}=X^{xy}$. The CNOTs on $y$ cancel.
Decomposing the controlled-$V$ gates further, the standard circuit uses 6 CNOTs and 7 $T$ or $T^\dagger$ gates; 6 CNOTs is optimal. On real hardware with nearest-neighbour connectivity,
extra SWAPs increase this count (\cref{sec:halfadder}).

\section{Universal gate sets}\label{sec:universal}
\begin{definition}
A set of gates is \emph{universal} if any unitary on any number of qubits can be implemented \emph{exactly} by a finite circuit of these gates, or \emph{approximately universal}
if any unitary can be approximated to any accuracy $\varepsilon$ in operator norm.
\end{definition}
\begin{theorem}
Single-qubit gates and CNOT are universal.
\end{theorem}
\begin{proof}[Sketch of proof (\cite{NielsenChuang}, \S4.5)]
\emph{Step 1: two-level unitaries.} Any unitary $U$ on a $d$-dimensional space is a product of at most $d(d-1)/2$ \emph{two-level} unitaries, i.e.\ unitaries that act non-trivially only
on two basis vectors. This is the complex version of the Givens-rotation (or Jacobi) method: one multiplies $U$ by two-level unitaries chosen to make the entries of the first column zero,
one by one, below the diagonal; unitarity then also makes the first row trivial, and one proceeds by induction on $d$.

\emph{Step 2: two-level unitaries from controlled gates.} A two-level unitary acting on $\ket s$ and $\ket t$ ($n$-bit strings) is implemented by a sequence of CNOT-type gates (a Gray
code) that maps $\ket s$ and $\ket t$ onto two strings differing in a single bit, then a single-qubit gate on that bit controlled by all the other bits, then the inverse Gray code.

\emph{Step 3: multi-controlled gates.} A $C^{n-1}(U)$ gate is decomposed into Toffoli gates (with ancillas) and a $C(U)$; Toffoli and $C(U)$ are made of CNOTs and single-qubit gates as above.
\end{proof}
The construction is not efficient: a generic $n$-qubit unitary needs $\order(4^n)$ gates. This is unavoidable. A unitary on $n$ qubits has $4^n$ real parameters, and a circuit of $G$ gates from
a finite set of $g$ gates can produce at most $(gn^2)^G$ distinct circuits; covering all unitaries to accuracy $\varepsilon$ requires $G=\Omega(4^n\log(1/\varepsilon)/\log n)$.
\emph{Most unitaries cannot be implemented efficiently}; the art of quantum algorithms is to find useful ones that can.

\paragraph{A finite universal set.} Hardware and error-correcting codes offer only a finite set of gates. The set $\{H,T,\CNOT\}$ (often written ``Clifford $+T$'') is approximately universal.
The key observation is that $HTHT$ is (up to phase) a rotation of the Bloch sphere by an angle $\theta$ with $\cos\frac\theta2=\cos^2\frac\pi8$, which is an irrational multiple of $\pi$.
Its powers are therefore dense in the rotations about its axis. Combined with a second rotation $THTH$ about a non-parallel axis, any rotation can be approximated (Boykin et al.\ 1999).
\begin{theorem}[Solovay--Kitaev]
If a finite set of single-qubit gates, closed under inverses, generates a dense subgroup of $SU(2)$, then any single-qubit gate can be approximated to accuracy $\varepsilon$ by a sequence
of $\order(\log^c(1/\varepsilon))$ gates of the set, with $c\approx2$--$4$ depending on the construction.
\end{theorem}
Approximation is thus cheap: the overhead is polylogarithmic. Without $T$, the gates $\{H,S,\CNOT\}$ generate only the finite \emph{Clifford group}, which maps Pauli operators to Pauli
operators and is classically simulable (Gottesman--Knill). The $T$ gate is the ``non-Clifford'' resource that makes quantum computers powerful, and in fault-tolerant architectures it is
the expensive gate (it requires magic-state distillation).

\section{Physical origin of two-qubit gates}
Two-qubit gates require an interaction between qubits. The three generic interactions found in nature generate the three standard entangling gates.
\begin{origin}[: from interactions to gates]
\textbf{Ising ($ZZ$) interaction} (nuclear spins in NMR with $J$-coupling, dispersively coupled superconducting qubits, Rydberg atoms): $H=\frac{\hbar J}{4}Z_1Z_2$. Since
$\proj{11}=\frac14(I-Z_1-Z_2+Z_1Z_2)$,
\begin{equation}
  \mathrm{CZ}=\e^{\ii\pi\proj{11}}=\e^{\ii\pi/4}\,\e^{-\ii\frac\pi4Z_1}\e^{-\ii\frac\pi4Z_2}\,\e^{\ii\frac\pi4Z_1Z_2}.
\end{equation}
All terms commute, so CZ is obtained by letting the qubits interact for the time $t$ with $Jt/4=-\pi/4$ (or $Jt=\pi$ with the opposite sign of $J$) together with single-qubit $Z$ rotations
by $\pi/2$. Adding Hadamards on the target gives a CNOT.

\textbf{Heisenberg exchange interaction} (electron spins in neighbouring quantum dots, \cref{sec:spinqubits}): $H=J\,\vec S_1\cdot\vec S_2=\frac{J\hbar^2}4\vec\sigma_1\cdot\vec\sigma_2$. The
identity $\vec\sigma_1\cdot\vec\sigma_2=2\,\mathrm{SWAP}-I$ (check it on the triplet, eigenvalue $+1$, and on the singlet, eigenvalue $-3$) gives
$U(t)=\e^{\ii\phi}\e^{-\ii\frac{J\hbar t}2\mathrm{SWAP}}=\e^{\ii\phi}\left(\cos\frac{J\hbar t}2-\ii\sin\frac{J\hbar t}2\,\mathrm{SWAP}\right)$, using $\mathrm{SWAP}^2=I$. For $J\hbar t=\pi$ this is the
SWAP gate (up to phase), and for half this time the entangling $\sqrt{\mathrm{SWAP}}$ gate. Two $\sqrt{\mathrm{SWAP}}$ and single-qubit rotations make a CNOT (Loss and DiVincenzo, 1998).

\textbf{$XY$ (flip-flop) interaction} (superconducting qubits coupled by a capacitor or a resonator, \cref{sec:transmon}): $H=\hbar g(\sigma_1^+\sigma_2^-+\sigma_1^-\sigma_2^+)=\frac{\hbar g}2(X_1X_2+Y_1Y_2)$.
It exchanges an excitation, $\ket{01}\leftrightarrow\ket{10}$, and acts as $\hbar g\,\sigma_x$ in this subspace. After $gt=\pi/2$ it gives the iSWAP gate (up to the sign of $\ii$),
$\ket{01}\to-\ii\ket{10}$, $\ket{10}\to-\ii\ket{01}$; after $gt=\pi/4$, $\sqrt{\mathrm{iSWAP}}$. Resonance is controlled by tuning the qubit frequencies.
\end{origin}
Any entangling two-qubit gate, together with arbitrary single-qubit gates, is universal (Brylinski 2002). Hardware therefore chooses whatever gate its interaction makes natural: CZ
(Google, Rydberg atoms), cross-resonance/ECR (IBM), Mølmer--Sørensen $XX$ (trapped ions, \cref{sec:ions}), and transpilers convert circuits into the native gate set.

\section{Exercises}
\begin{exercise}
Verify the operation of the CNOT gate: write the four basis vectors as column vectors using the tensor product, multiply them by \cref{eq:cnotmatrix}, and check the truth table.
\end{exercise}
\begin{exercise}
Show that $(H\otimes H)\CNOT_{1\to2}(H\otimes H)=\CNOT_{2\to1}$. Interpret it with the Bloch sphere.
\end{exercise}
\begin{exercise}
Show that $\CNOT$ cannot be written as $A\otimes B$. Show that it can create a Bell state from a product state, while $A\otimes B$ cannot change the Schmidt coefficients (\cref{app:schmidt}).
\end{exercise}
\begin{exercise}
Using the formula of the CZ derivation, find the evolution time and the single-qubit corrections needed to realize CZ with the NMR coupling $H=\frac{\hbar\pi J}{2}Z_1Z_2$, $J=200$\,Hz.
\end{exercise}

\chapter{Programming quantum computers with Qiskit}\label{ch:qiskit}

\section{General remarks}
It is instructive to compare present-day quantum computers with the first general-purpose electronic computer, ENIAC (1945).\footnote{Before electronic computers, calculations relied on
analog devices such as the slide rule, used from the 17th century up to the 1970s.} ENIAC worked in decimal arithmetic and needed 36 vacuum tubes to store a single decimal digit. Its memory
consisted of 20 accumulators, the equivalent of modern registers, each holding a signed 10-digit number. Some of them were combined for more complex operations: four accumulators were used
for multiplication and five for division. Since $10^{10}\approx2^{33}$, one accumulator corresponds roughly to a 32-bit binary register, and the whole machine stored only a few hundred bits.
Even so, the accumulators alone required of the order of $20\times10\times36\approx10^4$ vacuum tubes.

Current quantum processors have a comparable number of (noisy) qubits. In terms of programming, quantum computing is therefore still at the stage of ENIAC, or even before it: a computer is in
effect designed with a specific hardware wiring for each problem. Programming is done at the lowest level, gate by gate, even though some higher-level codes are already available.

Several software frameworks exist for programming quantum computers (Qiskit, Cirq, PennyLane, Q\#, CUDA-Q, \dots). We use Qiskit, an open-source Python library developed by IBM. It can run
circuits both on simulators and on real quantum hardware. The examples below have been tested with Qiskit 2.x and Qiskit Aer 0.18.

\section{Building a circuit}
A circuit is built from registers:
\begin{lstlisting}
from qiskit import QuantumRegister, ClassicalRegister, QuantumCircuit

qreg = QuantumRegister(2, 'q')    # one quantum register with 2 qubits q0, q1
creg = ClassicalRegister(2, 'c')  # one classical register with 2 bits
qc = QuantumCircuit(qreg, creg)   # a circuit acting on these registers

qc.h(qreg[0])                     # Hadamard gate on q0
qc.cx(qreg[0], qreg[1])           # controlled-X (CNOT): control q0, target q1
qc.draw('mpl')                    # draw the circuit with matplotlib
\end{lstlisting}
The method \texttt{cx} adds a controlled-$X$ gate, i.e.\ the CNOT gate, with \texttt{qreg[0]} as control and \texttt{qreg[1]} as target. This circuit prepares the Bell state
$\ket{\Phi^+}$ (\cref{ch:entanglement}). The option \texttt{'mpl'} of the \texttt{draw} method produces a figure using the Python plotting library matplotlib.

\paragraph{Exercise: matrices by hand.} Before using a simulator, it is useful to compute the matrices of small circuits by hand, for example with the symbolic library SymPy. Given the CNOT
matrix \texttt{Cx\_sp}, construct the matrices \texttt{H\_sp}, \texttt{I\_sp}, \texttt{IH\_sp} $=I\otimes H$, and $U=(I\otimes H)\,\CNOT\,(I\otimes H)$:
\begin{lstlisting}
import sympy as sp
from sympy.physics.quantum import TensorProduct

H_sp = sp.Matrix([[1, 1], [1, -1]]) / sp.sqrt(2)
I_sp = sp.eye(2)
Cx_sp = sp.Matrix([[1, 0, 0, 0], [0, 1, 0, 0],
                   [0, 0, 0, 1], [0, 0, 1, 0]])
IH_sp = TensorProduct(I_sp, H_sp)
U_sp = sp.simplify(IH_sp * Cx_sp * IH_sp)
\end{lstlisting}
The result is $U=\diag(1,1,1,-1)$, the controlled-$Z$ gate. Since $HXH=Z$, the Hadamard gates on the target convert the controlled-$X$ into a controlled-$Z$. By hand: $I\otimes H$ is
block diagonal with two blocks $H$, and $\CNOT=\diag(I,X)$, so $U=\diag(HIH,HXH)=\diag(I,Z)$.

\section{Simulators}
Aer is the Qiskit package for high-performance simulation of quantum circuits, including noisy ones. It provides several simulator backends:
\begin{lstlisting}
from qiskit_aer import Aer
simulator = Aer.get_backend('qasm_simulator')
\end{lstlisting}
\begin{itemize}
  \item \texttt{qasm\_simulator}: simulates measurements shot by shot, as a real device would. It is ideal by default, but noise models can be added.
  \item \texttt{statevector\_simulator}: returns the ideal final state vector.
  \item \texttt{unitary\_simulator}: returns the ideal unitary matrix of the whole circuit.
\end{itemize}
In recent versions all three are special cases of the general \texttt{AerSimulator}, whose simulation method can be \texttt{'statevector'} (on CPU or GPU), \texttt{'density\_matrix'} (also on
GPU), \texttt{'unitary'}, \texttt{'matrix\_product\_state'}, \texttt{'stabilizer'}, and others.

\begin{keypoint}[: the cost of classical simulation]
A state vector of $n$ qubits has $2^n$ complex amplitudes, i.e.\ $16\times2^n$ bytes in double precision: 16\,GB for $n=30$, 16\,PB for $n=50$. A density matrix needs the square of
that. Brute-force simulation is therefore limited to about 45--50 qubits on the largest supercomputers. Matrix-product-state and stabilizer simulators can handle many more qubits for weakly
entangled or Clifford circuits, in agreement with the remark at the end of \cref{sec:parallelism}.
\end{keypoint}

\section{Transpiling and running a circuit}
Before running, the circuit is \emph{transpiled}. Transpilation translates it into the gates actually available on the backend, maps the logical qubits onto physical qubits respecting the
connectivity (adding SWAPs if needed), and optimizes it. Unlike compilation, which translates a program to a lower level of abstraction, this is a conversion at the same level, from a circuit
to an equivalent circuit:
\begin{lstlisting}
from qiskit import transpile
from qiskit.visualization import plot_histogram

qc.barrier()            # visual separator; also blocks optimizations across it
qc.measure(qreg, creg)  # measure the qubits into the classical bits

qc_t = transpile(qc, simulator,
                 basis_gates=['id', 'rz', 'sx', 'x', 'cx'])
job = simulator.run(qc_t, shots=10000)
counts = job.result().get_counts(qc_t)
plot_histogram(counts)
\end{lstlisting}
A barrier is a visual marker in the circuit drawing. It also prevents the transpiler from moving or merging gates across it. The result \texttt{counts} is a Python dictionary that gives the number
of times each outcome was obtained. For the Bell state prepared above, it is approximately \texttt{\{'00': 5000, '11': 5000\}}. The statistical fluctuation of each count is $\sqrt{Np(1-p)}=50$.
The basis gates \texttt{rz}, \texttt{sx}$=\sqrt X$, \texttt{x} suffice for any single-qubit gate (Euler decomposition, \cref{thm:euler}), e.g.\ $H=\e^{\ii\pi/4}R_z(\pi/2)\sqrt XR_z(\pi/2)$.

\section{State vector and unitary matrix}
With the state-vector simulator, applied to a circuit without the final measurements, the state of the circuit can be retrieved and displayed as a \LaTeX{} formula:
\begin{lstlisting}
from qiskit.visualization import array_to_latex

simulator = Aer.get_backend('statevector_simulator')
psi = simulator.run(qc_t).result().get_statevector(qc_t)
array_to_latex(psi)
\end{lstlisting}
Similarly, the unitary simulator gives the matrix of the whole circuit, which can be checked numerically, for example its determinant:
\begin{lstlisting}
import numpy as np
from numpy import linalg as LA

simulator = Aer.get_backend('unitary_simulator')
U = simulator.run(qc_t).result().get_unitary(qc_t, decimals=3)
U = np.asarray(U)      # convert to a NumPy array
print(LA.det(U))       # |det U| = 1, since U is unitary
\end{lstlisting}

\section{Qubit ordering: big-endians and little-endians}\label{sec:endian}
Qiskit orders the qubits from right to left: in a basis state $\ket{q_{n-1}\dots q_1q_0}$, qubit $q_0$ is the least significant bit. This is the opposite of the convention used in these notes.
As a result, the matrices returned by Qiskit may appear with the roles of the qubits exchanged.

\paragraph{A little history.} In Jonathan Swift's \emph{Gulliver's Travels} (1726), the empire of Lilliput is torn by a bitter conflict about how a boiled egg should be opened. Traditionally, eggs
were broken at the larger end. However, the Emperor's grandfather cut his finger while doing so as a child, and his father then ordered all his subjects, under heavy penalties, to break their eggs
at the smaller end. The resistance of the Big-Endians led to several rebellions. The exiled rebels found refuge in the neighbouring empire of Blefuscu, and a long war followed. Swift's satire
targets the religious quarrels of his time. The point is that the dispute is passionate, but the choice itself is completely arbitrary.

In 1980, the computer scientist Danny Cohen borrowed Swift's terms in a famous note, \emph{On Holy Wars and a Plea for Peace} \cite{Cohen1981}. Cohen was writing about the order in which computers
store and transmit the bytes of a number. A 32-bit integer such as \texttt{0x12345678} occupies four bytes, and there are two natural ways to arrange them at increasing addresses:
\begin{center}
\begin{tabular}{@{}lcccc@{}}\toprule
Address & $n$ & $n+1$ & $n+2$ & $n+3$\\ \midrule
Big-endian & 12 & 34 & 56 & 78\\ Little-endian & 78 & 56 & 34 & 12\\ \bottomrule
\end{tabular}
\end{center}
A big-endian machine stores the most significant byte first, the way we write numbers on paper. A little-endian machine stores the least significant byte first. Little-endian ordering simplifies some
arithmetic operations, since addition proceeds from the least significant digit; big-endian ordering is easier for humans to read in a memory dump. As in Lilliput, neither is fundamentally better,
and both have been used: the Intel x86 family and, in practice, the ARM processors of today's smartphones and Apple silicon Macs are little-endian; the Motorola 68000, the IBM mainframes and the
PowerPC of classic Macintosh computers are big-endian (when Apple moved to Intel processors in 2006, it had to deal with the byte order of existing files and programs); ARM, PowerPC and MIPS are
bi-endian. The Internet protocols transmit numbers in big-endian order (the \emph{network byte order}). Many file formats therefore announce their byte order: a TIFF image starts with
\texttt{II} (Intel) or \texttt{MM} (Motorola), and a UTF-16 text file may begin with a byte-order mark.

\paragraph{Endianness of qubits.} Exactly the same ambiguity arises for qubits. A basis state of $n$ qubits is a string of bits, which can also be read as an integer $k$, the index of the corresponding
component of the state vector. The question is which qubit is the most significant bit of $k$.
\begin{itemize}
  \item \emph{Big-endian} (most textbooks, including Nielsen and Chuang \cite{NielsenChuang}, and these notes):
  \begin{equation}
    \ket{q_0q_1\dots q_{n-1}}=\ket{q_0}\otimes\ket{q_1}\otimes\dots\otimes\ket{q_{n-1}},\qquad k=\sum_{i=0}^{n-1}q_i2^{n-1-i}.
  \end{equation}
  \item \emph{Little-endian} (Qiskit):
  \begin{equation}
    \ket{q_{n-1}\dots q_1q_0}=\ket{q_{n-1}}\otimes\dots\otimes\ket{q_1}\otimes\ket{q_0},\qquad k=\sum_{i=0}^{n-1}q_i2^i .
  \end{equation}
  Qubit $q_i$ then has the weight $2^i$, as bit $i$ of an ordinary binary integer. This makes arithmetic circuits natural (and is consistent with the operation \texttt{1 << i} used in the
  Deutsch--Jozsa example), but it reverses the order of the tensor product.
\end{itemize}
\emph{Example 1: a single flipped qubit.} Apply an $X$ gate to qubit $q_0$ of a three-qubit register initially in $\ket{000}$:
\begin{lstlisting}
from qiskit import QuantumCircuit
from qiskit.quantum_info import Statevector

qc = QuantumCircuit(3)
qc.x(0)
print(Statevector(qc).probabilities_dict())   # {'001': 1.0}
\end{lstlisting}
Qiskit reports the outcome \texttt{'001'}, with $q_0$ as the rightmost bit, and the non-zero component of the state vector has index $k=1$. In the big-endian convention, the same state is written
$\ket{100}$ and has index $k=4$.

\emph{Example 2: the CNOT matrix.} For a CNOT with control $q_0$ and target $q_1$:
\begin{equation}
  \underbrace{\begin{pmatrix}1&0&0&0\\0&1&0&0\\0&0&0&1\\0&0&1&0\end{pmatrix}}_{\text{big-endian (these notes)}}\qquad
  \underbrace{\begin{pmatrix}1&0&0&0\\0&0&0&1\\0&0&1&0\\0&1&0&0\end{pmatrix}}_{\text{little-endian (Qiskit)}}
\end{equation}
In the Qiskit matrix, the control $q_0$ is the least significant bit, so the gate exchanges the basis states with indices $1=\ket{01}$ and $3=\ket{11}$ instead of 2 and 3. Both matrices describe
the same physical gate; only the labelling of the basis has changed. Formally, the two matrices are related by the SWAP permutation: $M_{\rm Qiskit}=\mathrm{SWAP}\,M\,\mathrm{SWAP}$.

\emph{Converting.} To compare Qiskit results with calculations done by hand, reverse the order of the qubits, with \texttt{QuantumCircuit.reverse\_bits()} or the method \texttt{reverse\_qargs()} of
the \texttt{Statevector} and \texttt{Operator} classes. For example, \texttt{Operator(qc.reverse\_bits())} gives the CNOT matrix in the big-endian form above. As in Lilliput, there is no right end
at which to break the egg. What matters is to know which convention a given book, article or program uses, and to convert carefully.

\section{Noisy simulation and fake backends}
To test a program under realistic conditions without access to a real quantum computer, Qiskit provides \emph{fake backends}: simulators that reproduce the qubit connectivity, the basis gates and the
noise of specific IBM devices, for example \texttt{FakeManilaV2} or \texttt{FakeMarrakesh} from \texttt{qiskit\_ibm\_runtime.fake\_provider}. The noise model contains, for each qubit and gate,
depolarizing errors, thermal relaxation with the measured $T_1$, $T_2$ (\cref{ch:open}), and readout errors.

\begin{keypoint}[: a back-of-the-envelope error estimate]
If a transpiled circuit has $N_1$ single-qubit gates with error $\varepsilon_1$, $N_2$ two-qubit gates with error $\varepsilon_2$ and $n$ measurements with error $\varepsilon_m$, the probability
that no error occurs is roughly
\begin{equation}
  P_{\rm success}\approx(1-\varepsilon_1)^{N_1}(1-\varepsilon_2)^{N_2}(1-\varepsilon_m)^{n}\approx\e^{-N_1\varepsilon_1-N_2\varepsilon_2-n\varepsilon_m}.
\end{equation}
With $\varepsilon_2\sim5\times10^{-3}$, a circuit becomes useless after a few hundred two-qubit gates. This is the first estimate to make when answering ``can it work in our case?''.
\end{keypoint}

\section{Higher-level tools}
Besides individual gates, Qiskit offers many higher-level building blocks:
\begin{itemize}
  \item \texttt{qiskit.circuit.library}: a library of ready-made circuits. It contains standard and generalized gates, arithmetic operations, Boolean logic, data-encoding circuits for quantum
  machine learning (e.g.\ the Pauli expansion \texttt{PauliFeatureMap} for supervised learning), state preparation (\texttt{Initialize}), phase and bit-flip oracles, the Grover operator,
  the QFT, and Hamiltonian evolution gates (\texttt{HamiltonianGate}, \texttt{PauliEvolutionGate}, which implements Trotter formulas \cref{eq:trotter}).
  \item \texttt{qiskit.quantum\_info}: utilities for quantum information. It covers quantum states and operators (e.g.\ the \texttt{Statevector} and \texttt{DensityMatrix} classes), partial traces,
  entropies, distances between states, fidelity (\cref{app:fidelity}), quantum channels (Kraus representation), and purity.
\end{itemize}

\section{The half adder}\label{sec:halfadder}
\subsection{The classical half adder}
The simplest arithmetic operation is the addition of two bits $x$ and $y$. The result can be $0$, $1$ or $2=10_2$, so it needs two output bits: the sum $S$ (weight 1) and the carry $C$ (weight 2):
\begin{center}
\begin{tabular}{@{}cc|cc@{}}\toprule
$x$&$y$&$C$&$S$\\ \midrule 0&0&0&0\\0&1&0&1\\1&0&0&1\\1&1&1&0\\ \bottomrule
\end{tabular}
\end{center}
The sum is the exclusive OR of the inputs and the carry is their AND:
\begin{equation}
  S=x\oplus y,\qquad C=x\wedge y.
\end{equation}
The circuit that implements these two functions is called a \emph{half adder}. It is only ``half'' of what is needed to add multi-bit numbers. When adding two binary numbers column by column,
each column must add three bits: the two digits and the carry coming from the previous column. A circuit that adds three bits is called a \emph{full adder}. It can be built from two half adders:
the first adds $x$ and $y$, the second adds the incoming carry to their sum, and the two partial carries are combined with an OR gate.

Adders are among the most important circuits of a classical computer. The arithmetic logic unit (ALU) of every processor is built around them. Through two's complement arithmetic, the same
adder also performs subtraction. Multiplication reduces to repeated shifts and additions, and adders are used constantly to compute memory addresses and increment the program counter.

\subsection{A quantum half adder}
A quantum circuit must be reversible. The classical half adder is not: the outputs $(C,S)=(0,1)$ do not tell us whether the input was $(0,1)$ or $(1,0)$. We therefore keep the inputs and write the
results into two additional qubits $s_0$ (sum) and $s_1$ (carry), both initialized to $\ket0$:
\begin{equation}
  \ket x\ket y\ket0\ket0\longrightarrow\ket x\ket y\ket{x\oplus y}\ket{x\wedge y}.
\end{equation}
\begin{itemize}
  \item The XOR is obtained with two CNOT gates targeting $s_0$: the first adds $x$, the second adds $y$ modulo 2, giving $s_0=x\oplus y$.
  \item The AND is obtained with a Toffoli gate (\texttt{ccx}), which flips the target $s_1$ only if both controls are $1$, giving $s_1=x\wedge y$.
\end{itemize}
\begin{lstlisting}
from qiskit import QuantumRegister, QuantumCircuit, transpile
from qiskit_aer import Aer

x = QuantumRegister(1, 'x')
y = QuantumRegister(1, 'y')
s = QuantumRegister(2, 's')   # s[0]: sum, s[1]: carry
qc = QuantumCircuit(x, y, s)

qc.h(x)                       # superposition of inputs 0 and 1
qc.h(y)
qc.barrier()
qc.cx(x, s[0])                # s0 = x
qc.cx(y, s[0])                # s0 = x XOR y
qc.ccx(x, y, s[1])            # s1 = x AND y
qc.barrier()
qc.measure_all()

simulator = Aer.get_backend('qasm_simulator')
counts = simulator.run(transpile(qc, simulator), shots=1000).result().get_counts()
\end{lstlisting}
\begin{figure}[h]\centering
\begin{quantikz}
\lstick{$x$} & \gate{H} & \slice{} & \ctrl{2} & \qw & \ctrl{3} & \slice{} & \meter{}\\
\lstick{$y$} & \gate{H} & \qw & \qw & \ctrl{1} & \ctrl{2} & \qw & \meter{}\\
\lstick{$s_0$} & \qw & \qw & \targ{} & \targ{} & \qw & \qw & \meter{}\\
\lstick{$s_1$} & \qw & \qw & \qw & \qw & \targ{} & \qw & \meter{}
\end{quantikz}
\caption{Quantum half adder. The Hadamard gates prepare the inputs in a superposition of all four values. Two CNOT gates compute the sum $s_0=x\oplus y$, and a Toffoli gate computes the carry $s_1=x\wedge y$.}
\label{fig:halfadder}
\end{figure}
Instead of a single input value, the Hadamard gates prepare the superposition of all four inputs. The circuit then evaluates the four additions simultaneously:
\begin{equation}
  \frac12\left(\ket{00}+\ket{01}+\ket{10}+\ket{11}\right)_{xy}\ket{00}_s\longrightarrow\frac12\sum_{x,y=0}^1\ket x\ket y\ket{x\oplus y}\ket{x\wedge y}.
\end{equation}
A run on the ideal simulator gives, for 1000 shots, approximately \texttt{\{'0000': 250, '0110': 250, '0101': 250, '1011': 250\}}. In Qiskit's little-endian convention the bit strings read
$s_1s_0yx$. For example, \texttt{'1011'} means $x=1$, $y=1$, sum $s_0=0$ and carry $s_1=1$, i.e.\ $1+1=10_2$. All four lines of the truth table appear, each with probability $1/4$. This illustrates
both the power and the limitation of quantum parallelism: the four additions are performed by a single run, but a measurement reveals only one of them, chosen at random. Useful quantum algorithms
must use interference to concentrate the probability on the desired answer before measuring.

\paragraph{Running on (simulated) hardware.} With a fake backend reproducing a real IBM processor, for example \texttt{FakeMarrakesh}, a few percent of the shots give outcomes that are not in the truth
table, such as \texttt{'1111'} or \texttt{'0111'}. The reason becomes clear after transpilation. The device has only one native two-qubit gate and limited connectivity, so the Toffoli gate alone is
decomposed into six two-qubit gates plus single-qubit rotations (\cref{sec:toffoli}), and SWAPs are added to bring the qubits together. The three logical gates of \cref{fig:halfadder} become about
15 two-qubit gates and nearly 60 single-qubit operations. With $\varepsilon_2\approx3\times10^{-3}$ and $\varepsilon_1\approx3\times10^{-4}$, the estimate of the previous section gives
$1-P_{\rm success}\approx15\times0.003+60\times0.0003\approx6\%$, in agreement with the observed error rate.

\subsection{Towards multi-bit addition: the ripple-carry adder}
To add two $n$-bit numbers $a$ and $b$, $n$ adders are chained as in addition by hand. The carry produced by the adder of bit $i$ is passed as an input to the adder of bit $i+1$, so the carry
``ripples'' from the least to the most significant bit. This \emph{ripple-carry adder} handles the incoming carries and adds numbers of arbitrary size, which a half adder cannot do. Its cost grows
only linearly with $n$.

Quantum versions of the ripple-carry adder were among the first quantum arithmetic circuits to be designed \cite{VBE1996}. They are better than a naive chain of half adders in an important respect:
the number of qubits. A naive reversible construction must keep every intermediate result, which consumes many ancilla qubits. The circuit of Cuccaro et al.\ \cite{Cuccaro2004} computes the sum in
place, $\ket a\ket b\to\ket a\ket{a+b}$, using a single ancilla qubit, by uncomputing the intermediate carries once they are no longer needed. Since qubits are the scarcest resource of present-day
quantum computers, this economy is essential. Both circuits are available in Qiskit as \texttt{VBERippleCarryAdder} and \texttt{CDKMRippleCarryAdder}, together with the single-bit
\texttt{HalfAdderGate} and \texttt{FullAdderGate}.

In classical processors, the ripple-carry adder is the simplest but also the slowest multi-bit adder, because the last bit must wait for the carry to propagate through all the previous ones. Faster
designs, such as carry-lookahead adders, compute the carries in parallel. The same idea leads to quantum adders of logarithmic depth, and to adders based on the quantum Fourier transform
(\texttt{DraperQFTAdder}): in the Fourier basis, adding $a$ is a set of single-qubit phase rotations (\cref{ch:shor}). Adders are the building blocks of modular exponentiation, the most expensive part
of Shor's algorithm.

\chapter[Oracle algorithms]{Oracle algorithms: Deutsch--Jozsa, Bernstein--Vazirani, Simon}\label{ch:dj}

The first quantum algorithms were formulated in the \emph{query} (or oracle) model: a function $f$ is given as a black box, and we count how many times it must be evaluated to learn some
property of $f$. In this model speed-ups can be proved rigorously, and the algorithms show the basic mechanism of quantum computation: parallelism, phase kickback and interference.

\section{The oracle and phase kickback}
A function $f:\{0,1\}^n\to\{0,1\}$ is accessed through the unitary oracle \cref{eq:Uf}, $U_f\ket{x,y}=\ket{x,y\oplus f(x)}$. If the output qubit is prepared in $\ket-=(\ket0-\ket1)/\sqrt2$,
\begin{equation}\label{eq:kickback}
  U_f\ket x\ket-=\ket x\frac{\ket{f(x)}-\ket{1\oplus f(x)}}{\sqrt2}=(-1)^{f(x)}\ket x\ket-,
\end{equation}
since flipping $\ket-$ just changes its sign. The output qubit is unchanged and the value of the function appears as a \emph{phase} of the input register (\emph{phase kickback}). We can forget
the output qubit and describe the oracle as the \emph{phase oracle} $O_f\ket x=(-1)^{f(x)}\ket x$.

\section{The Hadamard transform}
\begin{lemma}\label{lem:hadamard}
For a single qubit, $H\ket x=\frac1{\sqrt2}\sum_{z=0}^1(-1)^{xz}\ket z$. For $n$ qubits,
\begin{equation}\label{eq:Hn}
  H^{\otimes n}\ket x=\frac1{\sqrt{2^n}}\sum_{z\in\{0,1\}^n}(-1)^{x\cdot z}\ket z,\qquad x\cdot z=x_1z_1\oplus\dots\oplus x_nz_n.
\end{equation}
\end{lemma}
\begin{proof}
For one qubit: $H\ket0=(\ket0+\ket1)/\sqrt2$ and $H\ket1=(\ket0-\ket1)/\sqrt2$; the sign is negative only for $x=z=1$. For $n$ qubits, $H^{\otimes n}\ket{x_1\dots x_n}=\bigotimes_iH\ket{x_i}
=\bigotimes_i\frac1{\sqrt2}\sum_{z_i}(-1)^{x_iz_i}\ket{z_i}$; expanding the product gives the sum over all strings $z$ with the sign $(-1)^{\sum_ix_iz_i}$, and only the parity of the exponent matters.
\end{proof}
In particular $H^{\otimes n}\ket{0\dots0}=2^{-n/2}\sum_z\ket z$, the uniform superposition, and $H^{\otimes n}$ is its own inverse. $H^{\otimes n}$ is the Fourier transform over the group $\Z_2^n$.

\section{Deutsch's algorithm ($n=1$)}
Deutsch (1985) asked whether a function $f:\{0,1\}\to\{0,1\}$ is constant ($f(0)=f(1)$) or balanced ($f(0)\neq f(1)$). Classically, two evaluations are needed. Quantum mechanically:
\begin{align}
  \ket0\ket1&\xrightarrow{H\otimes H}\frac{\ket0+\ket1}{\sqrt2}\ket-\xrightarrow{U_f}\frac{(-1)^{f(0)}\ket0+(-1)^{f(1)}\ket1}{\sqrt2}\ket-\nonumber\\
  &\quad=\pm\frac{\ket0+(-1)^{f(0)\oplus f(1)}\ket1}{\sqrt2}\ket-\xrightarrow{H\otimes I}\pm\ket{f(0)\oplus f(1)}\ket- .
\end{align}
Measuring the first qubit gives $f(0)\oplus f(1)$, a \emph{global} property of $f$, with a single evaluation. Note that neither $f(0)$ nor $f(1)$ is learned.

\section{The Deutsch--Jozsa algorithm}
The Deutsch--Jozsa algorithm \cite{DeutschJozsa1992} was one of the first examples of a problem that a quantum computer solves dramatically faster than a classical one.

\paragraph{The problem.} We are given a function $f$ of $n$ bits (an integer $x\in\{0,\dots,2^n-1\}$) with values $f(x)\in\{0,1\}$. We are \emph{promised} that $f$ is either constant ($0$ for all $x$,
or $1$ for all $x$) or balanced ($0$ for exactly half of the inputs and $1$ for the other half). The task is to decide which. A classical deterministic algorithm must, in the worst case, check one more
than half of the inputs, i.e.\ $2^{n-1}+1$ queries, exponentially many. The Deutsch--Jozsa algorithm requires a single query.\footnote{A probabilistic classical algorithm does much better than the worst
case: after $k$ random queries giving equal values, the probability that a balanced function would produce this outcome is about $2^{1-k}$. The exponential separation holds for an exact, deterministic answer.}

\paragraph{The algorithm.} It uses $n$ input qubits and one output qubit, and starts in $\ket{\psi_0}=\ket0^{\otimes n}\ket1$.
\begin{enumerate}
  \item \emph{Superposition.} A Hadamard gate on each qubit gives
  \begin{equation}
    \ket{\psi_1}=\sum_{x=0}^{2^n-1}\frac{\ket x}{\sqrt{2^n}}\left[\frac{\ket0-\ket1}{\sqrt2}\right].
  \end{equation}
  \item \emph{Oracle.} By phase kickback \cref{eq:kickback},
  \begin{equation}
    \ket{\psi_2}=\sum_x\frac{(-1)^{f(x)}\ket x}{\sqrt{2^n}}\left[\frac{\ket0-\ket1}{\sqrt2}\right].
  \end{equation}
  Thanks to the superposition, a single application of the oracle has evaluated $f$ for all $2^n$ inputs at once (quantum parallelism).
  \item \emph{Interference.} Hadamard gates are applied again to the input qubits. By \cref{lem:hadamard},
  \begin{equation}
    \ket{\psi_3}=\sum_z\sum_x\frac{(-1)^{x\cdot z+f(x)}}{2^n}\ket z\left[\frac{\ket0-\ket1}{\sqrt2}\right].
  \end{equation}
  \item \emph{Measurement.} The input register is measured. The amplitude of $\ket{0}^{\otimes n}$ ($z=0$, for which $x\cdot z=0$) is
  \begin{equation}
    a_0=\frac1{2^n}\sum_x(-1)^{f(x)}.
  \end{equation}
  \begin{itemize}
    \item If $f$ is constant, all the terms are equal and $a_0=\pm1$. Since the state is normalized, all the other amplitudes vanish, and the measurement gives $0\dots0$ with certainty.
    \item If $f$ is balanced, half of the terms are $+1$ and half are $-1$. They cancel, $a_0=0$, and the result $0\dots0$ never occurs.
  \end{itemize}
\end{enumerate}
If $0\dots0$ is observed, the function is constant; otherwise, it is balanced. The answer is obtained with certainty after a single query.
\begin{center}
\begin{quantikz}
\lstick{$\ket0^{\otimes n}$} & \qwbundle{n} & \gate{H^{\otimes n}} & \gate[2]{U_f} & \gate{H^{\otimes n}} & \meter{}\\
\lstick{$\ket1$} & \qw & \gate{H} & & \qw & \qw
\end{quantikz}
\end{center}
The algorithm illustrates the two ingredients of a quantum speed-up. Quantum parallelism alone is not enough, since measuring $\ket{\psi_2}$ would give a single random value of $x$. It is the final
Hadamard transform that makes the $2^n$ contributions interfere, constructively or destructively depending on the global property of $f$ we are interested in. The individual values $f(x)$ are never
learned, only whether the function is constant or balanced. For a function that is neither (no promise), $|a_0|^2=(1-2w/2^n)^2$, where $w$ is the number of ones of $f$.

\section{The Bernstein--Vazirani algorithm}
Suppose now that $f_a(x)=a\cdot x\bmod2$ for an unknown string $a$. Classically, $n$ queries are needed (query $x=0\dots010\dots0$ to learn each bit of $a$). The Deutsch--Jozsa circuit gives
\begin{equation}
  H^{\otimes n}\left(\frac1{\sqrt{2^n}}\sum_x(-1)^{a\cdot x}\ket x\right)=H^{\otimes n}H^{\otimes n}\ket a=\ket a,
\end{equation}
using \cref{eq:Hn} and $(H^{\otimes n})^2=I$. A single query reveals $a$ with certainty.

\section{Implementation in Qiskit}
\paragraph{An oracle.} Constructing oracles often requires bitwise operations on integers. In Python, \texttt{x << y} shifts \texttt{x} by \texttt{y} places to the left ($x\cdot2^y$) and \texttt{x \& y}
is the bitwise AND. Thus \texttt{1 << i} is the binary number with a single 1 in position $i$, and \texttt{a \& (1 << i)} is non-zero iff bit $i$ of $a$ is 1. These operations give a simple family of
balanced functions on $n$-bit numbers $k<2^n$:
\begin{equation}
  f_a(k)=\left(\text{number of 1 bits in }a\,\&\,k\right)\bmod2,\qquad0<a<2^n,
\end{equation}
which is exactly $a\cdot k$. The number $a$ selects which bits of $k$ are kept. For $a=1$, $f_a(k)$ is the parity of $k$: half of the numbers are even and half are odd, so $f_a$ is balanced. The same holds
for any $a\neq0$: if bit $j$ of $a$ is set, flipping bit $j$ of $k$ flips $f_a(k)$, which pairs the inputs with $f=0$ and $f=1$. The corresponding oracle applies a CNOT from each selected input qubit to the output qubit:
\begin{lstlisting}
def dj_oracle(a, n):
    oracle = QuantumCircuit(n + 1)
    for i in range(n):
        if a & (1 << i):      # bit i of a is set
            oracle.cx(i, n)   # input qubit i controls the output qubit n
    return oracle
\end{lstlisting}
For $a=0$ the oracle is empty and the function is constant.

\paragraph{The circuit.}
\begin{lstlisting}
from qiskit import QuantumCircuit, transpile
from qiskit_aer import AerSimulator

n, a = 3, 5
qc = QuantumCircuit(n + 1, n)
qc.x(n)                                      # output qubit in |1>
qc.h(range(n + 1))                           # step 1: superposition
qc.compose(dj_oracle(a, n), inplace=True)    # step 2: oracle
qc.h(range(n))                               # step 3: interference
qc.measure(range(n), range(n))               # step 4: measurement

sim = AerSimulator()
counts = sim.run(transpile(qc, sim), shots=1000).result().get_counts()
\end{lstlisting}
For $a=0$ (constant function) all 1000 shots give \texttt{'000'}. For $a=5$ (balanced function) all shots give \texttt{'101'}, a non-zero result, as expected. For this family of oracles the measured
string is the number $a$ itself (Bernstein--Vazirani); with Qiskit's little-endian printing, the string \texttt{'101'} is directly the binary representation of $a=5$.

\section{Simon's algorithm\starred}
The Deutsch--Jozsa speed-up disappears if a small probability of error is allowed. Simon (1994) found a problem with an exponential separation even against randomized classical algorithms,
which inspired Shor.

\paragraph{Problem.} $f:\{0,1\}^n\to\{0,1\}^n$ satisfies $f(x)=f(y)\iff y=x\oplus s$ for an unknown non-zero ``period'' $s$. Find $s$. Classically, one must find a collision $f(x)=f(y)$; by the birthday
paradox this needs $\Omega(2^{n/2})$ queries.

\paragraph{Algorithm.} Prepare $2^{-n/2}\sum_x\ket x\ket0$, apply $U_f$ to get $2^{-n/2}\sum_x\ket x\ket{f(x)}$, and measure the second register. The result $f(x_0)$ leaves the first register in
$(\ket{x_0}+\ket{x_0\oplus s})/\sqrt2$. Applying $H^{\otimes n}$,
\begin{equation}
  \frac1{\sqrt{2^{n+1}}}\sum_z\left[(-1)^{x_0\cdot z}+(-1)^{(x_0\oplus s)\cdot z}\right]\ket z=\frac1{\sqrt{2^{n+1}}}\sum_z(-1)^{x_0\cdot z}\left[1+(-1)^{s\cdot z}\right]\ket z .
\end{equation}
Only strings with $s\cdot z=0$ survive, each with probability $2^{1-n}$. Each run gives a random $z$ orthogonal to $s$ (mod 2). After $\order(n)$ runs we have $n-1$ linearly independent equations
$s\cdot z_k=0$, which Gaussian elimination over $\Z_2$ solves for $s$ in $\order(n^3)$ classical operations. The quantum algorithm uses $\order(n)$ queries: an exponential speed-up. Shor's algorithm
(\cref{ch:shor}) is the same idea with the group $\Z_2^n$ replaced by $\Z_N$ and $H^{\otimes n}$ by the quantum Fourier transform.

\section{Exercises}
\begin{exercise}
Write the four possible oracles $U_f$ for $n=1$ as $4\times4$ matrices and check Deutsch's algorithm for each of them.
\end{exercise}
\begin{exercise}
Show that for the Deutsch--Jozsa algorithm the final output qubit is always in $\ket-$ and is never entangled with the input register. Why is entanglement nevertheless present in the algorithm
when $U_f$ is decomposed into gates? (Hint: consider the oracle $\CNOT$ alone with a control in $\ket+$ and a target in $\ket0$ instead of $\ket-$.)
\end{exercise}
\begin{exercise}
Run Simon's algorithm by hand for $n=2$, $s=11$, $f(00)=f(11)=00$, $f(01)=f(10)=01$.
\end{exercise}

\chapter{Grover's search algorithm}\label{ch:grover}

Grover's algorithm \cite{Grover1996} solves the problem of searching an unordered database, that is, a list of $N$ elements without any index or sorting. Classically, there is no better strategy than
checking the elements one by one. On average, $N/2$ checks are needed to find the desired element with probability $1/2$, so the cost grows as $N$. Grover's quantum algorithm needs only of the order of
$\sqrt N$ queries: a \emph{quadratic} speed-up. Since many problems (e.g.\ NP problems) can be phrased as a search for an input that satisfies an efficiently checkable condition, the algorithm is very general.

\section{The oracle}
The element to find, $w$, is characterized by the function
\begin{equation}
  f_w(x)=\begin{cases}0,&x\neq w,\\1,&x=w.\end{cases}
\end{equation}
This function is implemented as a quantum oracle acting on an $n$-qubit register $\ket x$ (with $N=2^n$) and one additional qubit $\ket y$: $U\ket x\ket y=\ket x\ket{y\oplus f_w(x)}$, where $\oplus$ denotes
addition modulo 2, $a\oplus b=(a+b)\bmod2$, so that $1\oplus1=0$. For $x=w$ the oracle flips the qubit $\ket y$: $\ket w\ket0\to\ket w\ket1$, $\ket w\ket1\to\ket w\ket0$, while for $x\neq w$ nothing changes.

The best choice for the additional qubit is $\ket y=\ket-=(\ket0-\ket1)/\sqrt2$. Flipping this qubit only changes its sign (phase kickback, \cref{eq:kickback}), so we can forget it and describe the oracle as the operator
\begin{equation}
  O:\ket x\to(-1)^{f_w(x)}\ket x,\qquad O=I-2\proj w .
\end{equation}
It is important to understand what the oracle is. It does not ``know'' the answer; it \emph{recognizes} it. For a search for the solution of a problem (e.g.\ a satisfying assignment of a Boolean formula),
$O$ is a circuit that evaluates the condition, which is usually easy, while finding the input that satisfies it is hard.

\section{The algorithm}
The register is first prepared in an equal superposition of all $N$ basis states by applying a Hadamard gate to each qubit of $\ket{0\dots0}$:
\begin{equation}
  \ket\psi=\frac1{\sqrt N}\sum_{x=0}^{N-1}\ket x=\frac1{\sqrt N}\left(\ket{0\dots00}+\ket{0\dots01}+\dots+\ket{1\dots11}\right).
\end{equation}
The \emph{Grover iteration} $G$ is then repeated several times. It consists of four steps:
\begin{enumerate}
  \item apply the oracle $O$;
  \item apply Hadamard gates to all qubits;
  \item apply a conditional phase shift, which changes the sign of every basis state except $\ket0$: $\ket x\to-(-1)^{\delta_{x0}}\ket x$, where $\delta_{x0}$ is the Kronecker delta. This operation is $2\proj0-I$;
  \item apply Hadamard gates to all qubits again.
\end{enumerate}
Since the Hadamard transform maps $\ket0$ onto $\ket\psi$ and is its own inverse, steps 2--4 together give
\begin{equation}
  H^{\otimes n}(2\proj0-I)H^{\otimes n}=2H^{\otimes n}\proj0H^{\otimes n}-H^{\otimes n}H^{\otimes n}=2\proj\psi-I.
\end{equation}
The Grover iteration is therefore
\begin{equation}\label{eq:grover}
  G=\left(2\proj\psi-I\right)O.
\end{equation}
The operator $D=2\proj\psi-I$ is called the \emph{diffusion} operator, or ``inversion about the mean'': for $\ket\phi=\sum_xa_x\ket x$ with mean amplitude $\bar a=\frac1N\sum_xa_x$,
$D\ket\phi=\sum_x(2\bar a-a_x)\ket x$. Each amplitude is reflected about the mean. After the oracle has flipped the sign of the solution, its amplitude is far below the mean, and the reflection
makes it grow; this is the intuition of Grover's original paper.

\section{Geometric interpretation}
Let us generalize slightly and assume that $M$ of the $N$ elements are solutions ($1\le M<N$). We define two normalized states, the superpositions of the non-solutions and of the solutions:
\begin{equation}
  \ket\alpha=\frac1{\sqrt{N-M}}\sum_{x\notin\text{sol.}}\ket x,\qquad\ket\beta=\frac1{\sqrt M}\sum_{x\in\text{sol.}}\ket x.
\end{equation}
The initial state lies in the plane spanned by $\ket\alpha$ and $\ket\beta$:
\begin{equation}
  \ket\psi=\sqrt{\frac{N-M}N}\ket\alpha+\sqrt{\frac MN}\ket\beta=\cos\frac\theta2\ket\alpha+\sin\frac\theta2\ket\beta,\qquad\sin\frac\theta2=\sqrt{\frac MN}.
\end{equation}
The oracle changes the sign of all the basis states that are solutions: $O\ket\alpha=\ket\alpha$, $O\ket\beta=-\ket\beta$. Both $O$ and $D$ map the plane onto itself, so the whole evolution stays in
this two-dimensional plane, however large $N$ is. In the basis $(\ket\alpha,\ket\beta)$,
\begin{equation}
  O=\begin{pmatrix}1&0\\0&-1\end{pmatrix},\qquad
  D=2\begin{pmatrix}\cos\frac\theta2\\ \sin\frac\theta2\end{pmatrix}\begin{pmatrix}\cos\frac\theta2&\sin\frac\theta2\end{pmatrix}-I=\begin{pmatrix}\cos\theta&\sin\theta\\ \sin\theta&-\cos\theta\end{pmatrix},
\end{equation}
\begin{equation}
  G=DO=\begin{pmatrix}\cos\theta&-\sin\theta\\ \sin\theta&\cos\theta\end{pmatrix}.
\end{equation}
In the $(\ket\alpha,\ket\beta)$ plane, $O$ is a reflection about $\ket\alpha$, and $D$ is a reflection about $\ket\psi$: it leaves the component along $\ket\psi$ unchanged and reverses the orthogonal component.
The product of two reflections is a rotation by twice the angle between the two mirrors. The angle between $\ket\alpha$ and $\ket\psi$ is $\theta/2$, so $G$ rotates the state by $\theta$ towards $\ket\beta$ (\cref{fig:grover}):
\begin{equation}
  G\ket\psi=\cos\frac{3\theta}2\ket\alpha+\sin\frac{3\theta}2\ket\beta,\qquad
  G^k\ket\psi=\cos\frac{(2k+1)\theta}2\ket\alpha+\sin\frac{(2k+1)\theta}2\ket\beta.
\end{equation}
The probability to find a solution after $k$ iterations is exactly
\begin{equation}\label{eq:groverprob}
  P_k=\sin^2\frac{(2k+1)\theta}2.
\end{equation}

\begin{figure}[h]\centering
\begin{tikzpicture}[scale=3]
  \draw[->] (0,0)--(1.25,0) node[right]{$\ket\alpha$ (not solutions)};
  \draw[->] (0,0)--(0,1.2) node[above]{$\ket\beta$ (solutions)};
  \draw[->,thick,blue] (0,0)--({cos(20)},{sin(20)}) node[right]{$\ket\psi$};
  \draw[->,thick,red] (0,0)--({cos(-20)},{sin(-20)}) node[right]{$O\ket\psi$};
  \draw[->,thick,green!50!black] (0,0)--({cos(60)},{sin(60)}) node[above right]{$G\ket\psi$};
  \draw (0.3,0) arc (0:20:0.3); \node at (0.38,0.06) {\small$\frac\theta2$};
  \draw ({0.5*cos(20)},{0.5*sin(20)}) arc (20:60:0.5); \node at (0.42,0.4) {\small$\theta$};
  \draw[dashed,gray] ({-0.3*cos(20)},{-0.3*sin(20)})--({1.2*cos(20)},{1.2*sin(20)});
\end{tikzpicture}
\caption{Geometric picture of one Grover iteration. The oracle reflects $\ket\psi$ about $\ket\alpha$, and the operator $2\proj\psi-I$ reflects the result about $\ket\psi$ (dashed). The net effect is a rotation by $\theta$ towards the solutions.}
\label{fig:grover}
\end{figure}

\section{Number of iterations}
\paragraph{Example: $N=4$, $M=1$.} $\sin\frac\theta2=\frac12$, so $\frac\theta2=30^\circ$ and $\theta=60^\circ$. After a single iteration the state is at the angle $3\theta/2=90^\circ$, exactly along
$\ket\beta$. One query is enough to find the solution with certainty. A classical search needs on average $2.25$ queries if the last element is inferred without checking ($\frac14\cdot1+\frac14\cdot2+\frac12\cdot3$),
or $2.5$ if all elements are checked in a random order until the solution is found.

\paragraph{General case.} We want $(2k+1)\theta/2$ to be as close as possible to $\pi/2$. Since $\frac\pi2-\frac\theta2=\arccos\sqrt{M/N}$, the optimal number of iterations is the integer closest to
\begin{equation}
  k_{\rm opt}=\frac{\arccos\sqrt{M/N}}\theta.
\end{equation}
This requires knowing the number of solutions $M$. If we iterate too long, the state rotates past $\ket\beta$ and the success probability decreases again: \cref{eq:groverprob} is periodic in $k$
(``souffl\'e problem''). Since the final state is within an angle $\theta/2$ of $\ket\beta$, the probability of success is at least $\cos^2(\theta/2)=1-M/N$.

The algorithm works best when $M\ll N$. Then, using $\sin\theta=2\sin\frac\theta2\cos\frac\theta2$,
\begin{equation}
  \theta\approx\sin\theta=2\sqrt{\frac MN\left(1-\frac MN\right)}\approx2\sqrt{\frac MN},
\end{equation}
and the number of iterations is
\begin{equation}
  k\approx\frac\pi{2\theta}\approx\frac\pi4\sqrt{\frac NM}\sim\sqrt N.
\end{equation}
This is the square-root efficiency of Grover's algorithm. For $N=10^6$ and $M=1$, about $785$ iterations instead of $5\times10^5$ classical checks.

\paragraph{Unknown $M$.} If $M$ is unknown, one can choose $k$ at random in $[0,m)$ for increasing $m=1,\lambda,\lambda^2,\dots$ (Boyer, Brassard, Høyer, Tapp 1998), which still finds a solution in
$\order(\sqrt{N/M})$ expected queries. Alternatively, \emph{quantum counting} estimates $\theta$, and therefore $M$, by phase estimation of the eigenvalues $\e^{\pm\ii\theta}$ of $G$ (\cref{ch:shor}).

\section{Optimality}
\begin{theorem}[Bennett, Bernstein, Brassard, Vazirani \cite{BBBV1997}]
Any quantum algorithm that finds a marked element among $N$ with probability at least $1/2$ must query the oracle $\Omega(\sqrt N)$ times.
\end{theorem}
\begin{proof}[Idea of the proof (hybrid argument)]
Let $\ket{\psi_k}$ be the state after $k$ queries when no element is marked, and $\ket{\psi_k^w}$ when $w$ is marked. Each query can move the state by at most $2|a_w|$, where $a_w$ is the amplitude on
$\ket w$: $\norm{\psi_k^w-\psi_k}\le\sum_{j<k}2|a_{w,j}|$. Summing over $w$ and using Cauchy--Schwarz and $\sum_w|a_{w,j}|^2=1$, $\sum_w\norm{\psi_k^w-\psi_k}^2\le4k^2$. But to identify $w$ with high
probability, each $\ket{\psi_k^w}$ must be far from the common state $\ket{\psi_k}$, which requires $\sum_w\norm{\psi_k^w-\psi_k}^2=\Omega(N)$. Hence $k=\Omega(\sqrt N)$.
\end{proof}
Grover's algorithm is therefore optimal: for \emph{unstructured} search, quantum computers give only a quadratic speed-up. Applied to an exponentially large search space $N=2^n$, it reduces $2^n$ to $2^{n/2}$,
which is still exponential. NP-complete problems are not expected to become easy.

\section{Amplitude amplification}
Grover's algorithm generalizes: if an algorithm $\mathcal A$ (any circuit) produces a ``good'' outcome with probability $p$, then the operator $Q=-\mathcal AS_0\mathcal A^{-1}S_{\rm good}$, where $S_0$ and
$S_{\rm good}$ flip the sign of $\ket0$ and of the good states, amplifies the success probability to $\approx1$ in $\order(1/\sqrt p)$ applications, instead of the $\order(1/p)$ repetitions needed classically
(Brassard, Høyer, Mosca, Tapp 2002). Grover is the special case $\mathcal A=H^{\otimes n}$, $p=M/N$. Amplitude amplification gives quadratic speed-ups for many problems, including Monte Carlo estimation (amplitude estimation).

\section{Implementation for $N=4$}
For two qubits and the marked element $w=11$, the phase oracle is simply the CZ gate. The diffusion operator, up to a global phase $-1$, is $H^{\otimes2}(X\otimes X)\,\mathrm{CZ}\,(X\otimes X)H^{\otimes2}$, since
$(X\otimes X)\mathrm{CZ}(X\otimes X)=\diag(-1,1,1,1)=-(2\proj{00}-I)$.
\begin{center}
\begin{quantikz}
\lstick{$\ket0$} & \gate{H} & \ctrl{1}\gategroup[2,steps=1,style={dashed,rounded corners,inner sep=2pt},label style={label position=below,anchor=north,yshift=-0.2cm}]{oracle} & \gate{H} & \gate{X} & \ctrl{1} & \gate{X} & \gate{H} & \meter{}\\
\lstick{$\ket0$} & \gate{H} & \control{} & \gate{H} & \gate{X} & \control{} & \gate{X} & \gate{H} & \meter{}
\end{quantikz}
\end{center}
\begin{lstlisting}
from qiskit import QuantumCircuit
from qiskit.quantum_info import Statevector

qc = QuantumCircuit(2)
qc.h([0, 1])
qc.cz(0, 1)                         # oracle marking |11>
qc.h([0, 1]); qc.x([0, 1])          # diffusion operator
qc.cz(0, 1)
qc.x([0, 1]); qc.h([0, 1])
print(Statevector(qc).probabilities_dict())   # {'11': 1.0}
\end{lstlisting}
For larger $n$, the conditional phase shift $2\proj0-I$ requires a multi-controlled $Z$ gate on $n$ qubits, which is decomposed into $\order(n)$ Toffoli gates with ancillas. The library function
\texttt{qiskit.circuit.library.grover\_operator} builds $G$ from any oracle circuit. On present hardware, the depth of the multi-controlled gates limits useful demonstrations to a few qubits.

\section{Exercises}
\begin{exercise}
Show that for $N=8$, $M=1$, the optimal number of iterations is 2, and compute the success probability. (Answer: $\sin^2(5\theta/2)$ with $\sin(\theta/2)=1/\sqrt8$, i.e.\ $\approx0.945$.)
\end{exercise}
\begin{exercise}
Show that for $M=N/4$ a single Grover iteration finds a solution with certainty, and that for $M=N/2$ Grover's algorithm does not help at all. What is the classical success probability with one query in each case?
\end{exercise}
\begin{exercise}
Show that $G$ has eigenvalues $\e^{\pm\ii\theta}$ in the $(\ket\alpha,\ket\beta)$ plane, with eigenvectors $(\ket\alpha\pm\ii\ket\beta)/\sqrt2$.
\end{exercise}

\chapter[QFT and Shor's algorithm]{The quantum Fourier transform and Shor's algorithm}\label{ch:shor}

In 1994 Peter Shor showed that a quantum computer can factor an integer $N$ in a time that grows only polynomially with the number of its digits \cite{Shor1997}. No polynomial classical algorithm is known
for this problem. The difficulty of factoring large numbers is the basis of the RSA public-key cryptosystem, which is why Shor's algorithm had such a strong impact. The algorithm combines a classical
reduction of factoring to the problem of finding the period of a function, and a quantum subroutine that finds this period efficiently using the quantum Fourier transform (QFT). We first construct
the QFT and the phase-estimation algorithm, which are the core of most exponential quantum speed-ups.

\section{The quantum Fourier transform}
\subsection{Definition}
For $Q=2^t$, the QFT is defined on the basis states of $t$ qubits by
\begin{equation}\label{eq:qft}
  U_{\rm QFT}\ket x=\frac1{\sqrt Q}\sum_{y=0}^{Q-1}\e^{2\pi\ii xy/Q}\ket y .
\end{equation}
On a general state it maps the amplitudes $a_x$ to their discrete Fourier transform $\tilde a_y=\frac1{\sqrt Q}\sum_x\e^{2\pi\ii xy/Q}a_x$. It is unitary, since
$\mel{x'}{U^\dagger U}x=\frac1Q\sum_y\e^{2\pi\ii(x-x')y/Q}=\delta_{xx'}$
(a geometric sum of $Q$-th roots of unity). For a single qubit ($Q=2$) it is the Hadamard gate.

\subsection{Product form and circuit}
Write $x=x_1x_2\dots x_t$ in binary ($x=\sum_jx_j2^{t-j}$) and use the binary fraction notation $0.x_lx_{l+1}\dots x_t=\sum_{j\ge l}x_j2^{l-1-j}$. Writing also $y=\sum_ly_l2^{t-l}$,
\begin{equation}
  \e^{2\pi\ii xy/Q}=\prod_{l=1}^t\e^{2\pi\ii xy_l2^{-l}},\qquad\text{and}\qquad\e^{2\pi\ii x2^{-l}}=\e^{2\pi\ii\,0.x_{t-l+1}\dots x_t},
\end{equation}
since the integer part of $x2^{-l}$ does not contribute. The sum over $y$ factorizes into a sum over each bit $y_l$:
\begin{equation}\label{eq:qftproduct}
  U_{\rm QFT}\ket{x_1\dots x_t}=\frac{\left(\ket0+\e^{2\pi\ii\,0.x_t}\ket1\right)\otimes\left(\ket0+\e^{2\pi\ii\,0.x_{t-1}x_t}\ket1\right)\otimes\dots\otimes\left(\ket0+\e^{2\pi\ii\,0.x_1x_2\dots x_t}\ket1\right)}{2^{t/2}}.
\end{equation}
This remarkable formula shows that the QFT of a basis state is a \emph{product state}, and gives the circuit. The phase $\e^{2\pi\ii\,0.x_1x_2\dots x_t}$ of the last factor is obtained by applying $H$ to $\ket{x_1}$
(which gives $\ket0+\e^{2\pi\ii\,0.x_1}\ket1$, since $\e^{\ii\pi x_1}=(-1)^{x_1}$) and then the controlled phase rotations
\begin{equation}
  R_k=\begin{pmatrix}1&0\\0&\e^{2\pi\ii/2^k}\end{pmatrix},\qquad k=2,\dots,t,
\end{equation}
controlled by $x_k$, which add $x_k2^{-k}$ to the binary fraction. Repeating for $x_2,\dots,x_t$ and reversing the order of the qubits at the end:
\begin{center}
\begin{quantikz}[column sep=0.25cm]
\lstick{$x_1$} & \gate{H} & \gate{R_2} & \gate{R_3} & \qw & \qw & \qw & \qw & \qw & \swap{2} & \qw\\
\lstick{$x_2$} & \qw & \ctrl{-1} & \qw & \gate{H} & \gate{R_2} & \qw & \qw & \qw & \qw & \qw\\
\lstick{$x_3$} & \qw & \qw & \ctrl{-2} & \qw & \ctrl{-1} & \gate{H} & \qw & \qw & \targX{} & \qw
\end{quantikz}
\end{center}
The number of gates is $t$ Hadamards, $t(t-1)/2$ controlled rotations and $\lfloor t/2\rfloor$ swaps: $\order(t^2)$, compared with $\order(tQ)=\order(t2^t)$ operations for the classical fast Fourier
transform of a vector of $2^t$ amplitudes. The rotations $R_k$ with $k\gtrsim\log_2t$ can be dropped with negligible error (approximate QFT, $\order(t\log t)$ gates). The exponential gain must be
interpreted with care: the QFT does not give access to the Fourier coefficients, only to samples from $|\tilde a_y|^2$. It is useful when the transformed state has a few sharp peaks that encode the
answer, as in phase estimation and period finding.

\section{Phase estimation}
\paragraph{Problem.} Given a unitary $U$ (as a circuit, together with its controlled powers $U^{2^j}$) and an eigenstate $\ket u$ with $U\ket u=\e^{2\pi\ii\varphi}\ket u$, estimate the phase $\varphi\in[0,1)$.

\paragraph{Algorithm.} Use $t$ ``counting'' qubits in $\ket0$ and a register in $\ket u$.
\begin{enumerate}
  \item Apply $H^{\otimes t}$ to the counting register: $2^{-t/2}\sum_{y=0}^{2^t-1}\ket y\ket u$.
  \item Apply controlled-$U^{2^j}$ from counting qubit $j$ (of weight $2^j$). Since $\ket u$ is an eigenstate, the phases are kicked back:
  \begin{equation}
    \frac1{2^{t/2}}\sum_y\ket y\,U^y\ket u=\frac1{2^{t/2}}\sum_y\e^{2\pi\ii\varphi y}\ket y\otimes\ket u .
  \end{equation}
  \item Apply the inverse QFT. If $\varphi=x/2^t$ exactly, the state above is $U_{\rm QFT}\ket x$ by \cref{eq:qft}, and the inverse QFT gives $\ket x$: the measurement returns the $t$ bits of $\varphi$ exactly.
\end{enumerate}
\paragraph{Accuracy.} If $\varphi$ is not a $t$-bit fraction, let $b$ be the integer closest to $2^t\varphi$ and $\delta=\varphi-b/2^t$, $|\delta|\le2^{-t-1}$. The amplitude of $\ket b$ is
\begin{equation}
  \frac1{2^t}\sum_{y=0}^{2^t-1}\e^{2\pi\ii\delta y}=\frac1{2^t}\frac{1-\e^{2\pi\ii\delta2^t}}{1-\e^{2\pi\ii\delta}},\qquad
  P(b)=\frac{\sin^2(\pi\delta2^t)}{2^{2t}\sin^2(\pi\delta)}\ge\frac{(2\cdot2^t\delta)^2}{2^{2t}(\pi\delta)^2}=\frac4{\pi^2}\approx0.405,
\end{equation}
using $|\sin\pi u|\ge2|u|$ for $|u|\le\frac12$ and $|\sin\pi\delta|\le\pi|\delta|$. With $t=m+\lceil\log_2(2+1/2\epsilon)\rceil$ qubits, $\varphi$ is obtained to $m$ bits with probability at least $1-\epsilon$.
Each additional bit of precision doubles the number of applications of $U$: the cost is $\order(1/\text{precision})$, the ``Heisenberg limit'', compared with $\order(1/\text{precision}^2)$ for repeated
measurements of $\expval U$ (shot noise). Phase estimation is used in Shor's algorithm, quantum counting, the HHL algorithm (\cref{ch:qml}) and quantum chemistry (energies of molecules as phases of $\e^{-\ii Ht}$).

\section{From factoring to period finding}
\subsection{The classical reduction}
\begin{enumerate}
  \item Let $N$ be the integer to factor (odd, and not a prime power; both cases are easily detected classically).
  \item Choose a random integer $1<a<N$.
  \item Compute $\gcdf(a,N)$, for example with Euclid's algorithm.
  \item If $\gcdf(a,N)\neq1$, it is a non-trivial factor of $N$, and we are done.
  \item Otherwise, use the quantum period-finding subroutine to find the period $r$ of the function
  \begin{equation}
    f(x)=a^x\bmod N,
  \end{equation}
  i.e.\ the smallest $r>0$ such that $a^r\equiv1\pmod N$ (the \emph{order} of $a$ modulo $N$; it exists because $a$ is invertible modulo $N$ and the invertible residues form a finite group).
  \item If $r$ is odd, or if $a^{r/2}\equiv-1\pmod N$, go back to step 2.
  \item Otherwise, the factors of $N$ are $\gcdf(a^{r/2}+1,N)$ and $\gcdf(a^{r/2}-1,N)$.
\end{enumerate}
All steps except step 5 are easy for a classical computer: Euclid's algorithm computes a greatest common divisor in a time polynomial in the number of digits, and $a^{r/2}\bmod N$ is computed by repeated squaring.

\paragraph{Why does it work?} By definition of the period, $a^r-1$ is a multiple of $N$. If $r$ is even, this can be factored as
\begin{equation}
  a^r-1=\left(a^{r/2}-1\right)\left(a^{r/2}+1\right)\equiv0\pmod N.
\end{equation}
The first factor is not a multiple of $N$, since $r$ is the smallest exponent with $a^r\equiv1$. If the second factor is not a multiple of $N$ either, i.e.\ if $a^{r/2}\not\equiv-1$, then $N$ divides the product
without dividing either factor. Each factor must therefore share a non-trivial divisor with $N$, which the gcd extracts. For an odd $N$ with at least two distinct prime factors, a random $a$ satisfies these
conditions with probability at least $1/2$ (the proof uses the Chinese remainder theorem: $r$ is the lcm of the orders modulo each prime power $p_i^{\alpha_i}$, and the conditions fail only if all these orders
contain the same power of 2), so only a few repetitions are needed on average.

\paragraph{Numerical example.} Let $N=15$ and choose $a=11$. Then $\gcdf(11,15)=1$, and we continue. The values of the function are
\begin{center}
\begin{tabular}{@{}lcccc@{}}\toprule
$x$ & 0&1&2&3\\ $a^x$ & 1&11&121&1331\\ $f(x)=a^x\bmod15$ & 1&11&1&11\\ \bottomrule
\end{tabular}
\end{center}
so the period is $r=2$. It is even, and $11^{2/2}\bmod15=11\neq-1\equiv14$, so we continue: $\gcdf(11+1,15)=3$, $\gcdf(11-1,15)=5$, and $15=3\times5$. For a number this small, the period is found by simply computing
the function. For a number with hundreds of digits, the period can be astronomically large (up to $N$), and this direct approach is hopeless. This is where the quantum computer comes in.

\section{Quantum period finding}
\subsection{Direct formulation}
The quantum subroutine uses two registers. The first, input register has $t$ qubits and dimension $Q=2^t$, chosen such that $N^2\le Q<2N^2$. The second, output register is large enough to hold the values of $f$ ($\lceil\log_2N\rceil$ qubits).
\begin{enumerate}
  \item \emph{Superposition.} Hadamard gates on the input register prepare $\frac1{\sqrt Q}\sum_{x=0}^{Q-1}\ket x\ket0$.
  \item \emph{Evaluation of the function.} A unitary circuit computing $f$ (modular exponentiation) evaluates the function for all inputs at once: $\frac1{\sqrt Q}\sum_x\ket x\ket{f(x)}$.
  \item \emph{Quantum Fourier transform} on the input register:
  \begin{equation}
    \frac1Q\sum_x\sum_y\e^{2\pi\ii xy/Q}\ket y\ket{f(x)}.
  \end{equation}
  \item \emph{Measurement.} Both registers are measured. The output register gives some value $f(x_0)$, and the input register gives some $y$. Since $f(x)=f(x_0)$ exactly for $x=x_0+rb$, $b=0,1,2,\dots$,
  the probability of a given pair of results is
  \begin{equation}
    P(y,f(x_0))=\abs{\frac1Q\sum_{x:\,f(x)=f(x_0)}\e^{2\pi\ii xy/Q}}^2=\abs{\frac1Q\sum_{b=0}^{m-1}\e^{2\pi\ii(x_0+rb)y/Q}}^2,\qquad m\approx Q/r.
  \end{equation}
\end{enumerate}
The phase $\e^{2\pi\ii x_0y/Q}$ factors out, and what remains is a geometric sum of the phases $\e^{2\pi\ii\,bry/Q}$. These add up constructively only if $ry/Q$ is close to an integer $k$, and cancel otherwise.
The measured $y$ is therefore, with high probability, close to a multiple of $Q/r$:
\begin{equation}
  \frac yQ\approx\frac kr .
\end{equation}
\paragraph{Peak probability.} Let $y$ be the integer closest to $kQ/r$, so that $ry=kQ+\epsilon$ with $|\epsilon|\le r/2$. The phases are $2\pi b\epsilon/Q$ and, exactly as in phase estimation,
$|\sum_{b<m}\e^{2\pi\ii b\epsilon/Q}|\ge\frac2\pi m$ because $m|\epsilon|/Q\le\frac12$. Hence $P(y,f(x_0))\ge\frac4{\pi^2}\frac{m^2}{Q^2}\approx\frac4{\pi^2r^2}$. Summing over the $r$ possible values of $f(x_0)$ and
the $r$ values of $k$, the probability to obtain one of the peaks is at least $4/\pi^2\approx40\%$.

\paragraph{Continued fractions.} Since $|y/Q-k/r|\le1/(2Q)\le1/(2N^2)<1/(2r^2)$, a classical theorem guarantees that $k/r$ is one of the \emph{convergents} of the continued-fraction expansion of $y/Q$, which
are computed in polynomial time. This gives $k/r$ in lowest terms, i.e.\ $r$ itself if $\gcdf(k,r)=1$ (which happens with probability $\Omega(1/\log\log r)$), or a divisor of $r$, in which case the procedure is
repeated and the least common multiple of the results is taken. The choice $Q\ge N^2$ is exactly what makes the continued-fraction step work.

\subsection{The phase-estimation formulation}
Equivalently, period finding is phase estimation for the unitary $U_a\ket z=\ket{az\bmod N}$ (a permutation of the residues coprime with $N$). Its eigenstates are
\begin{equation}
  \ket{u_s}=\frac1{\sqrt r}\sum_{k=0}^{r-1}\e^{-2\pi\ii sk/r}\ket{a^k\bmod N},\qquad U_a\ket{u_s}=\e^{2\pi\ii s/r}\ket{u_s},\quad s=0,\dots,r-1,
\end{equation}
as one checks by shifting the summation index. We cannot prepare a given $\ket{u_s}$ without knowing $r$, but $\frac1{\sqrt r}\sum_s\ket{u_s}=\ket1$: starting the target register in $\ket1$, phase estimation returns
$s/r$ for a random $s$. The controlled-$U_a^{2^j}$ gates are controlled multiplications by $a^{2^j}\bmod N$, whose constants are precomputed classically by repeated squaring. This is the \emph{modular
exponentiation} circuit, built from quantum adders (\cref{sec:halfadder}); it dominates the cost, with $\order(n^3)$ elementary gates for an $n$-bit number (or $\order(n^2\log n)$ with fast multiplication).
The QFT costs only $\order(n^2)$.

\begin{keypoint}
The QFT is the key element. It converts the periodicity of $f$, hidden in a superposition that cannot be read out directly, into sharp peaks in the probability distribution of $y$, through interference.
This is an exponential speed-up compared with the best known classical algorithm, the number field sieve, whose running time is $\exp\left(\order(n^{1/3}(\log n)^{2/3})\right)$.
\end{keypoint}

\section{Example: factoring 15}
The algorithm was first demonstrated in 2001 by factoring $N=15$ on a seven-qubit NMR quantum computer \cite{Vandersypen2001}. The input register had only $t=3$ qubits ($Q=8$). This is less than the $N^2$
required in general, but it is sufficient here because the period divides $Q$.

\paragraph{The easy case, $a=11$.} The period is $r=2$, and $Q/r=4$. If the output register is measured in the state $11$, the input register is left in $\frac12(\ket1+\ket3+\ket5+\ket7)$, and its Fourier transform
has amplitudes $\frac1{4\sqrt2}\sum_{x\text{ odd}}\e^{2\pi\ii xy/8}$, which vanish unless $y$ is a multiple of 4. The input register is found in a mixture of $\ket{000}$ and $\ket{100}$, i.e.\ $y=0$ or $y=4$, each
with probability $1/2$. The period of the distribution of $y$ is 4, hence $r=Q/4=2^3/4=2$.

\paragraph{The more complicated case, $a=7$.} Now $7^x\bmod15=1,7,4,13,1,7,\dots$, so that $r=4$. If the output is $4$, the input register is left in $(\ket2+\ket6)/\sqrt2$, with Fourier amplitudes
$\frac14\e^{\ii\pi y/2}(1+\e^{\ii\pi y})$, non-zero for even $y$ only. The input register is found in a mixture of $\ket{000}$, $\ket{010}$, $\ket{100}$, $\ket{110}$, i.e.\ $y=0,2,4,6$, each with probability $1/4$.
The spacing of the peaks is $Q/r=2$, hence $r=Q/2=2^3/2=4$. (In terms of continued fractions: $y/Q=2/8=1/4$ or $6/8=3/4$ give $r=4$; $4/8=1/2$ gives only the divisor $2$, and $0$ gives nothing.) Then
$7^{4/2}\bmod15=4\neq-1$, and $\gcdf(4+1,15)=5$, $\gcdf(4-1,15)=3$.

\paragraph{Outlook.} Factoring 15 is of course trivial classically. The importance of these experiments is to demonstrate all the elements of the algorithm (and many ``demonstrations'' simplify the circuit by using
knowledge of the answer, which must be kept in mind). Factoring numbers large enough to threaten RSA would require thousands of error-corrected logical qubits; resource estimates for RSA-2048 have decreased from
about $10^9$ physical qubits (2012) to about $2\times10^7$ (Gidney and Ekerå, 2019) and below $10^6$ noisy qubits running for about a week (Gidney, 2025), still far beyond present technology. This prospect has
nevertheless motivated the development and standardization (NIST, 2024) of post-quantum cryptographic methods, based on mathematical problems such as lattice problems, which are believed to resist quantum attacks.
Shor's algorithm also solves the discrete logarithm problem, which breaks elliptic-curve cryptography. Both are instances of the \emph{hidden subgroup problem} for abelian groups, which quantum computers solve efficiently.

\section{Exercises}
\begin{exercise}
Write the matrix of the QFT for $t=2$ and check that it is unitary. Verify the product formula \cref{eq:qftproduct} for $\ket{x}=\ket{11}$.
\end{exercise}
\begin{exercise}
Apply Shor's reduction to $N=21$ with $a=2$. Find the period, and the factors. What happens with $a=13$? (Answer: $r=6$, $2^3=8$, $\gcdf(9,21)=3$, $\gcdf(7,21)=7$; for $a=13$, $r=2$ and $13\equiv-8$, $\gcdf(14,21)=7$.)
\end{exercise}
\begin{exercise}
Compute the continued-fraction convergents of $y/Q=427/512$ and find the period if $N=21$. (Answer: $427/512\approx0.834$; the expansion is $[0;1,5,42,\dots]$ with convergents $0,\,1,\,5/6,\,211/253,\dots$; $5/6$ is the last one with denominator smaller than $N$ and gives $r=6$.)
\end{exercise}

\chapter{Efficiency of classical and quantum algorithms}\label{ch:complexity}

\section{Measuring the cost of an algorithm}
The usefulness of an algorithm depends on how its cost grows with the size $n$ of the problem: the number of digits of a number to factorize, the number of cities of a route to optimize, or the number
of qubits of a system to simulate. The cost is measured by the number of elementary steps (gates, for a quantum circuit) or, for oracle problems, the number of queries to the oracle. Only the leading behaviour
for large $n$ matters, and it is written with the ``big-O'' notation: $f(n)=\order(g(n))$ means that $f(n)\le Cg(n)$ for some constant $C$ and all large $n$; $\Omega$ denotes a lower bound in the same sense.
For example, $\order(n^2)$ means that the number of steps grows at most as $n^2$, up to a constant factor. Four broad classes of behaviour can be distinguished.
\begin{itemize}
  \item \textbf{Constant}, $\order(1)$: the cost does not depend on $n$ at all.
  \item \textbf{Polynomial}, $\order(n^k)$: the cost grows as a fixed power of $n$. Problems that admit polynomial algorithms are considered \emph{tractable}.
  \item \textbf{Exponential}, $\order(a^n)$ with $a>1$: adding a single unit to the problem size multiplies the cost by $a$. Such problems quickly become intractable.
  \item \textbf{Factorial}, $\order(n!)$: the cost grows even faster than any exponential, since $n!\approx\sqrt{2\pi n}(n/e)^n$ by Stirling's formula. This is typical of a brute-force search over all orderings of $n$ objects.
\end{itemize}

\section{Complexity classes}
A \emph{decision problem} has a yes/no answer. The main classes are:
\begin{itemize}
  \item \textbf{P}: problems solvable by a deterministic classical algorithm in polynomial time.
  \item \textbf{BPP}: problems solvable by a randomized classical algorithm in polynomial time, with error probability at most $1/3$ (repeating and taking the majority makes the error exponentially small).
  \item \textbf{NP}: problems whose ``yes'' answers can be \emph{verified} in polynomial time given a certificate (e.g.\ the factors of a number, a satisfying assignment, a route shorter than $L$).
  A problem is \textbf{NP-complete} if it is in NP and every NP problem reduces to it in polynomial time (SAT, Cook 1971 \cite{Cook1971}; travelling salesman, graph colouring, Ising spin glasses, exact cover, \dots).
  \item \textbf{BQP} (bounded-error quantum polynomial time): problems solvable by a uniform family of polynomial-size quantum circuits with error probability at most $1/3$. This is the class of problems
  that quantum computers solve efficiently.
  \item \textbf{PSPACE}: problems solvable with polynomial memory (and any time).
\end{itemize}
Known relations: $\mathrm P\subseteq\mathrm{BPP}\subseteq\mathrm{BQP}\subseteq\mathrm{PSPACE}$. The last inclusion holds because an amplitude of a quantum circuit is a sum over exponentially many paths, which can be
computed one by one with polynomial memory (Feynman path integral). Factoring is in BQP (Shor) and in NP, but it is \emph{not} believed to be NP-complete; NP-complete problems are not believed to be in BQP.
None of the strict separations is proved: we do not even know that $\mathrm P\neq\mathrm{PSPACE}$. In the oracle model, however, separations are proved (Deutsch--Jozsa, Simon, Grover's lower bound).

\section{Examples}
\Cref{tab:complexity} gives examples of each behaviour for classical and quantum algorithms. Two observations stand out. First, quantum computers do not speed up everything. Their advantage is specific to
problems with a structure that interference can exploit, such as the periodicity used by Shor's algorithm. Second, the speed-up can be of very different kinds. Shor's algorithm turns a problem for which no
polynomial classical algorithm is known into a polynomial one: an exponential speed-up. Grover's algorithm gives only a quadratic speed-up, $N\to\sqrt N$, and this is provably the best possible for unstructured
search \cite{BBBV1997}. Applied to an exponentially large search space $N=2^n$, it reduces $2^n$ to $2^{n/2}$, which is still exponential.

\begin{table}[h]\centering\small
\begin{tabularx}{\textwidth}{@{}l>{\raggedright\arraybackslash}X>{\raggedright\arraybackslash}X@{}}\toprule
Behaviour & Classical algorithms & Quantum algorithms\\ \midrule
Constant $\order(1)$ & Reading an element of an array from its index; testing whether a binary number is even (last bit only). &
Deutsch--Jozsa: deciding whether a function of $n$ bits is constant or balanced with a single oracle query. A deterministic classical algorithm needs $2^{n-1}+1$ queries in the worst case.\\[1ex]
Polynomial $\order(n^k)$ & Sorting a list, $\order(n\log n)$; matrix multiplication, $\order(n^3)$; testing whether an $n$-digit number is prime (AKS, 2002). &
Shor's factoring of an $n$-digit number, about $\order(n^3)$; Grover's search among $N$ elements, $\order(\sqrt N)$; simulation of the dynamics of $n$ interacting quantum particles \cite{Lloyd1996}.\\[1ex]
Exponential $\order(a^n)$ & Brute-force search of an $n$-bit key, $\order(2^n)$; exact simulation of $n$ qubits, which requires storing $2^n$ amplitudes; factoring with the best known classical algorithm
(sub-exponential, $\sim\e^{cn^{1/3}(\log n)^{2/3}}$). & Grover's search of an $n$-bit key, $\order(2^{n/2})$; NP-complete problems (e.g.\ Boolean satisfiability), for which no polynomial quantum algorithm is known or expected.\\[1ex]
Factorial $\order(n!)$ & Travelling salesman problem by checking all $(n-1)!/2$ routes; determinant or permanent of an $n\times n$ matrix computed directly from its definition, a sum of $n!$ terms. &
Grover's search over all routes of the travelling salesman problem, $\order(\sqrt{n!})$: still faster than any exponential.\\ \bottomrule
\end{tabularx}
\caption{Examples of algorithms with constant, polynomial, exponential and factorial numbers of steps, as a function of the problem size $n$.}\label{tab:complexity}
\end{table}

A brute-force algorithm is often much worse than the best available one. The determinant of an $n\times n$ matrix requires $n!$ terms if computed from its definition, but only $\order(n^3)$ operations by Gaussian
elimination. The permanent, which differs from the determinant only by the absence of signs, has no such shortcut: the best known algorithms remain exponential (it is \#P-hard, Valiant 1979). Small changes in a
problem can thus move it from one class to another. The permanent reappears in quantum physics: the amplitudes of non-interacting bosons scattered by a linear interferometer are permanents, which is the basis of
\emph{boson sampling} (Aaronson and Arkhipov, 2011), a proposal for demonstrating quantum computational advantage.

\Cref{fig:growth} shows how dramatic the differences between these classes are. Assume that each step takes one nanosecond. A polynomial algorithm with $n^3$ steps solves a problem of size $n=1000$ in one second.
An exponential algorithm with $2^n$ steps exceeds the age of the universe ($\approx4\times10^{17}$\,s, i.e.\ $\approx4\times10^{26}$ steps) already for $n\approx90$. A factorial algorithm does so for $n\approx27$. For such
problems, faster hardware does not help: a computer a million times faster only increases the accessible size by about 20 for $2^n$ (since $2^{20}\approx10^6$), and by about 4 for $n!$ (since $28\cdot29\cdot30\cdot31\approx10^6$).
Only a better algorithm, classical or quantum, can change the picture.

\begin{figure}[h]\centering
\begin{tikzpicture}
\begin{axis}[width=12cm,height=7cm,ymode=log,xmin=1,xmax=100,ymin=1,ymax=1e36,xlabel={Problem size $n$},ylabel={Number of steps},
  legend style={at={(0.5,-0.22)},anchor=north,legend columns=3,font=\small},ytick={1,1e5,1e10,1e15,1e20,1e25,1e30,1e35}]
\addplot[black,thick,domain=1:100] {1}; \addlegendentry{constant, $1$}
\addplot[blue,thick,domain=1:100,samples=100] {x^3}; \addlegendentry{polynomial, $n^3$}
\addplot[green!50!black,thick,domain=1:100,samples=100] {2^(x/2)}; \addlegendentry{$2^{n/2}$ (Grover)}
\addplot[red,thick,domain=1:100,samples=100] {2^x}; \addlegendentry{exponential, $2^n$}
\addplot[orange,thick] coordinates {(1,1.0000e+00) (2,2.0000e+00) (3,6.0000e+00) (4,2.4000e+01) (5,1.2000e+02) (6,7.2000e+02) (7,5.0400e+03) (8,4.0320e+04) (9,3.6288e+05) (10,3.6288e+06) (11,3.9917e+07) (12,4.7900e+08) (13,6.2270e+09) (14,8.7178e+10) (15,1.3077e+12) (16,2.0923e+13) (17,3.5569e+14) (18,6.4024e+15) (19,1.2165e+17) (20,2.4329e+18) (21,5.1091e+19) (22,1.1240e+21) (23,2.5852e+22) (24,6.2045e+23) (25,1.5511e+25) (26,4.0329e+26) (27,1.0889e+28) (28,3.0489e+29) (29,8.8418e+30) (30,2.6525e+32) (31,8.2228e+33) (32,2.6313e+35)}; \addlegendentry{factorial, $n!$}
\addplot[gray,dashed,domain=1:100] {4e26}; \addlegendentry{age of the universe (1\,ns/step)}
\end{axis}
\end{tikzpicture}
\caption{Number of steps as a function of the problem size $n$, on a logarithmic scale. Polynomial growth appears as a slowly rising curve, exponential growth as a straight line, and factorial growth as an
upward-curving line. The dashed line marks the number of nanoseconds since the Big Bang.}\label{fig:growth}
\end{figure}

\section{Quantum advantage in practice}
The asymptotic classes do not tell the whole story. Constants matter: a quantum gate is $10^3$--$10^6$ times slower than a classical operation, and error correction multiplies the cost further. A quadratic
speed-up therefore pays off only for very long computations: if a classical step takes $1$\,ns and a (logical) quantum step $1\,\mu$s, Grover beats brute force only when $\sqrt N>10^3$, i.e.\ beyond $N\sim10^6$,
and with the overheads of error correction, the cross-over moves to days of computation (Babbush et al., 2021). This is why the most promising applications are those with exponential or large polynomial speed-ups:
factoring and discrete logarithms, and the simulation of quantum systems (chemistry, materials). Experiments claiming ``quantum advantage'' (Google's random-circuit sampling with 53 qubits in 2019, boson sampling with
photons in 2020) address artificial sampling problems, chosen to be hard classically rather than useful, and improved classical algorithms have repeatedly narrowed the gap.

\section{Exercises}
\begin{exercise}
An algorithm needs $n^5$ steps and another $1.1^n$. For which $n$ does the exponential one become slower? Discuss why polynomial is ``tractable'' even if the exponent is large.
\end{exercise}
\begin{exercise}
Estimate the number of Grover iterations needed to search a 128-bit key (AES-128). If each iteration takes $1\,\mu$s, how long does it take? Why is AES-256 considered quantum-safe?
\end{exercise}

\chapter{Quantum error correction}\label{ch:qec}

Quantum states are fragile. The coupling to the environment (\cref{ch:open}) and imperfect gates continuously introduce errors, and a quantum computation of any useful length is impossible without correcting them.
With error rates $\sim10^{-3}$ per gate, a circuit of $10^3$ gates already fails most of the time, while Shor's algorithm for RSA-2048 needs $\sim10^{10}$ gates.

\section{Classical redundancy and the three quantum obstacles}
Classical computers protect information by redundancy: a bit is copied several times, $0\to000$, $1\to111$, and a majority vote recovers the correct value if one copy is corrupted. If each copy is flipped independently
with probability $p$, the vote fails only when at least two copies are flipped:
\begin{equation}
  p_{\rm fail}=3p^2(1-p)+p^3=3p^2-2p^3,
\end{equation}
which is smaller than $p$ when $p<1/2$. In the quantum case, this idea seems to face three obstacles:
\begin{enumerate}
  \item the no-cloning theorem forbids copying an unknown state (\cref{sec:nocloning});
  \item measuring the qubits to check them would destroy the superposition;
  \item quantum errors are continuous, since a qubit can be rotated by any small angle, and it seems impossible to correct a continuum of errors with finite resources.
\end{enumerate}
Quantum error correction overcomes all three obstacles \cite{Shor1995}. We illustrate it with the simplest example, the three-qubit bit-flip code.

\section{The three-qubit bit-flip code}
\subsection{Encoding}
The logical qubit is encoded in three physical qubits, $\ket0\to\ket{000}$, $\ket1\to\ket{111}$, so that a general state is encoded as
\begin{equation}
  a\ket0+b\ket1\to\ket\psi=a\ket{000}+b\ket{111}.
\end{equation}
This is not a copy of the state, which would be $(a\ket0+b\ket1)^{\otimes3}$, and it does not contradict the no-cloning theorem. It is an entangled state, prepared by two CNOT gates controlled by the first qubit
and acting on two additional qubits initialized to $\ket0$:
\begin{center}
\begin{quantikz}
\lstick{$a\ket0+b\ket1$} & \ctrl{1} & \ctrl{2} & \qw\\
\lstick{$\ket0$} & \targ{} & \qw & \qw\\
\lstick{$\ket0$} & \qw & \targ{} & \qw
\end{quantikz}
\end{center}

\subsection{Errors}
Suppose that the second qubit undergoes a small unwanted rotation about the $x$ axis, $E_2=\e^{-\ii\epsilon_2\sigma_x^{(2)}}=\cos\epsilon_2-\ii\sin\epsilon_2\,\sigma_x^{(2)}\approx1-\ii\epsilon_2\sigma_x^{(2)}$. Since $\sigma_x$ flips
$\ket0\leftrightarrow\ket1$, to first order in $\epsilon_2$
\begin{equation}
  E_2\ket\psi=\left[a\ket{000}+b\ket{111}\right]-\ii\epsilon_2\left[a\ket{010}+b\ket{101}\right].
\end{equation}
The state is now a superposition of the correct state and of the state with the second qubit flipped.

\subsection{Syndrome extraction}
To find out which qubit was flipped without measuring the qubits themselves, we measure \emph{parities}. Add ancilla qubits in the state $\ket{000}$:
\begin{equation}
  E_2\ket\psi\ket{000}=\left[a\ket{000}+b\ket{111}\right]\ket{000}-\ii\epsilon_2\left[a\ket{010}+b\ket{101}\right]\ket{000}.
\end{equation}
A circuit $S$ of CNOT gates writes the \emph{syndrome} of the error into the ancillas: the parity $q_1\oplus q_2$ of the first two code qubits into the first ancilla, and the parity $q_2\oplus q_3$ into the third ancilla
(two CNOTs each; the second ancilla is not needed and stays in $\ket0$, it is kept here only to follow the original notes). On the code basis states,
\begin{align}
  S\ket{000}\ket{000}&=\ket{000}\ket{000},& S\ket{111}\ket{000}&=\ket{111}\ket{000},\nonumber\\
  S\ket{100}\ket{000}&=\ket{100}\ket{100},& S\ket{011}\ket{000}&=\ket{011}\ket{100},\nonumber\\
  S\ket{010}\ket{000}&=\ket{010}\ket{101},& S\ket{101}\ket{000}&=\ket{101}\ket{101},\nonumber\\
  S\ket{001}\ket{000}&=\ket{001}\ket{001},& S\ket{110}\ket{000}&=\ket{110}\ket{001}.
\end{align}
The syndrome depends only on \emph{which qubit has been flipped}, not on the encoded state: both $\ket{000}$ and $\ket{111}$ give $000$, both $\ket{010}$ and $\ket{101}$ give $101$, and so on. Applied to the erroneous state,
\begin{equation}
  SE_2\ket\psi\ket{000}=\left[a\ket{000}+b\ket{111}\right]\ket{000}-\ii\epsilon_2\left[a\ket{010}+b\ket{101}\right]\ket{101}.
\end{equation}
The error is now entangled with the ancillas, while the coefficients $a$ and $b$ appear identically in both terms.

\subsection{Measurement and correction}
We now measure the ancilla qubits.
\begin{itemize}
  \item With probability $\cos^2\epsilon_2\approx1$, the result is $000$. The state collapses onto $a\ket{000}+b\ket{111}$, and nothing needs to be done.
  \item With probability $\sin^2\epsilon_2\approx\epsilon_2^2$, the result is $101$. The state collapses onto $a\ket{010}+b\ket{101}$, and applying $\sigma_x$ to the second qubit restores $a\ket{000}+b\ket{111}$.
\end{itemize}
In the same way, the syndromes $100$ and $001$ indicate a flip of the first and of the third qubit, which is corrected by applying $\sigma_x$ to that qubit. In all cases the original state is recovered \emph{exactly}.

This simple example shows how the three obstacles are overcome. Redundancy is obtained by entanglement, not by copying. The measurement of the syndrome reveals only which error occurred, never the values of $a$
and $b$, so the superposition is preserved. Finally, the measurement \emph{discretizes} the error: a continuous rotation by an arbitrary small angle is projected either onto ``no error'' or onto a full bit flip,
both of which can be handled. Correcting a discrete set of errors is therefore enough to correct a continuum of them.

\subsection{Stabilizer description}
The parity measurements are measurements of the commuting observables $Z_1Z_2$ and $Z_2Z_3$ (eigenvalue $+1$ for even parity, $-1$ for odd). The code space is the joint $+1$ eigenspace of these two operators,
which are called the \emph{stabilizers} of the code. Each single bit flip $X_j$ anticommutes with a different subset of the stabilizers and therefore produces a distinct pattern of $-1$ outcomes:
$X_1\to(-,+)$, $X_2\to(-,-)$, $X_3\to(+,-)$. Measuring the stabilizers does not disturb the encoded information because the logical operators $\bar X=X_1X_2X_3$ and $\bar Z=Z_1$ commute with them.

\subsection{Limitations}
The three-qubit code corrects a flip of a single qubit. If two qubits are flipped, the syndrome points to the wrong qubit and the correction produces a logical error. For independent errors with probability $p$ per qubit,
the failure probability is reduced from $p$ to $3p^2-2p^3\approx3p^2$, an improvement when $p<1/2$. The code does not protect against phase errors ($\sigma_z$): $Z_j$ commutes with the stabilizers and acts on the code
space as the logical $\bar Z$, flipping the relative sign of $a$ and $b$ undetected.

\section{The phase-flip code}
Phase flips become bit flips in the Hadamard basis $\ket\pm=(\ket0\pm\ket1)/\sqrt2$, since $HZH=X$. The phase-flip code encodes $\ket0\to\ket{{+}{+}{+}}$, $\ket1\to\ket{{-}{-}{-}}$ (the bit-flip encoder followed by $H^{\otimes3}$),
and its stabilizers are $X_1X_2$ and $X_2X_3$. It corrects one phase flip but no bit flips.

\paragraph{Experimental demonstration.} Quantum error correction was first demonstrated in 1998 with nuclear spins in liquid-state NMR \cite{Cory1998}. That experiment implemented the three-qubit code for phase errors
on the carbon and hydrogen nuclei of alanine or trichloroethylene molecules. The fidelity of the stored qubit, as a function of the delay during which it was exposed to natural decoherence, decayed with error correction
as $1-\order(t^2)$ instead of $1-\order(t)$, the signature of first-order protection.

\section{Shor's nine-qubit code}
Concatenating the two codes, Shor obtained the first code able to correct an arbitrary error on any single qubit \cite{Shor1995}:
\begin{gather}
  \ket{0_L}=\frac{(\ket{000}+\ket{111})(\ket{000}+\ket{111})(\ket{000}+\ket{111})}{2\sqrt2},\\
  \ket{1_L}=\frac{(\ket{000}-\ket{111})(\ket{000}-\ket{111})(\ket{000}-\ket{111})}{2\sqrt2}.
\end{gather}
Each block of three is a bit-flip code (stabilizers $Z_1Z_2$, $Z_2Z_3$, $Z_4Z_5$, \dots), protecting against one $X$ error per block. A $Z$ error on any qubit of a block flips the sign of that block; the signs of the three
blocks form a phase-flip code, with stabilizers $X_1X_2X_3X_4X_5X_6$ and $X_4\cdots X_9$.
\begin{theorem}[Discretization of errors]
If a code corrects the errors $X_j$, $Z_j$ and $X_jZ_j\propto Y_j$ on qubit $j$, it corrects \emph{any} error $E_j$ acting on qubit $j$ (any unitary, or any Kraus operator of a channel).
\end{theorem}
\begin{proof}
Expand $E_j=e_0I+e_1X_j+e_2Y_j+e_3Z_j$ on the Pauli basis. The syndrome measurement projects $E_j\ket\psi$ onto one of the four terms, each of which is then corrected; since the recovery is linear, it corrects the sum.
For a channel, apply this to each Kraus operator.
\end{proof}

\section{General theory\starred}
\paragraph{Knill--Laflamme conditions.} A code with projector $P$ onto the code space can correct a set of errors $\{E_a\}$ if and only if
\begin{equation}
  PE_a^\dagger E_bP=c_{ab}P
\end{equation}
for a Hermitian matrix $c_{ab}$ independent of the code state. Intuitively: different errors either map the code space to orthogonal subspaces (they can be distinguished) or act identically on it (they need not be),
and no error reveals information about the encoded state.

\paragraph{Stabilizer codes.} An $[[n,k,d]]$ stabilizer code encodes $k$ logical qubits in $n$ physical qubits as the common $+1$ eigenspace of $n-k$ commuting Pauli operators (the stabilizer generators). It
corrects errors on up to $\lfloor(d-1)/2\rfloor$ qubits, where the \emph{distance} $d$ is the minimal weight of a logical operator. Shor's code is $[[9,1,3]]$, Steane's code (built from the classical Hamming code)
is $[[7,1,3]]$ and allows transversal Clifford gates, and the five-qubit code $[[5,1,3]]$ is the smallest code that corrects an arbitrary single-qubit error: by the quantum Hamming bound, the $1+3n$ possible errors (no error,
or one of 3 Paulis on one of $n$ qubits) need $2(1+3n)\le2^n$ dimensions, which requires $n\ge5$.

\paragraph{Fault tolerance and the threshold theorem.} Error correction is itself performed with noisy gates. Fault-tolerant circuits are designed so that a single fault causes at most one error per code block
(e.g.\ by transversal gates, which act qubit by qubit between blocks). The \emph{threshold theorem} (Aharonov, Ben-Or; Kitaev; Knill, Laflamme, Zurek, 1996--97) states that if the physical error rate $p$ is below a
threshold $p_{\rm th}$, arbitrarily long computations can be performed with an overhead polylogarithmic in the size of the computation. For concatenated codes, each level of encoding maps $p\to cp^2$, so that after $k$ levels
$p_L=p_{\rm th}(p/p_{\rm th})^{2^k}$ with $p_{\rm th}=1/c$: the logical error decreases doubly exponentially.

\paragraph{The surface code.} The most promising code for hardware with nearest-neighbour connectivity is Kitaev's surface code: qubits on a $d\times d$ grid, with 4-body $X$ and $Z$ stabilizers on alternating plaquettes,
measured with ancilla qubits. Its threshold is about $1\%$, and its logical error rate scales as
\begin{equation}
  p_L\approx A\left(\frac p{p_{\rm th}}\right)^{(d+1)/2},
\end{equation}
with about $2d^2$ physical qubits per logical qubit. In 2024, Google demonstrated a surface-code memory below threshold, with the logical error rate decreasing by a factor $\Lambda\approx2$ each time the distance was
increased from $3$ to $5$ to $7$. For $p/p_{\rm th}\approx0.1$ and $p_L\sim10^{-12}$, $d\approx21$--$23$, i.e.\ about $10^3$ physical qubits per logical qubit. Clifford gates are relatively cheap in the surface code; the $T$ gate
requires \emph{magic-state distillation}, which dominates the overhead of fault-tolerant algorithms.

\section{Exercises}
\begin{exercise}
Draw the syndrome-extraction circuit of the bit-flip code with two ancillas. Show that a bit flip on an ancilla during the extraction leads to a wrong correction, and explain why syndromes are measured repeatedly in practice.
\end{exercise}
\begin{exercise}
Show that the error $Z_1Z_2$ is not detected by the Shor code and acts trivially on the code space. (Such errors are harmless: the Shor code is \emph{degenerate}.)
\end{exercise}
\begin{exercise}
With $p=10^{-3}$, $p_{\rm th}=10^{-2}$ and $A=0.1$, what surface-code distance gives $p_L<10^{-12}$? How many physical qubits are needed for $1000$ logical qubits?
(Answer: $(d+1)/2\ge11$, $d=21$; about $2\times21^2\times1000\approx9\times10^5$.)
\end{exercise}


\part{Optimization, annealing and learning}
\chapter[Annealing and adiabatic computation]{Optimization, annealing and adiabatic quantum computation}\label{ch:annealing}

Many practical problems can be formulated as the minimization of a cost function: fitting a model to data, scheduling, routing, or finding the ground state of a physical system. When the variables are discrete,
the number of possible configurations grows exponentially with the number of variables, and many such problems are NP-hard. In this chapter we show how an optimization problem can be cast in a form suitable for
a special kind of quantum computer, the \emph{quantum annealer}, how quantum annealing compares with its classical counterpart, simulated annealing, and how both relate to \emph{adiabatic quantum computation}.

\section{Linear regression as a binary optimization problem}
\subsection{Classical linear regression}
Consider a dataset of $N$ entries, each described by $d$ features $x_1,\dots,x_d$, together with $N$ labels $y$. Linear regression looks for the weights $w_0,\dots,w_d$ of the linear model
\begin{equation}
  y=w_0+w_1x_1+w_2x_2+\dots+w_dx_d .
\end{equation}
Adding a constant feature $x_0=1$ for the bias $w_0$, the dataset becomes an $N\times D$ matrix $X$ with $D=d+1$ columns, and the model reads $y=Xw$. The standard classical method, \emph{ordinary least squares},
minimizes the squared error
\begin{equation}
  E(w)=\norm{Xw-y}^2=w^TX^TXw-2w^TX^Ty+y^Ty.
\end{equation}
The minimum satisfies $\nabla E=2X^T(Xw-y)=0$ (the \emph{normal equations}), which gives the closed-form solution
\begin{equation}
  w=(X^TX)^{-1}X^Ty.
\end{equation}
The bottleneck is the inversion of the $D\times D$ matrix $X^TX$, which costs $\order(D^3)$ operations (forming $X^TX$ itself costs $\order(ND^2)$). For large problems, approximate iterative methods are used instead,
such as gradient descent, $w\to w-\alpha\nabla E(w)$, where $\alpha$ is the learning rate. Since $E$ is a convex quadratic function, gradient descent converges to the global minimum for small enough $\alpha$.

\subsection{The QUBO formulation}
Quantum annealers solve a specific class of problems, \emph{quadratic unconstrained binary optimization} (QUBO):
\begin{equation}\label{eq:qubo}
  \min_{z\in\{0,1\}^M}z^TAz+z^Tb,
\end{equation}
where $z$ is a vector of $M$ binary variables, $A\in\R^{M\times M}$ is the QUBO matrix and $b\in\R^M$ the QUBO vector. (Since $z_i^2=z_i$, the linear term can be absorbed in the diagonal of $A$.)
To bring linear regression into this form \cite{DatePotok2021}, note that $E(w)$ is already quadratic, but the weights are real numbers, not binary variables. We therefore \emph{binarize} them: each weight is
approximated by a combination of $K$ binary variables $\hat w_{ik}\in\{0,1\}$,
\begin{equation}
  w_i=\sum_{k=1}^Kp_k\hat w_{ik},
\end{equation}
where the precision vector $p=(p_1,\dots,p_K)^T$ fixes the values that can be represented. For example, $p=(0.25,0.5)^T$ allows $w_i\in\{0,0.25,0.5,0.75\}$; negative values are obtained with negative entries of $p$
(e.g.\ $p=(-1,0.5,0.25)$). Collecting all binary variables into a single vector of length $DK$, $\hat w=(\hat w_{11},\dots,\hat w_{1K},\dots,\hat w_{D1},\dots,\hat w_{DK})^T$, and defining the $D\times DK$ precision
matrix $P=I_D\otimes p^T$, we have $w=P\hat w$. Substituting and discarding the constant $y^Ty$,
\begin{equation}
  \min_{\hat w\in\{0,1\}^{DK}}E(\hat w)=\hat w^TP^TX^TXP\hat w-2\hat w^TP^TX^Ty.
\end{equation}
This is exactly the QUBO problem \cref{eq:qubo} with $z=\hat w$, $A=P^TX^TXP$, $b=-2P^TX^Ty$. The number of binary variables, $DK$, does not depend on the number of data points $N$: the dataset enters only through
$X^TX$ and $X^Ty$.

\paragraph{Performance.} Date and Potok \cite{DatePotok2021} implemented this approach on a D-Wave 2000Q quantum annealer and compared it with the Scikit-learn library on an 8-core Intel i9 workstation. They obtained a
speed-up of up to 2.8 for the largest dataset, with about $16.8\times10^6$ data points, and comparable regression errors. The size of the problem was limited by the hardware: the D-Wave 2000Q can handle at most 64
binary variables with all-to-all connectivity, so with $K=2$ bits per weight the number of features was limited to $D\le32$. Such results depend strongly on the specific hardware and software used, and should not be
read as a general quantum advantage: the regression problem is convex and is solved exactly in polynomial time classically.

\section{Simulated annealing}
QUBO problems are hard in general, and heuristic methods are widely used. One of the most successful is \emph{simulated annealing} (SA) \cite{Kirkpatrick1983}, based on the Metropolis algorithm of statistical physics.
It mimics the slow cooling (annealing) of a material, which allows the atoms to settle into a low-energy crystal instead of a disordered glass.

To minimize a function $f$ of $n$ variables, SA performs a random walk in configuration space, $x_{\rm init}\to x_1\to x_2\to\cdots$. At each step a random change of the configuration is proposed (e.g.\ flipping one bit).
Let $\Delta=f(x_{i+1})-f(x_i)$ be the resulting change of the cost function. The move is accepted with the \emph{Metropolis probability}
\begin{equation}\label{eq:metropolis}
  \mathrm{Prob(accept)}=\min\left(1,\e^{-\Delta/T}\right).
\end{equation}
A move that lowers the cost ($\Delta<0$) is always accepted. A move that raises it is accepted with probability $\e^{-\Delta/T}<1$, which allows the walker to climb out of local minima (\cref{fig:annealing}).

\begin{origin}[: why the Metropolis rule?]
The rule is designed so that the random walk samples the Boltzmann distribution $\pi(x)\propto\e^{-f(x)/T}$ of a physical system with energy $f$ at temperature $T$. A Markov chain with transition probabilities $W(x\to y)$
has $\pi$ as stationary distribution if it satisfies \emph{detailed balance}, $\pi(x)W(x\to y)=\pi(y)W(y\to x)$. With symmetric proposals, \cref{eq:metropolis} gives
$W(x\to y)/W(y\to x)=\e^{-(f(y)-f(x))/T}=\pi(y)/\pi(x)$. As $T\to0$, $\pi$ concentrates on the global minima. The difficulty is that at low $T$ the walk becomes trapped behind barriers, for a time
$\sim\e^{\Delta_{\rm barrier}/T}$ (Arrhenius law).
\end{origin}
The parameter $T$ plays the role of a temperature. It is slowly decreased during the search, so that the walker first explores the landscape broadly and then settles into a deep minimum. If the temperature is lowered slowly
enough, for instance as $T(i)=a/\ln i$ with a sufficiently large constant $a$ (of the order of the largest barrier), the algorithm is guaranteed to find the global minimum with probability approaching one \cite{Geman1984}.
Such a slow schedule is impractical, however, and faster schedules are used in practice, without guarantee.

\begin{figure}[h]\centering
\begin{tikzpicture}[xscale=1.2]
  \draw[->] (0,0)--(8.3,0) node[right]{$x$};
  \draw[->] (0,0)--(0,3.4) node[above]{$f(x)$};
  \draw[thick,blue] plot[smooth] coordinates {(0.2,2.8) (1.2,1.3) (2.0,1.0) (2.8,1.4) (3.6,2.6) (4.3,2.4) (5.1,0.9) (6.2,0.3) (7.2,1.2) (8,2.7)};
  \fill[red] (2.0,1.0) circle (2pt);
  \draw[->,red,thick] (2.05,1.1) .. controls (3,3.6) and (4,3.3) .. (5.0,1.1);
  \node[red,align=center,font=\small] at (3.3,3.2) {SA: over the barrier, prob.\ $\e^{-\Delta/T}$};
  \draw[->,green!50!black,thick,dashed] (2.2,1.05)--(5.0,1.05);
  \node[green!50!black,font=\small,align=center] at (3.6,0.55) {QA: tunnelling\\through the barrier};
  \node[font=\small] at (2.0,0.7) {local min.};
  \node[font=\small] at (6.3,0.05) {global min.};
\end{tikzpicture}
\caption{Escaping from a local minimum. In simulated annealing the walker must climb over the barrier by thermal activation. In quantum annealing it can also tunnel through it.}\label{fig:annealing}
\end{figure}

\section{Quantum annealing}
In quantum annealing (QA) \cite{Kadowaki1998}, the thermal fluctuations of simulated annealing are replaced by \emph{quantum fluctuations}. The cost function is encoded in a Hamiltonian $H_F$, the final Hamiltonian, whose
ground state is the solution. A second, ``disordering'' Hamiltonian $H_D$ is added, with a time-dependent strength $\Gamma(t)$:
\begin{equation}
  H(t)=H_F+\Gamma(t)H_D .
\end{equation}
The role of the temperature $T$ is now played by $\Gamma$, which is slowly decreased to zero. The essential requirement is that the two Hamiltonians do not commute, $[H_D,H_F]\neq0$, just as the position $x$ and momentum $p$
of a particle do not commute. The disordering term plays the role of a kinetic energy. It makes the system move in configuration space, and allows it to tunnel through energy barriers instead of climbing over them.
In practice $H_D$ is usually a transverse field, $H_D=-\sum_i\sigma^x_i$, which does not commute with the $\sigma^z$ operators of the cost Hamiltonian. Its ground state, the equal superposition of all configurations, is
the starting point of the evolution.

\paragraph{Tunnelling vs thermal activation.} For a barrier of height $\Delta$ and width $w$, thermal activation has a rate $\propto\e^{-\Delta/T}$, independent of $w$. Quantum tunnelling of a particle of mass $m$ has the
WKB rate $\propto\exp(-\frac{2}{\hbar}\int\sqrt{2m(V-E)}\dd x)\sim\exp(-w\sqrt{2m\Delta}/\hbar)$, which depends on $\sqrt\Delta$ and on $w$. Tunnelling therefore wins for \emph{tall but narrow} barriers, and loses
for wide ones. In a spin system, the ``width'' is the number of spins that must flip together (Hamming distance), and the transverse field must couple configurations through many orders of perturbation theory.

Convergence can also be guaranteed for quantum annealing. If the transverse field decreases as $\Gamma(i)\ge b/(i+1)^{c/n}$, where $b$ and $c$ are constants and $n$ is the number of variables, the system reaches the ground
state with probability approaching one \cite{MoritaNishimori2008}. This power-law schedule is much faster than the logarithmic schedule required by simulated annealing. In practice, however, little is known about the best
choice of $H_D$ or of the schedule $\Gamma(t)$, and whether quantum annealing offers a general speed-up over the best classical methods remains an open question.

\section{The Ising model}
\begin{origin}[: the Ising model]
The Ising model was proposed by W.~Lenz (1920) and solved in one dimension by his student E.~Ising (1925) as a model of ferromagnetism. Each site carries a magnetic moment that can point up or down, $s_i=\pm1$.
Neighbouring moments interact through the exchange interaction, a consequence of the Pauli principle and the Coulomb repulsion between electrons (\cref{sec:spinqubits}), which favours parallel ($J<0$ in the convention below) or
antiparallel ($J>0$) alignment; a magnetic field acts on each moment. In a strongly anisotropic magnet only the $z$ components matter, and the Heisenberg interaction $J\vec S_i\cdot\vec S_j$ reduces to $JS_i^zS_j^z$.
With random couplings of both signs, the model describes \emph{spin glasses}, whose rugged energy landscape, with exponentially many local minima, is the archetype of hard optimization problems.
\end{origin}
The cost Hamiltonian of a quantum annealer is that of an Ising model, a set of spins $s_i=\pm1$ with energy
\begin{equation}
  H=\sum_ih_is_i+\sum_{i<j}J_{ij}s_is_j,
\end{equation}
where $h_i$ are local fields and $J_{ij}$ are couplings ($J_{ij}<0$ favours parallel spins, i.e.\ ferromagnetic order). The task is to find the ground state, the spin configuration of lowest energy. For spins on a
$d$-dimensional grid this is easy in special cases (e.g.\ planar graphs without fields can be solved in polynomial time) but NP-hard in general, already for a two-dimensional grid with non-zero fields $h_i$ \cite{Barahona1982}.
On a quantum annealer, the spins are qubits and the Hamiltonian is
\begin{equation}
  H_{\rm Ising}=\sum_ih_i\sigma_i^z+\sum_{i<j}J_{ij}\sigma_i^z\sigma_j^z .
\end{equation}
It is diagonal in the computational basis: each basis state $\ket{s_1\dots s_n}$ is an eigenstate with eigenvalue $H(s)$.

\section{From optimization problems to the Ising model}
Many optimization problems can be mapped onto the Ising model \cite{Lucas2014}. The most direct is QUBO: given weights $Q_{ij}$, assign binary values $x_i\in\{0,1\}$ to minimize
\begin{equation}
  f_{\rm QUBO}(x)=\sum_{i\le j}Q_{ij}x_ix_j .
\end{equation}
The change of variables $s_i=1-2x_i$, i.e.\ $x_i=(1-s_i)/2$, transforms this into an Ising energy. For the diagonal terms, $x_i^2=x_i$ and $Q_{ii}x_i=\frac{Q_{ii}}2-\frac{Q_{ii}}2s_i$. For $i<j$,
\begin{equation}
  Q_{ij}x_ix_j=\frac{Q_{ij}}4\left(1-s_i-s_j+s_is_j\right).
\end{equation}
Collecting the terms,
\begin{equation}
  J_{ij}=\frac{Q_{ij}}4,\qquad h_i=-\frac{Q_{ii}}2-\frac14\sum_{j\neq i}Q_{ij},
\end{equation}
up to an irrelevant constant, where $Q_{ij}$ with $j<i$ stands for $Q_{ji}$. The linear regression problem of the previous section can thus be solved on a quantum annealer.

\paragraph{MAX-2-SAT.} Another example is weighted MAX-2-SAT, a satisfiability problem. We are given $m$ clauses, each containing two Boolean variables or their negations, for example $C_i=(l_{i1}\vee l_{i2})$, and each with a
weight $w_i$. The goal is to find the assignment that maximizes the total weight of the satisfied clauses. A clause is violated only if both of its literals are false. Writing the literals as binary variables (a negated variable
$\bar x$ is $1-x$), its penalty $w_i(1-l_{i1})(1-l_{i2})$ is quadratic, so the problem is again a QUBO, and with $s_i=1-2x_i$ an Ising problem.

\paragraph{Constraints.} Constraints are added as quadratic penalties. For example, ``exactly one of $x_1,\dots,x_k$ is 1'' becomes the penalty $\lambda(1-\sum_ix_i)^2$, with $\lambda$ larger than the scale of the cost function
(this is how the travelling salesman problem is mapped onto QUBO, with $n^2$ variables $x_{c,t}$ = ``city $c$ visited at time $t$'').

\section{Quantum annealers in practice}
The best-known quantum annealers are built by D-Wave with superconducting flux qubits (\cref{sec:transmon}), in which the two directions of a persistent current play the role of $s=\pm1$, and the transverse field is a tunnelling
between them. The D-Wave Advantage system has more than 5000 qubits, each coupled to up to 15 others, with more than 35\,000 couplers in total. Since the hardware graph is not fully connected, a general problem must first be
\emph{embedded}: a logical variable that interacts with many others is represented by a chain of physical qubits, strongly coupled ferromagnetically so that they behave as one. A fully connected problem of $n$ variables needs
$\order(n^2)$ physical qubits. This limits the size of problems with dense interactions, as in the linear regression example above. The annealer operates at about $15$\,mK, but the annealing time ($\mu$s) is long compared with
$T_2$ (ns), so it is an open quantum system at finite temperature, not a coherent adiabatic computer.

As an example of a realistic application, quantum annealing has been applied to optimization problems for constellations of low-Earth-orbit satellites. The problem is written as a QUBO, converted into an Ising model and solved
on the D-Wave Advantage, and the results are compared with simulated annealing on a classical laptop processor. A notable observation in comparisons of this kind is that quantum annealing finds many local valleys of the cost
function that are missed by simulated annealing \cite{Koshka2019}. These are not the global optimum, but they can be useful near-optimal solutions, or useful samples for applications such as machine learning.

\section{Adiabatic quantum computation}
\subsection{Optimization problems}
Adiabatic quantum computation (AQC) \cite{Farhi2001} is designed to solve optimization problems of the type introduced above. We are given an objective function $f:D^n\to\R$ of $n$ variables taking values in a discrete domain $D$,
and we must find the configuration $\vec x$ that minimizes $f(\vec x)$. We consider only unconstrained problems. Two examples illustrate how general this formulation is.

\paragraph{Boolean satisfiability (SAT).} SAT was the first problem proven to be NP-complete \cite{Cook1971}. A Boolean formula of $n$ variables is given as a conjunction of $m$ clauses, $B(x)=\bigwedge_{c=1}^mC_c$, and we look for
an assignment $x$ for which $B(x)$ is true. To turn this into an optimization problem, we associate with each clause a function $f_c(x)$ with $f_c(x)=0$ if $C_c$ is satisfied by $x$, and $f_c(x)\ge1$ otherwise, and minimize
\begin{equation}
  f_{\rm CNF}(x)=\sum_{c=1}^mf_c(x).
\end{equation}
The logical product $\wedge$ has become a sum because each $f_c$ measures the violation of its clause. The formula is satisfiable if and only if the minimum of $f_{\rm CNF}$ is zero.

\paragraph{Factoring.} Given an integer $N$, find integers $x,y>1$ that minimize $f_N(x,y)=(N-xy)^2$. The minimum is zero iff $N=xy$. (Writing $x,y$ in binary gives a polynomial of degree 4 in the bits, which can be reduced to
a quadratic one with auxiliary variables. This approach has no known speed-up, unlike Shor's algorithm.)

\subsection{The adiabatic algorithm}
An AQC algorithm uses a register of $n$ qubits in a state $\ket{\phi_t}$ that evolves under a time-dependent Hamiltonian $H(t)$ with the following properties.
\begin{enumerate}
  \item The initial Hamiltonian $H_I$ has a ground state that is easy to find and to prepare.
  \item The final, or problem, Hamiltonian $H_F$ is such that its ground state, $H_F\ket{S_{\rm ground}}=E_F\ket{S_{\rm ground}}$ with the minimal eigenvalue, is the basis state $\ket x$ that minimizes $f$.
  \item An evolution path $s(t)$ goes from 1 to 0 as $t$ goes from 0 to the final time $t_f$, for example linearly, $s(t)=1-t/t_f$, with
  \begin{equation}
    H(t)=s(t)H_I+\left(1-s(t)\right)H_F .
  \end{equation}
\end{enumerate}
The system is prepared in the ground state of $H_I$, and the Hamiltonian is slowly changed into $H_F$. If the change is slow enough, the system remains in the instantaneous ground state, and at $t=t_f$ it is in the ground state of $H_F$.
A measurement in the computational basis then gives the solution.

The linear interpolation is not necessarily the best choice. A curved path in the space of Hamiltonians, or a schedule $s(t)$ that slows down where needed, may perform much better. The initial Hamiltonian should preferably be
``orthogonal'' to $H_F$, i.e.\ not commute with it. If $[H_I,H_F]=0$, the two Hamiltonians share their eigenstates; the system simply stays in the ground state of $H_I$, which is not the solution, while the levels cross.
This is the same requirement as for the disordering Hamiltonian of quantum annealing.

\subsection{Example: the exact cover problem}
A family $C=\{U_\alpha,\alpha\in A\}$ of subsets $U_\alpha\subset X$ is a \emph{cover} of the set $X$ if $X\subseteq\bigcup_{\alpha\in A}U_\alpha$, i.e.\ if each element of $X$ belongs to some $U_\alpha$. It is an \emph{exact cover} if each
element of $X$ belongs to exactly one $U_\alpha$. Determining whether an exact cover exists is NP-complete, although checking a proposed solution takes only polynomial time. In the variant \emph{exact cover by 3-sets} (X3C), each $U_\alpha$
contains exactly three elements.

Farhi et al.\ \cite{Farhi2001} used the following formulation, as a satisfiability problem $B(x)=\bigwedge_{c=1}^mC_c$ of $n$ bits $x_i\in\{0,1\}$. Each of the $m$ clauses involves three bits $x_{c_1},x_{c_2},x_{c_3}$ and is satisfied if
exactly one of them is equal to 1. (In the set language, $x_i=1$ if the subset $U_i$ is chosen, and each element of $X$ gives a clause involving the subsets that contain it.) A suitable cost function for clause $c$ is
\begin{equation}
  f_c(x)=\left(1-x_{c_1}-x_{c_2}-x_{c_3}\right)^2,
\end{equation}
which vanishes if exactly one bit is 1, and is otherwise equal to 1 (no bit or two bits equal to 1) or 4 (three bits). The total cost $f_{\rm EC}(x)=\sum_cf_c(x)$ is non-negative, and equal to 0 when all clauses are satisfied.
Expanding the square, $(1-\sum_kx_{c_k})^2=1-2\sum_kx_{c_k}+\sum_kx_{c_k}^2+2\sum_{k<l}x_{c_k}x_{c_l}$, and using $x_i^2=x_i$,
\begin{equation}
  f_{\rm EC}(x)=m-\sum_iB_ix_i+2\sum_{i<j}C_{ij}x_ix_j,
\end{equation}
where $B_i$ is the number of clauses that contain the bit $x_i$, and $C_{ij}$ is the number of clauses that contain the pair $x_i,x_j$.

\paragraph{Final Hamiltonian.} The binary variables are mapped onto spins, $x_i=(1-s_i)/2$, $s_i\in\{-1,+1\}$. Since each clause contains three bits and three pairs, $\sum_iB_i=\sum_{i<j}C_{ij}=3m$; and since each clause containing
$x_i$ contains two pairs involving $x_i$, $\sum_{j\neq i}C_{ij}=2B_i$. Substituting,
\begin{align}
  f_{\rm EC}&=m-\sum_iB_i\frac{1-s_i}2+2\sum_{i<j}C_{ij}\frac{(1-s_i)(1-s_j)}4\nonumber\\
  &=m-\frac{3m}2+\frac{3m}2+\sum_i\frac{B_i}2s_i-\frac12\sum_i\Big(\sum_{j\neq i}C_{ij}\Big)s_i+\frac12\sum_{i<j}C_{ij}s_is_j\nonumber\\
  &=m-\frac12\sum_iB_is_i+\frac12\sum_{i<j}C_{ij}s_is_j,
\end{align}
an Ising energy. Replacing the spins by Pauli operators gives the problem Hamiltonian
\begin{equation}
  H_{\rm EC}=m-\frac12\sum_iB_i\sigma_i^z+\frac12\sum_{i<j}C_{ij}\sigma_i^z\sigma_j^z,
\end{equation}
which is diagonal in the computational basis, with eigenvalue $f_{\rm EC}(x)$ on the basis state $\ket x$. Recall that a term such as $\sigma_i^z\sigma_j^z$ means $I\otimes\dots\otimes\sigma_i^z\otimes\dots\otimes\sigma_j^z\otimes\dots\otimes I$,
an operator acting in the full $2^n$-dimensional space.

\paragraph{Initial Hamiltonian.} A convenient choice is
\begin{equation}
  H_I=\sum_{i=1}^n\frac{1-\sigma_i^x}2,
\end{equation}
whose ground state (with energy 0) has all qubits in the eigenstate of $\sigma^x$ with eigenvalue $+1$, $\ket{\phi_i}=(\ket0+\ket1)/\sqrt2$, i.e.\ the equal superposition
\begin{equation}
  \ket\phi=\ket{\phi_1}\ket{\phi_2}\cdots\ket{\phi_n}=\frac1{2^{n/2}}\sum_{x=0}^{2^n-1}\ket x.
\end{equation}
For two qubits, $\ket\phi=\frac12(\ket{00}+\ket{01}+\ket{10}+\ket{11})$. The path is $H(t)=s(t)H_I+(1-s(t))H_{\rm EC}$ with $s(t)=1-t/t_f$.

\subsection{The adiabatic theorem and the running time}
\begin{origin}[: where the adiabatic condition comes from]
Let $\ket{n(t)}$ and $E_n(t)$ be the instantaneous eigenstates and eigenvalues of $H(t)$. Expand $\ket{\psi(t)}=\sum_nc_n(t)\e^{-\ii\theta_n(t)}\ket{n(t)}$ with the dynamical phases $\theta_n=\frac1\hbar\int^tE_n\dd t'$. Inserting
into the Schrödinger equation and projecting on $\bra m$ gives
\begin{equation}
  \dot c_m=-\sum_nc_n\braket m{\dot n}\e^{\ii(\theta_m-\theta_n)},\qquad\braket m{\dot n}=\frac{\mel m{\dot H}n}{E_n-E_m}\quad(m\neq n),
\end{equation}
where the second relation follows from differentiating $H\ket n=E_n\ket n$. Starting in the ground state, $c_0\approx1$, the amplitude of an excited state $m$ is driven by a term $\mel m{\dot H}0/(E_0-E_m)$ multiplied by a phase
oscillating at the frequency $(E_m-E_0)/\hbar$. Integrating by parts, the oscillation averages this term out, leaving $|c_m|\sim\hbar|\mel m{\dot H}0|/(E_m-E_0)^2$. With $H=H(s)$ and $\dot H=\dot s\,\dd H/\dd s\sim(\dd H/\dd s)/t_f$,
the excitation is small if $t_f\gg\hbar\norm{\dd H/\dd s}/\delta^2$, where $\delta$ is the gap. (The term $m=n$ gives the Berry phase.)
\end{origin}
The adiabatic theorem guarantees that the system remains in the instantaneous ground state if three conditions hold:
\begin{enumerate}
  \item the system is initially in the ground state of $H(0)$;
  \item the spectral gap $\delta(s)$ between the ground state and the first excited state remains positive along the whole path, with minimum $\delta_m>0$;
  \item the evolution is slow enough:
  \begin{equation}\label{eq:adiabatic}
    t_f\gg\frac{\hbar\max_s\norm{\frac{\dd}{\dd s}H(s)}}{\delta_m^2}.
  \end{equation}
\end{enumerate}
Here $\norm A=\sup_x\norm{Ax}/\norm x$ is the matrix norm. For the linear path, $\dd H/\dd s=H_I-H_F$, and $\norm{H_I-H_F}\le\norm{H_I}+\norm{H_F}$ is bounded by the largest eigenvalues involved. For exact cover, the spectrum of $H_I$ lies in
$[0,n]$ and that of $H_{\rm EC}$ in $[0,4m]$. The number of clauses is at most the number of triplets of variables, $m\lesssim n^3$, so $\norm{\dd H/\dd s}$ is bounded by a polynomial in $n$. The running time of the algorithm is therefore
polynomial in $n$ if and only if the minimum gap $\delta_m$ decreases at most polynomially. For example, a gap closing as $\delta_m\sim n^{-2}$ gives $t_f=\order(n^4)$ (times the polynomial bound on the norm).

\paragraph{What about the spectral gap?} Unfortunately, the gap cannot be computed in general for large systems. Determining its scaling for NP-complete problems is equivalent to deciding whether quantum computers can solve them
efficiently in this way. Physical insight comes from random instances, characterized by the ratio $r=m/n$ of clauses to variables. For random 3-SAT there is a critical value $r_c\approx4.26$: below it almost all instances are
satisfiable, above it almost all are unsatisfiable, and the hardest instances lie close to $r_c$, where even classical algorithms need a long computation time (a phase transition in computational complexity).
\begin{itemize}
  \item For $r<r_c$, instances are typically easy and solved in polynomial time.
  \item Near $r_c$, the first numerical simulations of AQC for exact cover, limited to $n\le20$ bits, suggested a running time growing only polynomially \cite{Farhi2001}. Quantum Monte Carlo studies extended this to $n\le128$, still with a
  polynomial decrease of the median gap, in a range of sizes where a classical exact algorithm already shows exponential complexity \cite{Young2008}.
  \item For larger sizes, however, the gap of typical instances becomes exponentially small \cite{HenYoung2011}. The algorithm then requires an exponential time, unless the path $s(t)$ is modified.
\end{itemize}

\paragraph{Avoided crossings and the Landau--Zener formula.} The origin of these small gaps is avoided crossings between the ground state and an excited state \cite{Altshuler2010}, often near the end of the evolution, between the solution and
a competing local minimum. Near such a crossing, the system behaves like a two-level system with Hamiltonian
\begin{equation}
  H=\frac{\alpha t}2\sigma_z+V\sigma_x,
\end{equation}
whose levels $\pm\sqrt{(\alpha t/2)^2+V^2}$ approach each other to a minimal distance $2V$. If the crossing is traversed at the rate $\alpha$, the probability to leave the ground state (i.e.\ to follow the diabatic state) is given by the
\emph{Landau--Zener formula} (Landau, Zener, Stückelberg, Majorana, 1932)
\begin{equation}
  P=\exp\left(-\frac{2\pi V^2}{\hbar|\alpha|}\right).
\end{equation}
(It is obtained from the exact solution of the two coupled amplitude equations in terms of parabolic cylinder functions.) Adiabaticity requires $\hbar|\alpha|\ll V^2$, i.e.\ a sweep rate quadratically small in the minimum gap, as in
\cref{eq:adiabatic}. An exponentially small $V$ thus requires an exponentially long time.

\subsection{Grover's search as an adiabatic algorithm}
For unstructured search, take $H_I=I-\proj\psi$ (with $\ket\psi$ the uniform superposition) and $H_F=I-\proj w$, and $H(s)=(1-s)H_I+sH_F$ with $s$ from 0 to 1. The evolution stays in the plane spanned by $\ket w$ and $\ket\psi$, and the
$2\times2$ problem gives the gap
\begin{equation}
  g(s)=\sqrt{1-4\left(1-\frac1N\right)s(1-s)},\qquad g_{\min}=g(\tfrac12)=\frac1{\sqrt N}.
\end{equation}
With the linear path, \cref{eq:adiabatic} gives $t_f=\order(1/g_{\min}^2)=\order(N)$: no better than classical search. But if the schedule is adapted locally to the gap (Roland and Cerf \cite{RolandCerf2002}), $\dd s/\dd t=\epsilon g^2(s)$,
the total time is
\begin{equation}
  t_f=\frac1\epsilon\int_0^1\frac{\dd s}{g^2(s)}=\frac1\epsilon\sqrt{\frac N{1-1/N}}\arctan\sqrt{N-1}\approx\frac\pi{2\epsilon}\sqrt N,
\end{equation}
which recovers Grover's quadratic speed-up. The evolution must be slow only near the minimum gap and can be fast elsewhere.

\subsection{Power of adiabatic quantum computation}
More generally, adiabatic quantum computation is equivalent to the gate model of quantum computation. AQC can simulate any quantum circuit, and vice versa, with at most a polynomial slowdown \cite{Aharonov2007}. Any quantum algorithm can
therefore in principle be formulated as an adiabatic evolution (the circuit-to-AQC direction uses Feynman's ``history state'' construction, with a gap that closes polynomially). The difficulty is not in the power of the model, but in finding
paths along which the spectral gap remains large. Conversely, a gate-based computer can simulate an adiabatic evolution by Trotterization (\cref{eq:trotter}); the \emph{quantum approximate optimization algorithm} (QAOA, Farhi, Goldstone,
Gutmann 2014) alternates $p$ layers $\e^{-\ii\gamma_kH_F}\e^{-\ii\beta_kH_I}$ with angles optimized classically, and becomes adiabatic for large $p$ (\cref{ch:qml}).

\section{Exercises}
\begin{exercise}
Write the QUBO for the problem ``choose exactly two of four items, minimizing the total cost $c_1x_1+\dots+c_4x_4$'' with a penalty term, and convert it to Ising form.
\end{exercise}
\begin{exercise}
Derive the gap $g(s)$ of the adiabatic Grover problem by writing $H(s)$ in the orthonormal basis $\{\ket w,\ket{w^\perp}\}$ of the plane, with $\ket{w^\perp}\propto\ket\psi-N^{-1/2}\ket w$.
\end{exercise}
\begin{exercise}
For $n=3$ bits and a single exact-cover clause on bits $1,2,3$, compute $B_i$ and $C_{ij}$, write the diagonal of $H_{\rm EC}$ (an $8\times8$ matrix) and identify its three degenerate ground states. Compute numerically the gap of $H(s)=sH_I+(1-s)H_{\rm EC}$ along the linear path. Why does it close at $s=0$?
\end{exercise}

\chapter[Quantum machine learning]{Quantum machine learning and variational algorithms}\label{ch:qml}

Machine learning consists largely of optimization and linear algebra on large datasets: training a model means minimizing a cost function, often by manipulating large matrices. Quantum machine learning (QML) asks
whether quantum computers can perform these tasks faster or better \cite{Biamonte2017}. Linear regression, discussed in \cref{ch:annealing}, is a good example, since it can be approached in three different ways.

\section{Three approaches to linear regression}
\paragraph{Optimization with quantum annealers.} As we have seen, training a linear regression model can be written as a QUBO problem by binarizing the weights, and solved on a quantum annealer \cite{DatePotok2021}.
The same strategy applies to other models whose training reduces to a quadratic cost function, such as support vector machines or clustering. The limitations are those of quantum annealing itself: the number and connectivity of
the qubits, the finite precision of the binarized weights, and the absence of a guaranteed speed-up.

\paragraph{Quantum linear algebra.} Ordinary least squares requires solving the linear system $X^TXw=X^Ty$. The HHL algorithm \cite{HHL2009} solves a linear system $A\vec x=\vec b$ in a time that grows only logarithmically with the
size of the matrix, for sparse and well-conditioned matrices. Its application to least-squares fitting \cite{Wiebe2012} promises an exponential speed-up over the classical $\order(D^3)$ inversion. The algorithm works as follows
($A$ Hermitian, with eigenvectors $\ket{u_j}$ and eigenvalues $\lambda_j$):
\begin{enumerate}
  \item Encode $\vec b$ as a quantum state $\ket b=\sum_j\beta_j\ket{u_j}$ (amplitude encoding: $\log_2N$ qubits for $N$ numbers).
  \item Phase estimation of $\e^{\ii At}$ (\cref{ch:shor}) gives $\sum_j\beta_j\ket{u_j}\ket{\lambda_j}$.
  \item A rotation of an ancilla controlled by the eigenvalue register,
  $\ket{\lambda_j}\ket0\to\ket{\lambda_j}\big(\sqrt{1-C^2/\lambda_j^2}\,\ket0+\frac C{\lambda_j}\ket1\big)$,
  then uncomputation of $\ket{\lambda_j}$.
  \item Post-selecting the ancilla in $\ket1$ leaves $\propto\sum_j\frac{\beta_j}{\lambda_j}\ket{u_j}=A^{-1}\ket b$, the normalized solution $\ket x$.
\end{enumerate}
The cost is $\tilde\order(\log N\,s^2\kappa^2/\epsilon)$ for an $s$-sparse matrix with condition number $\kappa$ and precision $\epsilon$. These promises come with important caveats \cite{Aaronson2015}. The data must first be loaded into a quantum
state, which may itself take a time proportional to their size (or require a ``quantum RAM'' that does not exist). The solution is obtained as a quantum state $\ket x$, from which only a few properties (e.g.\ $\mel x Mx$) can be read out
efficiently, not all the weights. The condition number enters polynomially. And for several of these problems, ``quantum-inspired'' classical algorithms with a comparable scaling have since been found \cite{Tang2019}, when the classical
data structure allows the same sampling access.

\paragraph{Variational quantum models.} The approach best suited to present-day noisy hardware uses \emph{parameterized quantum circuits}. Classical data are encoded into a quantum state by a \emph{feature map}, such as angle encoding
$\ket{\phi(x)}=\bigotimes_iR_y(x_i)\ket0$ or the Pauli expansion circuit \texttt{PauliFeatureMap} of Qiskit's circuit library. A circuit $U(\vec\theta)$ with adjustable rotation angles then acts on this state, and the expectation value
of a measured observable, $f_{\vec\theta}(x)=\mel{\phi(x)}{U^\dagger(\vec\theta)OU(\vec\theta)}{\phi(x)}$, serves as the prediction of the model. The angles are trained by a classical optimizer that minimizes the squared error, exactly
as in classical regression, with the quantum computer evaluating the model and the classical computer updating its parameters. Alternatively, the quantum state can be used to compute a \emph{kernel}, a similarity measure between data points in
an exponentially large Hilbert space, $k(x,x')=|\braket{\phi(x)}{\phi(x')}|^2$, which is then used by a classical algorithm such as a support vector machine \cite{Havlicek2019}.

\section{Variational algorithms}
\subsection{The variational principle and VQE}
The \emph{variational quantum eigensolver} (VQE, Peruzzo et al.\ 2014) estimates the ground-state energy of a Hamiltonian $H$ (a molecule, a spin model) by minimizing $E(\vec\theta)=\mel{\psi(\vec\theta)}H{\psi(\vec\theta)}$ over a family of
trial states prepared by a circuit. It rests on the \emph{variational principle}:
\begin{equation}
  E(\vec\theta)\ge E_0\quad\text{for all }\vec\theta,
\end{equation}
since, expanding $\ket{\psi(\vec\theta)}=\sum_nc_n\ket{n}$ on the eigenstates of $H$, $E(\vec\theta)=\sum_n|c_n|^2E_n\ge E_0\sum_n|c_n|^2=E_0$. The Hamiltonian is written as a sum of Pauli strings, $H=\sum_kh_kP_k$ (for molecules, via the
Jordan--Wigner mapping of fermions to qubits), and each $\expval{P_k}$ is measured separately, with the shot-noise cost discussed in \cref{ch:measurement}.

\subsection{QAOA}
The quantum approximate optimization algorithm uses the trial state
\begin{equation}
  \ket{\vec\gamma,\vec\beta}=\e^{-\ii\beta_pH_I}\e^{-\ii\gamma_pH_F}\cdots\e^{-\ii\beta_1H_I}\e^{-\ii\gamma_1H_F}\ket+^{\otimes n},
\end{equation}
with the Ising cost Hamiltonian $H_F$ of \cref{ch:annealing} and the mixer $H_I=\sum_iX_i$, and minimizes $\expval{H_F}$ over the $2p$ angles. It is a Trotterized, optimized version of quantum annealing. For MaxCut, $H_F=\sum_{(ij)\in E}Z_iZ_j$.

\subsection{Gradients: the parameter-shift rule}
Training requires gradients. Finite differences are very sensitive to shot noise, but for gates of the form $\e^{-\ii\theta P/2}$ with $P^2=I$ (all Pauli rotations) the gradient is given exactly by two evaluations of the circuit.
\begin{proposition}[Parameter-shift rule]
If $f(\theta)=\mel\psi{U^\dagger(\theta)AU(\theta)}\psi$ with $U(\theta)=\e^{-\ii\theta P/2}$, $P^2=I$ (other gates before and after are absorbed in $\ket\psi$ and $A$), then
\begin{equation}
  \frac{\dd f}{\dd\theta}=\frac12\left[f\left(\theta+\frac\pi2\right)-f\left(\theta-\frac\pi2\right)\right].
\end{equation}
\end{proposition}
\begin{proof}
By \cref{prop:involution}, $U(\theta)=\cos\frac\theta2-\ii\sin\frac\theta2P$, so $f$ is a quadratic form in $(\cos\frac\theta2,\sin\frac\theta2)$, i.e.\ a trigonometric polynomial of degree one in $\theta$: $f(\theta)=a+b\cos\theta+c\sin\theta$.
Then $f(\theta+\frac\pi2)-f(\theta-\frac\pi2)=-2b\sin\theta+2c\cos\theta=2f'(\theta)$.
\end{proof}

\subsection{Difficulties}
Variational algorithms face serious obstacles: (i) \emph{barren plateaus}: for random, deep, or highly expressive circuits, the gradients vanish exponentially with the number of qubits (McClean et al.\ 2018), so that an exponential number of
shots is needed to train them; (ii) noise also flattens the cost landscape; (iii) the classical optimization is non-convex with many local minima; (iv) the number of measurements needed to estimate $\expval H$ to chemical accuracy can be huge.
No practical quantum advantage has been demonstrated with variational algorithms so far.

\section{Useful subroutines}
\paragraph{The swap test.} The overlap of two states, needed for kernels and fidelity estimation, can be measured with one ancilla:
\begin{center}
\begin{quantikz}
\lstick{$\ket0$} & \gate{H} & \ctrl{1} & \gate{H} & \meter{}\\
\lstick{$\ket\psi$} & \qw & \swap{1} & \qw & \qw\\
\lstick{$\ket\phi$} & \qw & \targX{} & \qw & \qw
\end{quantikz}
\end{center}
After the controlled SWAP and the second Hadamard, the state is $\frac12\ket0(\ket{\psi\phi}+\ket{\phi\psi})+\frac12\ket1(\ket{\psi\phi}-\ket{\phi\psi})$, and
\begin{equation}
  P(0)=\frac14\norm{\ket{\psi\phi}+\ket{\phi\psi}}^2=\frac{1+|\braket\psi\phi|^2}2 .
\end{equation}
(For a kernel, it is simpler to apply $U_\phi(x')^\dagger U_\phi(x)$ to $\ket0$ and measure the probability of $\ket{0\dots0}$.)

\section{Outlook}
Quantum machine learning is a very active field, but no practical advantage over classical machine learning has been demonstrated so far. Quantum computers are most likely to help for data that are themselves quantum, such as the outputs
of quantum sensors or quantum simulations, rather than for classical datasets. There are rigorous results showing exponential advantages in learning from quantum experiments with a quantum memory (Huang et al.\ 2022).

\section{Exercises}
\begin{exercise}
Compute $f(\theta)=\mel0{R_y(\theta)^\dagger ZR_y(\theta)}0$ and check the parameter-shift rule.
\end{exercise}
\begin{exercise}
For angle encoding of a single feature, $\ket{\phi(x)}=R_y(x)\ket0$, compute the kernel $k(x,x')$. Which classical kernel is it?
(Answer: $\cos^2\frac{x-x'}2$.)
\end{exercise}


\part{Physical implementations}
\chapter{Physical implementations of qubits}\label{ch:implementations}

\section{Requirements and conflicting constraints}
To build a quantum computer, a physical system must satisfy a number of requirements, formulated by DiVincenzo \cite{DiVincenzo2000}:
\begin{enumerate}
  \item robust qubits, i.e.\ well-defined two-level systems, in a scalable architecture, that keep their quantum coherence long enough;
  \item the possibility to prepare the qubits in a known initial state (e.g.\ by cooling to the ground state, \cref{sec:thermal}, or by optical pumping);
  \item the possibility to manipulate them with single-qubit gates;
  \item a controllable coupling between qubits, to perform two-qubit gates such as the CNOT (together, 3 and 4 give a universal gate set, \cref{sec:universal});
  \item an efficient, qubit-specific measurement of the final state;
\end{enumerate}
and, for quantum communication, (6) the ability to convert stationary qubits into flying qubits (photons) and (7) to transmit them faithfully.

These requirements are partly contradictory. A robust qubit must be well isolated from its environment, so that it does not decohere. At the same time it must be accessible: it must interact strongly with the control fields, with the
other qubits and with the measuring apparatus. \emph{Every interaction channel used for control is also a channel through which noise can enter.} All physical implementations are compromises between isolation and accessibility.

A useful figure of merit is the ratio
\begin{equation}
  n_{\rm op}=\frac{\tau_Q}{\tau_{\rm op}}
\end{equation}
between the coherence time $\tau_Q$ of the qubits ($T_2$, \cref{ch:open}) and the time $\tau_{\rm op}$ needed to perform an elementary operation. It gives the number of operations that can be carried out before the quantum information
is lost, and its inverse is a rough lower bound on the error per gate. \Cref{tab:platforms} gives order-of-magnitude estimates.
\begin{table}[h]\centering
\begin{tabular}{@{}lccc@{}}\toprule
System & $\tau_Q$ (s) & $\tau_{\rm op}$ (s) & $n_{\rm op}$\\ \midrule
Nuclear spin (NMR) & $10^{-2}$--$10^{8}$ & $10^{-3}$--$10^{-6}$ & $10^{5}$--$10^{14}$\\
Optical cavity & $10^{-5}$ & $10^{-14}$ & $10^{9}$\\
Superconducting (transmon) & $10^{-4}$ & $10^{-8}$ & $10^{4}$\\
Ion trap & $10^{-1}$ & $10^{-14}$ & $10^{13}$\\
Neutral atoms (Rydberg) & $1$ & $10^{-7}$--$10^{-6}$ & $10^{6}$\\
Electron spin in Si quantum dot & $10^{-3}$--$10^{-1}$ & $10^{-8}$--$10^{-7}$ & $10^{5}$\\ \bottomrule
\end{tabular}
\caption{Coherence time $\tau_Q$, operation time $\tau_{\rm op}$ and number of operations $n_{\rm op}=\tau_Q/\tau_{\rm op}$ (orders of magnitude; first rows adapted from \cite{NielsenChuang}, with added estimates for superconducting qubits,
neutral atoms and spins). The values of $\tau_{\rm op}$ for optical systems and ion traps are idealized lower limits; practical two-qubit gates are much slower ($\sim10$--$100\,\mu$s for ions).}\label{tab:platforms}
\end{table}

\paragraph{The harmonic oscillator and the need for nonlinearity.} Many of these systems are built from harmonic oscillators: modes of the electromagnetic field, vibrational modes of trapped ions, or LC circuits. A harmonic oscillator alone is
not a good qubit. Its levels are equally spaced (\cref{sec:oscillator}), so a field that drives the transition $\ket0\to\ket1$ also drives $\ket1\to\ket2$, and the system leaves the qubit subspace. Moreover, a harmonic oscillator driven by a
classical field stays in a coherent state, which behaves classically: linear systems cannot produce the non-classical states needed for computation. The two-level truncation of the damped oscillator of \cref{ch:open} is justified only when
this leakage is suppressed. \emph{A nonlinearity (anharmonicity) is therefore needed}, both to isolate two levels and to make qubits interact. The various implementations differ mainly in how they provide it: atoms are naturally anharmonic;
superconducting circuits use the Josephson junction; photons need atoms or nonlinear crystals.

\section{Superconducting circuits}\label{sec:transmon}
The qubits are electrical circuits made of superconducting metals (aluminium, niobium, tantalum) and cooled to about $10$\,mK. They are macroscopic (micrometres to millimetres) yet quantum: superconductivity removes dissipation, and the
collective variables of the circuit (charge and flux) obey quantum mechanics.

\subsection{Quantization of an LC circuit}\label{sec:lc}
\begin{origin}[: circuit quantum mechanics]
Consider an LC circuit and choose as coordinate the node flux $\Phi(t)=\int^tV\dd t'$ (the time integral of the voltage across the capacitor). The energy stored in the capacitor, $\frac12CV^2=\frac12C\dot\Phi^2$, plays the role of a kinetic
energy, and the energy of the inductor, $\frac{\Phi^2}{2L}$ (since $\Phi=LI$), that of a potential energy. The Lagrangian
\begin{equation}
  \mathcal L=\frac C2\dot\Phi^2-\frac{\Phi^2}{2L}
\end{equation}
gives the equation of motion $C\ddot\Phi=-\Phi/L$, i.e.\ Kirchhoff's law. The momentum conjugate to $\Phi$ is $\partial\mathcal L/\partial\dot\Phi=C\dot\Phi=Q$, the charge on the capacitor, and the Hamiltonian is
\begin{equation}
  H=\frac{Q^2}{2C}+\frac{\Phi^2}{2L}.
\end{equation}
Canonical quantization, $[\hat\Phi,\hat Q]=\ii\hbar$, gives a harmonic oscillator (mass $C$, frequency $\omega_r=1/\sqrt{LC}$), with $\hat\Phi=\Phi_{\rm zpf}(\hat a+\hat a^\dagger)$, $\Phi_{\rm zpf}=\sqrt{\hbar Z/2}$, $Z=\sqrt{L/C}$. For
$\omega_r/2\pi=5$\,GHz, $\hbar\omega_r/k_B=0.24$\,K: at $10$\,mK the circuit is in its quantum ground state ($\bar n\approx10^{-11}$). Microwave resonators (coplanar waveguides) are LC oscillators of this kind.
\end{origin}

\subsection{The Josephson junction}
\begin{origin}[: the Josephson relations]
A Josephson junction is a thin insulating barrier (about 1\,nm of AlO$_x$) between two superconductors. Each superconductor is described by a macroscopic wave function of the Cooper pairs, $\psi_k=\sqrt{n_k}\e^{\ii\theta_k}$ ($k=1,2$).
Cooper pairs tunnel through the barrier. Following Feynman (Lectures, vol.~III, chap.~21), the two amplitudes obey coupled equations exactly like the double well of \cref{sec:doublewell},
$\ii\hbar\dot\psi_1=U_1\psi_1+K\psi_2$, $\ii\hbar\dot\psi_2=U_2\psi_2+K\psi_1$, with $U_1-U_2=qV$, $q=2e$. Separating real and imaginary parts and writing $\varphi=\theta_2-\theta_1$, one obtains
\begin{equation}
  I=I_c\sin\varphi,\qquad\dot\varphi=\frac{2e}\hbar V,
\end{equation}
the \emph{DC and AC Josephson relations} (Josephson, 1962). The second relation means $\varphi=2\pi\Phi/\Phi_0$, with the flux quantum $\Phi_0=h/2e$. The energy stored in the junction is $\int IV\dd t=-E_J\cos\varphi$, with the
\emph{Josephson energy} $E_J=\hbar I_c/2e$. For small $\varphi$, $I\approx I_c\varphi$ is the relation of an inductor with $L_J=\hbar/(2eI_c)$; for larger $\varphi$ the inductance $L_J/\cos\varphi$ depends on the current:
\emph{the Josephson junction is a non-linear, non-dissipative inductor}, the only such element available in electronics.
\end{origin}

\subsection{The transmon}
Replacing the inductor of the LC circuit by a junction (shunted by a capacitance $C$) gives, with the number of Cooper pairs $\hat n=\hat Q/2e$ and the phase $\hat\varphi$, conjugate variables $[\hat\varphi,\hat n]=\ii$,
\begin{equation}\label{eq:transmon}
  H=4E_C(\hat n-n_g)^2-E_J\cos\hat\varphi,\qquad E_C=\frac{e^2}{2C},
\end{equation}
where $n_g$ is an offset charge set by a gate voltage or by stray charges. For $E_J\lesssim E_C$ this is the \emph{Cooper-pair box} (Nakamura et al., 1999, the first superconducting qubit): the qubit states are $n=0$ and $n=1$ near $n_g=1/2$,
coupled by $E_J/2$, a realization of the double-well model. Its energy depends strongly on $n_g$, so charge noise dephases it within nanoseconds.

The \emph{transmon} (``transmission-line shunted plasma oscillation'' qubit) \cite{Koch2007} works in the opposite regime $E_J/E_C\sim50$. The phase is then localized near $\varphi=0$ and we may expand the cosine,
$-E_J\cos\varphi\approx-E_J+\frac{E_J}2\varphi^2-\frac{E_J}{24}\varphi^4$. The quadratic part is an oscillator with
\begin{equation}
  \hbar\omega_p=\sqrt{8E_JE_C},\qquad\hat\varphi=\left(\frac{2E_C}{E_J}\right)^{1/4}(\hat b+\hat b^\dagger),\qquad\hat n=\frac\ii2\left(\frac{E_J}{2E_C}\right)^{1/4}(\hat b^\dagger-\hat b).
\end{equation}
(Check: $4E_C\hat n^2+\frac{E_J}2\hat\varphi^2$ has the form $\frac{p^2}{2m}+\frac12m\omega^2x^2$ with ``mass'' $1/8E_C$ and $m\omega^2=E_J$.) The quartic term is a perturbation. Using $\mel m{(\hat b+\hat b^\dagger)^4}m=6m^2+6m+3$ and
$\frac{E_J}{24}\left(\frac{2E_C}{E_J}\right)=\frac{E_C}{12}$, first-order perturbation theory gives
\begin{equation}
  E_m\approx-E_J+\sqrt{8E_JE_C}\left(m+\tfrac12\right)-\frac{E_C}{12}\left(6m^2+6m+3\right).
\end{equation}
The transition frequencies are $E_{01}=\sqrt{8E_JE_C}-E_C$ and $E_{12}=\sqrt{8E_JE_C}-2E_C$: the \emph{anharmonicity} is
\begin{equation}
  \alpha=E_{12}-E_{01}=-E_C.
\end{equation}
Typical values are $E_{01}/h\approx4$--$6$\,GHz and $E_C/h\approx200$--$300$\,MHz. The anharmonicity is what makes the two lowest levels a qubit: a microwave pulse resonant with $E_{01}$ is detuned by $\alpha$ from $E_{12}$, and pulses longer than
$\sim1/|\alpha|\approx5$\,ns do not cause leakage (the DRAG technique allows even shorter pulses). The transmon is, in essence, the anharmonic oscillator of the beginning of these notes. Its decisive advantage is that the dependence on $n_g$
(``charge dispersion'') decreases exponentially, as $\e^{-\sqrt{8E_J/E_C}}$, because tunnelling of the phase between $\varphi$ and $\varphi+2\pi$ (which is what feels $n_g$) is exponentially suppressed. At $E_J/E_C=50$ the charge noise is
reduced by $\sim10^{-9}$ while the anharmonicity is reduced only polynomially.

Replacing the single junction by two junctions in a loop (a SQUID) makes the effective $E_J=2E_J^{(0)}|\cos(\pi\Phi_{\rm ext}/\Phi_0)|$ tunable by an external magnetic flux, and therefore the qubit frequency. This is used for tunable couplers and
two-qubit gates, at the price of sensitivity to flux noise. \emph{Flux qubits} (used in D-Wave annealers) instead have a double-well potential in $\varphi$, with clockwise and anticlockwise persistent currents as the two states.

\paragraph{Gates and coupling.} Single-qubit gates are microwave pulses (Rabi oscillations, \cref{sec:rabi}), with $20$--$50$\,ns duration. The qubits are coupled through capacitors or inductances (the latter acting like a transformer), or through
a common microwave resonator, giving the $XY$ or $ZZ$ interactions of \cref{ch:gates}. Two-qubit gates take $30$--$300$\,ns with fidelities up to $99.9\%$.

\subsection{Circuit QED: the Jaynes--Cummings model and dispersive readout}\label{sec:jc}
\begin{origin}[: the Jaynes--Cummings Hamiltonian]
Consider a two-level system (an atom, or a transmon) of frequency $\omega_q$ coupled to a single mode of frequency $\omega_r$ of a cavity or resonator. The electric field of the mode at the position of the atom is $\hat E=E_{\rm zpf}(\hat a+\hat a^\dagger)$,
where $E_{\rm zpf}=\sqrt{\hbar\omega_r/2\varepsilon_0V}$ is the vacuum fluctuation of the field in the mode volume $V$ (obtained by quantizing the mode as an oscillator, \cref{sec:oscillator}, with energy $\int\varepsilon_0E^2\dd V$). The dipole coupling
$-\hat d\cdot\hat E$ with $\hat d=d(\sigma^++\sigma^-)$ gives
\begin{equation}
  H=\hbar\omega_r\hat a^\dagger\hat a+\frac{\hbar\omega_q}2\sigma_z+\hbar g(\sigma^++\sigma^-)(\hat a+\hat a^\dagger),\qquad\hbar g=-dE_{\rm zpf},
\end{equation}
with $\sigma_z=\proj e-\proj g$ here. In the interaction picture, $\sigma^+\hat a$ and $\sigma^-\hat a^\dagger$ oscillate at $\omega_q-\omega_r$, while $\sigma^+\hat a^\dagger$ and $\sigma^-\hat a$ oscillate at $\omega_q+\omega_r$. Near resonance and for
$g\ll\omega$, the fast terms are dropped (rotating-wave approximation, as in \cref{sec:rabi}), which gives the \emph{Jaynes--Cummings Hamiltonian} (1963)
\begin{equation}\label{eq:jc}
  H_{\rm JC}=\hbar\omega_r\hat a^\dagger\hat a+\frac{\hbar\omega_q}2\sigma_z+\hbar g\left(\sigma^+\hat a+\sigma^-\hat a^\dagger\right).
\end{equation}
The coupling exchanges one excitation between the atom and the field: an absorbed photon excites the atom, an emitted photon de-excites it. For a transmon the same Hamiltonian follows from a capacitive coupling between the charge $\hat n$ of the
transmon and the voltage of the resonator.
\end{origin}
$H_{\rm JC}$ conserves the number of excitations $\hat a^\dagger\hat a+\proj e$, so it is block diagonal in the two-dimensional subspaces $\{\ket{e,n},\ket{g,n+1}\}$, in which it reads (up to a constant)
\begin{equation}
  H_n=\hbar\begin{pmatrix}\Delta/2&g\sqrt{n+1}\\g\sqrt{n+1}&-\Delta/2\end{pmatrix},\qquad\Delta=\omega_q-\omega_r.
\end{equation}
This is the two-level problem of \cref{ch:twolevel}: the eigenenergies are split by $\hbar\sqrt{\Delta^2+4g^2(n+1)}$.
\begin{itemize}
  \item \emph{Resonance} ($\Delta=0$): the eigenstates are the ``dressed states'' $(\ket{e,n}\pm\ket{g,n+1})/\sqrt2$, split by $2\hbar g\sqrt{n+1}$. An excited atom in the empty cavity exchanges its excitation with the field at the
  \emph{vacuum Rabi frequency} $2g$. Strong coupling requires $g$ larger than the cavity and atom decay rates ($g\gg\kappa,\gamma$). The $\sqrt{n+1}$ dependence makes the spectrum anharmonic: \emph{the atom acts as a giant optical nonlinearity}.
  \item \emph{Dispersive regime} ($|\Delta|\gg g$): there is no real exchange of excitations, but second-order perturbation theory gives energy shifts. The state $\ket{e,n}$ couples to $\ket{g,n+1}$ (energy difference $\hbar\Delta$), and is shifted
  by $+\hbar g^2(n+1)/\Delta$; the state $\ket{g,n}$ couples to $\ket{e,n-1}$ (energy difference $-\hbar\Delta$), and is shifted by $-\hbar g^2n/\Delta$. These shifts are reproduced, up to a constant, by
  \begin{equation}\label{eq:dispersive}
    H_{\rm disp}=\hbar\left(\omega_r+\chi\sigma_z\right)\hat a^\dagger\hat a+\frac{\hbar(\omega_q+\chi)}2\sigma_z,\qquad\chi=\frac{g^2}\Delta .
  \end{equation}
  The resonator frequency depends on the qubit state: $\omega_r\pm\chi$. Sending a microwave pulse through the resonator and measuring the phase of the transmitted or reflected signal measures $\sigma_z$ without exchanging energy with the qubit (a
  quantum non-demolition measurement). This \emph{dispersive readout} is used in all superconducting quantum computers (with $\chi/2\pi\sim1$\,MHz, measurement times of $100$--$500$\,ns and quantum-limited parametric amplifiers). Historically,
  the readout was first performed with SQUIDs (superconducting quantum interference devices). The term $\chi\sigma_z\hat a^\dagger\hat a$ can also be read the other way round: each photon in the resonator shifts the qubit frequency by $2\chi$ (AC Stark shift).
  For a transmon, the multilevel structure modifies $\chi$ to $\chi=\frac{g^2}{\Delta}\frac{\alpha}{\Delta+\alpha}$.
\end{itemize}
The circuits are fabricated with lithographic techniques, which makes them easy to scale, and their parameters can be designed at will. The price is a strong coupling to the many degrees of freedom of the solid (two-level defects in amorphous oxides,
quasiparticles, flux noise), so the main problem is noise and decoherence: $T_1,T_2\sim100\,\mu$s--$1$\,ms today, compared with $1$\,ns in 1999.

\section{Optical cavity quantum electrodynamics}
The Jaynes--Cummings model was first realized with atoms in optical and microwave cavities (Haroche, Kimble, Rempe; Nobel prize 2012 for Haroche and Wineland). The qubit is a single photon in two modes, and gates are made with phase shifters,
beam splitters and the atom-induced nonlinearity. The strength of the coupling is set by $g\propto E_{\rm zpf}\propto1/\sqrt V$: small mode volumes and long photon lifetimes (high finesse) are needed. The drawback is an intrinsic conflict.
A strong light--matter coupling requires the photon to stay in the cavity for a long time, whereas getting the photon in and out of the cavity requires a strong coupling to the outside. Cavity QED is now mostly used for quantum networks
(atom--photon interfaces) rather than for computing.

\section{Trapped ions}\label{sec:ions}
\paragraph{Trapping.} By Earnshaw's theorem, no static electric field configuration can trap a charge in three dimensions ($\nabla^2\phi=0$ forbids a minimum). The \emph{Paul trap} uses an oscillating quadrupole field
$\phi\propto(x^2-y^2)\cos\Omega t$. A particle in a rapidly oscillating field $E_0(\vec r)\cos\Omega t$ performs a fast micromotion whose averaged kinetic energy acts as a \emph{ponderomotive} (pseudo-)potential,
\begin{equation}
  U_{\rm eff}(\vec r)=\frac{q^2E_0(\vec r)^2}{4m\Omega^2},
\end{equation}
which is minimal where the field vanishes, at the centre of the quadrupole. (Derivation: write $\vec r=\vec R+\vec\xi$ with the fast oscillation $\vec\xi=-\frac{qE_0}{m\Omega^2}\cos\Omega t$; average the force over one period to second order in $\xi$.)
The ions are confined in vacuum, cooled by lasers to the motional ground state, and form a linear crystal (a chain) because of their Coulomb repulsion. The collective vibrations of the chain are harmonic modes of frequency $\nu\sim1$--$5$\,MHz.

\paragraph{Qubits.} The qubit is encoded in two long-lived internal levels of each ion, typically hyperfine levels of the ground state (where the nuclear spin plays a role, e.g.\ $^{171}$Yb$^+$ at $12.6$\,GHz), or a ground and a metastable level
(optical qubits, e.g.\ $^{40}$Ca$^+$). Hyperfine qubits have coherence times of seconds to an hour. Single-qubit gates are performed with laser pulses (or microwaves): Rabi oscillations again.

\paragraph{Coupling internal and motional states.} A laser of wave vector $k$ and frequency $\omega$ acting on an ion at position $\hat x=x_0(\hat a+\hat a^\dagger)$, $x_0=\sqrt{\hbar/2m\nu}$, gives in the RWA
\begin{equation}
  H=\frac{\hbar\Omega}2\left(\sigma^+\e^{\ii(k\hat x-\omega t)}+\text{h.c.}\right),\qquad k\hat x=\eta(\hat a+\hat a^\dagger),
\end{equation}
with the \emph{Lamb--Dicke parameter} $\eta=kx_0\sim0.1$. Expanding $\e^{\ii\eta(\hat a+\hat a^\dagger)}\approx1+\ii\eta(\hat a+\hat a^\dagger)$ and going to the interaction picture with respect to $\hbar\omega_0\sigma_z/2+\hbar\nu\hat a^\dagger\hat a$, the terms
oscillate at $\omega-\omega_0$ (carrier), $\omega-\omega_0+\nu$ and $\omega-\omega_0-\nu$ (sidebands). Tuning the laser to the \emph{red sideband} $\omega=\omega_0-\nu$ and keeping only the resonant terms,
\begin{equation}
  H_{\rm red}=\frac{\ii\hbar\eta\Omega}2\left(\sigma^+\hat a-\sigma^-\hat a^\dagger\right),
\end{equation}
which is again a Jaynes--Cummings Hamiltonian, with the phonon playing the role of the photon. The blue sideband $\omega=\omega_0+\nu$ gives the ``anti-JC'' coupling $\sigma^+\hat a^\dagger+\sigma^-\hat a$. The ions interact through their Coulomb
repulsion, so their collective vibrations, the phonons of the harmonic trap, serve as a \emph{bus} between the qubits: in the Cirac--Zoller gate \cite{CiracZoller1995} (1995), a red-sideband $\pi$ pulse maps the state of the control ion onto the
common motional mode, a conditional $2\pi$ rotation on the target ion through an auxiliary level gives a phase $-1$ only if both the phonon and the target are excited, and the mapping is reversed: a CZ gate.

\paragraph{Mølmer--Sørensen gate.} Today's ion-trap computers use the Mølmer--Sørensen gate (1999): both ions are illuminated by two frequencies $\omega_0\pm(\nu+\delta)$. The phonons are only virtually excited, and second-order perturbation theory gives the
effective interaction
\begin{equation}
  H_{\rm MS}=\hbar\frac{(\eta\Omega)^2}{\delta}\,\sigma_x^{(1)}\sigma_x^{(2)},
\end{equation}
which is insensitive to the temperature of the motion to first order. After the time $t$ with $(\eta\Omega)^2t/\delta=\pi/4$ it produces $\e^{-\ii\frac\pi4X_1X_2}$, which maps $\ket{00}\to(\ket{00}-\ii\ket{11})/\sqrt2$, a maximally entangling gate.
Because all ions in a chain share the motional modes, any pair can be coupled: ion traps have all-to-all connectivity.

\paragraph{Readout.} State-dependent fluorescence: a laser drives a cycling transition from one qubit state only. An ion in that state scatters millions of photons and appears bright; in the other state it stays dark. Fidelities exceed $99.9\%$.

Ion traps offer excellent coherence times and the best gate fidelities ($99.9\%$ two-qubit, $99.9999\%$ single-qubit). The main drawbacks are the slowness of the gates ($10$--$100\,\mu$s), which becomes more severe as the number of ions grows
(the motional spectrum becomes crowded), heating of the motional modes, and the difficulty of scaling a single chain beyond $\sim50$ ions. Modular architectures shuttle ions between zones of segmented traps (``quantum CCD'') or connect traps with photons.

\section{Nuclear magnetic resonance}\label{sec:nmr}
The qubits are the spins of atomic nuclei (${}^1$H, ${}^{13}$C, ${}^{19}$F, \dots) in a molecule placed in a strong magnetic field $B_0\sim10$\,T. Each spin has the Zeeman Hamiltonian $-\hbar\gamma B_0I_z$ (\cref{sec:spin}), with Larmor frequencies of hundreds of MHz,
shifted slightly by the electronic environment (\emph{chemical shift}): different nuclei of the molecule have different frequencies and can be addressed individually. Gates are performed by radio-frequency pulses, exactly as the Rabi problem of
\cref{sec:rabi} (in fact NMR is where Rabi oscillations, rotating frames, $T_1$ and $T_2$ (Bloch, 1946) and spin echoes (Hahn, 1950) were invented). The qubits are coupled through the chemical bonds of the molecule: the \emph{$J$-coupling}
$H_J=2\pi\hbar J_{12}I_z^{(1)}I_z^{(2)}$ ($J\sim10$--$200$\,Hz), mediated by the electrons (hyperfine interaction with the nucleus, transmitted along the bond), which is an Ising interaction and gives the CZ gate of \cref{ch:gates}. The readout detects the voltage induced
in a coil by the precessing spins (free induction decay). NMR provided the first demonstrations of many quantum algorithms, including Grover's (1998), quantum error correction (1998) \cite{Cory1998} and Shor's algorithm (2001) \cite{Vandersypen2001}.

Its fundamental drawback is that the experiment is performed on a macroscopic ensemble of $\sim10^{18}$ molecules at room temperature. The nuclear polarization is tiny, $\epsilon=\hbar\gamma B_0/2k_BT\sim10^{-5}$ (\cref{sec:thermal}), so the state is
$\hat\rho\approx\frac{I}{2^n}+\epsilon\,\Delta\hat\rho$, almost maximally mixed. One prepares a \emph{pseudo-pure} state $\hat\rho=(1-p)\frac I{2^n}+p\proj{0\dots0}$ by averaging over several experiments; the identity part gives no signal, and the
deviation behaves like a pure state. But the weight $p\sim n\epsilon/2^n$ of the effective pure state decreases exponentially with the number of qubits, which limits NMR to about ten qubits. (The states are also so mixed that they are separable:
NMR experiments do not involve entanglement, which raised a debate about their ``quantumness''.)

\section{Electron spins in quantum dots}\label{sec:spinqubits}
In the proposal of Loss and DiVincenzo (1998), the qubit is the spin of a single electron confined in a semiconductor quantum dot (an ``artificial atom'' defined by gate electrodes in a GaAs or Si/SiGe heterostructure). Single-qubit gates use
electron spin resonance (an oscillating magnetic field) or electric-dipole spin resonance with a micromagnet. The two-qubit interaction is the \emph{exchange interaction}, controlled by the voltage of the gate between two dots.

\begin{origin}[: the exchange interaction from the Hubbard model]
Consider two dots, each with one orbital, and two electrons. The Hubbard model is
\begin{equation}
  H=-t\sum_{\sigma=\uparrow,\downarrow}\left(c_{1\sigma}^\dagger c_{2\sigma}+c_{2\sigma}^\dagger c_{1\sigma}\right)+U\sum_{i=1,2}n_{i\uparrow}n_{i\downarrow},
\end{equation}
where $t$ is the tunnelling amplitude between the dots and $U$ the charging energy of a doubly occupied dot. For $U\gg t$, the low-energy states have one electron per dot: four spin states, the singlet $S$ and the triplet $T_{0,\pm}$. Tunnelling
couples the singlet (and only the singlet) to the doubly occupied states, at energy $U$. Indeed the Pauli principle forbids two electrons with the same spin in the same orbital, so the triplets cannot hop at all. The coupling of the singlet
$\frac1{\sqrt2}(\ket{\uparrow,\downarrow}-\ket{\downarrow,\uparrow})$ to the symmetric combination of doubly occupied states $\frac1{\sqrt2}(\ket{\uparrow\downarrow,0}+\ket{0,\uparrow\downarrow})$ has the matrix element $-2t$. Second-order perturbation
theory therefore lowers the singlet by $(2t)^2/U$ and leaves the triplets unchanged:
\begin{equation}
  E_T-E_S=J=\frac{4t^2}U,\qquad H_{\rm eff}=J\left(\vec S_1\cdot\vec S_2-\tfrac14\right)
\end{equation}
(with $\vec S$ in units of $\hbar$; $\vec S_1\cdot\vec S_2=-\frac34$ for the singlet and $+\frac14$ for the triplets). The same mechanism (kinetic exchange) is responsible for antiferromagnetism in Mott insulators, and the Ising model of \cref{ch:annealing}
is its anisotropic limit. Since $t$ depends exponentially on the barrier between the dots, $J$ can be switched on and off by a gate voltage over many orders of magnitude.
\end{origin}
Pulsing $J$ for a time $\tau$ with $J\tau/\hbar=\pi/2$ gives the $\sqrt{\mathrm{SWAP}}$ gate (\cref{ch:gates}). The readout uses spin-to-charge conversion: a spin-dependent tunnelling event changes the charge of the dot, which is detected by a nearby
charge sensor (Elzerman et al.\ 2004). The main decoherence source in GaAs is the hyperfine interaction with $\sim10^6$ nuclear spins; silicon, isotopically purified in $^{28}$Si (nuclear spin zero), gives coherence times of milliseconds.
Spin qubits are tiny ($\sim100$\,nm) and compatible with CMOS fabrication, which makes them attractive for scaling; the challenge is the variability between dots and the wiring of many gates.

\section{Neutral atoms and Rydberg blockade}
Neutral atoms (Rb, Cs, Sr, Yb) are trapped individually in arrays of optical tweezers (tightly focused laser beams, which attract the atom by the AC Stark shift towards the intensity maximum), with hundreds to thousands of sites, rearranged by
moving tweezers. The qubit is a pair of hyperfine or nuclear-spin levels of the ground state, with seconds of coherence. Atoms in the ground state hardly interact. Two-qubit gates use \emph{Rydberg states}, highly excited states with principal quantum
number $n\sim50$--$100$, whose huge size ($\propto n^2$) gives a strong van der Waals interaction $V(r)=C_6/r^6$, with $C_6\propto n^{11}$.

\paragraph{Rydberg blockade.} If two atoms closer than the \emph{blockade radius} $R_b=(C_6/\hbar\Omega)^{1/6}$ (several $\mu$m) are driven to the Rydberg state $\ket r$ with Rabi frequency $\Omega$, the doubly excited state $\ket{rr}$ is shifted out of resonance by
$V\gg\hbar\Omega$: only one atom can be excited. A CZ gate follows (Jaksch et al., 2000). Let only $\ket1$ be coupled to $\ket r$. Apply a $\pi$ pulse on the control atom, a $2\pi$ pulse on the target, and a $\pi$ pulse on the control. A $2\pi$ Rabi
rotation multiplies a state by $-1$ (\cref{sec:simplest}: $U(2\pi/\omega)=-I$). Then $\ket{00}\to\ket{00}$; $\ket{01}\to-\ket{01}$ (target $2\pi$); $\ket{10}\to-\ket{10}$ (control $\pi+\pi=2\pi$); $\ket{11}\to-\ket{11}$ (the control goes to $\ket r$, the
target is blocked and does nothing, and the control returns with a sign $-1$). The result $\diag(1,-1,-1,-1)=\mathrm{CZ}\cdot(Z\otimes Z)$ is a CZ up to single-qubit $Z$ gates. Neutral-atom processors have demonstrated two-qubit fidelities above
$99.5\%$, parallel gates on hundreds of atoms, and the first logical-qubit algorithms (2023--2024). They are also natural analog simulators and annealers of Ising models, since the blockade implements the constraint ``no two neighbours excited''.

\section{Photons}
\paragraph{Linear optics.} A qubit can be encoded in a single photon, either in its polarization (\cref{sec:photonic}) or in which of two modes it occupies ($\ket{01}$ or $\ket{10}$, the \emph{dual-rail} encoding). Single-qubit gates are easy and very accurate:
phase shifters and beam splitters, realized for example with integrated waveguides. A beam splitter mixes two modes,
\begin{equation}
  \begin{pmatrix}\hat a_{\rm out}\\ \hat b_{\rm out}\end{pmatrix}=\begin{pmatrix}\cos\vartheta&\ii\sin\vartheta\\ \ii\sin\vartheta&\cos\vartheta\end{pmatrix}\begin{pmatrix}\hat a_{\rm in}\\ \hat b_{\rm in}\end{pmatrix},
\end{equation}
a unitary matrix (energy conservation), i.e.\ the rotation $R_x(-2\vartheta)$ of a dual-rail qubit; the 50:50 splitter has $\vartheta=\pi/4$. The readout uses single-photon detectors (superconducting nanowires reach $>95\%$ efficiency). Photons hardly
decohere and travel at the speed of light, which makes them ideal for communication.

\paragraph{Two-photon interference.} Two identical photons entering the two ports of a 50:50 beam splitter always exit together (Hong--Ou--Mandel effect, 1987): with $\hat a^\dagger\hat b^\dagger\to\frac12(\hat a^\dagger+\ii\hat b^\dagger)(\ii\hat a^\dagger+\hat b^\dagger)
=\frac\ii2(\hat a^{\dagger2}+\hat b^{\dagger2})$, the amplitudes for one photon in each output cancel. This bosonic interference is the basic resource of linear optical quantum computing and of boson sampling.

\paragraph{The interaction problem.} But for the same reason that they hardly decohere, photons hardly interact with each other. A two-qubit gate requires an optical nonlinearity, such as the Kerr effect in a nonlinear crystal: the refractive index
depends on the intensity, $n=n_0+n_2I$, which for two modes gives the cross-Kerr Hamiltonian $H=\hbar\chi\hat n_a\hat n_b$ and, after a time $t$ with $\chi t=\pi$, a phase $-1$ only when both modes contain a photon: a CZ gate in dual-rail encoding.
Available Kerr nonlinearities are very weak, and materials with a strong nonlinearity also absorb strongly (the real and imaginary parts of the susceptibility are related by the Kramers--Kronig relations). A conditional phase shift of $\pi$ between
two single photons would be accompanied by the absorption of about 50 photons for each photon transmitted. These huge losses make this direct approach impractical.

\paragraph{Measurement-induced nonlinearity.} Knill, Laflamme and Milburn (KLM, 2001) showed that linear optics, single-photon sources, detectors and feed-forward suffice for universal quantum computing: the nonlinearity is provided by the measurement
of ancilla photons, and gates succeed with some probability, heralded by the detection pattern; teleportation (\cref{ch:teleportation}) boosts the success probability. Modern photonic architectures (fusion-based quantum computing) build large entangled
cluster states from small resource states and probabilistic ``fusion'' measurements. The main challenges are photon losses and deterministic single-photon sources (quantum dots in cavities are approaching the requirements).

\section{Current status}
Superconducting circuits and trapped ions are at present the most advanced platforms for gate-based quantum computing, and both are available through cloud services. Neutral atoms held in optical tweezers, photonic circuits and spin qubits in
semiconductors are also developing rapidly. Other approaches include topological qubits (Majorana modes, still at the stage of demonstrating the basic physics) and bosonic codes in microwave cavities (cat qubits, GKP states), which encode a qubit in an
oscillator and use its large Hilbert space for error correction. No platform yet satisfies all the requirements at the scale needed for fault-tolerant quantum computation. Achieving it requires quantum error correction (\cref{ch:qec}), which multiplies the
number of physical qubits needed for each logical qubit by $10^2$--$10^3$. The first demonstrations of logical qubits below threshold (2023--2024) mark the beginning of this new stage.

\section{Exercises}
\begin{exercise}
A transmon has $E_J/h=20$\,GHz and $E_C/h=0.25$\,GHz. Compute $E_J/E_C$, $f_{01}$ and the anharmonicity. (Answer: $80$; $f_{01}\approx\sqrt{8\cdot20\cdot0.25}-0.25=6.07$\,GHz; $-250$\,MHz.)
\end{exercise}
\begin{exercise}
Show that $H_{\rm JC}$ commutes with $\hat N=\hat a^\dagger\hat a+\proj e$. Compute the time evolution of $\ket{e,0}$ at resonance, and the probability to find the photon in the cavity.
\end{exercise}
\begin{exercise}
For a $^{40}$Ca$^+$ ion with $\lambda=729$\,nm and trap frequency $\nu/2\pi=1$\,MHz, compute the Lamb--Dicke parameter (mass $40\times1.66\times10^{-27}$\,kg). (Answer: $\eta\approx0.1$.)
\end{exercise}
\begin{exercise}
Verify the matrix element $-2t$ in the Hubbard derivation by writing the states with creation operators, e.g.\ $\ket{\uparrow,\downarrow}=c^\dagger_{1\uparrow}c^\dagger_{2\downarrow}\ket{\rm vac}$, and keeping track of the fermionic signs.
\end{exercise}


\appendix
\part{Appendices}
\chapter{The Schmidt decomposition and purification}\label{app:schmidt}

\section{Derivation}
Consider a pure state of a bipartite system, $\ket\psi_{AB}\in\Ha_A\otimes\Ha_B$. In arbitrary orthonormal bases $\{\ket i_A\}$ and $\{\ket\mu_B\}$ of the two subsystems, any such state can be written
\begin{equation}
  \ket\psi_{AB}=\sum_{i,\mu}\psi_{i\mu}\ket i_A\otimes\ket\mu_B .
\end{equation}
Grouping the terms with the same $\ket i_A$, we can also write
\begin{equation}
  \ket\psi_{AB}=\sum_i\ket i_A\otimes\ket{\tilde\imath}_B,\qquad\text{with}\qquad\ket{\tilde\imath}_B=\sum_\mu\psi_{i\mu}\ket\mu_B .
\end{equation}
In general the states $\ket{\tilde\imath}_B$ are neither orthogonal nor normalized. Let us show that they are orthogonal if the basis of $A$ is suitably chosen. We choose the basis $\{\ket i_A\}$ in which the reduced density matrix of $A$ is diagonal:
\begin{equation}
  \hat\rho_A=\sum_ip_i\proj i .
\end{equation}
On the other hand, by definition of the reduced density matrix,
\begin{equation}
  \hat\rho_A=\tr_B\proj\psi=\tr_B\Big(\sum_{i,j}\ketbra ij\otimes\ketbra{\tilde\imath}{\tilde\jmath}\Big)=\sum_{i,j}\braket{\tilde\jmath}{\tilde\imath}\ketbra ij,
\end{equation}
because, using the completeness relation $\sum_k\proj k=I$ in $\Ha_B$, $\tr_B\ketbra{\tilde\imath}{\tilde\jmath}=\sum_k\braket k{\tilde\imath}\braket{\tilde\jmath}k=\braket{\tilde\jmath}{\tilde\imath}$. Comparing the two expressions of $\hat\rho_A$,
\begin{equation}
  \braket{\tilde\jmath}{\tilde\imath}_B=p_i\delta_{ij}:
\end{equation}
the states $\ket{\tilde\imath}_B$ are orthogonal. We can therefore normalize them, $\ket{i'}_B=p_i^{-1/2}\ket{\tilde\imath}_B$ (for $p_i\neq0$), and obtain
\begin{equation}\label{eq:schmidt}
  \ket\psi_{AB}=\sum_i\sqrt{p_i}\ket i_A\otimes\ket{i'}_B .
\end{equation}
This is the \emph{Schmidt decomposition}: any pure state of a bipartite system can be written as a single sum over pairs of orthonormal states $\ket i_A$, $\ket{i'}_B$ of the two subsystems, with non-negative coefficients $\sqrt{p_i}$, the
\emph{Schmidt coefficients}. The bases depend on the state. Compared with the general expansion, which has $d_Ad_B$ terms, the Schmidt form has at most $\min(d_A,d_B)$ terms.

\paragraph{Singular value decomposition.} In fixed bases $\{\ket a_A\}$ and $\{\ket\mu_B\}$, the Schmidt decomposition reads
\begin{equation}
  \psi_{a\mu}=\sum_i(U_A)_{ai}\sqrt{p_i}\,(U_B^T)_{i\mu},
\end{equation}
where the unitary matrices $U_A$ and $U_B$ transform the fixed bases into the Schmidt bases. This is the singular value decomposition (SVD) $\Psi=U_A\Sigma U_B^T$ of the matrix $\Psi=(\psi_{a\mu})$ of coefficients. Conversely, the SVD theorem
(any complex matrix is $U\Sigma V^\dagger$ with $U,V$ unitary and $\Sigma\ge0$ diagonal) gives a second proof of the Schmidt decomposition. The Schmidt coefficients are the singular values of $\Psi$, which gives a practical way to compute them
numerically (\texttt{numpy.linalg.svd}). The reduced density matrices are $\hat\rho_A=\Psi\Psi^\dagger$ and $\hat\rho_B=(\Psi^\dagger\Psi)^T$ in these bases.

\section{Consequences}
\paragraph{Equal spectra.} Tracing out $A$ instead of $B$ gives
\begin{equation}
  \hat\rho_B=\tr_A\proj\psi=\sum_ip_i\proj{i'}.
\end{equation}
The reduced density matrices $\hat\rho_A$ and $\hat\rho_B$ therefore have the same non-zero eigenvalues $p_i$, even if the two subsystems have different dimensions. In particular they have the same purity and the same von Neumann entropy,
which justifies the definition of the entanglement entropy $S(\hat\rho_A)=S(\hat\rho_B)=-\sum_ip_i\log_2p_i$.

\paragraph{Schmidt number and entanglement.} The number of non-zero Schmidt coefficients is the \emph{Schmidt number} (the rank of $\Psi$). A state is entangled if and only if its Schmidt number is greater than 1. For a separable state,
$\ket\psi_{AB}=\ket\varphi_A\otimes\ket\chi_B$, there is a single term, and the reduced density matrices $\hat\rho_A=\proj\varphi$, $\hat\rho_B=\proj\chi$ describe pure states. Otherwise, they are mixed, and subsystem $A$ alone has no wave function.
For example, the Bell state $(\ket{00}+\ket{11})/\sqrt2$ is already in Schmidt form, with two coefficients $1/\sqrt2$: its Schmidt number is 2, and $\hat\rho_A=\hat\rho_B=I/2$, in agreement with the partial trace computed in \cref{sec:reduced}.
For two qubits, $\det\Psi=\sqrt{p_1p_2}\det U_A\det U_B$, so $|\det\Psi|=\sqrt{p_1p_2}$, which explains the criterion of \cref{prop:det} and the concurrence $C=2\sqrt{p_1p_2}$.

\paragraph{Entanglement cannot be created locally.} A unitary operation $U_A\otimes U_B$ acting separately on the two subsystems only rotates the Schmidt bases,
\begin{equation}
  (U_A\otimes U_B)\ket\psi_{AB}=\sum_i\sqrt{p_i}\left(U_A\ket i_A\right)\otimes\left(U_B\ket{i'}_B\right),
\end{equation}
and leaves the Schmidt coefficients unchanged. In particular, local operations cannot increase the Schmidt number. Entanglement can only be created by an interaction between the subsystems, such as a CNOT gate. (More generally, entanglement cannot
increase on average under local operations and classical communication; Nielsen's theorem (1999) states that $\ket\psi$ can be converted into $\ket\phi$ by LOCC iff the Schmidt coefficients of $\ket\psi$ are majorized by those of $\ket\phi$.)

\section{Purification}\label{app:purification}
Conversely, any mixed state can be seen as part of a larger pure state. Suppose that $\hat\rho_A$ does not describe a pure state. We can introduce a fictitious \emph{reference system} $R$ such that the joint state of $A$ and $R$ is pure and
\begin{equation}
  \hat\rho_A=\tr_R\proj{AR}.
\end{equation}
\begin{proof}
Write $\hat\rho_A$ in its orthonormal eigenbasis, $\hat\rho_A=\sum_ip_i\proj{i^A}$. Let $R$ have the same state space as $A$, with an orthonormal basis $\{\ket{i^R}\}$, and define
\begin{equation}
  \ket{AR}=\sum_i\sqrt{p_i}\ket{i^A}\ket{i^R}.
\end{equation}
Then
\begin{equation}
  \tr_R\proj{AR}=\sum_{i,j}\sqrt{p_ip_j}\ketbra{i^A}{j^A}\tr\left(\ketbra{i^R}{j^R}\right)=\sum_{i,j}\sqrt{p_ip_j}\ketbra{i^A}{j^A}\delta_{ij}=\sum_ip_i\proj{i^A}=\hat\rho_A.\qedhere
\end{equation}
\end{proof}
The state $\ket{AR}$ is written directly in Schmidt form, which shows the close relation between purification and the Schmidt decomposition. Physically, purification says that mixedness can always be attributed to entanglement with
something else (the ``church of the larger Hilbert space''). A compact formula: with the unnormalized maximally entangled vector $\ket\Omega=\sum_i\ket i\ket i$, every purification of $\hat\rho$ is $(\sqrt{\hat\rho}\otimes V)\ket\Omega$
for some unitary (or isometry) $V$ on $R$. This is used in the proof of Uhlmann's theorem (\cref{app:fidelity}).

\section{The Hughston--Jozsa--Wootters theorem}\label{app:hjw}
A given density matrix can be realized by many different ensembles of pure states \cite{HJW1993}:
\begin{equation}
  \hat\rho_A=\sum_ip_i\proj{\varphi_i}=\sum_\mu q_\mu\proj{\psi_\mu},
\end{equation}
where neither the $\ket{\varphi_i}$ nor the $\ket{\psi_\mu}$ need to be orthogonal. Each ensemble corresponds to a purification,
\begin{equation}
  \ket{\Phi_1}_{AB}=\sum_i\sqrt{p_i}\ket{\varphi_i}_A\ket{\alpha_i}_B,\qquad\ket{\Phi_2}_{AB}=\sum_\mu\sqrt{q_\mu}\ket{\psi_\mu}_A\ket{\beta_\mu}_B,
\end{equation}
where $\{\ket{\alpha_i}\}$ and $\{\ket{\beta_\mu}\}$ are orthonormal states of a system $B$ of sufficiently large dimension. (Check: $\tr_B\proj{\Phi_1}=\sum_{ij}\sqrt{p_ip_j}\ketbra{\varphi_i}{\varphi_j}\braket{\alpha_j}{\alpha_i}=\hat\rho_A$.)
\begin{theorem}[HJW]
Two purifications of the same $\hat\rho_A$ on the same space $B$ are related by a unitary change of basis in $B$ alone:
\begin{equation}
  \ket{\Phi_1}_{AB}=(I_A\otimes U_B)\ket{\Phi_2}_{AB}.
\end{equation}
\end{theorem}
\begin{proof}
Both states give the same reduced density matrix $\hat\rho_A$, so their Schmidt decompositions can be written with the same coefficients and the same eigenbasis $\{\ket k_A\}$ of $\hat\rho_A$ (if some eigenvalues are degenerate, we choose the same basis
of each eigenspace for both; this is possible because \cref{eq:schmidt} holds for any eigenbasis of $\hat\rho_A$):
\begin{equation}
  \ket{\Phi_1}_{AB}=\sum_k\sqrt{\lambda_k}\ket k_A\ket{k_1}_B,\qquad\ket{\Phi_2}_{AB}=\sum_k\sqrt{\lambda_k}\ket k_A\ket{k_2}_B .
\end{equation}
The linear map defined by $\ket{k_1}_B=U_B\ket{k_2}_B$ maps one orthonormal set onto another; it can be extended to a unitary on the whole of $\Ha_B$, which maps one state onto the other.
\end{proof}
An equivalent algebraic form (\cite{NielsenChuang}, Thm.~2.6): $\sum_ip_i\proj{\varphi_i}=\sum_\mu q_\mu\proj{\psi_\mu}$ iff $\sqrt{p_i}\ket{\varphi_i}=\sum_\mu u_{i\mu}\sqrt{q_\mu}\ket{\psi_\mu}$ for some unitary matrix $u$ (padding the
smaller ensemble with zero-probability states).

The physical meaning is remarkable. If $A$ and $B$ share the state $\ket{\Phi_1}$, then by choosing in which basis to measure $B$, an observer can decide which ensemble describes $A$: measuring $B$ in the basis $\{\ket{\alpha_i}\}$ leaves $A$ in
$\ket{\varphi_i}$ with probability $p_i$; measuring in the basis $\{U_B\ket{\beta_\mu}\}$ leaves $A$ in $\ket{\psi_\mu}$ with probability $q_\mu$. Since both ensembles give the same $\hat\rho_A$, an observer of $A$ alone cannot tell which choice was made,
and no information is transmitted. This is the precise content of ``steering'' (Schrödinger, 1935) and of the EPR argument (\cref{ch:bell}), and again of the no-signalling principle.

\section{Exercises}
\begin{exercise}
Find the Schmidt decomposition of $\frac12(\ket{00}+\ket{01}+\ket{10}-\ket{11})$. (Hint: compute the SVD of $\frac12\begin{pmatrix}1&1\\1&-1\end{pmatrix}$.) Is the state entangled?
\end{exercise}
\begin{exercise}
Show that the two ensembles $\{\frac12,\ket0;\frac12,\ket1\}$ and $\{\frac12,\ket+;\frac12,\ket-\}$ are related by the unitary $u=H$ in the HJW theorem, and find the corresponding measurements on $B$ for the purification $\ket{\Phi^+}$.
\end{exercise}

\chapter{Fidelity, Uhlmann's theorem and the quantum metric}\label{app:fidelity}

How close are two quantum states? The question is central for experiments (how good is a prepared state or a gate?), for error correction (how much does an error damage the state?) and for quantum metrology
(how well can a parameter encoded in a state be estimated?). This appendix completes the section announced but left empty in the original notes.

\section{Distinguishability of pure states}
For pure states the natural measure is the \emph{fidelity} $F(\psi,\phi)=|\braket\psi\phi|^2$: the probability that $\ket\phi$ passes a test ``is it $\ket\psi$?'' (the projective measurement $\{\proj\psi,I-\proj\psi\}$). It is $1$ for equal states (up to phase)
and $0$ for orthogonal ones. For qubits, $F=\cos^2(\Theta/2)$ where $\Theta$ is the angle between the Bloch vectors (\cref{eq:overlapbloch}). The angle $\arccos\sqrt F$ is a distance on the space of rays (the Fubini--Study distance, see below).

\section{Trace distance and the Helstrom bound}
The \emph{trace distance} between density matrices is
\begin{equation}
  D(\hat\rho,\hat\sigma)=\frac12\norm{\hat\rho-\hat\sigma}_1=\frac12\tr\abs{\hat\rho-\hat\sigma}=\frac12\sum_k|\lambda_k|,
\end{equation}
where $\lambda_k$ are the eigenvalues of the Hermitian operator $\hat\rho-\hat\sigma$. For qubits, $D=\frac12|\vec r-\vec s|$: half the Euclidean distance in the Bloch ball.
\begin{theorem}[Helstrom, 1969]
If a system is prepared in $\hat\rho$ or $\hat\sigma$ with equal probabilities, the best possible measurement identifies the state correctly with probability
\begin{equation}
  P_{\rm succ}=\frac12\left(1+D(\hat\rho,\hat\sigma)\right).
\end{equation}
\end{theorem}
\begin{proof}
A two-outcome measurement is a POVM $\{E,I-E\}$, $0\le E\le I$, where $E$ means ``guess $\hat\rho$''. Then $P_{\rm succ}=\frac12\tr(E\hat\rho)+\frac12\tr((I-E)\hat\sigma)=\frac12+\frac12\tr(E(\hat\rho-\hat\sigma))$. Writing
$\hat\rho-\hat\sigma=Q_+-Q_-$ with $Q_\pm\ge0$ supported on the positive and negative eigenspaces, $\tr(E(\hat\rho-\hat\sigma))\le\tr(EQ_+)\le\tr Q_+=D$, since $\tr(\hat\rho-\hat\sigma)=0$ implies $\tr Q_+=\tr Q_-=D$. Equality holds for $E$ the projector onto the
positive eigenspace.
\end{proof}
The trace distance is therefore the operational measure of distinguishability. It is contractive: $D(\mathcal E(\hat\rho),\mathcal E(\hat\sigma))\le D(\hat\rho,\hat\sigma)$ for any channel $\mathcal E$: no physical processing can make states more distinguishable.
For pure states, $D=\sqrt{1-|\braket\psi\phi|^2}$.

\section{Fidelity of mixed states and Uhlmann's theorem}
\begin{definition}
The fidelity of two density matrices is
\begin{equation}
  F(\hat\rho,\hat\sigma)=\left(\tr\sqrt{\sqrt{\hat\rho}\,\hat\sigma\sqrt{\hat\rho}}\right)^2=\norm{\sqrt{\hat\rho}\sqrt{\hat\sigma}}_1^2 .
\end{equation}
\end{definition}
(Some authors use the square root of this quantity.) When $\hat\sigma=\proj\phi$ is pure, $\sqrt{\hat\rho}\proj\phi\sqrt{\hat\rho}=\proj{\tilde\phi}$ with $\ket{\tilde\phi}=\sqrt{\hat\rho}\ket\phi$, whose square root is $\proj{\tilde\phi}/\norm{\tilde\phi}$, so
$F=\norm{\tilde\phi}^2=\mel\phi{\hat\rho}\phi$: the probability that $\hat\rho$ passes the test ``is it $\ket\phi$?''. This is the quantity measured in experiments as the ``fidelity of the prepared state''. For commuting states (classical probability
distributions $p_i$, $q_i$) it reduces to the Bhattacharyya overlap $(\sum_i\sqrt{p_iq_i})^2$.

\begin{theorem}[Uhlmann, 1976]
\begin{equation}
  F(\hat\rho,\hat\sigma)=\max_{\ket{\psi_\rho},\ket{\psi_\sigma}}\abs{\braket{\psi_\rho}{\psi_\sigma}}^2,
\end{equation}
where the maximum is over all purifications of $\hat\rho$ and $\hat\sigma$ on the same space $A\otimes R$ (and it suffices to vary one of them).
\end{theorem}
\begin{proof}
By \cref{app:purification}, all purifications have the form $\ket{\psi_\rho}=(\sqrt{\hat\rho}\otimes U)\ket\Omega$, $\ket{\psi_\sigma}=(\sqrt{\hat\sigma}\otimes V)\ket\Omega$, with $\ket\Omega=\sum_i\ket i\ket i$ and unitaries $U,V$ on $R$. Using the ``transpose trick''
$(I\otimes W)\ket\Omega=(W^T\otimes I)\ket\Omega$ and $\mel\Omega{X\otimes I}\Omega=\tr X$,
\begin{equation}
  \braket{\psi_\rho}{\psi_\sigma}=\mel\Omega{\sqrt{\hat\rho}\sqrt{\hat\sigma}\otimes U^\dagger V}\Omega=\tr\left(\sqrt{\hat\rho}\sqrt{\hat\sigma}\,W\right),\qquad W=(U^\dagger V)^T .
\end{equation}
For any matrix $X$ with polar decomposition $X=|X^\dagger|V_0$, we have $|\tr(XW)|\le\tr|X|$ for all unitaries $W$, with equality for $W=V_0^\dagger$ (Cauchy--Schwarz for the Hilbert--Schmidt product:
$|\tr(|X^\dagger|^{1/2}\,|X^\dagger|^{1/2}V_0W)|\le\sqrt{\tr|X^\dagger|}\sqrt{\tr(W^\dagger V_0^\dagger|X^\dagger|V_0W)}=\tr|X|$). With $X=\sqrt{\hat\rho}\sqrt{\hat\sigma}$, $\tr|X|=\tr\sqrt{X^\dagger X}$, and the operators $X^\dagger X=\sqrt{\hat\sigma}\hat\rho\sqrt{\hat\sigma}$ and
$XX^\dagger=\sqrt{\hat\rho}\hat\sigma\sqrt{\hat\rho}$ have the same non-zero eigenvalues. Hence $\max|\braket{\psi_\rho}{\psi_\sigma}|=\tr\sqrt{\sqrt{\hat\rho}\hat\sigma\sqrt{\hat\rho}}$.
\end{proof}
Uhlmann's theorem makes many properties of the fidelity obvious:
\begin{itemize}
  \item $0\le F\le1$, and $F=1$ iff $\hat\rho=\hat\sigma$; it is symmetric, $F(\hat\rho,\hat\sigma)=F(\hat\sigma,\hat\rho)$, which is not evident from the definition.
  \item It is invariant under unitaries, $F(U\hat\rho U^\dagger,U\hat\sigma U^\dagger)=F(\hat\rho,\hat\sigma)$.
  \item It can only increase under a quantum channel (monotonicity): purify the channel by a unitary on a larger space (\cref{ch:open}); the purifications of the outputs include the images of the purifications of the inputs, over which the maximum is taken.
  \item It is related to the trace distance by the Fuchs--van de Graaf inequalities $1-\sqrt F\le D\le\sqrt{1-F}$.
\end{itemize}
For two qubits with Bloch vectors $\vec r$ and $\vec s$, a direct computation gives
\begin{equation}
  F(\hat\rho,\hat\sigma)=\frac12\left[1+\vec r\cdot\vec s+\sqrt{(1-|\vec r|^2)(1-|\vec s|^2)}\right].
\end{equation}

\paragraph{Gate fidelity.} To characterize a gate, one compares the implemented channel $\mathcal E$ with the ideal unitary $U$. The \emph{average gate fidelity} $\bar F=\int\dd\psi\,\mel\psi{U^\dagger\mathcal E(\proj\psi)U}\psi$, averaged over all pure states,
is related to the \emph{process} (entanglement) fidelity $F_{\rm pro}$, the fidelity of the Choi state $(\mathcal E\otimes I)(\proj{\Phi})$ with $(U\otimes I)\ket\Phi$, by $\bar F=\frac{dF_{\rm pro}+1}{d+1}$ for dimension $d$. For a depolarizing channel
$\mathcal E(\hat\rho)=(1-p)\hat\rho+pI/d$, $\bar F=1-p\frac{d-1}d$. The ``99.9\% two-qubit gate fidelities'' quoted for hardware are average gate fidelities, usually measured by \emph{randomized benchmarking}: sequences of random Clifford gates of increasing length,
whose survival probability decays exponentially at a rate given by the average error, independently of state-preparation and measurement errors.

\section{The quantum metric}
\subsection{Fubini--Study metric}
For two neighbouring pure states $\ket{\psi(\vec\lambda)}$ and $\ket{\psi(\vec\lambda+\dd\vec\lambda)}$ depending on parameters $\vec\lambda$, expand the infidelity to second order. Using the normalization, $1-|\braket{\psi}{\psi+\dd\psi}|^2=\dd s^2$ with
\begin{equation}
  \dd s^2=\sum_{\mu\nu}g_{\mu\nu}\dd\lambda_\mu\dd\lambda_\nu,\qquad
  g_{\mu\nu}=\operatorname{Re}\left[\braket{\partial_\mu\psi}{\partial_\nu\psi}-\braket{\partial_\mu\psi}\psi\braket\psi{\partial_\nu\psi}\right].
\end{equation}
The subtraction makes $g$ independent of the phase convention (it is invariant under $\ket\psi\to\e^{\ii\chi(\vec\lambda)}\ket\psi$). This is the \emph{Fubini--Study metric}, the natural Riemannian metric on the space of quantum states.
For a qubit, $\ket\psi=(\cos\frac\theta2,\e^{\ii\varphi}\sin\frac\theta2)$ gives $\braket{\partial_\theta\psi}{\partial_\theta\psi}=\frac14$, $\braket\psi{\partial_\theta\psi}=0$, $\braket{\partial_\varphi\psi}{\partial_\varphi\psi}=\sin^2\frac\theta2$,
$\braket\psi{\partial_\varphi\psi}=\ii\sin^2\frac\theta2$, hence
\begin{equation}
  \dd s^2=\frac14\left(\dd\theta^2+\sin^2\theta\,\dd\varphi^2\right):
\end{equation}
the space of pure qubit states, with the Fubini--Study metric, is exactly a sphere of radius $1/2$, the Bloch sphere. The geodesic distance between two states is $\arccos|\braket\psi\phi|=\Theta/2$.

\subsection{Quantum geometric tensor and Berry phase}
The full Hermitian tensor $Q_{\mu\nu}=\braket{\partial_\mu\psi}{\partial_\nu\psi}-\braket{\partial_\mu\psi}\psi\braket\psi{\partial_\nu\psi}$ is the \emph{quantum geometric tensor}. Its real part is the metric $g_{\mu\nu}$; its imaginary part is
$-\frac12$ the \emph{Berry curvature} $\Omega_{\mu\nu}=\partial_\mu\mathcal A_\nu-\partial_\nu\mathcal A_\mu$, with the Berry connection $\mathcal A_\mu=\ii\braket\psi{\partial_\mu\psi}$. When the parameters are slowly varied around a closed loop $C$, an eigenstate of a
Hamiltonian $H(\vec\lambda)$ acquires, in addition to the dynamical phase, the geometric \emph{Berry phase} (1984)
\begin{equation}
  \gamma_C=\oint_C\mathcal A\cdot\dd\vec\lambda=\iint_S\Omega\,\dd S .
\end{equation}
(It is the term $m=n$ dropped in the derivation of the adiabatic theorem in \cref{ch:annealing}.) For the qubit, $\operatorname{Im}Q_{\theta\varphi}=\frac14\sin\theta$ gives a Berry curvature $\Omega_{\theta\varphi}=-\frac12\sin\theta$: a spin $\frac12$ that follows a magnetic
field around a loop acquires a phase equal to $\mp\frac12$ the solid angle enclosed by the field direction. The metric and the curvature are the two halves of quantum geometry; the integral of the Berry curvature over a closed surface is $2\pi$ times an
integer, the Chern number, which underlies topological phases (quantum Hall effect, topological insulators and topological photonics). The quantum metric is measurable, for example from the response to periodic modulations or from the
anomalous trajectories of wave packets, and has been measured for photons in microcavities.

\subsection{Mixed states: Bures metric and quantum Fisher information}
For mixed states, the Bures distance $d_B^2=2(1-\sqrt F)$ defines, to second order, the Bures metric $\dd s_B^2=\frac14F_Q\,\dd\lambda^2$ for a single parameter. The quantity $F_Q$ is the \emph{quantum Fisher information}; for pure states $F_Q=4g_{\lambda\lambda}$.
It sets the ultimate precision of the estimation of $\lambda$ from $\nu$ repetitions of a measurement (quantum Cramér--Rao bound):
\begin{equation}
  \Delta\lambda\ge\frac1{\sqrt{\nu F_Q}} .
\end{equation}
The interpretation is geometric: a parameter can be estimated precisely if the states move fast (in the Fubini--Study sense) when it changes. For a phase $\lambda$ imprinted by $\e^{-\ii\lambda G}$, $F_Q=4\Delta G^2$. With $n$ independent qubits and
$G=\frac12\sum_iZ_i$, $\Delta G^2\le n/4$ and $\Delta\lambda\ge1/\sqrt{\nu n}$ (the standard quantum limit). With the entangled GHZ state $(\ket{0\dots0}+\ket{1\dots1})/\sqrt2$, $\Delta G^2=n^2/4$ and $\Delta\lambda\ge1/(n\sqrt\nu)$ (the
\emph{Heisenberg limit}): entanglement improves metrology by a factor $\sqrt n$, the principle behind quantum-enhanced atomic clocks and interferometers. Phase estimation (\cref{ch:shor}) reaches the same scaling.

\section{Exercises}
\begin{exercise}
Compute $F$ and $D$ for $\hat\rho=\proj0$ and $\hat\sigma=I/2$, and check the Fuchs--van de Graaf inequalities.
\end{exercise}
\begin{exercise}
Show that the fidelity of a Werner state $\hat\rho_W$ (\cref{ex:werner}) with the singlet is $F=\frac{1+3p}4$. Deduce that $F>1/2$ iff $p>1/3$, i.e.\ iff the state is entangled.
\end{exercise}
\begin{exercise}
Compute the Berry phase of the ground state of $H=-\vec B\cdot\vec\sigma$ when $\vec B$ rotates once around the $z$ axis at fixed polar angle $\theta$. (Answer: $\pm\pi(1-\cos\theta)$, half the solid angle.)
\end{exercise}

\chapter{Formula sheet}\label{app:formulas}

\section*{Pauli matrices and rotations}
\begin{align*}
  &\sigma_j\sigma_k=\delta_{jk}I+\ii\varepsilon_{jkl}\sigma_l,\qquad(\vec a\cdot\vec\sigma)(\vec b\cdot\vec\sigma)=\vec a\cdot\vec b\,I+\ii(\vec a\times\vec b)\cdot\vec\sigma,\\
  &\e^{-\ii\alpha\vec n\cdot\vec\sigma/2}=\cos\tfrac\alpha2\,I-\ii\sin\tfrac\alpha2\,\vec n\cdot\vec\sigma,\qquad HXH=Z,\quad HZH=X,\quad HYH=-Y,\\
  &\hat\rho=\tfrac12(I+\vec r\cdot\vec\sigma),\qquad\tr\hat\rho^2=\tfrac12(1+|\vec r|^2),\qquad P_\pm^{(\vec n)}=\tfrac12(I\pm\vec n\cdot\vec\sigma).
\end{align*}

\section*{Tensor products}
\begin{align*}
  &(A\otimes B)(C\otimes D)=AC\otimes BD,\qquad\tr(A\otimes B)=\tr A\tr B,\\ &\tr_B(\ketbra{a_1}{a_2}\otimes\ketbra{b_1}{b_2})=\ketbra{a_1}{a_2}\braket{b_2}{b_1},\\
  &\CNOT=\proj0\otimes I+\proj1\otimes X,\qquad\mathrm{CZ}=\proj0\otimes I+\proj1\otimes Z,\\ &\vec\sigma_1\cdot\vec\sigma_2=2\,\mathrm{SWAP}-I.
\end{align*}

\section*{Bell states}
\[
  \ket{\Phi^\pm}=\tfrac1{\sqrt2}(\ket{00}\pm\ket{11}),\quad\ket{\Psi^\pm}=\tfrac1{\sqrt2}(\ket{01}\pm\ket{10}),\quad
  \mel{\Psi^-}{(\vec a\cdot\vec\sigma)\otimes(\vec b\cdot\vec\sigma)}{\Psi^-}=-\vec a\cdot\vec b .
\]

\section*{Hadamard and Fourier transforms}
\[
  H^{\otimes n}\ket x=\frac1{\sqrt{2^n}}\sum_z(-1)^{x\cdot z}\ket z,\qquad U_{\rm QFT}\ket x=\frac1{\sqrt Q}\sum_y\e^{2\pi\ii xy/Q}\ket y,\]
\[
  \sum_{y=0}^{Q-1}\e^{2\pi\ii ky/Q}=Q\,\delta_{k\equiv0\ (\mathrm{mod}\ Q)} .
\]

\section*{Harmonic oscillator and thermal states}
\[
  \hat a\ket n=\sqrt n\ket{n-1},\quad\hat a^\dagger\ket n=\sqrt{n+1}\ket{n+1},\quad\bar n=\frac1{\e^{\hbar\omega/k_BT}-1},\quad p_n=\frac{\bar n^n}{(\bar n+1)^{n+1}} .
\]

\section*{Master equation}
\[
  \dot{\hat\rho}=-\tfrac\ii\hbar[H,\hat\rho]+\sum_k\left(L_k\hat\rho L_k^\dagger-\tfrac12\{L_k^\dagger L_k,\hat\rho\}\right),\qquad\frac1{T_2}=\frac1{2T_1}+\frac1{T_\varphi},\]
\[
  \frac{\dd\expval A}{\dd t}=\tfrac\ii\hbar\expval{[H,A]}+\sum_k\expval{L_k^\dagger AL_k-\tfrac12\{L_k^\dagger L_k,A\}}.
\]

\section*{Algorithms}
\begin{align*}
  &\text{Grover: }\sin\tfrac\theta2=\sqrt{M/N},\quad P_k=\sin^2\tfrac{(2k+1)\theta}2,\quad k_{\rm opt}\approx\tfrac\pi4\sqrt{N/M}.\\
  &\text{Shor: }a^r\equiv1\ (\mathrm{mod}\ N),\quad\text{factors }\gcdf(a^{r/2}\pm1,N),\quad N^2\le Q<2N^2,\quad y/Q\approx k/r.\\
  &\text{Phase estimation: }P(\text{best }t\text{-bit estimate})\ge4/\pi^2.\\
  &\text{Adiabatic: }t_f\gg\hbar\max\norm{\dd H/\dd s}/\delta_{\min}^2,\qquad\text{Landau--Zener: }P=\exp(-2\pi V^2/\hbar|\alpha|).\\
  &\text{QUBO}\to\text{Ising: }x_i=\tfrac{1-s_i}2,\quad J_{ij}=\tfrac{Q_{ij}}4,\quad h_i=-\tfrac{Q_{ii}}2-\tfrac14\textstyle\sum_{j\neq i}Q_{ij}.
\end{align*}

\section*{Euclid's algorithm and continued fractions}
$\gcdf(a,b)=\gcdf(b,a\bmod b)$, $\gcdf(a,0)=a$; e.g.\ $\gcdf(12,15)=\gcdf(15,12)=\gcdf(12,3)=\gcdf(3,0)=3$. The continued fraction of $x\in(0,1)$ is obtained by iterating $a_k=\lfloor1/x_k\rfloor$, $x_{k+1}=1/x_k-a_k$; the
convergents $p_k/q_k$ satisfy $p_k=a_kp_{k-1}+p_{k-2}$, $q_k=a_kq_{k-1}+q_{k-2}$ (with the convention $x=[0;a_1,a_2,\dots]$, $p_{-1}=1,q_{-1}=0,p_0=0,q_0=1$). If $|x-p/q|<1/(2q^2)$, then $p/q$ is a convergent of $x$.

\section*{Physical constants and scales}
\[
  h=6.626\times10^{-34}\,\mathrm{J\,s},\quad k_B=1.381\times10^{-23}\,\mathrm{J/K},\quad\Phi_0=\frac h{2e}=2.07\times10^{-15}\,\mathrm{Wb},\]
\[
  h\times1\,\mathrm{GHz}=k_B\times48\,\mathrm{mK}=4.14\,\mu\mathrm{eV}.
\]


\backmatter
\begin{thebibliography}{99}
\addcontentsline{toc}{chapter}{References}

\bibitem{Aaronson2015} S.~Aaronson. Read the fine print. \emph{Nature Physics}, 11:291--293, 2015.
\bibitem{Aharonov2007} D.~Aharonov, W.~van Dam, J.~Kempe, Z.~Landau, S.~Lloyd, and O.~Regev. Adiabatic quantum computation is equivalent to standard quantum computation. \emph{SIAM Journal on Computing}, 37(1):166--194, 2007.
\bibitem{Altshuler2010} B.~Altshuler, H.~Krovi, and J.~Roland. Anderson localization makes adiabatic quantum optimization fail. \emph{PNAS}, 107(28):12446--12450, 2010.
\bibitem{Aspect1982} A.~Aspect, P.~Grangier, and G.~Roger. Experimental realization of Einstein--Podolsky--Rosen--Bohm Gedankenexperiment: a new violation of Bell's inequalities. \emph{Physical Review Letters}, 49(2):91--94, 1982.
\bibitem{Barahona1982} F.~Barahona. On the computational complexity of Ising spin glass models. \emph{Journal of Physics A}, 15(10):3241--3253, 1982.
\bibitem{BB84} C.~H.~Bennett and G.~Brassard. Quantum cryptography: public key distribution and coin tossing. In \emph{Proceedings of the IEEE International Conference on Computers, Systems and Signal Processing}, pages 175--179, Bangalore, 1984. Reprinted in \emph{Theor.\ Comput.\ Sci.} 560:7--11, 2014.
\bibitem{BBBV1997} C.~H.~Bennett, E.~Bernstein, G.~Brassard, and U.~Vazirani. Strengths and weaknesses of quantum computing. \emph{SIAM Journal on Computing}, 26(5):1510--1523, 1997.
\bibitem{Bell1964} J.~S.~Bell. On the Einstein Podolsky Rosen paradox. \emph{Physics Physique Fizika}, 1(3):195--200, 1964.
\bibitem{Bennett1993} C.~H.~Bennett, G.~Brassard, C.~Crépeau, R.~Jozsa, A.~Peres, and W.~K.~Wootters. Teleporting an unknown quantum state via dual classical and Einstein--Podolsky--Rosen channels. \emph{Physical Review Letters}, 70(13):1895--1899, 1993.
\bibitem{Biamonte2017} J.~Biamonte, P.~Wittek, N.~Pancotti, P.~Rebentrost, N.~Wiebe, and S.~Lloyd. Quantum machine learning. \emph{Nature}, 549:195--202, 2017.
\bibitem{Bouwmeester1997} D.~Bouwmeester, J.-W.~Pan, K.~Mattle, M.~Eibl, H.~Weinfurter, and A.~Zeilinger. Experimental quantum teleportation. \emph{Nature}, 390:575--579, 1997.
\bibitem{CHSH1969} J.~F.~Clauser, M.~A.~Horne, A.~Shimony, and R.~A.~Holt. Proposed experiment to test local hidden-variable theories. \emph{Physical Review Letters}, 23(15):880--884, 1969.
\bibitem{CiracZoller1995} J.~I.~Cirac and P.~Zoller. Quantum computations with cold trapped ions. \emph{Physical Review Letters}, 74(20):4091--4094, 1995.
\bibitem{Cohen1981} D.~Cohen. On holy wars and a plea for peace. \emph{Computer}, 14(10):48--54, 1981. First circulated as Internet Experiment Note IEN 137, 1980.
\bibitem{Cook1971} S.~A.~Cook. The complexity of theorem-proving procedures. In \emph{Proceedings of the Third Annual ACM Symposium on Theory of Computing (STOC '71)}, pages 151--158, 1971.
\bibitem{Cory1998} D.~G.~Cory, M.~D.~Price, W.~Maas, E.~Knill, R.~Laflamme, W.~H.~Zurek, T.~F.~Havel, and S.~S.~Somaroo. Experimental quantum error correction. \emph{Physical Review Letters}, 81(10):2152--2155, 1998.
\bibitem{Cuccaro2004} S.~A.~Cuccaro, T.~G.~Draper, S.~A.~Kutin, and D.~P.~Moulton. A new quantum ripple-carry addition circuit. arXiv:quant-ph/0410184, 2004.
\bibitem{DatePotok2021} P.~Date and T.~Potok. Adiabatic quantum linear regression. \emph{Scientific Reports}, 11:21905, 2021.
\bibitem{DeutschJozsa1992} D.~Deutsch and R.~Jozsa. Rapid solution of problems by quantum computation. \emph{Proceedings of the Royal Society of London A}, 439:553--558, 1992.
\bibitem{Dieks1982} D.~Dieks. Communication by EPR devices. \emph{Physics Letters A}, 92(6):271--272, 1982.
\bibitem{DiVincenzo2000} D.~P.~DiVincenzo. The physical implementation of quantum computation. \emph{Fortschritte der Physik}, 48(9--11):771--783, 2000.
\bibitem{EPR1935} A.~Einstein, B.~Podolsky, and N.~Rosen. Can quantum-mechanical description of physical reality be considered complete? \emph{Physical Review}, 47(10):777--780, 1935.
\bibitem{Ekert1991} A.~K.~Ekert. Quantum cryptography based on Bell's theorem. \emph{Physical Review Letters}, 67(6):661--663, 1991.
\bibitem{Farhi2001} E.~Farhi, J.~Goldstone, S.~Gutmann, J.~Lapan, A.~Lundgren, and D.~Preda. A quantum adiabatic evolution algorithm applied to random instances of an NP-complete problem. \emph{Science}, 292:472--475, 2001.
\bibitem{FVH1957} R.~P.~Feynman, F.~L.~Vernon, and R.~W.~Hellwarth. Geometrical representation of the Schrödinger equation for solving maser problems. \emph{Journal of Applied Physics}, 28(1):49--52, 1957.
\bibitem{Geman1984} S.~Geman and D.~Geman. Stochastic relaxation, Gibbs distributions, and the Bayesian restoration of images. \emph{IEEE Trans.\ Pattern Analysis and Machine Intelligence}, PAMI-6(6):721--741, 1984.
\bibitem{Grover1996} L.~K.~Grover. A fast quantum mechanical algorithm for database search. In \emph{Proceedings of the 28th Annual ACM Symposium on Theory of Computing (STOC '96)}, pages 212--219, 1996.
\bibitem{HHL2009} A.~W.~Harrow, A.~Hassidim, and S.~Lloyd. Quantum algorithm for linear systems of equations. \emph{Physical Review Letters}, 103:150502, 2009.
\bibitem{Havlicek2019} V.~Havlíček, A.~D.~Córcoles, K.~Temme, A.~W.~Harrow, A.~Kandala, J.~M.~Chow, and J.~M.~Gambetta. Supervised learning with quantum-enhanced feature spaces. \emph{Nature}, 567:209--212, 2019.
\bibitem{HenYoung2011} I.~Hen and A.~P.~Young. Exponential complexity of the quantum adiabatic algorithm for certain satisfiability problems. \emph{Physical Review E}, 84:061152, 2011.
\bibitem{Horodecki1996} M.~Horodecki, P.~Horodecki, and R.~Horodecki. Separability of mixed states: necessary and sufficient conditions. \emph{Physics Letters A}, 223:1--8, 1996.
\bibitem{HJW1993} L.~P.~Hughston, R.~Jozsa, and W.~K.~Wootters. A complete classification of quantum ensembles having a given density matrix. \emph{Physics Letters A}, 183:14--18, 1993.
\bibitem{Jacak2021} M.~M.~Jacak, P.~Jóźwiak, J.~Niemczuk, and J.~E.~Jacak. Quantum generators of random numbers. \emph{Scientific Reports}, 11:16108, 2021.
\bibitem{Kadowaki1998} T.~Kadowaki and H.~Nishimori. Quantum annealing in the transverse Ising model. \emph{Physical Review E}, 58(5):5355--5363, 1998.
\bibitem{Kirkpatrick1983} S.~Kirkpatrick, C.~D.~Gelatt, and M.~P.~Vecchi. Optimization by simulated annealing. \emph{Science}, 220:671--680, 1983.
\bibitem{Koch2007} J.~Koch, T.~M.~Yu, J.~Gambetta, A.~A.~Houck, D.~I.~Schuster, J.~Majer, A.~Blais, M.~H.~Devoret, S.~M.~Girvin, and R.~J.~Schoelkopf. Charge-insensitive qubit design derived from the Cooper pair box. \emph{Physical Review A}, 76:042319, 2007.
\bibitem{Koshka2019} Y.~Koshka and M.~A.~Novotny. Comparison of D-Wave quantum annealing and classical simulated annealing for local minima determination. arXiv:1911.03338, 2019.
\bibitem{Lloyd1996} S.~Lloyd. Universal quantum simulators. \emph{Science}, 273:1073--1078, 1996.
\bibitem{Lucas2014} A.~Lucas. Ising formulations of many NP problems. \emph{Frontiers in Physics}, 2:5, 2014.
\bibitem{MoritaNishimori2008} S.~Morita and H.~Nishimori. Mathematical foundation of quantum annealing. \emph{Journal of Mathematical Physics}, 49:125210, 2008.
\bibitem{NielsenChuang} M.~A.~Nielsen and I.~L.~Chuang. \emph{Quantum Computation and Quantum Information}. Cambridge University Press, 10th anniversary edition, 2010.
\bibitem{Peres1996} A.~Peres. Separability criterion for density matrices. \emph{Physical Review Letters}, 77(8):1413--1415, 1996.
\bibitem{RolandCerf2002} J.~Roland and N.~J.~Cerf. Quantum search by local adiabatic evolution. \emph{Physical Review A}, 65:042308, 2002.
\bibitem{Shor1995} P.~W.~Shor. Scheme for reducing decoherence in quantum computer memory. \emph{Physical Review A}, 52(4):R2493--R2496, 1995.
\bibitem{Shor1997} P.~W.~Shor. Polynomial-time algorithms for prime factorization and discrete logarithms on a quantum computer. \emph{SIAM Journal on Computing}, 26(5):1484--1509, 1997.
\bibitem{Tang2019} E.~Tang. A quantum-inspired classical algorithm for recommendation systems. In \emph{Proceedings of the 51st Annual ACM SIGACT Symposium on Theory of Computing (STOC 2019)}, pages 217--228, 2019.
\bibitem{Vandersypen2001} L.~M.~K.~Vandersypen, M.~Steffen, G.~Breyta, C.~S.~Yannoni, M.~H.~Sherwood, and I.~L.~Chuang. Experimental realization of Shor's quantum factoring algorithm using nuclear magnetic resonance. \emph{Nature}, 414:883--887, 2001.
\bibitem{VBE1996} V.~Vedral, A.~Barenco, and A.~Ekert. Quantum networks for elementary arithmetic operations. \emph{Physical Review A}, 54(1):147--153, 1996.
\bibitem{vonNeumann1951} J.~von Neumann. Various techniques used in connection with random digits. In \emph{Monte Carlo Method}, National Bureau of Standards Applied Mathematics Series 12, pages 36--38, 1951.
\bibitem{Wiebe2012} N.~Wiebe, D.~Braun, and S.~Lloyd. Quantum algorithm for data fitting. \emph{Physical Review Letters}, 109:050505, 2012.
\bibitem{WoottersZurek1982} W.~K.~Wootters and W.~H.~Zurek. A single quantum cannot be cloned. \emph{Nature}, 299:802--803, 1982.
\bibitem{Young2008} A.~P.~Young, S.~Knysh, and V.~N.~Smelyanskiy. Size dependence of the minimum excitation gap in the quantum adiabatic algorithm. \emph{Physical Review Letters}, 101:170503, 2008.

\paragraph{Further reading.} Besides Nielsen and Chuang \cite{NielsenChuang}: J.~Preskill, \emph{Lecture notes on quantum computation} (Caltech, Physics 219); R.~P.~Feynman, \emph{The Feynman Lectures on Physics}, vol.~III (two-state systems, ammonia maser, Josephson junction);
S.~Haroche and J.-M.~Raimond, \emph{Exploring the Quantum} (Oxford, 2006) for cavity QED and decoherence; A.~Blais et al., ``Circuit quantum electrodynamics'', \emph{Rev.\ Mod.\ Phys.}\ 93, 025005 (2021); H.-P.~Breuer and F.~Petruccione,
\emph{The Theory of Open Quantum Systems} (Oxford, 2002); T.~Albash and D.~A.~Lidar, ``Adiabatic quantum computation'', \emph{Rev.\ Mod.\ Phys.}\ 90, 015002 (2018).
\end{thebibliography}


\end{document}
